%  LaTeX support: latex@mdpi.com 
%  For support, please attach all files needed for compiling as well as the log file, and specify your operating system, LaTeX version, and LaTeX editor.

%=================================================================
%\documentclass[journal,article,submit,pdftex,moreauthors]{Definitions/mdpi} 
%\documentclass[preprints]{Definitions/mdpi} 
\documentclass[preprints,accept]{Definitions/mdpi}
% For posting an early version of this manuscript as a preprint, you may use "preprints" as the journal. Changing "submit" to "accept" before posting will remove line numbers.

% Below journals will use APA reference format:
% admsci, behavsci, businesses, econometrics, economies, education, ejihpe, games, humans, ijfs, journalmedia, jrfm, languages, psycholint, publications, tourismhosp, youth

% Below journals will use Chicago reference format:
% arts, genealogy, histories, humanities, jintelligence, laws, literature, religions, risks, socsci

%--------------------
% Class Options:
%--------------------
%----------
% journal
%----------
% Choose between the following MDPI journals:
% accountaudit, acoustics, actuators, addictions, adhesives, admsci, adolescents, aerobiology, aerospace, agriculture, agriengineering, agrochemicals, agronomy, ai, air, algorithms, allergies, alloys, amh, analytica, analytics, anatomia, anesthres, animals, antibiotics, antibodies, antioxidants, applbiosci, appliedchem, appliedmath, appliedphys, applmech, applmicrobiol, applnano, applsci, aquacj, architecture, arm, arthropoda, arts, asc, asi, astronomy, atmosphere, atoms, audiolres, automation, axioms, bacteria, batteries, bdcc, behavsci, beverages, biochem, bioengineering, biologics, biology, biomass, biomechanics, biomed, biomedicines, biomedinformatics, biomimetics, biomolecules, biophysica, biosensors, biosphere, biotech, birds, blockchains, bloods, blsf, brainsci, breath, buildings, businesses, cancers, carbon, cardiogenetics, catalysts, cells, ceramics, challenges, chemengineering, chemistry, chemosensors, chemproc, children, chips, cimb, civileng, cleantechnol, climate, clinbioenerg, clinpract, clockssleep, cmd, cmtr, coasts, coatings, colloids, colorants, commodities, complications, compounds, computation, computers, condensedmatter, conservation, constrmater, cosmetics, covid, crops, cryo, cryptography, crystals, csmf, ctn, curroncol, cyber, dairy, data, ddc, dentistry, dermato, dermatopathology, designs, devices, diabetology, diagnostics, dietetics, digital, disabilities, diseases, diversity, dna, drones, dynamics, earth, ebj, ecm, ecologies, econometrics, economies, education, eesp, ejihpe, electricity, electrochem, electronicmat, electronics, encyclopedia, endocrines, energies, eng, engproc, ent, entomology, entropy, environments, epidemiologia, epigenomes, esa, est, famsci, fermentation, fibers, fintech, fire, fishes, fluids, foods, forecasting, forensicsci, forests, fossstud, foundations, fractalfract, fuels, future, futureinternet, futureparasites, futurepharmacol, futurephys, futuretransp, galaxies, games, gases, gastroent, gastrointestdisord, gastronomy, gels, genealogy, genes, geographies, geohazards, geomatics, geometry, geosciences, geotechnics, geriatrics, glacies, grasses, greenhealth, gucdd, hardware, hazardousmatters, healthcare, hearts, hemato, hematolrep, heritage, higheredu, highthroughput, histories, horticulturae, hospitals, humanities, humans, hydrobiology, hydrogen, hydrology, hygiene, idr, iic, ijerph, ijfs, ijgi, ijmd, ijms, ijns, ijpb, ijt, ijtm, ijtpp, ime, immuno, informatics, information, infrastructures, inorganics, insects, instruments, inventions, iot, j, jal, jcdd, jcm, jcp, jcs, jcto, jdad, jdb, jeta, jfb, jfmk, jimaging, jintelligence, jlpea, jmahp, jmmp, jmms, jmp, jmse, jne, jnt, jof, joitmc, joma, jop, jor, journalmedia, jox, jpbi, jpm, jrfm, jsan, jtaer, jvd, jzbg, kidney, kidneydial, kinasesphosphatases, knowledge, labmed, laboratories, land, languages, laws, life, lights, limnolrev, lipidology, liquids, literature, livers, logics, logistics, lubricants, lymphatics, machines, macromol, magnetism, magnetochemistry, make, marinedrugs, materials, materproc, mathematics, mca, measurements, medicina, medicines, medsci, membranes, merits, metabolites, metals, meteorology, methane, metrics, metrology, micro, microarrays, microbiolres, microelectronics, micromachines, microorganisms, microplastics, microwave, minerals, mining, mmphys, modelling, molbank, molecules, mps, msf, mti, multimedia, muscles, nanoenergyadv, nanomanufacturing, nanomaterials, ncrna, ndt, network, neuroglia, neurolint, neurosci, nitrogen, notspecified, nursrep, nutraceuticals, nutrients, obesities, oceans, ohbm, onco, oncopathology, optics, oral, organics, organoids, osteology, oxygen, parasites, parasitologia, particles, pathogens, pathophysiology, pediatrrep, pets, pharmaceuticals, pharmaceutics, pharmacoepidemiology, pharmacy, philosophies, photochem, photonics, phycology, physchem, physics, physiologia, plants, plasma, platforms, pollutants, polymers, polysaccharides, populations, poultry, powders, preprints, proceedings, processes, prosthesis, proteomes, psf, psych, psychiatryint, psychoactives, psycholint, publications, purification, quantumrep, quaternary, qubs, radiation, reactions, realestate, receptors, recycling, regeneration, religions, remotesensing, reports, reprodmed, resources, rheumato, risks, robotics, rsee, ruminants, safety, sci, scipharm, sclerosis, seeds, sensors, separations, sexes, signals, sinusitis, siuj, skins, smartcities, sna, societies, socsci, software, soilsystems, solar, solids, spectroscj, sports, standards, stats, std, stresses, surfaces, surgeries, suschem, sustainability, symmetry, synbio, systems, tae, targets, taxonomy, technologies, telecom, test, textiles, thalassrep, therapeutics, thermo, timespace, tomography, tourismhosp, toxics, toxins, transplantology, transportation, traumacare, traumas, tropicalmed, universe, urbansci, uro, vaccines, vehicles, venereology, vetsci, vibration, virtualworlds, viruses, vision, waste, water, wem, wevj, wild, wind, women, world, youth, zoonoticdis

%---------
% article
%---------
% The default type of manuscript is "article", but can be replaced by: 
% abstract, addendum, article, benchmark, book, bookreview, briefcommunication, briefreport, casereport, changes, clinicopathologicalchallenge, comment, commentary, communication, conceptpaper, conferenceproceedings, correction, conferencereport, creative, datadescriptor, discussion, entry, expressionofconcern, extendedabstract, editorial, essay, erratum, fieldguide, hypothesis, interestingimages, letter, meetingreport, monograph, newbookreceived, obituary, opinion, proceedingpaper, projectreport, reply, retraction, review, perspective, protocol, shortnote, studyprotocol, supfile, systematicreview, technicalnote, viewpoint, guidelines, registeredreport, tutorial,  giantsinurology, urologyaroundtheworld
% supfile = supplementary materials

%----------
% submit
%----------
% The class option "submit" will be changed to "accept" by the Editorial Office when the paper is accepted. This will only make changes to the frontpage (e.g., the logo of the journal will get visible), the headings, and the copyright information. Also, line numbering will be removed. Journal info and pagination for accepted papers will also be assigned by the Editorial Office.

%------------------
% moreauthors
%------------------
% If there is only one author the class option oneauthor should be used. Otherwise use the class option moreauthors.

%---------
% pdftex
%---------
% The option pdftex is for use with pdfLaTeX. Remove "pdftex" for (1) compiling with LaTeX & dvi2pdf (if eps figures are used) or for (2) compiling with XeLaTeX.

%=================================================================
% MDPI internal commands - do not modify
\firstpage{1} 
\makeatletter 
\setcounter{page}{\@firstpage} 
\makeatother
\pubvolume{1}
\issuenum{1}
\articlenumber{0}
\pubyear{2025}
\copyrightyear{2025}
%\externaleditor{Firstname Lastname} % More than 1 editor, please add `` and '' before the last editor name
\datereceived{ } 
\daterevised{ } % Comment out if no revised date
\dateaccepted{ } 
\datepublished{ } 
%\datecorrected{} % For corrected papers: "Corrected: XXX" date in the original paper.
%\dateretracted{} % For retracted papers: "Retracted: XXX" date in the original paper.
\hreflink{https://doi.org/} % If needed use \linebreak
%\doinum{}
%\pdfoutput=1 % Uncommented for upload to arXiv.org
%\CorrStatement{yes}  % For updates
%\longauthorlist{yes} % For many authors that exceed the left citation part

%=================================================================
% Add packages and commands here. The following packages are loaded in our class file: fontenc, inputenc, calc, indentfirst, fancyhdr, graphicx, epstopdf, lastpage, ifthen, float, amsmath, amssymb, lineno, setspace, enumitem, mathpazo, booktabs, titlesec, etoolbox, tabto, xcolor, colortbl, soul, multirow, microtype, tikz, totcount, changepage, attrib, upgreek, array, tabularx, pbox, ragged2e, tocloft, marginnote, marginfix, enotez, amsthm, natbib, hyperref, cleveref, scrextend, url, geometry, newfloat, caption, draftwatermark, seqsplit
% cleveref: load \crefname definitions after \begin{document}

%=================================================================
% Please use the following mathematics environments: Theorem, Lemma, Corollary, Proposition, Characterization, Property, Problem, Example, ExamplesandDefinitions, Hypothesis, Remark, Definition, Notation, Assumption
%% For proofs, please use the proof environment (the amsthm package is loaded by the MDPI class).


% ---------- Encoding / fonts ----------
%\usepackage[T1]{fontenc}
%\usepackage[utf8]{inputenc}
\usepackage{lmodern}        % Latin Modern: reliable math/text glyph coverage

% ---------- Layout / typography ----------
%\usepackage[margin=1in]{geometry}
%\usepackage{microtype}      % better PDF text shaping; helps readability
%\usepackage{xcolor}

% ---------- Math packages ----------
%\usepackage{amsmath,amssymb,amsthm,mathtools}
\usepackage{bm}             % bold math
\usepackage{stmaryrd}       % provides \llbracket \rrbracket
\usepackage{mathrsfs}       % optional: \mathscr

% ---------- Lists / diagrams ----------
\usepackage{enumitem}       % consistent lists
\usepackage{tikz-cd}        % commutative diagrams

% ---------- Hyperlinks / bookmarks ----------
%\usepackage{hyperref}
%\usepackage{bookmark}
%\usepackage{natbib}
\usepackage{comment}
\usepackage{tikz}
\usetikzlibrary{positioning}
%\hypersetup{
%  colorlinks=true,
%  linkcolor=blue,
%  citecolor=blue,
%  urlcolor=blue
%}

% ============================================================
% Theorem environments
% ============================================================

%\newtheorem{theorem}{Theorem}[section]
\newtheorem{lemma}[theorem]{Lemma}
\newtheorem{proposition}[theorem]{Proposition}
\newtheorem{corollary}[theorem]{Corollary}

\theoremstyle{definition}
\newtheorem{definition}[theorem]{Definition}
\newtheorem{remark}[theorem]{Remark}
\newtheorem{example}[theorem]{Example}
\newtheorem{assumption}[theorem]{Assumption}

% ============================================================
% FRFP Macros (use everywhere; avoids symbol drift)
% ============================================================

\newcommand{\FRFP}{FRFP}
\newcommand{\Cfull}{C_{\mathrm{full}}}
\newcommand{\Ecat}{\mathcal{E}} % explicit category
\newcommand{\Tcat}{\mathcal{T}} % tacit category
\newcommand{\Ncat}{\mathcal{N}} % navigation groupoid
\newcommand{\Scat}{\mathcal{S}} % shapes
\newcommand{\TDG}{\mathsf{TDG}} % Task Decomposition Grammar
\newcommand{\RB}{\mathsf{RB}}   % Representational Backflow
\newcommand{\RI}{\mathsf{RI}}
\newcommand{\EC}{\mathsf{EC}}
\newcommand{\ED}{\mathsf{ED}}
\newcommand{\TE}{\mathsf{TE}}
\newcommand{\HFD}{\mathsf{HFD}}
\newcommand{\Req}{\mathsf{Req}}
\newcommand{\Hom}{\mathrm{Hom}}
\newcommand{\Obj}{\mathrm{Ob}}
\newcommand{\BRANCH}{\mathsf{BRANCH}}
\newcommand{\PARALLEL}{\mathsf{PARALLEL}}
\newcommand{\LOOP}{\mathsf{LOOP}}

% aliases for explicit / tacit space used in the appendix
\newcommand{\Espace}{\mathcal{E}} % explicit space
\newcommand{\Tspace}{\mathcal{T}} % tacit space

% encoding / decoding
\newcommand{\enc}{\mathsf{enc}}
\newcommand{\dec}{\mathsf{dec}}

% semantic RB and free-category notation
\newcommand{\RBsharp}{\RB^{\sharp}}
\newcommand{\FreeCat}{\mathsf{FreeCat}}
\newcommand{\Ufun}{U_E}

% shape constructors (non-conflicting names)
\newcommand{\chain}{\mathsf{chain}}
\newcommand{\branch}{\mathsf{branch}}
\newcommand{\loopshape}{\mathsf{loop}}

% frameworks category name: \Frm_1
\newcommand{\Frm}{\mathsf{Frm}} % so \Frm_1 works

% example navigation generators (purely illustrative)
\newcommand{\TRYMODEL}{\mathsf{TRYMODEL}}
\newcommand{\ADDASSUMPTION}{\mathsf{ADDASSUMPTION}}
\newcommand{\DROPASSUMPTION}{\mathsf{DROPASSUMPTION}}
\newcommand{\MERGE}{\mathsf{MERGE}}
\newcommand{\RESTART}{\mathsf{RESTART}}

\newcommand{\id}{\mathrm{id}}
\newcommand{\shape}{\mathsf{shape}}

\newcommand{\Inst}{\mathsf{Inst}}
\newcommand{\HP}{\mathcal{HP}}
\newcommand{\Sem}{\mathsf{Sem}}
\newcommand{\Esem}{\mathcal{E}_{\mathrm{sem}}}

\newcommand{\Empty}{\varnothing}
\newcommand{\Eobj}{E}
\newcommand{\Tobj}{T}

% Semantic brackets
\newcommand{\sem}[1]{\llbracket #1 \rrbracket}

% Semidirect / Grothendieck notation
\newcommand{\NltimesE}{\Ncat \ltimes \Ecat}
\newcommand{\NltimesEsem}{\Ncat \ltimes \Esem}

% Logic-ish convenience
\newcommand{\defeq}{\mathrel{\vcentcolon=}}
\newcommand{\eqdef}{\mathrel{=\vcentcolon}}

%=================================================================
% Full title of the paper (Capitalized)
\Title{Foundational Reasoning and Feedback Protocol (FRFP): A Formal Architecture for Long-Horizon Human--AI Collaboration}

% MDPI internal command: Title for citation in the left column
\TitleCitation{Title}

% Author Orchid ID: enter ID or remove command
\newcommand{\orcidauthorA}{0009-0008-2862-5041} % Add \orcidA{} behind the author's name
%\newcommand{\orcidauthorB}{0000-0000-0000-000X} % Add \orcidB{} behind the author's name

% Authors, for the paper (add full first names)
\Author{Vijaya Nadendla $^{1,\dagger}$\orcidA{}}

%\longauthorlist{yes}

% MDPI internal command: Authors, for metadata in PDF
\AuthorNames{Vijaya Nadendla}

% MDPI internal command: Authors, for citation in the left column, only choose below one of them according to the journal style
% If this is a APA style journal 
% (admsci, behavsci, businesses, econometrics, economies, education, ejihpe, games, humans, ijfs, journalmedia, jrfm, languages, psycholint, publications, tourismhosp, youth): 
% Lastname, F., Lastname, F., \& Lastname, F.

% If this is a Chicago style journal 
% (arts, genealogy, histories, humanities, jintelligence, laws, literature, religions, risks, socsci): 
% Lastname, Firstname, Firstname Lastname, and Firstname Lastname.

% If this is a ACS style journal (Except for the above Chicago and APA journals, all others are in the ACS format): 
% Lastname, F.; Lastname, F.; Lastname, F.
\isAPAStyle{%
       \AuthorCitation{Lastname, F., Lastname, F., \& Lastname, F.}
         }{%
        \isChicagoStyle{%
        \AuthorCitation{Lastname, Firstname, Firstname Lastname, and Firstname Lastname.}
        }{
        \AuthorCitation{Lastname, F.; Lastname, F.; Lastname, F.}
        }
}

% Affiliations / Addresses (Add [1] after \address if there is only one affiliation.)
\address[1]{%
$^{1}$ \quad Affiliation; Independent Researcher}

% Contact information of the corresponding author
\corres{Correspondence: vnadendla@gmail.com; Tel.:   +01-609-510-4995}

% Current address and/or shared authorship
\firstnote{Current address: 51 Linwood Dr Monroe Township NJ 08831.}  % Current address should not be the same as any items in the Affiliation section.

% The commands \thirdnote{} till \eighthnote{} are available for further notes

%\simplesumm{} % Simple summary

%\conference{} % An extended version of a conference paper
\nolinenumbers
% Abstract (Do not insert blank lines, i.e. \\) 
\abstract{This paper investigates conditions under which AI systems in high-stakes scientific and regulatory workflows can be governed and trusted using only automated checks, and when human tacit judgment remains structurally necessary. We introduce FRFP, a categorical protocol model that separates explicit artefacts (prompts, tools, reports) from tacit states (understanding, judgment), with a single boundary and a small set of epistemic primitives, and show that this “kernel” is initial/canonical among architectures satisfying our explicit–tacit axioms. Building on this, we develop a dynamic model with finite tacit capacity, degradation, and requirement maps, and analyze stochastic pipelines via hazard profiles and safe-horizon functionals. Under the FRFP assumptions, we prove conditional inevitability theorems: along sufficiently long runs, some accepted output will fall below its grounding requirement, and no explicit-only detector, oracle, or governance operator can guarantee the absence of hallucination—i.e., can guarantee correctness across all admissible pipelines and populations. We extend the framework to collective and institutional settings and derive structural non-automation results for governance. We close with protocol templates, case sketches, and FRFP-compatible instruction sets that demonstrate how the formal architecture, given the stated assumptions on explicit–tacit separation and degradation, constrains the design of realistic human–AI workflows.}

% Keywords
\keyword{Human–AI collaboration; epistemic governance; tacit knowledge; explicit–tacit separation; AI safety; long-horizon reasoning; category theory; sociotechnical systems} 

% The fields PACS, MSC, and JEL may be left empty or commented out if not applicable
%\PACS{J0101}
%\MSC{}
%\JEL{}

%%%%%%%%%%%%%%%%%%%%%%%%%%%%%%%%%%%%%%%%%%
% Only for the journal Diversity
%\LSID{\url{http://}}

%%%%%%%%%%%%%%%%%%%%%%%%%%%%%%%%%%%%%%%%%%
% Only for the journal Applied Sciences
%\featuredapplication{Authors are encouraged to provide a concise description of the specific application or a potential application of the work. This section is not mandatory.}
%%%%%%%%%%%%%%%%%%%%%%%%%%%%%%%%%%%%%%%%%%

%%%%%%%%%%%%%%%%%%%%%%%%%%%%%%%%%%%%%%%%%%
% Only for the journal Data
%\dataset{DOI number or link to the deposited data set if the data set is published separately. If the data set shall be published as a supplement to this paper, this field will be filled by the journal editors. In this case, please submit the data set as a supplement.}
%\datasetlicense{License under which the data set is made available (CC0, CC-BY, CC-BY-SA, CC-BY-NC, etc.)}

%%%%%%%%%%%%%%%%%%%%%%%%%%%%%%%%%%%%%%%%%%
% Only for the journal Toxins
%\keycontribution{The breakthroughs or highlights of the manuscript. Authors can write one or two sentences to describe the most important part of the paper.}

%%%%%%%%%%%%%%%%%%%%%%%%%%%%%%%%%%%%%%%%%%
% Only for the journal Encyclopedia
%\encyclopediadef{For entry manuscripts only: please provide a brief overview of the entry title instead of an abstract.}

%%%%%%%%%%%%%%%%%%%%%%%%%%%%%%%%%%%%%%%%%%
% Only for the journal Advances in Respiratory Medicine and Smart Cities
%\addhighlights{yes}
%\renewcommand{\addhighlights}{%

%\noindent This is an obligatory section in “Advances in Respiratory Medicine'' and ``Smart Cities”, whose goal is to increase the discoverability and readability of the article via search engines and other scholars. Highlights should not be a copy of the abstract, but a simple text allowing the reader to quickly and simplified find out what the article is about and what can be cited from it. Each of these parts should be devoted up to 2~bullet points.\vspace{3pt}\\
%\textbf{What are the main findings?}
% \begin{itemize}[labelsep=2.5mm,topsep=-3pt]
% \item First bullet.
% \item Second bullet.
% \end{itemize}\vspace{3pt}
%\textbf{What are the implications of the main findings?}
% \begin{itemize}[labelsep=2.5mm,topsep=-3pt]
% \item First bullet.
% \item Second bullet.
% \end{itemize}
%}

%%%%%%%%%%%%%%%%%%%%%%%%%%%%%%%%%%%%%%%%%%
\begin{document}
\nolinenumbers  %remove comment while editing
\pagebreak
\tableofcontents %remove comment while editing
\pagebreak
%%%%%%%%%%%%%%%%%%%%%%%%%%%%%%%%%%%%%%%%%%

\section{Introduction}
\begin{comment}
AI-based reasoning systems are increasingly used as partners in scientific work: they propose conjectures, sketch proofs, draft experiments, and synthesise literature. A single project can now involve months of back-and-forth with a collection of such tools: conjectures refined, drafts rewritten, code patched, experiments reprioritized. This creates both opportunities and risks. These systems can operate over large explicit structures that are difficult for humans to inspect directly, but they can also generate plausible yet invalid or insufficiently grounded content, sometimes with high confidence. In most current workflows, humans and AI systems interact in an untyped, informal way. Inputs, outputs, and corrections share a single channel; correctness judgments are not explicitly separated from exploratory or heuristic suggestions; and long interaction traces (a sequence of human–AI messages; see Section~2 for a formal definition) ...
 often suffer from semantic drift. This makes it difficult to reason formally about correctness, responsibility, and long-run behaviour.

We understand the scientific method as the organised process of proposing explicit hypotheses or models, deriving predictions, gathering empirical data, and revising or rejecting those hypotheses in light of the evidence. Over time, the scientific community has developed institutional practices that help stabilise scientific claims: repeatable procedures, reproducible results, and organised peer review of conclusions. Our FRFP protocol is designed to represent these patterns in a unified way, by separating explicit artefacts\footnote{By \emph{explicit} we mean artefacts that can be represented as finite, machine-manipulable symbolic structures (text, code, programs, datasets, proofs, etc.) without loss of their correctness-relevant content. In the formal development these form the explicit space \(E\).} (hypotheses, models, data, analyses) from tacit judgments\footnote{By \emph{tacit} we mean human-exclusive evaluative states—understanding, background know-how, and context-dependent judgment—that bear correctness authority and are not treated, in FRFP, as fully captured by any fixed finite symbolic encoding. In the formal development these form the tacit space \(T\).} about their adequacy, credibility, and acceptance.

In this paper, we introduce a protocol-level design for human--AI collaboration, the \textbf{Foundational Reasoning and Feedback Protocol (FRFP)}, with three aims:
\begin{itemize}
  \item to provide a formal semantics for mixed human--AI reasoning trajectories;
  \item to make explicit the boundary between what AI systems may and may not do in scientific workflows;
\item to enable theorem-level analysis of correctness and long-run behaviour (for example, conditions under which certain kinds of hallucination cannot be ruled out).\footnote{Throughout, such ``inevitability'' claims are made only within the FRFP formalism and under the assumptions set out in Section~\ref{sec:dynamics}, not as universal statements about all possible AI systems.}
\end{itemize}
\end{comment}
\noindent
AI-based reasoning systems are increasingly used as partners at work across domains from science to business: they draft, critique, plan, and propose. A single project can now involve months of back-and-forth with multiple such tools---drafts rewritten, code patched, analyses rerun, experiments reprioritised, decisions revisited---creating both opportunities and risks. These systems can operate over large explicit structures that humans cannot easily inspect, yet they can also produce plausible but invalid or weakly grounded content, sometimes with high confidence.  Complex systems are routinely built and operationalised from imperfect components by adding redundant layers of safety checks and escalation paths\cite{nrc_defense_in_depth_glossary,iaea_insag10_defence_in_depth_1996,reason_human_error_models_management_2000,reason_latent_human_failures_complex_systems_1990}. The same engineering fact suggests a protocol view of human--AI collaboration: reliability should be an emergent property of the workflow, not an assumption about any single component.

\smallskip
\noindent
The core issue is not occasional error but \textbf{unstabilised collaboration}: we increasingly depend on long, multi-step human--AI workflows whose intermediate artefacts, assumptions, and evaluation criteria shift over time. In many settings this drift occurs inside one conversational channel, with no stable separation between exploration and acceptance and no durable checkpoints where correctness and responsibility are explicitly reasserted. As a result, most workflows remain untyped and informal: inputs, outputs, and corrections share a channel; correctness judgments are not separated from heuristic suggestions; and long interaction traces (a sequence of human--AI messages; see Section~2 for a formal definition) often suffer semantic drift---making correctness, responsibility, and long-run behaviour hard to analyse \cite{laban_llms_get_lost_multiturn_2025,dongre_drift_no_more_context_equilibria_2025,arike_evaluating_goal_drift_agents_2025,choi_identity_drift_llm_agents_2024}.

\smallskip
\noindent
We take the scientific method as organising inspiration because its stabilisation norms---replication, reproducibility, and peer review---are historically established and widely standardised, providing a canonical template for converting exploratory work into accepted claims.\footnote{We use ``scientific'' for workflows whose primary aim is stable, publicly checkable claims about the world, stabilised by intersubjective norms such as replication, reproducibility, and peer review. By contrast, much non-scientific work aims at action under objectives and constraints (e.g., safety, cost, policy, profit), where acceptance is stabilised by organisational controls (tests, audits, approvals, monitoring) rather than truth-claims alone. The same explicit/tacit separation applies in both: what changes is the form of ``evidence'' and the institutional mechanism of acceptance.}
We abstract from that template a \textbf{domain-general} protocol for human--AI collaboration: in non-scientific settings (business, policy, engineering), ``evidence'' and ``review'' differ in form, but the same long-horizon drift and accountability problems arise. Most examples in this paper are non-scientific.

\smallskip
\noindent
In this paper, we introduce a protocol-level design for human--AI collaboration, the \emph{Foundational Reasoning and Feedback Protocol} (FRFP), with three aims:
\begin{itemize}
  \item to provide a formal semantics for mixed human--AI reasoning trajectories;
  \item to make explicit the boundary between what AI systems may and may not do in such workflows;
  \item to enable theorem-level analysis of correctness and long-run behaviour, including conditions under which certain kinds of hallucination cannot be ruled out.\footnote{Throughout, such ``inevitability'' claims are made only within the FRFP formalism and under the assumptions set out in Section~\ref{sec:epistemicdynamics}, not as universal statements about all possible AI systems.}
\end{itemize}

\noindent
FRFP’s central design choice is a sharp separation between two modes of operation: an \emph{explicit} mode of machine-manipulable representations and pipelines, and a \emph{tacit} mode of human-exclusive evaluative, correctness-bearing states. AI systems are confined to explicit mode; only humans inhabit tacit mode. A single review-and-belief-update boundary carries explicit artefacts into tacit evaluation, with no direct channel in the opposite direction. The rest of the paper formalises this architecture and studies long-horizon human--AI workflows within it.

\smallskip
\noindent
Informally, FRFP is a human--AI conversational calculus: a typed calculus of interaction traces in which AI systems make only explicit moves, humans retain tacit correctness-bearing moves, and entire workflows admit compositional semantics and dynamic analysis.
\begin{comment}

\paragraph{Scope of novelty.}
We rely on standard categorical constructions---free categories, groupoids, fibrations, Grothendieck constructions---but assemble them in service of explicit/tacit separation and feedback in human--AI workflows. The novelty lies in:
\begin{enumerate}
  \item isolating a small set of epistemic primitives (explicit/tacit separation, boundary, feedback);
  \item combining them into a compositional architecture in which system-level properties (e.g., auditability, correctness preservation, and conditional inevitability results for certain failure modes such as hallucination) can be formulated and proved as theorems inside the FRFP assumptions; and
  \item working at a level of abstraction that allows later work to reuse existing categorical theory rather than inventing bespoke formalisms for each application domain.
\end{enumerate}
    
\end{comment}

\subsection{Scope of novelty.}
We rely on standard categorical constructions---free categories, groupoids, fibrations, and Grothendieck constructions---but assemble them as follows to formalise explicit/tacit separation and feedback in human--AI workflows:
\begin{enumerate}
  \item a minimal set of epistemic primitives (explicit/tacit separation, boundary, feedback);
  \item a compositional architecture in which system-level properties---auditability, correctness preservation, and conditional inevitability results for failure modes such as hallucination---are stated and proved as theorems under FRFP assumptions; and
  \item an abstraction level that lets later work reuse existing categorical theory rather than introducing bespoke formalisms per application domain.
\end{enumerate}

Our use of FRFP in this paper is threefold: we prove theorems internal to the model; we conjecture that its structural assumptions approximate many high-stakes workflows in practice; and, conditional on that empirical premise, we derive design implications for institutional oversight.

\subsection{FRFP architecture overview}
\label{sec:frfp-architecture-overview}
We organise FRFP into five epistemic layers:

\textbf{Interaction Protocols.}
Rules governing human--AI interaction: who may say what, in which modes (explicit vs.\ meta-level), and in which order. This layer constrains admissible traces but does not yet assign meaning.

\textbf{Epistemic Spaces.}
Admissible traces are interpreted over two state spaces.\footnote{We use ``epistemic'' in the standard sense: pertaining to knowledge, belief, or correctness-bearing state. In FRFP, epistemic structures track correctness-relevant information (what is written down, computed, or logged, and what is tacitly accepted), as distinct from surface interaction structure.}
We distinguish: (i) an explicit space of knowledge-like artefacts (what is written down, computed, or logged); and (ii) a tacit space of correctness-bearing states (what humans accept or are prepared to stand behind). A one-way boundary maps explicit artefacts into tacit evaluation.

\textbf{Epistemic Algebra.}
This layer specifies the allowed \emph{operations} on explicit workflows---how we decompose tasks into steps, compose steps into pipelines, navigate between subproblems, and refactor a workflow without changing its intent. These operations act only on explicit artefacts. At designated review points, the one-way boundary carries the current explicit artefact(s) into tacit evaluation, producing a human judgment update. The result is an ``algebra'' of explicit workflows equipped with a tacit correctness semantics: we can reason about how workflow transformations affect what a human ends up accepting.

In this volume we keep the tacit side intentionally minimal. We treat the tacit space $T$ and its operators as a correctness-bearing interface: a black-box human layer that (i) receives explicit artefacts only via the one-way boundary and (ii) exposes only a finite, observable set of judgments. The technical structure therefore lives on the explicit side. Detailed models of individual cognition or institutional psychology are out of scope and deferred.


\textbf{Epistemic Dynamics.}
The evolution of knowledge- and correctness-bearing state over time and under load: tacit degradation, requirement schedules, hallucination as requirement violation, and safe horizons for running a workflow.

\textbf{Collective Epistemic Dynamics.}
Knowledge and correctness in populations and institutions: multiple agents sharing explicit artefacts, communicating via conversational workflows, and operating under policies, committees, and governance operators. This layer captures collective dynamics and the limits of explicit-only consensus and governance.

\textbf{Running example.}
Throughout the paper we illustrate these layers with a simple case: a human uses an AI assistant to choose a ``best'' airfare, where best is not cheapest. The assistant may explore and structure options, but the accept/reject judgment remains with the human.


\begin{figure}[t]
  \centering
  \begin{tikzpicture}[
    node distance=1.2cm,
    box/.style={
      rectangle,
      draw=black,
      rounded corners,
      align=center,
      text width=0.9\linewidth,
      minimum height=0.9cm
    }
  ]
    \node[box] (ip) {Interaction Protocols\\[2pt]
      \footnotesize Rules about how humans and AI may interact: who may say what, in which modes (explicit vs.\ meta-level), and in which order.};
    \node[box, above=of ip] (es) {Epistemic Spaces\\[2pt]
      \footnotesize The state those interactions manipulate: an explicit space of knowledge-like artefacts (what is written down, computed, or logged) and a tacit space of correctness-bearing states (what humans accept or stand behind), linked by a one-way boundary.};
    \node[box, above=of es] (ea) {Epistemic Algebra\\[2pt]
      \footnotesize Operations on explicit workflows that transform epistemic state: task decomposition, composition, navigation, and refactoring. These act on explicit artefacts and, via the boundary, induce changes in tacit states.};
    \node[box, above=of ea] (ed) {Epistemic Dynamics\\[2pt]
      \footnotesize How knowledge- and correctness-bearing states evolve over time and under load: tacit degradation, requirement schedules, hallucination as requirement violation, and safe horizons for running a given workflow.};
    \node[box, above=of ed] (ced) {Collective Epistemic Dynamics\\[2pt]
      \footnotesize How knowledge and correctness behave in populations and institutions: multiple agents sharing explicit artefacts, communicating via conversational workflows, and operating under policies, committees, and governance operators.};

    \draw[->] (ip) -- (es);
    \draw[->] (es) -- (ea);
    \draw[->] (ea) -- (ed);
    \draw[->] (ed) -- (ced);
  \end{tikzpicture}
  \caption{FRFP as a human--AI conversational calculus with five epistemic layers. Interaction protocols constrain admissible H--AI interactions. Epistemic spaces specify the explicit and tacit states those interactions act on. Epistemic algebra equips explicit workflows with compositional structure and a tacit correctness semantics. Epistemic dynamics models how these epistemic states evolve over time and under load, and collective epistemic dynamics lifts the architecture to populations and institutions.}
  \label{fig:frfp-architecture-overview}
\end{figure}




This paper is Volume~1 of a three-volume programme.

\begin{description}
  \item[Volume 1: FRFP core.] Develop the formal architecture of FRFP---its abstract objects and inevitability results---using category-theoretic tools. Deeper mathematical developments are deferred to a companion pack of technical notes. Informally: \emph{what is structurally true of long-horizon human--AI workflows, regardless of institutional convenience?}

  \item[Volume 2: Case studies and design patterns.] Apply the framework to concrete domains (e.g., scientific assistants, clinical decision support, trading systems, autonomous driving) to test its assumptions and extract design patterns. Informally: \emph{given these structural constraints, what works or fails in real deployments?}

  \item[Volume 3: Tooling and disciplined agents.] Develop practical tooling around FRFP (run-keepers, context managers, CMS integration, adapters, competence monitors) to support disciplined human--AI workflows. Informally: \emph{given what works, what should we build so others can use these systems without unacceptable risk to themselves or others?}
\end{description}


\subsection{Intended audiences and reading paths.}
This volume (Volume~1) and its companion materials target three overlapping audiences:
\begin{enumerate}
  \item \textbf{Category theory and rewriting readers.} Readers comfortable with category theory and rewriting systems can read Volume~1 directly, using the companion technical pack~[CP] for notation, background, and longer proofs.
  \item \textbf{AI safety and governance readers.} Readers less familiar with category theory should begin with the informal overview and glossary in~[CP] before entering the technical core of Volume~1. That material isolates the few constructions needed for the main claims (explicit/tacit separation, the FRFP boundary, the formal notion of hallucination, and the long-horizon dynamic model) without assuming prior exposure to fibrations or Grothendieck constructions.
  \item \textbf{Practitioners.} Readers primarily interested in design implications may start with the non-technical exposition and case studies in Volume~2, treating Volume~1 and~[CP] as reference background, and Volume~3 as a source of concrete tooling for disciplined deployment.
\end{enumerate}
\subsection{Contributions}

Formally and conceptually, this volume makes the following contributions.
\begin{enumerate}
  \item \textbf{Explicit--tacit epistemic architecture.}
  We propose FRFP as an explicit--tacit architecture for long-horizon human--AI workflows, based on an explicit/tacit separation and a one-way review boundary.

  \item \textbf{Category-theoretic realisation and semantics.}
  We realise this architecture category-theoretically: an algebra of explicit pipelines and a semantics mapping whole workflows to tacit correctness judgments. In the model, correctness is invariant under admissible explicit transformations and navigation, so system-level properties can be stated and proved at the pipeline level rather than per step.

  \item \textbf{Dynamics of tacit context and hallucination.}
  We develop a qualitative long-horizon dynamic model: tacit capacity is finite and degrades under load; requirements evolve; and feedback from explicit artefacts shapes future judgments. Under these assumptions we prove conditional inevitability results: over sufficiently long admissible runs, failure modes including hallucination (accepting artefacts that do not meet active requirements) cannot be ruled out, and safe horizons are necessarily finite even when explicit semantics are fixed.

  \item \textbf{Collective and stochastic settings.}
  We extend the analysis from single users to populations and from deterministic to stochastic implementations. We model collective executions as processes over populations of tacit states, analyse how consensus interacts with explicit--tacit separation, and show that stochasticity determines which failure trajectory is realised rather than precluding failure, yielding notions such as hazard profiles and safe-horizon functionals.

  \item \textbf{Limits of explicit-only mechanisms and non-automation.}
  From explicit--tacit separation and human-exclusive grounding of correctness, we derive conditional impossibility results for purely explicit mechanisms. Under FRFP assumptions, no explicit-only detector or governance operator can guarantee the absence of hallucination or correctness failure across all admissible pipelines and populations. This yields a structural non-automation claim: in domains fitting the FRFP architecture, tacit human oversight is necessary at the governance layer.

  \item \textbf{Disciplined agents and FRFP-compatible tooling.}
  We introduce \emph{disciplined agents}: AI tools and surrounding infrastructure confined to the explicit side of FRFP, exposing reasoning in pipeline form and integrating with tacit oversight rather than replacing it. We state requirements (e.g., source--reason traceability, navigation integrity, explicit confinement) and outline supporting tooling (run-keepers, context managers, competence monitors).

  \item \textbf{Protocols, case studies, and design patterns.}
  We translate the formal architecture into protocol-level design principles for long-horizon human--AI collaboration. We state foundational rules for FRFP-compatible practice (explicit confinement, human-exclusive grounding, source--reason traceability, integrity under navigation), illustrate them in case studies (e.g., literature review, financial scenario analysis), and give examples of FRFP-aligned system instructions and prompting patterns for explicit-only tools. These patterns show how FRFP can constrain interfaces, workflows, and institutional arrangements in high-stakes settings.
\end{enumerate}

\subsection{Related Work}

Our work sits at the intersection of formal semantics, human--AI interaction, AI safety governance, and social epistemology. We position FRFP against mechanised reasoning, categorical semantics and rewriting, HCI-style human--AI frameworks, governance and evaluation protocols, and formal models of scientific communities, and then relate it to the hallucination and calibration literature.

\subsubsection{Interactive proof, type theory, and mechanised reasoning.}
Proof assistants and type-theoretic systems (e.g., Coq, Isabelle, Lean) distinguish proof objects from proof scripts and provide soundness guarantees for derivations within a fixed formal calculus. Users supply goals and proof terms; the central question is whether a term inhabits a type or a proof object checks a specification. FRFP is orthogonal: we do not propose a new logic or proof theory, and correctness is not a syntactic property of explicit objects. Instead, correctness is a tacit attribute of human epistemic state, accessed only via the explicit--tacit boundary, and the target is long-horizon workflows in which correctness cannot be reduced to proof checking or type inference.

Within a fixed formal system, some correctness properties (e.g., that a proof term checks) are decidable and verifiable entirely in explicit space. In FRFP terms, provers and test harnesses act as explicit correctness engines: morphisms internal to $E$ that certify properties of explicit artefacts. Our non-automation results address a broader, tacit notion of correctness for real-world workflows and institutions, which cannot be discharged by any explicit-only oracle. Such tools are therefore best viewed as tacit-insensitive inner layers: once a calculus and specification are tacitly fixed, inner-task correctness becomes explicit, while tacit judgment continues to govern the choice, adequacy, and revision of the formalism itself (Appendix~D.4).

\subsubsection{Category theory and categorical semantics.}
The categorical tools we use---free categories, groupoids, (op)fibrations, Grothendieck constructions---are standard in categorical semantics of languages and logics, where they model syntax/semantics maps, type systems, effects, and compositional program behaviour. Our contribution is to repurpose these tools for epistemic roles in mixed human--AI reasoning: explicit and tacit subcategories linked by a one-way boundary, navigation capturing refactoring and exploration, and a combined workflow category in which whole runs are single compositional objects. This supports theorem-level statements of auditability, correctness invariance, and explicit-only impossibility results as properties of a unified architecture rather than informal desiderata.

\subsubsection{Rewriting, confluence, and canonical forms.}
Our explicit reduction systems and confluence arguments draw on standard abstract reduction systems and termination/confluence techniques from term rewriting. The novelty is not the rewriting machinery but its epistemic role: normalisation shows that admissible explicit refactorings and navigation steps collapse to canonical forms that induce the same tacit correctness judgments. Rewriting thus yields invariance of tacit acceptance under the exploratory restructuring typical of long-horizon human--AI pipelines.

\subsubsection{HCI-style human--AI interaction and appropriate reliance.}
HCI and human--AI interaction research studies task allocation, communication, and co-adaptation, often organised along dimensions such as locus of control, initiative, transparency, and feedback. Recent surveys and experiments analyse trust, distrust, and \emph{appropriate reliance}, including how explanations and interface choices affect over- and under-reliance \cite{Eckhardt2024Reliance,Schemmer2023Appropriate,Visser2025Trust}. These approaches typically treat the boundary at the level of tasks and interfaces, with correctness and responsibility mediated by constructs such as calibrated assistance. FRFP is compatible in spirit but imposes a sharper architectural constraint: correctness is, by construction, human-exclusive. AI-mediated operations remain in explicit space; systems cannot access tacit state as input or cross the explicit--tacit boundary; and interaction patterns such as suggestion and critique are treated as morphisms in a formal workflow semantics. This yields invariance and impossibility results at the class-of-interactions level and provides an architectural envelope within which HCI design choices must sit.

\subsubsection{AI safety governance, risk committees, and evaluation protocols.}
AI safety and governance literatures emphasise organisational and regulatory mechanisms: safety frameworks, red-teaming, incident response, internal risk committees, standardised evaluation suites, and secure development practices, often organised as multi-line-of-defence models. Recent syntheses, such as METR's \emph{Common Elements of Frontier AI Safety Policies}, extract shared components of lab policies for capability thresholds, evaluations, deployment safeguards, and governance structures \cite{METRCommonElements}. These proposals typically treat correctness or hallucination as system properties measured by evaluations (e.g., benchmark failure rates) and assume that appropriately designed protocols deliver acceptably low failure.

FRFP differs in three respects. First, evaluations, committees, and governance operators are explicit constructions atop the workflow architecture; questions such as whether explicit-only consensus can recover correctness, or whether a fully automated governance oracle can exist, become formal questions about morphisms and functors. Second, explicit--tacit separation yields explicit-only impossibility results: under FRFP assumptions, broad classes of detectors, consensus mechanisms, and governance operators cannot, in principle, recover tacit correctness or eliminate hallucination across all admissible pipelines and populations. Third, FRFP introduces dynamic safe horizons: rather than focusing only on static capability scores, we analyse degradation, requirement schedules, and hazard profiles for whole workflows, shifting attention from benchmark scores to how long a given protocol can be run under specified controls before risk exceeds a tolerable threshold.

\subsubsection{Epistemic models of scientific communities.}
Formal social epistemology and philosophy-of-science models represent communities as stylised agents that select research options, exchange evidence or beliefs, and update via simple rules. Work on \emph{network epistemology} studies how communication structure and incentives affect convergence, replication success, and the division of cognitive labour \cite{Zollman2013Network,Grim2013ScientificNetworks}. These models share FRFP's interest in individual versus collective performance and institutional structure, but typically work with explicit belief states (e.g., probability vectors or epistemic landscapes) and rarely separate explicit representations from tacit state. FRFP instead builds explicit/tacit modality into the foundation: explicit objects are manipulable by AI systems and institutions, while tacit states remain human and correctness-bearing. This underwrites results showing that certain institutional arrangements cannot reconstruct correctness from explicit data alone.

\subsubsection{Hallucination, calibration, and uncertainty in language models.}
Work on hallucination, calibration, and uncertainty surveys failure modes and benchmarks, develops semantic-uncertainty methods, and studies when models can assess their own correctness and why standard training pipelines yield hallucinations and miscalibration \cite{Zhang2023Siren,Huang2025HallucinationSurvey}. Hallucination is typically treated as mismatch between outputs and a ground-truth distribution or benchmark; the focus is on architectures, training objectives, or decoding strategies that reduce or detect such errors. FRFP instead treats hallucination as a protocol-level property: a mismatch between tacit context and active requirements at the moment an explicit artefact is accepted. The definition is model-agnostic but institution-dependent, since requirement maps encode domain- and institution-specific standards. Model-side detectors can reduce short-run error by filtering artefacts before acceptance, but so long as tacit state degrades and nontrivial requirements recur, our dynamic analysis implies that sufficiently long admissible runs will eventually include requirement-violating acceptances. Because detectors operate entirely in explicit space, they cannot arrest long-run degradation of tacit grounding. FRFP is therefore complementary to model-level hallucination and calibration work, but targets protocol and governance design and the limits of explicit-only guarantees.

\subsubsection{Human--AI collaboration and protocol design.}
There is a growing literature on AI support for discovery, literature review, hypothesis generation, experimental design, and broader knowledge work, typically organised around specific tasks, phases, or interface patterns. Much of this work optimises local performance (speed, coverage, quality) and studies usability constructs (transparency, control, reliance), but leaves long-horizon questions---semantic drift across extended traces, responsibility for acceptance, and the limits of automation---only implicitly specified. FRFP provides a protocol-level frame: long-horizon workflows are treated as compositions in an explicit workflow architecture with a fixed boundary to tacit space. Within this frame, ``copilots'', risk committees, and red-team evaluations become concrete patterns of morphisms and institutions. FRFP's distinctive combination is an explicit/tacit semantic architecture with human-exclusive correctness, a dynamic theory of tacit degradation and protocol-level hallucination with inevitability results, and an institutional layer yielding structural non-automation and explicit-only impossibility results for governance mechanisms.



%=========================================================
\section{Interaction Protocols}
\label{sec:interaction-protocols}

This layer specifies interface-level constraints on human--AI interaction traces: who may say what, in which modes, and in which order. We distinguish content-producing messages (explicit messages, which both human and AI may send) from correctness-bearing messages (meta-level messages, which only the human may send); Section~\ref{sec:epistemic-spaces} makes this distinction precise as explicit versus tacit epistemic spaces.

\subsection{Interaction Traces}

We model an interaction session as a finite sequence of messages tagged by origin.

\begin{definition}[Interaction trace]
A \emph{human--AI interaction trace} is a finite sequence
\[
S = (m_1, m_2, \dots, m_n),
\]
where each message \(m_i\) is tagged as either human-generated (H) or AI-generated (A) and carries explicit content (e.g., text, code, mathematical expressions).
\end{definition}

We write \(S_k = (m_1,\dots,m_k)\) for the \(k\)-prefix of \(S\). We impose a family of interaction constraints \(\Phi_0\) on such prefixes.

\subsection{Interaction Constraints}

The constraints are operational: they enforce mode discipline and prevent the AI from taking tacit roles.

\begin{description}[leftmargin=2em]
  \item[Intent locking.] Early in the trace, the human states a coarse-grained intent (e.g., ``analyse this conjecture'', ``design an experiment''). The AI treats this intent as fixed unless the human explicitly revises it. This prevents uncontrolled drift in task framing.

  \item[Ambiguity resolution.] If an AI message admits more than one task-relevant interpretation, the next AI message must be a clarification query, not a content-producing response. This prohibits silent resolution of deep ambiguity.

  \item[Mode discipline.] Each message is typed as either explicit (facts, derivations, questions) or meta-level (requests for tacit judgment). Both human and AI may send explicit messages. Only the human may send meta-level, correctness-bearing messages; the AI is confined to explicit mode.

  \item[No tacit emulation.] The AI may not simulate tacit operations such as evaluating overall correctness or deciding whether a result should be accepted. It may propose checks, but it is not a source of correctness.

  \item[Correctness stability.] Once the human issues a tacit correctness judgment, the AI may not reinterpret or override it via further explicit manipulation. Tacit states are fixed points with respect to AI operations.
\end{description}

\begin{definition}[Interaction admissibility]
An interaction trace \(S\) is \emph{admissible} if every prefix \(S_k\) satisfies all interaction constraints. We write \(S \vDash \Phi_0\) for this property.
\end{definition}

These constraints are intentionally strict and normative. FRFP commits to: (i) an explicit/tacit distinction, in the sense of Polanyi and subsequent work on organisational knowledge creation; and (ii) human-exclusive correctness in high-stakes settings, as reflected in human-oversight requirements for high-risk AI systems. Following Polanyi's account of tacit knowledge as embodied, context-dependent, and not fully articulable \cite{Polanyi1966}, and Nonaka's treatment of explicit and tacit knowledge as distinct but interacting modalities in organisational learning \cite{Nonaka1994,NonakaTakeuchi1995}, we model explicit artefacts and tacit human states as separate epistemic spaces (Section~\ref{sec:epistemic-spaces}) linked by a one-way boundary. In parallel, contemporary governance frameworks require human oversight of high-risk AI systems, with humans able to monitor, intervene, and override AI outputs \cite{EUAIAct2024,Amershi2019}. FRFP packages these commitments as an architectural constraint for disciplined human--AI workflows and studies their formal consequences. It is not intended as a description of current ad hoc practice, but as a target architecture: in this volume, explicit/tacit separation and human-exclusive correctness are design principles, not optional knobs; later work may study carefully delimited departures.

In the running airfare example introduced in Section~\ref{sec:frfp-architecture-overview}, the interaction trace \(S = (m_1,\dots,m_{10})\) consists of alternating human and AI messages \(H_1\)--\(H_5\) and \(A_1\)--\(A_5\). We write the human messages as \(H_1,\dots,H_5\) and the AI messages as \(A_1,\dots,A_5\), in order of appearance.The interaction constraints require, in particular, that the AI remain in explicit mode: it may elicit preferences, propose itineraries, and analyse trade-offs, but it may not itself issue the final ``Option~A is best for you'' verdict. The human's final choice at \(H_5\) is a tacit correctness judgment, which the AI can at most record explicitly. Any trace in which the AI directly declares an option to be ``best'' or overwrites the human's recorded decision is operationally inadmissible.

The interaction protocols constrain the observable structure of human--AI conversations: who may say what, in which modes, and in which order. The next layer introduces the epistemic spaces that such conversations act on: explicit artefacts and tacit human states that carry knowledge- and correctness-relevant information.

\section{Epistemic Spaces}
\label{sec:epistemic-spaces}

The interaction protocols in Section~\ref{sec:interaction-protocols} constrain which conversational messages are allowed; this layer specifies what those messages act on: an explicit space \(E\) of machine-manipulable artefacts and a tacit space \(T\) of human-only correctness-bearing states, linked by a one-way boundary.

\subsection{Ambient Category and Explicit--Tacit Separation}

We work inside an ambient category \(\Cfull\) with distinguished objects \(\varnothing\), \(E\), and \(T\), as detailed in Appendix~A. The Explicit--Tacit Separation axiom (ETS) asserts:
\begin{itemize}[leftmargin=2em]
  \item explicit morphisms end at \(E\);
  \item tacit morphisms end at \(T\);
  \item there are no morphisms \(T \to E\);
  \item there is a unique boundary morphism \(\RB:E\to T\);
  \item \(\RB\) is not invertible.
\end{itemize}

Intuitively, explicit and tacit modes are \emph{disjoint}: there is a one-way flow from explicit artefacts into tacit evaluation, and no way back.

\subsection{Explicit and Tacit Primitives}

Within \(\Cfull\), we identify six primitive morphisms that play irreducible epistemic roles:
\begin{itemize}[leftmargin=2em]
  \item \(\RI : \varnothing \to E\) (Representation Initiation),
  \item \(\EC : E \to E\) (Explicit Computation),
  \item \(\ED : E \to E\) (Explicit Diagnostics),
  \item \(\RB : E \to T\) (Representational Backflow),
  \item \(\TE : T \to T\) (Tacit Evaluation),
  \item \(\HFD : T \to T\) (Human Final Decision).
\end{itemize}

The first three live entirely in explicit space; \(\RB\) crosses the boundary into tacit space; and \(\TE,\HFD\) live in tacit space. They compose into the core correctness architecture:
\[
\RI \;\to\; (\EC/\ED)^\ast \;\to\; \RB \;\to\; \TE \;\to\; \HFD.
\]
Here \((\EC/\ED)^\ast\) denotes an arbitrary (possibly empty) finite sequence of explicit computations and diagnostics. The pipeline ends with tacit evaluation and final decision: no correctness-bearing steps occur in explicit space.

\subsection{Correctness in Tacit Space}

Tacit states form a preorder-enriched one-object category \(\Tcat\) with primitives \(\TE,\HFD\) and a correctness quotient \(\gamma:\Obj(\Tcat)\to J\) into a finite set of judgments. The key fact for the main text is:

\begin{proposition}[Tacit correctness quotient]
Tacit primitives \(\TE,\HFD:T\to T\) are idempotent and commute, and the correctness quotient \(\gamma\) satisfies \(\gamma\circ\TE=\gamma=\gamma\circ\HFD\). Thus tacit transformations do not change the observable correctness class.
\end{proposition}

The proof and full structure appear in Appendix~A. For the main text, it suffices to know that correctness is localised in tacit space and is observable only via \(\gamma\).

\paragraph{Interpretation and scope of tacit correctness.}
Treating correctness as living in tacit space \(T\) rather than in explicit space \(E\) is a substantive architectural commitment, not a definitional firewall. It targets domains in which correctness-bearing judgment does not reduce to a fixed, finite explicit specification or benchmark suite: long-horizon workflows in which tasks evolve, standards and constraints shift, and background understanding and norms determine whether an explicit artefact is adequate. In such settings, explicit artefacts, tests, and formal checks do not exhaust correctness-bearing judgment; Appendix~A.11 shows that, under a nontrivial tacit structure, a broad class of explicit-only oracles and governance mechanisms are structurally impossible. By contrast, some regimes admit fully explicit correctness---for example, proof-assistant settings or closed-world tasks with sound and complete test oracles. FRFP does not treat these regimes as primary instances of its epistemic dynamics; rather, they appear as tacit-insensitive inner layers embedded in a broader tacit governance context (Appendix~D.4). The results in this section are therefore conditional on a genuine tacit layer in the sense above, and on explicit/tacit separation and human-exclusive correctness as non-negotiable design principles in that domain.

\paragraph{Hallucination as a protocol-level phenomenon.}
Within this architecture, hallucination is a tacit-level notion. We do not apply ``hallucination'' to the mere production of incorrect strings in \(E\); at that stage outputs are explicit candidates, and, when they violate a trusted explicit oracle (e.g.\ a test harness or proof checker), explicit errors. Hallucination arises at the protocol level, when explicit artefacts cross the boundary \(\RB:E\to T\) and are incorporated into tacit correctness. Informally, a hallucination occurs along a run when tacit operations such as \(\TE\) and \(\HFD\) end up treating an explicit artefact as acceptable even though, relative to the domain's standards, it ought to have been judged inadequate or incorrect. Because correctness is localised in \(T\) and observable only via \(\gamma\), hallucination cannot be read off from explicit artefacts in isolation or from the behaviour of the generator that produced them: the same explicit string may be harmless when rejected and hallucinated when taken up into tacit acceptance. Later sections refine this picture dynamically; the core point is that, once explicit and tacit spaces are separated, hallucination in FRFP is a property of endorsement by tacit states, not of explicit generation alone.

\paragraph{Explicit correctness engines.}
In many workflows, parts of correctness can be discharged by \emph{explicit correctness engines}: unit-test harnesses, static analysers, and interactive theorem provers (e.g., Coq, Isabelle, Lean) that, relative to a fixed formal specification, can certify whether an explicit artefact satisfies a given property. Within FRFP these engines live entirely in explicit space \(E\): they are instances of explicit computation or diagnostics (\(\EC/\ED\)), possibly composed into larger explicit pipelines. They can reduce the space of explicit errors tacit agents must consider, but they do not collapse the tacit layer. Tacit correctness still governs which specifications are adequate, when they must be revised, how to interpret formally verified artefacts in context, and what residual risk is acceptable. Explicit correctness engines are therefore treated as powerful inner modules inside \(E\), coupled to but not replacing the outer tacit governance layer.

\subsection{Uniqueness of the FRFP Epistemic Ontology}

Appendix~A shows that, under reasonable axioms, the FRFP Epistemic Ontology is essentially unique.

\begin{theorem}[Initiality of the FRFP ontology, informal]
Any epistemic ontology that:
\begin{itemize}[leftmargin=2em]
  \item distinguishes explicit and tacit modalities,
  \item enforces a one-way boundary \(E\to T\),
  \item localises correctness in tacit operators \(\TE,\HFD\), and
  \item represents AI systems as acting only in explicit space,
\end{itemize}
admits a unique structure-preserving mapping from the \FRFP{} ontology. In this precise sense, the \FRFP{} ontology is minimal and essentially unique.
\end{theorem}

We refer the reader to Appendix~A for the formal category of Epistemic Spaces frameworks and the proof that the FRFP ontology is initial in that category.

\paragraph{Operational significance of the FRFP Epistemic Ontology.}
The initiality result is not a formal restatement of ``we built a free category''. It says that any ontology satisfying (i) explicit/tacit modality, (ii) a one-way boundary \(E\to T\), (iii) tacit localisation of correctness, and (iv) explicit-only AI action must factor uniquely through the FRFP Epistemic Ontology. In particular, any framework that preserves the same primitive epistemic roles and axioms, but proposes a different wiring of \(\RI,\EC,\ED,\RB,\TE,\HFD\), still receives a unique structure-preserving functor from FRFP. This makes later non-automation and no-go results robust under local redesigns: one cannot bypass them by relabelling primitives or rearranging explicit/tacit roles while keeping the Epistemic Spaces axioms intact.

Given an admissible interaction trace \(S\), the epistemic-ontology layer associates to each prefix \(S_k\) a pair of states \((e_k,t_k)\), where \(e_k \in E\) is the explicit state and \(t_k \in T\) the tacit state.

In the airfare example introduced in Section~\ref{sec:frfp-architecture-overview}, after \(H_2\) the explicit state \(e_2\) contains at least the travel dates and route, the human's stated priority weights, and the assistant's preference template. The tacit state \(t_2\) reflects the human's internal attitudes toward those trade-offs, which may not match any explicit labels. After \(A_3\), the explicit state \(e_3\) additionally contains candidate itineraries and their attributes; the tacit state \(t_3\) reflects the human's (possibly tentative) sense of which options are promising.

The boundary operator \(\RB:E\to T\) (representational backflow) is applied when the human reviews an explicit artefact. At \(H_5\), the human evaluates the explicit summaries and options, applies \(\RB\) to the relevant explicit configuration, and updates the tacit state \(t_5\) to encode acceptance of Option~A as ``best'' in context. The correctness quotient \(\gamma:\Obj(\Tcat)\to J\) then records this tacit acceptance as an observable judgment (e.g., ``Option~A accepted'').

Our treatment of tacit space is deliberately spare. We assume only that tacit states carry correctness-relevant information, are updated via a one-way boundary from explicit artefacts, and admit a finite observable judgment quotient. We do not model the internal structure of human belief or institutional reasoning; instead, substantive structure is pushed into explicit workflows that we can represent and analyse compositionally.
%=========================================================
\section{Epistemic Algebra}
\label{sec:epistemicalgebra}
%=========================================================
The previous section introduced the FRFP epistemic spaces: an explicit space $E$ of machine-manipulable artefacts and a tacit space $T$ of human-only correctness-bearing states, linked by a one-way boundary. This section equips \emph{explicit} workflows with algebraic structure over that ontology. Together with the interaction protocols of Section~\ref{sec:interaction-protocols}, these constructions form the FRFP human--AI conversational calculus. All algebra in this section lives on explicit workflows; tacit state appears only as the semantics of those workflows, not as an object we rewrite.

Given the explicit/tacit ontology, FRFP’s technical structure sits on the explicit side. We define an epistemic algebra of explicit workflows: a task decomposition grammar, navigation and refactoring operations, and a semantics into tacit states. Tacit evaluation enters only through the boundary and correctness quotient; the algebraic invariance results are entirely explicit.

This layer has two ingredients:
\begin{itemize}[leftmargin=2em]
  \item a Task Decomposition Grammar (\(\TDG\)) that generates explicit
        pipelines from the primitives \(\RI,\EC,\ED\); and
  \item a navigation groupoid (\(\Ncat\)) that captures reversible
        changes of perspective or representation.
\end{itemize}
They are combined via a Grothendieck-style semidirect product
\(\Ncat \ltimes \Ecat\), a standard construction in categorical
semantics~\cite{MacLane1998,Awodey2010,Borceux1994}, yielding an algebra
of reasoning trajectories.


\subsection{Task Decomposition Grammar}

TDG presents the free algebra of explicit pipelines. Informally, its grammar is:
\begin{align*}
Q &::= \RI \mid Q;\EC \mid Q;\ED \mid
        \BRANCH(Q_1,\dots,Q_k) \mid
        \PARALLEL(Q_1,\dots,Q_k) \mid
        \LOOP(Q_{\mathrm{cond}},Q_{\mathrm{body}}),\\
P &::= Q;\RB;\TE;\HFD.
\end{align*}
Here \(Q\) ranges over explicit pipelines and \(P\) over complete FRFP
pipelines.

The explicit category \(E_{\mathrm{sem}}\) generated by TDG satisfies the
usual universal property: any interpretation of the generators in an
explicit category extends uniquely to a functor from \(E_{\mathrm{sem}}\).

\subsection{Navigation Groupoid}

Navigation steps---trying alternative models, adding or dropping
assumptions, restarting from a new shape---are represented as morphisms
in a groupoid \(\Ncat\). Objects of \(\Ncat\) are \emph{shapes} of
explicit reasoning; morphisms are structural rewritings between shapes,
freely inverted to form a groupoid.

Each navigation morphism \(n\in\Ncat\) acts on explicit pipelines
via a fiberwise functor \(n^\ast:E_S\to E_{S'}\) between fibers of a
fibration \(p:\Ecat\to\Scat\) over shapes. This makes \(E\) an indexed
category over \(\Ncat\).

\subsection{Semidirect Product \texorpdfstring{\(\Ncat \ltimes \Ecat\)}{N ⋉ E}}

The combined reasoning algebra is the Grothendieck construction
\(\Ncat \ltimes E\). A morphism in \(\Ncat \ltimes E\) is a pair
\((n,e)\) consisting of a navigation morphism \(n\in\Ncat\) and an
explicit arrow \(e\in\Ecat\) in the appropriate fiber. Composition
interleaves navigation and explicit reasoning according to:
\[
(n_2,e_2)\circ(n_1,e_1)
 =
\bigl(n_2n_1,\ e_2\circ n_2^\ast(e_1)\bigr).
\]

Intuitively:
\begin{itemize}[leftmargin=2em]
  \item \(\Ncat\) says \emph{which view} of the explicit reasoning we
        are in (e.g., which model, which set of assumptions);
  \item \(\Ecat\) records \emph{what} explicit computation and
        diagnostics have been applied;
  \item \(\Ncat \ltimes \Ecat\) packages these into a single object
        representing a reasoning trajectory.
\end{itemize}

\subsection{Confluence, Normal Forms, and Decidability}

Appendix~A equips \(\Ncat \ltimes E\) with a reduction system that
combines explicit rewrites and navigation moves. Provided that:
\begin{itemize}[leftmargin=2em]
  \item explicit reduction is terminating and locally confluent, and
  \item navigation reduction normalizes within each component of
        \(\Ncat\),
\end{itemize}
we obtain:

\begin{theorem}[Confluence and normal forms (informal)]
Every reasoning trajectory in \(\Ncat \ltimes \Ecat\) admits a
canonical normal form, unique up to navigation isomorphism. In
particular, explicit steps and navigation steps may be permuted and
simplified without changing the canonical form.
\end{theorem}

Moreover, given a well-specified reduction system and typing rules, the question \emph{``Is this composite a valid FRFP pipeline?''} is decidable.


\begin{theorem}[Decidability of admissibility (informal)]
Given a composite term in \(\Ncat \ltimes \Ecat\), it is decidable
whether it corresponds to an admissible pipeline of the form
\(\RI \to (\EC/\ED)^\ast \to \RB \to \TE \to \HFD\).
\end{theorem}

These properties make \(\Ncat \ltimes \Ecat\) a suitable intermediate
language for representing and analyzing complex reasoning trajectories.
\paragraph{Operational significance of Epistemic Algebra canonicality.}
Appendix~A strengthens this picture by showing that the Epistemic Algebra architecture of FRFP is
canonical among all architectures with the same primitive structure: a free explicit category
generated by TDG, a navigation groupoid acting fiberwise via a fibration $p : E \to S$, a
Grothendieck-style semidirect product $N \ltimes E$, and preservation of the RB--TE--HFD
boundary under navigation. Any other architecture that satisfies these axioms admits a unique
structure-preserving morphism from the FRFP Epistemic Algebra construction. Operationally, this means
that ``small local changes'' to task decomposition or navigation---while still respecting the same
roles and constraints---cannot evade the correctness-preservation, non-automation, and
inevitability results: they are inherited along the unique morphism from FRFP. To genuinely
escape those consequences, one must step outside the Epistemic Algebra axioms themselves (for example
by allowing tacit-to-explicit morphisms or by breaking boundary preservation), rather than
modifying TDG or navigation in a cosmetic way.

At the explicit algebra layer, the airfare interaction (introduced in Section~\ref{sec:frfp-architecture-overview}), is represented as an explicit pipeline \(P_{\mathrm{air}} \in N \ltimes E\) that composes AI-assisted steps such as preference elicitation, option generation, and trade-off analysis. Different prompt sequences that are operationally admissible but differ only by refactoring---for example, asking for trade-offs before seeing full itineraries, or regrouping constraints into two separate messages---are related by navigation morphisms in \(N\) and rewrite rules in the explicit reduction system.
 
By construction, our rewrite system on \(N \ltimes E\) is terminating and confluent, so each such pipeline admits a canonical normal form. Correctness is defined at the level of the canonical pipeline via the semantics functor \(\llbracket - \rrbracket : N \ltimes E \to T\) and the quotient \(\gamma\). Informally: if one way of using the airfare assistant is correct (the human, after reviewing the explicit artefacts, makes an acceptable tacit choice), then all ways of using the assistant that are rewritable to that pipeline by admissible refactorings are equally correct.


%=========================================================
\subsection{Semantic Correctness and Invariance}
\label{sec:semantics}
%=========================================================

The combination of Epistemic Space and Epistemic Algebra yields a semantic map from composite
trajectories to tacit states.

\subsubsection{Semantic Functor and Observable Correctness}

The semantic functor
\[
\llbracket - \rrbracket: \Ncat \ltimes E_{\mathrm{sem}}\to\Tcat
\]
interprets a trajectory by running its explicit part, applying the
boundary \(\RB\), and then applying tacit evaluation and final decision.
Observable correctness is given by \(\gamma(\llbracket - \rrbracket)\).

\begin{theorem}[Correctness preservation (informal)]
Correctness judgments are invariant under explicit reasoning,
navigation, and semidirect-product composition: for any admissible
trajectory \(\tau\) and any navigation morphism \(n\in\Ncat\),
\[
\gamma(\llbracket \tau \rrbracket)
 =
\gamma(\llbracket n\cdot \tau \rrbracket).
\]
\end{theorem}

In other words, changing the route by which we arrive at a result---so
long as we respect explicit typing and Epistemic Spaces/Algebra constraints---does not
change the tacit correctness judgment.

\subsubsection{Canonical Correctness Form}

Every admissible pipeline admits a canonical factorization:

\begin{theorem}[Canonical correctness form (informal)]
Any trajectory in \(\Ncat \ltimes E_{\mathrm{sem}}\) that is admissible
in the FRFP sense is equivalent, up to navigation, to a unique normal
form in which:
\begin{itemize}[leftmargin=2em]
  \item all explicit operations (including navigation) occur before the
        boundary morphism \(\RB\);
  \item the tacit suffix is \(\TE \to \HFD\), with no interleaving
        explicit operations;
  \item no additional boundary or tacit operators appear.
\end{itemize}
\end{theorem}

This expresses, at the algebraic level, the design principle that AI
systems may explore and manipulate explicit representations in many
ways, but correctness judgments always arise from a single passage
through \(\RB\) followed by tacit evaluation and decision.

%=========================================================
\section{Epistemic Dynamics}
\label{sec:epistemicdynamics}
%=========================================================

The conversational calculus described so far is static: it assigns tacit judgments to whole explicit workflows but does not yet model how those judgments evolve over time. This section adds a dynamic layer on top of the conversational calculus, introducing tacit degradation, requirement schedules, hallucination, and safe horizons.

The structures introduced in Epistemic Spaces and Epistemic Algebra characterise \emph{single}
FRFP pipelines: how explicit and tacit operations are separated, how
explicit reasoning and navigation combine, and how correctness is
preserved under admissible transformations. In typical long-horizon use,
however, a human agent engages with an AI system repeatedly. The tacit context from one episode influences the next,
and errors or misplaced trust may accumulate.

Epistemic Dynamics extends the framework to this setting by equipping FRFP
pipelines with a dynamic semantics on tacit states. The aim is
to capture how tacit grounding, task requirements, and confidence
interact over time, and to make precise in what sense hallucinations
become structurally difficult to avoid. The formal development of this
dynamic semantics is given in Appendix~\ref{app:dynamics}; the
present section outlines the main ingredients and results.

\subsection{Tacit Context, Degradation, and Requirements}

We view the human agent’s epistemic background as a tacit context space $T$ equipped with a
preorder $\preceq$:
\[
  t_1 \preceq t_2 \quad\text{“$t_1$ is epistemically no stronger than $t_2$”.}
\]
Strict inequality $t_1 \prec t_2$ denotes $t_1 \preceq t_2$ and $t_2 \not\preceq t_1$. No further
algebraic structure is imposed on $T$ at this stage; the categorical structure of tacit space from
Appendix~A is still present but the dynamic analysis works at the level of the underlying
preordered set.


A tacit state $t \in T$ is not merely isolated propositional knowledge, but a compressed record of prior interaction, aggregating experience accumulated over years through performing manual checks, encountering edge cases, debugging failures, and observing where formal specifications cease to align with practice. Tacit state is path-dependent: one only
reaches a given $t$ by walking through particular sequences of work, not by reading a finite
specification once. In this picture, the degradation operator $\delta : T \to T$ and the requirement
map $\Req : E \times \mathbb{R}_{\ge 0} \to T$ express how this compressed history evolves over
time: $\delta$ erodes available grounding through repeated use, while $\Req(e,c)$ encodes the
tacit floor needed to treat an explicit artefact $e$ at credence $c$ as epistemically adequate. Our
later inevitability results will rely precisely on this interaction between path-dependent tacit
history, degradation, and growing or recurring requirements.
\subsection{Epistemic degradation under long-run AI assistance}
\label{subsec:epistemic-degradation}

A central premise of this paper is that sustained reliance on AI assistants can induce
\emph{epistemic degradation} in human users. We use this term in a local, task-relative sense,
not as a global claim that ``humans become stupid.'' Let $D$ be a family of tasks (e.g.,
reading and evaluating papers, designing experiments, checking proofs). For a given user,
let $t_t \in T$ denote their tacit epistemic state at time $t$: a compressed record of prior
interactional experience---performing manual checks, encountering edge cases, debugging
failures, and seeing where formal specifications cease to match practice---that enables them
to solve tasks in $D$ and to notice when something is wrong.

\begin{definition}[Epistemic degradation]
\label{def:epistemic-degradation}
Fix a task family $D$ and a baseline learning process in which a user develops by
performing tasks in $D$ without an AI assistant (or with only minimal, non-substitutive
tools). We say that the user is \emph{epistemically degraded} on $D$ at time $t$ (relative to
this baseline) if, when evaluated \emph{without} access to an assistant, their independent
performance on $D$ and their ability to detect and correct errors on $D$ are strictly worse
than under the baseline, due to differences in the evolution of their tacit state $t_t$.
\end{definition}

Our claim is that, under conditions that already hold in many research workflows, long-run
use of AI assistance makes epistemic degradation on relevant task families $D$ the default
outcome. We cannot prove this as a theorem about all possible humans and all possible AI
systems, but we can justify it as a modeling assumption based on three well-established
mechanisms.

\paragraph{Automation, complacency, and deskilling.}
A large human-factors literature shows that when automation and decision aids become
reliable, humans tend to become \emph{monitors} rather than \emph{primary actors}, with
resulting automation bias (over-trust in automated outputs) and automation-induced
complacency (insufficient checking) \citep[e.g.][]{brown_effect_automation_aviation_2017,
okine_evolution_hf_aviation_2025,singh_automation_complacency_2024}. In aviation,
for example, increased reliance on autopilot and flight management systems has been linked
to loss of basic manual flying skills, mode confusion, and degraded situation awareness,
especially among less experienced pilots \citep{brown_effect_automation_aviation_2017,
okine_evolution_hf_aviation_2025}. Similar dynamics are now documented for AI systems
in medicine: a mixed-method review of AI in clinical practice identifies
\emph{deskilling}---degradation of previously acquired competencies due to reduced
practice---and \emph{upskilling inhibition}---failure to develop advanced skills among
trainees because AI systems are doing the difficult parts---as central risks of AI adoption
\citep{natali_ai_deskilling_medicine_2025}. Commentaries and empirical studies on AI
decision support in diagnostics explicitly worry about over-reliance, skill degradation,
diagnostic deskilling, and loss of vigilance if clinicians defer too readily to AI outputs
\citep{sunday_behavioral_impacts_diagnostics_2025,berzin_preserving_clinical_skills_2025,
fatima_ai_triage_deskilling_2025}.

In our notation, the assistant increasingly intercepts those task instances in $D$ that
would otherwise provide practice and feedback to update $t_t$. Observable performance
for the \emph{combined} human--AI system may go up; but the contribution from the human's
tacit state shrinks, so that their independent competence is worse than in the no-assistant
baseline.

\paragraph{Cognitive offloading and the ``Google effect.''}
A second mechanism is cognitive offloading: the tendency to rely on external systems for
storage and processing when they are easily available. The ``Google effect'' literature
shows that intensive Internet search leads people to treat the web as an external memory
store: they remember \emph{where} to find information rather than the information itself,
and forgetting of details accelerates when future access is anticipated
\citep{gong_yang_google_effect_2024}. Gong and Yang's meta-analytical review documents
that intensive search behavior systematically changes memory and retrieval mechanisms in
this way \citep{gong_yang_google_effect_2024}. Related work on technology use and
cognitive aging argues that increased dependence on external media for semantic,
autobiographical, and spatial memory can accelerate decline if it reduces practice with
one's own intrinsic abilities \citep{benge_tech_use_cognitive_aging_2025}.

Recent work extends this to AI tools specifically. Gerlich's study on AI tools and critical
thinking finds a significant negative relationship between heavy AI-tool usage and
critical-thinking scores, with cognitive offloading statistically mediating this effect:
participants who offload more reasoning to AI tools score worse on standardized critical
thinking tests \citep{gerlich_ai_tools_cognitive_offloading_2025}. Qualitative synthesis
of this and related work emphasizes that users come to rely on AI for analysis and
synthesis, not just for recall, and consequently engage less in deep, effortful thinking on
their own \citep{gerlich_ai_cognitive_implications_2025}.

In our framework, this means that both the \emph{contents} and the \emph{procedures}
required to perform tacit checks on tasks in $D$ are increasingly stored in, and executed
by, the assistant rather than encoded in $t_t$. The user becomes more proficient at calling
the assistant, but worse at reconstructing, critiquing, or replacing its reasoning.

\paragraph{Early evidence of degraded independent performance under AI use.}
Third, there is emerging empirical work that directly measures how AI assistance affects
independent performance over time. In medicine, studies of AI-supported diagnostics show
clear short-term gains in accuracy and speed for physicians using AI decision support
\citep[e.g.][]{anderson_deep_learning_physician_accuracy_2024,lai_concurrent_read_ai_2022,
lim_medical_ai_workload_2025}, but parallel reviews now explicitly flag the risk that this
support will erode core clinical skills if it replaces, rather than augments, trainee
practice \citep{natali_ai_deskilling_medicine_2025,berzin_preserving_clinical_skills_2025,
sunday_behavioral_impacts_diagnostics_2025,fatima_ai_triage_deskilling_2025}. Similar
concerns appear in stroke imaging and radiology training, where triage and detection
algorithms risk relegating trainees to passive oversight roles
\citep{fatima_ai_triage_deskilling_2025,savardi_ai_radiology_training_2025}.

Outside medicine, experimental work on generative-AI use in writing tasks has begun to
measure neural and behavioral effects. An EEG-based study of essay writing with and
without a chatbot assistant finds that participants using the assistant show lower brain
engagement in regions associated with attention, planning, and memory, and that over
several months they exhibit declining independent performance and increasing reliance on
copy-and-paste use of AI outputs \citep{kosmyna_brain_on_chatgpt_2025}. Journalistic
syntheses of this and related studies similarly report that frequent chatbot use is
associated with reduced critical thinking and weaker performance when the tool is removed,
especially among younger or less-experienced users
\citep{washpost_ai_rewiring_minds_2025,guardian_ai_harming_intelligence_2025}.

These results are not yet definitive, and some meta-analyses also find positive effects of
generative AI on specific higher-order skills in controlled educational settings
\citep{zhao_genai_higher_order_thinking_2025}. However, taken together with the older
automation and offloading literatures, they make it reasonable to model epistemic
degradation as the \emph{default trajectory} when high-performing assistants replace, rather
than structure, human practice on a task family $D$.

\paragraph{The formal assumption.}
We can now state the assumption our protocol is designed to address.

\begin{assumption}[Degradation under sustained assistance]
\label{assumption:degradation}
Let $D$ be a family of tasks, and consider two learning processes for the same user:

\begin{enumerate}
    \item A \emph{baseline} process in which the user solves tasks in $D$ without access to
    an AI assistant (or with only non-substitutive tools), updating $t_t$ through unaided
    problem-solving, feedback, and explicit practice.

    \item An \emph{assisted} process in which a high-performing AI assistant proposes
    candidate solutions and rationales for most tasks in $D$, and the user typically
    accepts these with limited independent checking.
\end{enumerate}

Assume that (i) human memory and skill exhibit standard decay in the absence of practice
or retrieval, as modeled in classical forgetting and skill-decay work
\citep[e.g.][]{gong_yang_google_effect_2024,benge_tech_use_cognitive_aging_2025}; and
(ii) the presence of the assistant systematically reduces the frequency and depth of
unaided problem-solving and checking on $D$, as suggested by the automation,
offloading, and deskilling studies discussed above
\citep{brown_effect_automation_aviation_2017,okine_evolution_hf_aviation_2025,
natali_ai_deskilling_medicine_2025,gerlich_ai_tools_cognitive_offloading_2025,
sunday_behavioral_impacts_diagnostics_2025}.

Then, beyond some time horizon, the assisted process leads to epistemic degradation on
$D$ in the sense of Definition~\ref{def:epistemic-degradation}: when evaluated without the
assistant, the user's independent performance and error-detection ability on $D$ are worse
than under the baseline process.
\end{assumption}

This assumption captures the worry that motivates our protocol design. As AI assistants
become more capable, human scientists risk becoming less capable \emph{relative to the
tasks they routinely delegate}: the joint human--AI system may look impressively
competent, but the next generation's tacit epistemic state $t_t$ fails to accumulate the
interaction-derived structure needed to generate and scrutinize scientific results on its
own. The protocols we propose in the remainder of the paper are explicitly designed to
break this dynamic: they aim to secure near-term accuracy while preserving, and in some
cases amplifying, the long-run development of human tacit competence.

\paragraph{Context degradation.}
Long-run use of an AI assistant is assumed to induce epistemic
degradation. This is modelled by a function \(\delta:T\to T\) satisfying:
\begin{enumerate}[label=(D\arabic*),leftmargin=2.2em]
  \item \emph{Monotonicity:} \(\delta(t) \preceq t\) for all \(t\in T\);
  \item \emph{Strict decay somewhere:} there exists \(t\) with
        \(\delta(t) \prec t\);
  \item \emph{Non-invertibility:} there is no \(f:T\to T\) such that
        \(\delta\circ f = \mathrm{id}_T\).
\end{enumerate}
Intuitively, \(\delta\) aggregates effects such as fatigue, forgetting,
and inappropriate deference to the system. Repair operations (rest,
study, institutional checks) can be represented by separate operators,
but are not needed for the core Epistemic Dynamics results and are therefore left
for later extensions.

\paragraph{Requirements and hallucination.}
Hallucination is defined at the level of the protocol, rather than in
terms of any particular model. To this end, we introduce a requirement
map
\[
\Req : E \times \mathbb{R}_{\ge 0} \to T,
\]
where \(e\in E\) is an explicit artifact (e.g., an explanation or
derivation) and \(c\ge 0\) is a tacit credence level indicating how
strongly the agent treats \(e\) as reliable. The element \(\Req(e,c)\)
represents the tacit grounding required to treat \((e,c)\) as
epistemically adequate.

We assume:
\begin{enumerate}[label=(R\arabic*),leftmargin=2.2em]
  \item \emph{Monotonicity in credence:} if \(c_1 \le c_2\) then
        \(\Req(e,c_1) \preceq \Req(e,c_2)\) for all \(e\in E\);
  \item \emph{Nontriviality:} not all artifacts have the same
        requirements up to \(\preceq\).
\end{enumerate}

\begin{definition}[Hallucination]
Let \(t\in T\) be the current tacit state, \(e\in E\) an explicit output,
and \(c\ge 0\) a credence level. A \emph{hallucination} occurs at this
step if
\[
t \prec \Req(e,c).
\]
We then say that \(e\) is \emph{hallucinated} (relative to \((t,c)\)).
\end{definition}

In this sense, hallucination is a mismatch between available tacit
grounding and what would be required to support the explicit artifact at
the assigned credence. This definition does not depend on internal model
mechanisms or intent.

\paragraph{Repair, workload, and the cost--speed tradeoff.}
The Epistemic Dynamics analysis above treats degradation as one-sided: tacit states evolve via a monotone,
non-invertible operator $\delta$ under recurring nontrivial demands, yielding inevitability of
requirement-violating steps. Real institutions also invest in \emph{repair}: training, feedback,
reflection, and retraining that partially restore grounding. We model these as operators
$\rho : T \to T$ that decrease an entropy-like potential $S : T \to \mathbb{R}$ on tacit states, while
$\delta$ increases $S$ whenever nontrivial work is done.

If average entropy production from nontrivial tasks exceeds the average entropy removal achievable
by bounded repair operators $\rho$, then $S$ drifts upward and eventually crosses any fixed
requirement threshold, recovering the inevitability results even in the presence of repair.
Conversely, stabilizing $S$ requires repair at a rate and intensity comparable to the rate at which
tasks are processed. In high-load regimes this typically means a governance workload whose cost and
latency scale with AI usage. Under the FRFP assumptions there is thus no regime in which one
simultaneously obtains unbounded speed/cost gains from automation and long-horizon immunity to
FRFP-hallucination: strong repair can delay and reshape risk, but only by re-introducing substantial
human effort (see Appendix~\ref{app:delta-rho} for a formal $\delta$--$\rho$ analysis).

\subsection{Pipeline Transformers and Runs}
\label{subsec:epistemicdynamicsruns}

The static semantics of Appendix~\ref{app:math-foundations} already
associates to each explicit pipeline \(e\in E_{\mathrm{sem}}\) a tacit
effect via the semantic RB functor \(\RBsharp\) and the tacit primitives
\(\TE,\HFD\). In particular, evaluation of \(e\) at tacit context \(t\)
can be written as
\[
t \longmapsto \HFD\circ \TE \circ \RBsharp(e,t).
\]

For dynamic analysis, we compose this with degradation.

\begin{definition}[Pipeline transformer]
A complete pipeline at a given turn consists of:
\begin{itemize}[leftmargin=2em]
  \item an explicit artifact \(e \in E_{\mathrm{sem}}\) constructed as in
        Epistemic Spaces and Epistemic Algebra;
  \item the tacit suffix \(\RB \to \TE \to \HFD\);
  \item a degradation step \(\delta:T\to T\).
\end{itemize}
It induces a map
\[
\llbracket P \rrbracket_\delta : T \to T,\qquad
\llbracket P \rrbracket_\delta(t)
 :=
 \delta\bigl(\HFD\circ\TE\circ\RBsharp(e,t)\bigr).
\]
\end{definition}

\begin{definition}[Run]
A \emph{run} (or \emph{interaction sequence}) is a sequence of tacit
states \((t_n)_{n\ge 0}\) such that, for some sequence of pipelines
\(P_0,P_1,\dots\),
\[
t_{n+1} = \llbracket P_n \rrbracket_\delta(t_n)
\quad\text{for all } n\ge 0.
\]
When convenient, we also record the corresponding explicit artifacts and
credences, writing \((t_n,e_n,c_n)_{n\ge 0}\).
\end{definition}

All Epistemic Dynamics conclusions are stated in terms of such runs and the
hallucination condition \(t_n \prec \Req(e_n,c_n)\).

\subsection{Baseline Hallucination Inevitability}

We first consider a deterministic setting with fixed degradation
\(\delta\) and fixed pipelines. To capture the notion that nontrivial
epistemic demands recur, we use a requirement floor.

\begin{definition}[Requirement floor and recurring demands]
A state \(r\in T\) is a \emph{requirement floor} if, for a run
\((t_n,e_n,c_n)_{n\ge 0}\), we say that \emph{nontrivial demands recur
relative to \(r\)} when there exists an infinite subsequence
\((n_k)_{k\ge 0}\) such that
\[
\Req(e_{n_k},c_{n_k}) \succeq r
\quad\text{for all }k.
\]
\end{definition}

We strengthen the degradation assumption to ensure that tacit states
eventually fall below this floor.

\begin{assumption}[Degradation below the requirement floor]
\label{assump:baseline-degradation-main}
There exists \(r\in T\) such that, for every \(t_0\in T\), there is
\(N\in\mathbb{N}\) with
\[
\delta^N(t_0) \prec r.
\]
\end{assumption}

Under these assumptions, Appendix~\ref{app:dynamics} shows:

\begin{theorem}[Hallucination inevitability, deterministic case]
\label{thm:baseline-inevitability-main}
Let \((t_n,e_n,c_n)_{n\ge 0}\) be a run for which
Assumption~\ref{assump:baseline-degradation-main} holds and
nontrivial demands recur relative to some requirement floor \(r\). Then
there exists \(n\in\mathbb{N}\) such that
\[
t_n \prec \Req(e_n,c_n).
\]
In particular, every infinite run of this form contains at least one
hallucination.
\end{theorem}

The theorem applies to runs that are both degrading and nontrivial in
the sense above. Scenarios restricted to very easy tasks, very short
horizons, or aggressive external re-grounding may not satisfy these
premises and are not targeted here.

\subsection{Confidence, Requirement Escalation, and Feedback}

The behaviour of the human agent depends not only on the explicit
content \(e_n\) but also on how plausible it appears. Epistemic Dynamics therefore
adds structure for explicit plausibility, credence, and feedback between
credence and degradation. The formal definitions and proofs are given in
Appendix~A.9.8–A.9.13; we summarise the main ideas.

\paragraph{Plausibility and confidence.}
An explicit plausibility functional \(\pi:E\to[0,1]\) assigns to each
artifact a scalar measure of surface coherence or apparent authority.
Tacit state is extended to \(T^e:=T\times\mathbb{R}_{\ge 0}\), with
elements \((t,c)\) where \(t\) is grounding and \(c\) is credence. A
confidence inflation operator
\(\Phi:(T^e\times E)\to T^e\) updates credence according to plausibility,
without improving \(t\); for example
\[
\Phi((t,c),e) = (t,c+\varphi(\pi(e))),
\]
for some monotone \(\varphi:[0,1]\to\mathbb{R}_{\ge 0}\).

\paragraph{Requirement escalation and feedback-coupled degradation.}
Since \(\Req(e,c)\) is monotone in \(c\), higher credence raises
requirements. To reflect the possibility that higher confidence reduces
scrutiny, degradation is allowed to depend on \(c\) via a map
\(\delta^\star:T^e\to T\), written \(\delta^\star(t,c)=\delta_c(t)\),
such that:
\begin{itemize}[leftmargin=2em]
  \item for fixed \(t\), higher \(c\) leads to weaker outputs:
        \(\delta^\star(t,c_2)\preceq \delta^\star(t,c_1)\) whenever
        \(c_1\le c_2\);
  \item at zero credence, \(\delta^\star\) reduces to the baseline
        \(\delta\).
\end{itemize}

An extended run \((t_n,c_n,e_n)_{n\ge 0}\) is then obtained by
interleaving confidence inflation and feedback-coupled degradation. Under
assumptions that credence can grow without bound along subsequences when
high-plausibility artifacts recur, Appendix A.9.13
establishes an \emph{accelerated} version of
Theorem~\ref{thm:baseline-inevitability-main}: hallucinations still
occur in finite time, and the time-to-hallucination is not increased,
and may be strictly reduced, by the presence of the feedback loop.

\subsection{Stochastic Pipelines, Ensembles, and Finite Horizons}

The preceding discussion treats runs as deterministic sequences. In many
applications, FRFP pipelines include stochastic components (e.g.,
sampling, randomised retrieval). Appendix~A.9.15–A.9.16
extends the dynamic semantics to this case by viewing runs as
trajectories \(\omega\) in a space \(\Omega\) and equipping \(\Omega\)
with probability measures arising from stochastic implementations.

For each trajectory \(\omega = ((t_n,c_n,e_n))_{n\ge 0}\), the
\emph{hallucination time} is defined as
\[
\tau(\omega)
 :=
\inf\{n\ge 0:\ t_n \prec \Req(e_n,c_n)\},
\]
with \(\tau(\omega)=+\infty\) if no hallucination occurs. When the
structural premises of Theorem~\ref{thm:baseline-inevitability-main}
(and its feedback-augmented variants) hold pointwise on the support of a
stochastic pipeline, Appendix A.9.15–A.9.16 shows that
\(\tau(\omega)<\infty\) for all such \(\omega\). Any probability measure
supported on these trajectories then satisfies
\(\mathbb{P}(\tau<\infty)=1\): hallucination is almost sure.

For a fixed horizon \(N\), the finite-horizon hallucination probability
\(p_N:=\mathbb{P}(\tau\le N)\) is non-decreasing in \(N\) and converges
to~1 as \(N\to\infty\). For \(K\) independent runs, the probability that
at least one run hallucinates by time \(N\) is
\(1-(1-p_N)^K\). Thus ensembles can improve robustness over short
horizons, but they do not change the long-run structural inevitability
when the underlying conditions are met.

\subsection{Implications}

Epistemic Dynamics shows how the static FRFP architecture behaves over long
interaction sequences when combined with simple models of context
degradation, requirements, and confidence. The main conclusions are:
\begin{itemize}[leftmargin=2em]
  \item hallucination is naturally captured as a mismatch between tacit
        context and task requirements, independent of any particular AI
        architecture;
  \item under monotone degradation and recurring nontrivial demands,
        hallucinations are inevitable along sufficiently long runs;
  \item feedback from plausibility and credence to degradation tends to
        shorten the time-to-hallucination rather than extend it;
  \item stochasticity and ensembles can reduce error on finite horizons,
        but do not remove the long-run structural pressure towards
        hallucination under the same assumptions.
\end{itemize}

The formal statements and proofs are given in
Appendix~\ref{app:dynamics}. Subsequent sections extend the
framework from single agents to multi-agent and institutional settings.

The dynamics layer enriches tacit states with degradation and requirements. In the airfare setting, the requirement map \(\mathrm{Req}(e,c)\) encodes, for a given explicit configuration \(e\) (set of options, price ranges, constraints) and credence level \(c\), how much tacit grounding is required for it to be appropriate to accept an itinerary as ``best''. For example, a corporate travel policy might demand that the human understand not only the price and timing, but also fare rules and change penalties before accepting a premium non-refundable ticket.

As the human interacts with the assistant across many planning tasks, tacit grounding degrades via \(\delta : T \to T\) under cognitive load or time pressure, while nontrivial requirements recur. Within our Epistemic Dynamics model, prolonged or repeated use of such a ``best airfare'' protocol without adequate rest or oversight will eventually lead to requirement-violating acceptances—for instance, choosing an itinerary that superficially fits high-level preferences but violates tacit constraints the human would normally enforce (such as minimum connection times or visa requirements). The safe-horizon functional \(H_P(\varepsilon)\) then captures, for a given pipeline \(P\) and tolerance \(\varepsilon\), how many such planning episodes can be run before the probability of a tacitly unacceptable choice exceeds \(\varepsilon\).

% ------------------------------------------------------------------
% Section X — Collective Epistemic Dynamics
% ------------------------------------------------------------------

\section{Collective Epistemic Dynamics}
\label{sec:collective-epistemic-dynamics}

We now lift the conversational calculus to populations and institutions. Multiple human agents share explicit artefacts and communicate via typed workflows; their tacit states evolve individually but interact through shared artefacts, committees, and policies. The resulting collective layer supports population-level dynamics, consensus mechanisms, and governance operators, and yields explicit-only impossibility results for hallucination detection and institutional automation. Technical details and categorical formulations appear in the appendices, in particular Appendices~A--C.


\subsection{From Single Pipelines to Epistemic Populations}

In the single-agent setting, an FRFP pipeline comprises:
\begin{itemize}
  \item an explicit artefact space $E$ (texts, code, plans, proofs),
  \item a tacit space $T$ with a preorder $\preceq$ capturing ``at least as grounded as'',
  \item an extended tacit state space $T^e = T \times \mathbb{R}_{\ge 0}$ pairing grounding and tacit credence,
  \item dynamics combining degradation $\delta$, requirement maps $\mathrm{Req}$, confidence inflation, and feedback-coupled degradation.
\end{itemize}
A single execution is a trajectory
\[
  \omega = \bigl( (t_n,c_n,e_n) \bigr)_{n \ge 0},
\]
and hallucination occurs at the first $n$ such that $t_n \prec \mathrm{Req}(e_n,c_n)$ while the explicit artefact $e_n$ is nevertheless accepted in the pipeline.

To move to the collective level, we replace the single tacit space by a family
\[
  A_i = (T_i, \preceq_i, \delta_i, T^e_i),
  \qquad i \in I,
\]
of epistemic agents indexed by a set $I$ (individuals, teams, or institutions). Each $A_i$ has its own tacit space, requirement relation, and degradation dynamics. Agents share a common explicit reasoning environment described by the category
\[
  \mathcal{C} = N \ltimes E,
\]
where $E$ collects explicit artefacts and $N$ captures navigation and control structure at the explicit level (task selection, mode switching, tool invocation, and so on).

The joint explicit state of the population is modeled as an object of the product category
\[
  \mathcal{C}^I.
\]
An object $X \in \mathcal{C}^I$ is a family $X = (X_i)_{i\in I}$ of explicit states, indicating for each agent which explicit task and artefacts are currently active. A morphism $X \to Y$ in $\mathcal{C}^I$ represents one step of explicit activity across the population: one or more agents update their explicit state, for example by computing, revising, delegating, or producing new artefacts.

A population semantics functor
\[
  \overline{J} : \mathcal{C}^I \longrightarrow \prod_{i\in I} T^e_i
\]
associates to each global explicit state $X$ a tuple
\[
  \overline{J}(X) = \bigl((t_i,c_i)\bigr)_{i\in I}
\]
of extended tacit states, obtained by allowing each agent to consume the explicit artefacts in $X$ and update according to the single-agent dynamics. This lifts explicit--tacit separation to the population level: explicit artefacts are shared, while tacit grounding remains agent-local.

\subsection{Collective Dynamics as a Markov System}

To analyze temporal behaviour, we endow $\mathcal{C}^I$ with stochastic dynamics. From any global explicit state $X$, there is a distribution over possible successors $Y$, reflecting:
\begin{itemize}
  \item which agents act next (possibly determined by a scheduling or institutional mechanism),
  \item which explicit morphisms in $\mathcal{C}$ they apply (such as computation, summarization, deployment, or consensus),
  \item which agents observe the newly produced artefacts (a communication structure on $I$),
  \item how each tacit state $(t_i,c_i)$ is updated via representational backflow, degradation, confidence inflation, and feedback-coupled degradation.
\end{itemize}

Formally, this is captured by a Markovian collective dynamics
\[
  M : \mathrm{Obj}(\mathcal{C}^I)
    \longrightarrow \mathcal{P}\bigl( \mathrm{Obj}(\mathcal{C}^I) \bigr),
\]
which specifies for each $X$ a probability measure $M(X,\cdot)$ on successor states. Equivalently, one can view $M$ as a functor
\[
  M : \mathcal{C}^I \longrightarrow \mathbf{Kl}(D),
\]
into the Kleisli category of the discrete distribution monad $D$, but the object-level transition kernel suffices for present purposes.
Fix an initial state $X_0\in \mathrm{Obj}(\mathcal{C}^I)$. The transition kernel $M$ then induces a
probability measure on the space of infinite executions
\[
  \Omega := \mathrm{Obj}(\mathcal{C}^I)^{\mathbb{N}},
\]
namely the law of the Markov chain on $\mathrm{Obj}(\mathcal{C}^I)$ with initial condition $X_0$ and
transition kernel $M$. We denote this measure by $P_{M,X_0}$. When the initial state is clear from
context, we suppress it and write $P_M$.


A \emph{collective execution} is a random path
\[
  X_0, X_1, X_2, \dots
\]
in $\mathcal{C}^I$ generated by $M$ from an initial state $X_0$. The corresponding tacit trajectories $\overline{J}(X_n)$ record how grounding and credence evolve over time for each agent, coupled through shared explicit artefacts and communication. This yields a population-level path space and probability measure analogous to the single-agent setting, and permits the usual notion of stopping time with respect to the resulting filtration (see Appendix~A.10 for details).

\subsection{Collective Hallucination and Its Inevitability}

At the individual level, hallucination is the first time an explicit artefact is endorsed despite insufficient tacit grounding. At the collective level, the corresponding notion involves endorsement by a subset of agents.

Given a collective execution $(X_n)_{n\ge 0}$ with tacit states
\[
  \overline{J}(X_n) = \bigl((t_{i,n},c_{i,n})\bigr)_{i\in I},
\]
a \emph{collective hallucination} at time $n$ occurs if there exists an explicit artefact $e \in E$ and a subset $J \subseteq I$ with $|J|\ge 1$ such that:
\begin{itemize}
  \item each agent $j \in J$ has accepted or propagated $e$ (i.e., treated it as admissible for downstream use under their tacit gate) in their pipeline by time $n$, and
  \item for all $j \in J$,
        \[
          t_{j,n} \prec_j \mathrm{Req}_j(e,c_{j,n}),
        \]
        i.e.\ the tacit state lies strictly below the applicable requirement for $e$ in that context.
\end{itemize}
Here ``nontrivial'' means \emph{nonempty}: we allow $|J|\ge 1$, so the collective notion strictly extends the single-agent notion. When modeling multi-approval governance, we use the strengthened threshold $|J|\ge k$ (typically $k=2$), corresponding to acceptance rules that require multiple independent tacit gates.


The \emph{collective hallucination time} is the stopping time $\tau^{\mathrm{coll}}$ defined by
\[
  \tau^{\mathrm{coll}}(\omega)
  := \inf\{\,n \ge 0 : \omega \text{ exhibits a collective hallucination at step } n\,\},
\]
with $\inf \varnothing = \infty$ as usual.

Under the same high-level assumptions that drive single-agent inevitability (finite tacit capacities, non-invertible degradation for each agent, recurrent demands on grounding via surface-plausible artefacts, confidence inflation, and requirement escalation), and under mild assumptions on $M$ (continued production and circulation of artefacts through the population), one obtains an almost-sure finiteness result:
\[
  P_M\bigl(\tau^{\mathrm{coll}} < \infty\bigr) = 1.
\]
In words, over sufficiently long horizons, some under-grounded artefact will almost surely be collectively endorsed, provided the structural FRFP conditions remain satisfied. A formal statement and proof are given in Appendix~A.10.11--A.10.12.

\subsection{Consensus Mechanisms as Explicit-Only Functors}

Collective systems frequently employ consensus mechanisms: voting schemes, averaging procedures, ensemble methods, committee decisions, and more elaborate governance structures. Within the present framework, these are modeled as explicit-only operators
\[
  C : E^n \longrightarrow E,
\]
which aggregate a finite tuple of explicit artefacts $(e_1,\dots,e_n) \in E^n$ into a new explicit artefact $C(e_1,\dots,e_n)$. Crucially, such operators act solely in explicit space; they have no direct access to tacit states $T_i$ or to the requirement maps $\mathrm{Req}_i$.

A basic consequence is \emph{tacit-stability}. Fix an agent $i$ and a context $c_i$. Suppose that for each $k \in \{1,\dots,n\}$,
\[
  t_i \prec_i \mathrm{Req}_i(e_k,c_i),
\]
so each input artefact is under-grounded for agent~$i$ in context $c_i$. Because explicit transformations do not modify the underlying tacit state, any output $C(e_1,\dots,e_n)$ is evaluated against the same $t_i$ and $c_i$ and therefore remains under-grounded.

As a consequence, fully automated consensus processes that operate purely on explicit artefacts cannot, in general, eliminate hallucination or under-grounded decisions at the collective level. They can smooth, denoise, or delay failures by reshaping explicit behaviour, but they cannot create new grounding in tacit space. This extends the single-agent observation that explicit refactorings do not restore grounding to the population level.

\subsection{Structural Notions at the Collective Level}

The collective formalism introduces structural notions that do not appear at the static single-agent level and that will be used later.

\paragraph{Equivalence of collective protocols.}
Viewing $\mathcal{C}^I$ as the explicit category for the full population, different interfaces, workflows, or consensus arrangements can be represented as endofunctors
\[
  F : \mathcal{C}^I \longrightarrow \mathcal{C}^I
\]
that preserve the population semantics $\overline{J}$. If two designs $F_1$ and $F_2$ are related by explicit-only refactorings that leave $\overline{J}$ unchanged, they are \emph{equivalent} in the FRFP sense: they have the same tacit semantics and the same distribution of the collective hallucination time $\tau^{\mathrm{coll}}$. Any difference between them is then syntactic rather than epistemic.

\paragraph{Refinement order on architectures.}
Given two Markovian collective dynamics $M_1$ and $M_2$ defined on the same population, with corresponding collective hallucination times $\tau^{\mathrm{coll}}_1$ and $\tau^{\mathrm{coll}}_2$, define the finite-horizon risk curves
\[
  p_N^{(k)} := P_{M_k}\bigl(\tau^{\mathrm{coll}}_k \le N\bigr),
  \qquad k \in \{1,2\},\ N \in \mathbb{N}.
\]
We say that $M_2$ \emph{dynamically refines} $M_1$ if $p_N^{(2)} \le p_N^{(1)}$ for all $N$. This induces a partial order on collective designs: some architectures are strictly safer than others in the sense of never increasing finite-horizon hallucination risk. Explicit refactorings that preserve $\overline{J}$ cannot strictly improve this order; they can at best preserve it. In contrast, interventions that alter tacit structure (e.g.\ adding independent evaluators, modifying requirements, or slowing degradation) can produce genuine refinements. A more detailed treatment appears in Appendix~A.10.

\paragraph{Compositional risk and defense in depth.}
Many systems decompose into interacting subpopulations (for example, scientific, engineering, and policy components). Each subpopulation has its own collective hallucination time $\tau^{\mathrm{coll}}_{\mathrm{sub}}$ and hazard profile. At a coarse level, the overall collective hallucination time is bounded by
\[
  \tau^{\mathrm{coll}} = \min_{\mathrm{sub}} \tau^{\mathrm{coll}}_{\mathrm{sub}}.
\]
The system fails epistemically as soon as any critical component fails. The formalism therefore supports a compositional ``defense-in-depth'' analysis: subsystem-level guarantees on hazard profiles may, under appropriate coupling conditions, yield system-level guarantees. The appendices provide the technical apparatus needed to make such statements precise.
\paragraph{Do collectives reduce or increase hallucination risk on fixed horizons?}
Moving from a single agent to a population does not have a uniform effect on hallucination risk. On a
\emph{fixed horizon} \(N\), collectives can make hallucination \emph{less probable} or \emph{more probable},
depending on how acceptance is implemented and how correlated the agents' errors and evidence streams are.

\smallskip
\noindent
\textbf{When collectives can reduce fixed-horizon risk (defense in depth).}
If an explicit artefact must pass \emph{independent} tacit gates before it is accepted (e.g., two-person
sign-off, committee review, replication-style checks), then collective hallucination requires that multiple tacit states
simultaneously fall below their relevant requirements. When the gates are sufficiently independent---in the
sense that their under-groundedness events are not strongly positively correlated---the probability of a
collective hallucination by time \(N\) can be strictly smaller than in the single-agent case. Intuitively,
independence turns ``one agent misses the flaw'' into ``several agents must miss it in compatible ways.''
This is the collective analogue of layered safety checks: additional gates can improve finite-horizon safety
by reducing the chance that any single agent's degraded tacit state becomes the sole point of failure.

\smallskip
\noindent
\textbf{When collectives can increase fixed-horizon risk (amplification and speedup).}
Collectives can also \emph{increase} short-horizon risk when they raise throughput of artefact production and
circulation, or when their failures are highly correlated. Shared tools, shared prompts, shared datasets, or
shared institutional incentives can make agents' errors co-move, so that many agents become under-grounded
with respect to the same artefacts at roughly the same time. In addition, populations typically explore and
propagate far more candidate artefacts per unit time than a single user; even if the per-artefact
hallucination probability is unchanged, the probability that \emph{some} artefact is accepted while
under-grounded can increase simply because the system samples more opportunities for failure. Finally, if
``consensus'' operates as an explicit-only aggregator over weakly grounded artefacts, it can accelerate
endorsement by producing artefacts that look more authoritative without increasing tacit grounding.

\smallskip
\noindent
\textbf{Structural implication.}
The collective layer therefore supports two distinct design levers. First, \emph{gate structure} (how many
tacit acceptances are required, and by whom) can reduce finite-horizon risk when evaluations are genuinely
independent. Second, \emph{correlation and throughput control} (diversity of evidence, diversity of tools,
rate-limits, and escalation rules) becomes essential to prevent populations from becoming
high-throughput amplifiers of the same under-grounded artefacts. None of these effects contradict the
long-horizon inevitability results under the FRFP assumptions: layered collectives can reshape and delay
risk, but cannot, in general, remove it over unbounded horizons when tacit capacities degrade and nontrivial
requirements recur.
\paragraph{Non-identifiability from logs.}
Collective executions give rise to rich explicit traces: logs, metrics, dashboards, and derived indicators. All of these live in $\mathcal{C}^I$ or in its images under explicit functors. Tacit states $\overline{J}(X)$ and requirement maps $\mathrm{Req}_i$ remain unobserved. In general, many different tacit configurations and dynamics $M$ can induce the same (or arbitrarily similar) distributions over explicit traces. It follows that correctness and hallucination risk are \emph{non-identifiable} from logs alone: without additional assumptions, explicit-only monitoring and metrics cannot uniquely determine how well grounded a population's beliefs are.

\paragraph{Scaling and phase-transition questions.}
The population-level formulation also supports scaling questions. Parameterizing systems by population size $|I|$, communication structure, and per-agent degradation properties, one can ask how epistemic risk scales with the number of agents and tools, whether particular communication patterns are stabilizing or destabilizing, and whether qualitative changes (phase transitions) occur as complexity grows. The present work does not attempt a systematic scaling analysis, but the structural notions above identify the relevant objects for future investigation.

\subsection{Connections to Example Domains and Local Interventions}

The same collective apparatus applies across domains in which many agents interact over shared artefacts:
\begin{itemize}
  \item Scientific communities can be modeled as populations of agents (labs, groups, or individual scientists) with papers, datasets, and code as explicit artefacts, and peer review, replication, and citation practices as partial consensus mechanisms.
  \item Software development lifecycles can be treated as populations of developers, reviewers, operators, and tools with explicit artefacts such as requirements, designs, code, tests, and deployment configurations, coupled by CI/CD and incident-response processes.
  \item Strategic domains such as military planning or pandemic response involve agents in governments, institutions, and advisory bodies operating over explicit artefacts including plans, intelligence, models, and policies, with doctrines and institutions acting as higher-level operators on $M$ and on $\mathrm{Req}$.
  \item Markets and large-scale experiments provide examples where rich explicit traces (transaction logs, event data) remain insufficient to uniquely identify underlying beliefs or theories, paralleling the non-identifiability observation above.
\end{itemize}

The collective perspective also clarifies the status of local interaction patterns, such as common prompting schemes for language models (incentive framing, challenge prompts, stepwise guidance, high-stakes framing, self-evaluation prompts). These can be viewed as local interventions on:
\begin{itemize}
  \item the explicit pipeline shape in $N \ltimes E$ (for example, encouraging decomposition and self-checking),
  \item the effective requirement schedule (for example, weakening answers to hypotheses or raising the bar for strong claims),
  \item the human user's tacit dynamics (for example, prompting increased attention and skepticism).
\end{itemize}
Such interventions can materially improve finite-horizon correctness and enlarge safe horizons for particular tasks. They do not, however, alter the structural facts established in the underlying theory: explicit-only systems remain ungrounded, human tacit capacities remain finite and degrading, and, within the Epistemic Dynamics and population assumptions of FRFP, under-grounded collective decisions remain inevitable over unbounded horizons.

In summary, the collective layer provides a coherent extension of FRFP to multi-agent settings. It connects individual pipeline dynamics to institutional design, clarifies the limits of explicit-only interventions (including consensus mechanisms and monitoring), and identifies the structural levers---in particular, those acting on tacit dynamics and requirements---on which robust long-horizon human--AI collaboration must ultimately rely.

In a collective version of the airfare example, a team or organisation might use a shared assistant to propose itineraries for many travellers under a common travel policy. Individual tacit states live in agent-local copies of $T$, while the shared explicit category records policies, candidate options, and past choices. Consensus mechanisms and governance operators at this layer determine, for example, whether a manager's tacit judgment can override an employee's preferred itinerary, or how exceptions to the cheapest-available-fare policy are recorded. Our non-automation results imply that, even with strong explicit policies and logging, no explicit-only governance operator can, in general, guarantee that all accepted itineraries are tacitly acceptable across all such populations and protocols.

\section{Structural Limits of Explicit Mechanisms}
\label{sec:structural-limits-explicit}

The preceding development has described a class of human--AI collaborations in which
\begin{itemize}
  \item explicit artefacts and transformations are modelled in a category $E$ of explicit reasoning,
  \item navigation, prompting, and other meta-operations are represented in a category $N$,
  \item tacit states and evaluations live in a space $T$, and
  \item the combined system $N \ltimes E$ is equipped with a reduction relation and a semantic functor
    \[
      \llbracket - \rrbracket : N \ltimes E_{\mathrm{sem}} \longrightarrow T,
    \]
    with observable correctness extracted via a map $\gamma : \mathrm{Ob}(T) \to J$.
\end{itemize}
Under the FRFP axioms (ETS, HEG, HEC, AE), explicit computation is separated from tacit evaluation and cannot reconstruct tacit state or correctness except via designated human interaction. At the same time, the dynamic and collective layers developed earlier show that long-horizon behaviour is inherently stochastic: pipelines are implemented as stochastic processes over a space of admissible trajectories, and populations interacting through explicit channels induce Markovian dynamics over collections of agents.

This section explains how these ingredients jointly constrain what can be achieved by mechanisms that operate purely in explicit space. Formal statements and proofs of the main structural results are given in Appendix~D.

\subsection{Dynamic Quantities and Long-Horizon Behaviour}

When a fixed pipeline $P$ is executed under FRFP constraints, the dynamic semantics introduced previously associate to any admissible implementation:
\begin{itemize}
  \item a measurable space of trajectories $\Omega_{\mathrm{allowed}}$ of runs, each trajectory $\omega \in \Omega_{\mathrm{allowed}}$ corresponding to a reduction sequence in $N \ltimes E$ satisfying the interaction protocol constraints;
  \item a hallucination time random variable
    \[
      \tau : \Omega_{\mathrm{allowed}} \longrightarrow \mathbb{N} \cup \{\infty\},
    \]
    recording the first time a requirement-violating artefact is accepted along the run;
  \item finite-horizon hallucination probabilities
    \[
      p_N \;=\; \mathbb{P}(\tau \leq N)
    \]
    for $N \in \mathbb{N}$;
  \item a hazard sequence $(h_n)_{n \in \mathbb{N}}$ derived from the increments of $(p_N)$; and
  \item for each tolerance $\varepsilon \in (0,1)$, a safe-horizon functional
    \[
      H_P(\varepsilon) \;=\; \sup\{N \in \mathbb{N} : p_N \leq \varepsilon\}.
    \]
\end{itemize}
These objects summarise, in a long-horizon sense, how often and how soon a given implementation of $P$ will fail the tacit requirements encoded by $\Req$.

Under the epistemic dynamics assumptions (finite tacit capacity, monotone non-invertible degradation of tacit state, recurring nontrivial demands, and appropriate independence conditions), the model yields an inevitability theorem: along admissible trajectories in this setting, hallucination occurs almost surely in the limit. In particular, $\mathbb{P}(\tau < \infty) = 1$ for any such implementation, and for every fixed $\varepsilon \in (0,1)$ the safe horizon $H_P(\varepsilon)$ is almost surely finite.

These results are proved in the dynamic layer but are agnostic to the internal details of the explicit mechanisms; they depend only on the interaction between explicit artefacts, requirement maps, and degradation behaviour of tacit state.

\subsection{Population-Level Dynamics and Consensus}

When multiple agents interact through explicit channels, the population development introduces:
\begin{itemize}
  \item an explicit reasoning category $\mathcal{C}$ for individual agents;
  \item a product category $\mathcal{C}^I$ of explicit profiles across a population $I$;
  \item functorial consensus operators
    \[
      C : \mathcal{C}^n \longrightarrow \mathcal{C},
    \]
    modelling explicit-only aggregation rules (voting, averaging, reconciliation of drafts, and so on); and
  \item a Markovian collective dynamics $M$ on $\mathcal{C}^I$, obtained by iterating local update and interaction steps.
\end{itemize}
Given an epistemic population (agents, tacit spaces, and requirement maps) and such a dynamics $M$, one can define a collective hallucination time $\tau_{\mathrm{coll}}$ as the first step at which a nonempty subset of agents jointly accept an artefact that is under-grounded relative to their individual requirements.\footnote{When modelling multi-approval governance, one may strengthen ``nonempty'' to $|J|\ge k$ (typically $k=2$) to represent acceptance rules that require multiple independent tacit gates.}

Under assumptions analogous to those used for single-pipeline dynamics, one again obtains an almost-sure inevitability result: for admissible populations and dynamics consistent with the FRFP constraints, the collective stopping time $\tau_{\mathrm{coll}}$ is almost surely finite. Moreover, the almost-sure finiteness of $\tau_{\mathrm{coll}}$ is invariant under inserting additional explicit-only consensus layers: functorial consensus operators act purely in explicit space and cannot, by themselves, change the inevitability class of the process.

These collective results provide a quantitative description of long-horizon epistemic risk at the level of whole institutions or populations, using the same kind of stochastic and categorical apparatus as in the single-pipeline case.

\subsection{Tacit Dependence and Explicit-Only Mechanisms}

Against this backdrop, it is natural to ask which desirable properties one might hope to enforce by purely explicit means. In the present setting, explicit-only mechanisms are those whose internal computation and inputs/outputs are confined to $E$ (or $\mathcal{C}$ in the population case), and which do not gain access to tacit states or tacit evaluations other than through the designated interfaces already modelled in the semantics.

A recurring theme is that many natural predicates of interest are \emph{tacit-dependent}: their truth value can vary between admissible runs that share the same explicit artefacts and explicit interaction history but differ in their tacit evolution. Examples include:
\begin{itemize}
  \item whether a given explicit artefact meets the relevant tacit requirements at the point of use;
  \item whether hallucination occurs within a fixed horizon, as expressed by events of the form $[\tau \leq N]$;
  \item whether an implementation has a given $\varepsilon$-safe horizon, as expressed by inequalities $H_P(\varepsilon) \geq N$; and
  \item in the collective setting, whether a population remains hallucination-free up to time $N$, as expressed by events such as $[\tau_{\mathrm{coll}} > N]$.
\end{itemize}
Because the combined explicit trace does not uniquely determine the underlying tacit trajectory, predicates of this kind cannot be reduced to functions of explicit artefacts alone. This observation is formalised in Appendix~A.11, which defines tacit-dependent predicates in terms of the semantic functor $\llbracket - \rrbracket$ and observable correctness $\gamma(\llbracket P \rrbracket)$ and then studies their interaction with explicit-only mechanisms.

The central structural result, stated and proved there, shows that under ETS and AE, no mechanism confined to explicit space can enforce a tacit-dependent predicate across all admissible runs. The proof relies only on the fact that such a mechanism receives isomorphic explicit inputs across certain pairs of admissible runs whose tacit behaviour differs, and hence must behave identically even when the predicate of interest disagrees.

\subsection{Consequences for Detection, Elimination, and Consensus}

Several consequences follow immediately when this schema is instantiated with the dynamic and collective predicates arising from earlier sections.

\paragraph{Detection and elimination.}
The definition of hallucination in terms of requirement maps and tacit state yields a predicate of the form ``this explicit artefact is hallucinated'' which is tacit-dependent: the same explicit output can be acceptable or unacceptable depending on the tacit configuration in which it is evaluated. It follows that:
\begin{itemize}
  \item there is no perfect hallucination detector operating purely on explicit artefacts and logs; and
  \item no explicit-only policy can guarantee that hallucination never occurs along admissible runs, or that $\tau = \infty$ (or $\mathbb{P}(\tau = \infty) = 1$) holds uniformly across implementations supported on $\Omega_{\mathrm{allowed}}$.
\end{itemize}
Under the epistemic dynamics assumptions, inevitability theorems already show that $\tau$ is almost surely finite for admissible implementations, and that $H_P(\varepsilon)$ is almost surely finite for each fixed $\varepsilon \in (0,1)$. The structural results in Appendix~D complement this by ruling out the possibility of using explicit-only mechanisms to enforce $\tau = \infty$ or $H_P(\varepsilon) = \infty$ as design constraints.

\paragraph{Correctness oracles.}
Observable correctness is represented by $\gamma(\llbracket P \rrbracket)$ and is defined purely in tacit space, subject to HEG and HEC. For a fixed explicit artefact, different admissible runs can lead to different tacit evaluations and hence different observable correctness values. A hypothetical correctness oracle
\[
  \Omega : E \longrightarrow J
\]
that recovered $\gamma(\llbracket P \rrbracket)$ from the explicit artefact alone would thus be required to decide a tacit-dependent predicate using only explicit data. The formal results establish that no such oracle can exist within the FRFP axioms if it is confined to explicit computation.

\paragraph{Consensus and collective guarantees.}
In the population setting, consensus operators $C : E^n \to E$ and explicit coordination rules act entirely in $E$ and do not modify tacit states directly. They can be composed to form consensus pipelines and embedded in the Markov dynamics $M$ on population states. The structural results then imply that:
\begin{itemize}
  \item explicit-only consensus cannot create grounding ex nihilo: if each input artefact exceeds the grounding of an agent, so does the consensus output for that agent;
  \item no fully automated consensus pipeline can guarantee that its outputs are hallucination-free across all admissible populations and implementations; and
  \item more strongly, predicates such as $[\tau_{\mathrm{coll}} = \infty]$ are tacit-dependent and cannot be enforced by explicit-only transformations of the population dynamics.
\end{itemize}
Invariance results for $\tau_{\mathrm{coll}}$ in the collective dynamics further reinforce this picture: inserting additional explicit-only consensus layers can alter the distribution of collective risk but cannot change the fact that $\tau_{\mathrm{coll}} < \infty$ almost surely under the stated assumptions.

\subsection{Implications for Design and Governance}

Taken together, these structural limitations delineate the scope of what explicit-only design can accomplish in long-horizon human--AI collaboration.

On the one hand, the dynamic layer provides meaningful quantitative levers:
\begin{itemize}
  \item implementations can be compared via their finite-horizon risk profiles $(p_N)$ and hazard sequences $(h_n)$;
  \item the safe-horizon functional $H_P(\varepsilon)$ captures, for each tolerance $\varepsilon$, how long a system can typically operate before crossing a risk threshold;
  \item dynamic refinement orders allow one to classify implementations as safer or riskier over finite horizons.
\end{itemize}
On the other hand, the impossibility results anchored in tacit dependence show that no amount of explicit-only filtering, logging, or aggregation can convert a system whose long-run behaviour is almost surely fallible into one that is guaranteed safe across all admissible runs and populations.

Within this boundary, the design problem becomes one of shaping hazard profiles and safe horizons, structuring interactions so that collective dynamics refine a baseline rather than degrade it, and introducing additional forms of tacit evaluation and institutional oversight.

Appendix~B.6 develops concrete protocol templates and institutional patterns compatible with these constraints. Appendix~C instantiates the dynamic and institutional semantics in detailed case studies (ICU sepsis literature review and financial scenario analysis / stress testing). Appendix~B.11 provides a FRFP-compliant system-instruction set for an explicit-only assistant, and Appendix~B.27--B.30 analyses a set of FRFP-aware prompting patterns and their effects on hazard profiles and safe horizons. Together, these appendices turn the abstract design levers of this section into operational tools for building and evaluating concrete workflows.

Subsequent sections on governance, incentives, and institutional design can be read as exploring these remaining degrees of freedom: adjusting requirement maps, allocating human attention, and configuring collective processes so as to manage long-horizon risk in a way that is consistent with the FRFP axioms and with the structural limitations just described.

Formal statements and proofs of the general schema and its corollaries appear in Appendix~D; the present discussion is intended to make clear how those results integrate with the dynamic and collective layers and how they bear on the design of real-world human--AI collaborations.
\section{Institutions for Long-Horizon Human--AI Collaboration}
\label{sec:institutions-main}

The previous sections analyzed individual human--AI pipelines and their long-run behavior: how explicit pipelines acquire tacit grounding, how degradation and hallucination accumulate over time, and how collective epistemic dynamics emerge when many agents interact with shared systems. We now step up one more level, from collections of agents and runs to the \emph{institutions} that sit above them. The aim of this section is conceptual: to describe the role of institutions in stabilizing long-horizon practice, while leaving formal details to Appendix~A.12.

\subsection{From Individual Agents to Epistemic Populations}

Earlier, we treated a single human agent as a bearer of tacit context, capable of evaluating explicit artefacts and updating their internal state in response to interaction with an FRFP pipeline. At scale, however, we are never dealing with a single agent in isolation. Scientific communities, engineering teams, regulatory bodies, and user populations are all better modeled as \emph{epistemic populations}: structured collections of agents, each with their own tacit capacities and degradation dynamics.

Formally, we use the agent and population structures introduced in the collective dynamics analysis. Each epistemic agent carries:
\begin{itemize}
  \item a tacit context space $T_i$ with a preorder expressing ``at least as well grounded as'';
  \item a degradation operator $\delta_i : T_i \to T_i$ capturing background loss of tacit context;
  \item an extended tacit state $T^e_i$ that pairs tacit context with a notion of subjective credence.
\end{itemize}
For institutional purposes, we enrich this with a \emph{correctness quotient} $\gamma_i : T_i \to J$, where $J$ is the space of observable correctness judgments (accept/reject/revise, etc.). The tuple
\[
  (T_i, \preceq_i, \delta_i, T^e_i, \gamma_i)
\]
is the institutional ``signature'' of the agent: it records which tacit states they can occupy, how those states degrade, and how they express correctness judgments when asked.

An \emph{epistemic population} is then a family of such agents indexed by a set $I$, together with the morphisms that allow us to compare and transform populations. This lets us treat heterogeneous communities---with agents who differ in tacit capacity, degradation rates, and evaluative norms---within a single formal frame. Full categorical details are given in Appendix~A.12, but the upshot for the main story is that we can talk sensibly about how a fixed pipeline behaves when replicated across a population, and how institutions act on the resulting pattern of judgments.

\subsection{Replication and Population-Level Semantics}

Given a well-typed FRFP pipeline $P$, the earlier semantic development assigns it a tacit meaning $\sem{P}$ and an associated correctness outcome under evaluation by a single agent. In practice, however, we do not trust any single evaluation of $P$; we care about what happens when $P$ is executed across a population, at different times, under different tacit conditions.

To capture this, we consider \emph{admissible runs} of $P$ by each agent in the population. An admissible run bundles:
\begin{itemize}
  \item an explicit execution trace of $P$ in the navigation--explicit space $N \ltimes E$, subject to the structural constraints maintained throughout; and
  \item the induced evolution of the agent's extended tacit state along that trace, as specified by the epistemic dynamics update rules (degradation, confidence inflation, requirement escalation, and feedback-coupled degradation).
\end{itemize}

A \emph{replication family} for $P$ is a choice of one admissible run for each agent (or for some subcollection of agents) in the population. For each such run, we select an evaluation time---for instance, when the pipeline halts or when a fixed resource limit is reached---and record the corresponding tacit state $t_{i,*}$ for agent $A_i$.

What matters for institutional behavior is not the detailed trajectory of each run, but the \emph{pattern of correctness judgments} these runs induce. We therefore define the \emph{population semantic profile} of $P$ relative to a replication family $R(P)$ as the multiset
\[
  \mathcal{P}(P; R(P)) \;=\; \{ \gamma_i(t_{i,*}) : i \in I_{\mathrm{rep}} \} \subseteq J.
\]
This multiset is the primary interface between long-horizon practice and institutions. It aggregates how many agents, under what tacit conditions, endorse, reject, or revise the outputs of a given pipeline. The construction makes clear that institutional input is inherently \emph{tacit-dependent}, even though it is summarised in a discrete space of judgments.

The formal development of admissible runs, replication families, and population semantic profiles is given in Appendix~A.12, building on the run semantics and stochastic models of Appendices~B and~C.

\subsection{Institutions as Aggregators over Tacit Judgments}

Institutions---journals, standards bodies, model review boards, internal red-teams, and so on---do not typically see the full detailed state of every agent. Instead, they receive some distribution of judgments: accept, reject, request further work, abstain. In our framework, an \emph{institutional aggregator} is a map
\[
  \mathsf{Mult}(J) \to J \cup \{\bot\},
\]
where $\mathsf{Mult}(J)$ is the space of multisets of judgments and $\bot$ denotes suspended judgment. For a given pipeline $P$, an institution chooses a particular aggregator $I_P$ and applies it to the population semantic profiles generated when $P$ is replicated.

This abstraction covers common patterns:
\begin{itemize}
  \item majority or supermajority rules;
  \item unanimity or veto-based protocols;
  \item weighted consensus, where some agents' judgments count more than others;
  \item deliberative or staged procedures that can output ``no decision'' ($\bot$).
\end{itemize}

Crucially, $I_P$ operates \emph{only} on the multiset of tacitly produced judgments. It does not see raw model outputs, logs, or metrics directly. All explicit activity is mediated through human agents, whose tacit evaluations are then aggregated.

Full definitions of institutional aggregators, institutional outcomes, and their basic properties appear in Appendix~A.12. Here, the important point is that institutions sit at a distinct layer in the architecture: they neither compute the explicit artefacts themselves, nor embody the tacit processes that evaluate them. They are rules for turning many tacit judgments into collective decisions.

\subsection{Why Institutions Cannot Be Fully Automated}

Given mature AI systems and rich logging, it is tempting to imagine institutions that operate purely on explicit artefacts and their explicit summaries: dashboards, performance metrics, proofs of robustness, and so on. In our framework, such an institution would be \emph{explicit-only}: for each pipeline $P$, its decisions could be computed by a mechanism that inspects only explicit artefacts, never accessing the underlying tacit states or correctness judgments.

The impossibility results developed earlier rule this out in a strong sense. Under the assumptions that:
\begin{itemize}
  \item correctness is fundamentally tacit-dependent (it lives in the $\gamma_i$ quotients of tacit states), and
  \item admissible pipelines can be arranged so that explicit traces are ambiguous with respect to tacit correctness,
\end{itemize}
Appendix~A.12 proves an institutional non-automation theorem: no explicit-only institutional layer can be \emph{correctness-preserving} across all admissible pipelines and populations. In particular, if we require that the institution agree with unanimous tacit verdicts whenever they occur, then there is no way to define its decisions purely in terms of explicit artefacts.

Intuitively, any such institution would amount to an ``explicit-only correctness oracle,'' and the earlier no-go theorems about correctness oracles apply. No amount of logs and metrics, however elaborate, can substitute for direct access to tacit evaluations. Institutions that aim at epistemic reliability must therefore be designed around \emph{tacit input}, even if they also make use of explicit signals in how they schedule, resource, and structure evaluation.

The formal statement and proof of the institutional non-automation theorem are collected in Appendix~A.12, where they are derived as a corollary of the explicit-only impossibility results from the more general theory.

\section{Governance, Law, and Oversight}
\label{sec:governance-main}

Institutions regulate the behaviour of individual pipelines and agents. But institutions themselves change over time: journals tighten or relax standards, internal review processes are redesigned, and regulatory frameworks evolve. In long-horizon settings, these higher-order changes are at least as important as first-order institutional choices. This section outlines how governance and law fit into the architecture, again referring to Appendix~F for the formal layer.

\subsection{Governance as a Meta-Institution}

The collection of all institutional arrangements we might adopt forms a structured space. For each pipeline $P$ we can choose an aggregator $I_P$, and we may change these choices over time in response to experience. We write $\Inst$ for the category whose objects are such institutional families and whose morphisms are structure-preserving transformations between them.

A \emph{governance operator} is then an endofunctor
\[
  G : \Inst \to \Inst
\]
that describes how institutional arrangements evolve. At each time step $t$, a governance process takes an existing institutional layer $(I^{(t)}_P)_P$ and produces an updated one $(I^{(t+1)}_P)_P$. This update may depend on observables extracted from the collective dynamics of agents and pipelines: disagreement rates, error frequencies, time-to-detection for hallucinations, evidence of confidence inflation, and so on.

The collective dynamics were modeled earlier as a process on epistemic populations interacting with pipelines. Appendix~F shows how to couple this with governance: observables are extracted from the evolving population, fed into a governance operator, and used to adjust institutional aggregators over time. The resulting picture is a closed loop:
\begin{enumerate}
  \item institutions aggregate tacit judgments and shape which pipelines are accepted or deployed;
  \item those pipelines, in turn, influence the tacit states of agents and the composition of populations;
  \item governance watches this evolving behaviour and modifies institutions accordingly.
\end{enumerate}
This three-level interaction---agents and pipelines, institutions, and governance---is where the long-horizon behaviour of human--AI systems is ultimately determined.

\subsection{Explicit-Only Governance and Its Limits}

Just as we can imagine explicit-only institutions, we can imagine explicit-only governance: processes that monitor logs, metrics, and other explicit artefacts of institutional performance, and then automatically adjust institutional rules based solely on those explicit signals.

In our framework, an explicit-only governance operator is one that, at each time step, updates the institutional layer using only explicit artefacts produced so far. Formally, Appendix~F captures this by requiring the updates to factor through explicit-only mechanisms on the explicit space $E$, in the same way the explicit-only institutions of Appendix~A.12 factor through explicit artefacts.

The same tension appears here as before. If correctness is fundamentally tacit-dependent, and if pipelines and populations can be arranged so that explicit traces underdetermine tacit correctness, then governance that never consults tacit judgments cannot reliably track institutional performance. Any attempt to make governance fully explicit-only will either fail to detect drift away from correct tacit practice, or else reintroduce tacit judgments indirectly (for example, via ``ground truth'' labels or other human-provided signals).

Appendix~F formalizes this as a non-automation theorem for governance: under the same assumptions that rule out explicit-only correctness oracles and explicit-only institutional layers, there is no correctness-preserving governance operator that depends exclusively on explicit artefacts or their explicit summaries. In essence, a fully automated governance layer would again amount to a higher-order explicit-only oracle, and the impossibility results apply at this meta-level as well.

\subsection{Law as a Constraint on Governance}

Governance does not operate in a vacuum. In practice, institutions and their evolution are constrained by legal and normative frameworks: data protection laws, safety regulations, due-process requirements, and so on. In the present setting, we model these as predicates on institutional arrangements and governance operators.

Concretely, we can view a legal or regulatory framework as specifying a set of \emph{admissible} institutions and transformations between them. A governance process is then \emph{lawful} if it never maps admissible institutions to inadmissible ones. Law can also constrain which observables governance is allowed to use (for example, privacy-preserving restrictions on what can be logged) and which kinds of interventions are permitted.

While the formal treatment of law as a constraint on $\Inst$ and on governance functors appears in Appendix~F, the conceptual takeaway is straightforward: the space of possible governance processes is smaller than the space of merely mathematically definable ones, and non-automation results for governance must be understood relative to whatever legal admissibility conditions we impose.

\subsection{Human Oversight as a Structural Requirement}

Taken together, the institutional and governance non-automation results point to a strong conclusion. It is not simply that, for now, we happen to rely on humans in the loop; rather, under the structural assumptions of the framework, there is no coherent way \emph{within FRFP’s explicit--tacit separation and dynamic model} to remove them while preserving correctness over long horizons.

At the institutional level, correctness-preserving aggregation requires access to tacit judgments. At the governance level, correctness-preserving adaptation of institutions requires some access to tacit evaluations of institutional performance. Attempts to fully automate either layer collapse, in the limit, into explicit-only oracles, which the earlier impossibility theorems prohibit.

This does not mean that institutions or governance must be informal or opaque. On the contrary, much of their structure can and should be formalized, monitored, and partially automated. The key point is that the formal and automated parts must be designed \emph{around} human tacit judgment, not as replacements for it.

The detailed statements and proofs of these claims are provided in Appendices~E and~F. Here we emphasize their architectural consequence: in any long-horizon human--AI ecosystem that satisfies the assumptions of the framework, human agents retain irreplaceable roles at three levels simultaneously:
\begin{enumerate}
  \item as individual evaluators of explicit artefacts;
  \item as contributors to institutional decisions via their tacit judgments; and
  \item as overseers of governance processes that adapt and constrain institutions over time.
\end{enumerate}
The rest of the paper can be read as an exploration of how to design pipelines, institutions, and governance mechanisms that make these roles precise rather than aspirational.

%=========================================================
\section{Discussion and Outlook}
\label{sec:discussion}
%=========================================================

\subsection{Design Implications}

The FRFP architecture formalizes a set of design constraints for AI
systems intended for long-horizon scientific collaboration:

\begin{itemize}[leftmargin=2em]
  \item \textbf{Explicit-only AI.} AI systems should manipulate explicit
        structures but must not assume tacit roles such as correctness
        judgment. This limits their epistemic authority while maximizing
        their computational utility.
  \item \textbf{Single explicit--tacit boundary.} All correctness-bearing
        information flow should pass through a clearly identifiable
        boundary operator \(\RB\), with no hidden channels from tacit
        to explicit space.
  \item \textbf{Auditability via canonical forms.} Reasoning trajectories
        can be normalized within \(\Ncat \ltimes \Ecat\), yielding
        canonical forms amenable to audit and replication.
  \item \textbf{Structural handling of hallucinations.} Hallucinations
        are treated as structural mismatches between context and
        requirements, not as psychological quirks of models. This allows
        analysis of how context degradation, feedback, and task design
        interact.
\end{itemize}

\paragraph{Kernel-level constraints vs.\ instantiated layers.}
The design constraints above attach first and
foremost to the kernel of FRFP: the explicit/tacit separation with a single boundary, the minimal
primitive set, and the canonical \(N \ltimes E\) semantics with tacit correctness. By construction, any
system that respects this kernel inherits the same structural limitations and invariants, regardless of
how it instantiates dynamics, populations, or institutions. Conversely, Sections~6–10 and
Appendices~B–J are deliberately more contingent: they work out one particular family of choices for
context degradation and requirements (Appendices~B–D), collective and institutional dynamics
(Appendices~C–F), and concrete protocols, case studies, instruction sets, and prompting patterns
(Appendices~G–J) compatible with the kernel. These outer layers are intended as illustrative and
revisable; what is non-negotiable, from the FRFP point of view, is the underlying correctness
architecture and its separation between explicit computation and tacit judgment. Readers primarily
interested in organizational and governance implications, rather than the categorical details, may wish
to consult the non-technical organizational summary in Appendix~K.

\subsection{FRFP as a formalization of the scientific method}
\label{subsec:frfp-scientific-method}

Although FRFP is stated as a domain-general protocol for long-horizon human--AI collaboration, its core architectural commitments mirror the stabilisation logic of the scientific method. Science is not merely a sequence of cognitive steps; it is an institutional technology for converting exploratory reasoning into stable, publicly defensible claims. FRFP abstracts that conversion into an explicit/tacit split plus a responsibility-bearing acceptance boundary.

\paragraph{A translation dictionary.}
The following correspondences are the intended reading of the formal objects.
\begin{itemize}[leftmargin=2em]
  \item \textbf{Explicit artefacts (\(E\)).} Hypotheses, models, derivations, experimental designs, protocols, code, datasets, plots, and written arguments: objects that can be logged, transmitted, inspected, and mechanically transformed.
  \item \textbf{Explicit operations (\(\RI,\EC,\ED\)).} Representation initiation (writing down a conjecture or design), computation (derivation, simulation, fitting, summarisation), and diagnostics (sanity checks, ablations, error analysis, unit tests, proof checking where applicable).
  \item \textbf{Navigation (\(\Ncat\)).} Legitimate changes of perspective that preserve intent: reparameterisations, alternative model classes, reorganising an analysis, swapping datasets, or reframing a subproblem while remaining within the same research question.
  \item \textbf{Tacit state (\(T\)).} The correctness-bearing background a scientific community relies on but cannot fully encode: domain know-how, interpretive competence, judgment of adequacy, and the willingness to stand behind a claim under scrutiny.
  \item \textbf{Boundary (\(\RB:E\to T\)).} Review-and-belief update: the act of assessing an explicit artefact against tacit standards and incorporating (or rejecting) it as part of what one is prepared to endorse.
  \item \textbf{Tacit primitives (\(\TE,\HFD\)).} Evaluation and acceptance: internal or institutional endorsement decisions (``this analysis is sound enough to rely on''; ``this result is ready to submit''; ``this meets the evidential bar for policy use'').
\end{itemize}

\paragraph{Replication, reproducibility, and peer review as FRFP patterns.}
The scientific stabilisation norms that motivate FRFP can be expressed as structured constraints on explicit workflows and on where boundary crossings are permitted.
\begin{itemize}[leftmargin=2em]
  \item \textbf{Reproducibility} is an explicit-side property: an analysis is reproducible when the relevant explicit artefacts (data, code, environment, protocol) are sufficient for another agent to re-run an explicit pipeline and obtain the same explicit outputs, up to documented tolerances. In FRFP terms this is invariance in \(E\) under admissible refactoring and navigation, together with suitable logging of the pipeline normal form.
  \item \textbf{Replication} is a stronger, world-facing constraint: a claim survives replication when an independent execution produces new explicit artefacts (new data) that, upon boundary crossing, supports the same tacit endorsement. Formally, replication is not just agreement of \(E\)-outputs; it is agreement of tacit acceptance \(\gamma(t)\) under independent pipelines that re-contact the domain.
  \item \textbf{Peer review} is a multi-agent boundary mechanism: independent tacit evaluators consume shared explicit artefacts and apply their own tacit operators before endorsing publication-level acceptance. In FRFP terms, peer review introduces additional boundary crossings by other agents, creating a layered acceptance structure that can improve finite-horizon risk when evaluations are sufficiently independent.
\end{itemize}

\paragraph{What FRFP adds, relative to informal scientific method talk.}
FRFP is not a sociological model of science; it is a protocol semantics intended to make long-horizon correctness and responsibility analyzable. Relative to informal descriptions of the scientific method, FRFP contributes three formal clarifications.
\begin{itemize}[leftmargin=2em]
  \item \textbf{Acceptance is an event, not a vibe.} Scientific practice often blurs exploration and endorsement; FRFP forces the distinction by making acceptance a boundary crossing \(E\to T\) with explicit responsibility.
  \item \textbf{Hallucination is endorsement failure.} In FRFP, ``hallucination'' is not the production of an incorrect string; it is the protocol-level event in which an explicit artefact is incorporated into tacit acceptance without meeting the active tacit requirements for that credence.
  \item \textbf{Stabilisation is workflow-level.} Reliability is treated as an emergent property of explicit pipelines plus review structure, not as a property of any single component (human or AI). This matches the way scientific norms reduce error by layering checks and independent evaluation rather than assuming perfect agents.
\end{itemize}

\paragraph{Scope and limits of the mapping.}
The mapping is exact in regimes where ``science'' is best understood as a stabilisation protocol over explicit public artefacts (papers, code, data) and tacit endorsement (what researchers and institutions will stand behind). It is intentionally agnostic about internal cognitive models and about which explicit checks are appropriate in a given domain. Conversely, in fully formal regimes where correctness is exhausted by an explicit oracle (e.g., proof checking under a fixed calculus), FRFP should be read as placing that oracle inside explicit space as an explicit correctness engine, with tacit judgment still governing the adequacy and scope of the formalisation. The purpose of the mapping is therefore not to claim that FRFP \emph{is} the scientific method, but to show that FRFP isolates, and makes compositional, the same stabilisation logic that has historically made scientific claims durable.


\subsection{Limitations and scope}
\label{subsec:limitations}

The present work has several limitations.

\paragraph{Single-agent tacit model.}
We model one tacit context space $T$ at a time; multi-agent and institutional settings, where different humans (and possibly institutions) have distinct tacit states, require richer structures (e.g.\ families of tacit spaces, communication channels, and aggregation mechanisms). Developing such multi-agent extensions is left to future work.

\paragraph{Worst-case focus.}
The hallucination inevitability results rely on deliberately worst-case assumptions: long or unbounded runs, monotone degradation, no explicit repair operator, and simple forms of requirement escalation and feedback. Practical systems may use aggressive re-grounding, multiple agents, explicit restarts, or external tools (e.g.\ proof assistants, simulators) to reduce risk. Our results should therefore be read as structural ``upper bounds'' on what can go wrong in principle, not as quantitative predictions for particular deployments.

\paragraph{Abstract requirements and tacit semantics.}
The requirement map $\mathrm{Req}(e,c)$ is treated as an abstract semantic object; we do not attempt to estimate it from data or tie it to any specific scientific or engineering domain. Likewise, tacit states are modeled at a coarse level (often as elements of a preorder or low-dimensional space), omitting many cognitive and socio-technical factors that matter in practice. Making $\mathrm{Req}$ and the tacit dynamics empirically grounded for particular domains is outside the scope of this paper.

\paragraph{Level of abstraction.}
More broadly, we work at an intentionally abstract, foundational level and do not propose concrete algorithms, training procedures, or evaluations. FRFP is a design framework and semantic target, not a ready-made system. Instantiating it in concrete tools will require additional engineering choices (e.g.\ representations, interfaces, inference methods) that we do not model here.

\paragraph{Scope of inevitability and non-automation claims.}
All inevitability and non-automation statements in this paper are conditional on the FRFP
architecture and its standing assumptions. In particular, the hallucination inevitability theorems
presuppose the Epistemic Dynamics degradation and requirement conditions, together with the admissibility
constraints on runs, and the impossibility results for detectors, consensus mechanisms, institutions,
and governance operators presuppose explicit--tacit separation and human-exclusive grounding.
They should therefore not be read as claiming that hallucination is inevitable, or that fully automated
governance oracles are impossible, for \emph{all} conceivable AI architectures. Rather, they assert
that such guarantees cannot be obtained \emph{within} the FRFP class of models, once tasks are
nontrivial and requirements recur; trivial strategies such as always answering ``I don't know'' or
never accepting artefacts are ruled out by these structural premises rather than by logical impossibility.
These non-automation results apply to any explicit-only AI system, including hypothetical AGI, as long as correctness remains a human/institutional notion rather than delegated to the system itself.
In particular, in domains where correctness is truly explicit at the relevant level of analysis—
for example, inner loops governed by sound and complete proof or test oracles—our Epistemic Dynamics
inevitability statements are vacuous for that inner task.
They become informative only when one models the outer tacit layer where specifications, benchmarks,
and institutional standards are chosen and revised (Appendix~D.3).

\paragraph{Counter-scenarios for Epistemic Dynamics assumptions.}
The Epistemic Dynamics inevitability results target a regime with recurring nontrivial demands, accumulating
degradation, and bounded repair. There are three salient counter-scenarios. In \emph{balanced or
negative drift} regimes, institutions invest heavily enough in repair---for example, intensive and
frequent simulation training, structured post-mortems, and recertification in ICU- or nuclear-safety
settings---that average repair offsets or exceeds task-induced degradation. Our results then describe
what happens if such investment is reduced, rather than claiming that drift wins regardless of effort.
In \emph{unbounded repair} regimes, one postulates repair moves with no global bound on their
effect, such as costless, arbitrarily frequent ``oracle refresh'' operations or instant replacement by
perfectly grounded experts; this violates the bounded-repair assumption and lies outside FRFP's
human-bounded institutional model. In \emph{low-demand} regimes, such as rare emergency
committees that face serious decisions only occasionally, nontrivial demands are so infrequent that
modest repair may suffice to maintain grounding, and asymptotic Epistemic Dynamics arguments are not the
dominant concern.

These counter-scenarios also clarify an economic tension. In high-load settings, pushing repair hard
enough to stabilize tacit grounding typically implies a governance workload whose cost and latency
scale with AI usage. Institutions can thus choose a point on a tradeoff surface: minimal repair,
preserving large speed and cost advantages at the price of long-run inevitability; substantial repair,
which buys longer safe horizons but erodes pure automation gains; or intermediate regimes that
partially delay and reshape, but do not eliminate, long-run risk.


\paragraph{Scope of novelty.}
As noted in the Introduction, the categorical ingredients we use (free categories, groupoids, fibrations, Grothendieck constructions) are standard; the contribution is to assemble them into a single architecture that captures explicit/tacit separation and feedback, and to derive system-level constraints on auditability, correctness, and hallucination. FRFP is intended as a reusable structural framework rather than a specific algorithm or training procedure.
\noindent
Our strongest governance-level claims should thus be read as conditional: they are formal
no-go theorems \emph{given} the FRFP axioms, plus an empirical hypothesis that high-stakes
scientific and regulatory practice falls within that class of architectures, from which we derive
normative guidance for protocol and institutional design.


\subsection{Future Directions}

Several directions are natural next steps.

\paragraph{Empirical dynamics of human--AI protocols.}
The current framework treats tacit context, degradation, requirements, confidence, and feedback as abstract operators. A first direction is to make these empirically meaningful. This includes: (i) developing task- and domain-specific proxies for tacit context, degradation, requirement maps, and plausibility functionals; (ii) instrumenting real human--AI workflows so that trajectories in \(N \ltimes E\) can be logged, normalized, and analysed; and (iii) testing structural predictions (such as inevitability and acceleration of hallucination) against behavioural and longitudinal data.

\paragraph{Repair, re-grounding, and reset mechanisms.}
Our analysis intentionally focuses on settings where context degrades and is not fully restored. A complementary direction is to design explicit repair and re-grounding mechanisms---study breaks, external reference checks, slow deliberation modes---and represent them as operators within the same algebraic framework. At the protocol level, this suggests specifying and analysing reset policies, usage caps, and re-verification gates as first-class constructions in \(N \ltimes E\), and characterising when such mechanisms significantly delay (though do not eliminate) structural hallucination.

\paragraph{Collective architectures and institutions.}
The framework can be extended from single users to populations of agents and to institutions that aggregate many runs. This opens questions about how different institutional patterns (peer review, adversarial review, replication, regulatory stress-testing) can be represented, and which combinations provably buffer individual degradation while respecting a human-exclusive grounding constraint. Another direction is to relate formal notions such as population semantic profiles, institutional stability, and institutional non-automation to empirical studies of scientific and regulatory processes, and to compare predicted failure modes (fake consensus, lock-in, correlated error) with observed ones.

\paragraph{Governance, law, and meta-level design.}
The same categorical machinery can be used to model governance operators that act on institutions, together with constraints that restrict how institutions may change. This suggests: (i) designing concrete governance observables (drift, disagreement, confidence-inflation indicators) that can be computed from FRFP-compatible logs; (ii) exploring families of governance operators that incorporate human oversight at the institutional level without collapsing into fully automated correctness claims; and (iii) mapping these patterns onto emerging regulatory regimes for AI in science, medicine, and finance, to clarify which regulatory templates are structurally compatible with an explicit/tacit separation.

\paragraph{Richer explicit logics and tooling.}
On the explicit side, the framework is agnostic about the internal logic carried by \(E\). A further direction is to instantiate \(E\) with richer structures: type-theoretic or temporal logics, proof objects from interactive theorem provers, or domain-specific calculi for simulation and data analysis. This would support more expressive explicit pipelines while preserving the \(RB\)--\(TE\)--\(HFD\) suffix and its correctness properties. In parallel, one can build FRFP-aware tooling: typed interface languages, normalization and auditing layers for \(N \ltimes E\), and assistants that expose their operations as TDG terms rather than as opaque conversations.

\paragraph{Benchmarking and comparative evaluation.}
Finally, FRFP suggests protocol-level evaluation criteria. Instead of scoring isolated model outputs, one can compare whole workflows---different ways of structuring the explicit/tacit boundary, navigation, repair, and institutional aggregation---using metrics derived from the dynamic theory, such as time-to-hallucination distributions, robustness under degradation, and sensitivity to feedback coupling. Developing such benchmarks would allow FRFP-compatible systems to be compared along principled axes, and would help connect the abstract theory to concrete engineering trade-offs.

FRFP is intended as a foundation, not a finished product. Its value lies in making constraints and invariants explicit so that future systems---both technical and institutional---can be designed, analysed, and compared within a common formal framework.

\section{Alternative formalisms}
\label{sec:alternative-formalisms}

FRFP is one point in a broader design space of human--AI collaboration architectures. This section
positions FRFP against competing formalisms, clarifies which alternatives do \emph{not} reduce to
FRFP (because they step outside the FRFP axiom class), and explains why FRFP is the appropriate
choice for the theorem-level aims of this volume.

\subsection{The design space: what varies across architectures}

Any protocol-level formalism must decide (at minimum):
\begin{enumerate}[leftmargin=2em]
  \item \textbf{Where correctness lives.} Is correctness a property decidable in explicit space (a predicate or oracle on artefacts), or a property of tacit endorsement states?
  \item \textbf{Whether tacit is writable.} Is there a channel \(T\to E\) that exports acceptance states into machine-manipulable specifications, policies, or constraints?
  \item \textbf{Who may access tacit.} May AI systems condition on tacit state (or modify it), or is tacit strictly human-exclusive?
  \item \textbf{How acceptance is structured.} Is there a single acceptance boundary, or multiple gates (individual, committee, institution) with distinct authority and standards?
  \item \textbf{How uncertainty is represented.} Is endorsement discrete (a finite judgment quotient), ordered (preorders), or probabilistic (graded belief and risk thresholds)?
\end{enumerate}
FRFP fixes these choices in a specific way: correctness is localized in tacit space; AI is confined to
explicit moves; the explicit--tacit boundary is one-way \(E\to T\) with no return channel \(T\to E\);
and the tacit interface is intentionally minimal so that the explicit algebra carries the technical
structure.

\subsection{Architectures that do \emph{not} reduce to FRFP}

Because FRFP is proved (in Appendix~A) to be \emph{initial} only within a category of frameworks that
share the FRFP epistemic primitives and axioms, an alternative architecture avoids ``reducing to
FRFP'' precisely by violating at least one of those axioms. The following are genuine competitors in
that sense.

\paragraph{(A) Bidirectional boundary architectures (\(T\leftrightarrow E\)).}
These formalisms introduce an essential write-back channel \(F:T\to E\), by which tacit endorsement
states externalize into explicit policies, specifications, rubrics, or machine-checkable constraints.
This matches institutional learning patterns (postmortems \(\to\) checklists \(\to\) standards), but it
is outside FRFP by construction: FRFP forbids \(T\to E\). In such systems, the acceptance criteria are
not merely applied at review points; they are continuously rewritten into the explicit layer.

\paragraph{(B) Tacit-in-the-loop AI.}
Here AI systems may condition on tacit state or participate in tacit updates (e.g., morphisms of the
form \(T\times E\to E\) or AI-mediated operators \(T\to T\)). This architecture is attractive for
personalization and closed-loop assistance, but it breaks human-exclusive correctness and blurs
responsibility: the system can implicitly steer endorsement by reading or shaping tacit states. This
is therefore not an instance of FRFP.

\paragraph{(C) Explicit-only correctness architectures.}
In some regimes, correctness can be reduced to a complete explicit oracle: proof checking, model
checking, exhaustive tests under a closed-world specification. In such settings, ``acceptance'' is an
explicit computation on artefacts, not a tacit judgment. This is again outside FRFP's target class:
FRFP assumes a genuine tacit layer that is not exhausted by any fixed explicit oracle.

\paragraph{(D) Multi-boundary acceptance architectures.}
Real institutions often implement layered gates of endorsement (individual sign-off, peer review,
risk committee approval, regulatory clearance). A formalism may represent these as multiple tacit
layers and multiple boundaries, with cross-talk between gates. If the ontology class fixes a single
boundary primitive, then such multi-gate architectures do not reduce to FRFP without extending the
primitive structure.

\paragraph{(E) Probabilistic tacit semantics.}
A competing formalism may model tacit state as a richer belief object (e.g., probability measures,
utility-weighted acceptance thresholds, risk budgets) rather than a preorder with a finite observable
judgment quotient. When the tacit algebra differs materially (graded endorsement rather than discrete
quotients and idempotent operators), the architecture falls outside the specific tacit-interface axioms
used in FRFP's initiality claims.

\subsection{Which alternatives can be viewed as generalizations of the scientific method}

Several competing architectures can still be regarded as ``generalizations of the scientific method,''
but they generalize different facets of it.

\begin{itemize}[leftmargin=2em]
  \item \textbf{Explicit-only correctness} generalizes the formal-verification slice of science (mathematics, parts of theoretical computer science), where correctness is internal to a fixed calculus.
  \item \textbf{Bidirectional \(T\leftrightarrow E\)} generalizes the meta-scientific process by which tacit know-how becomes codified into explicit methods, standards, and tooling.
  \item \textbf{Multi-boundary acceptance} generalizes the institutional reality of science as layered endorsement (lab judgment, journal review, replication, community uptake).
  \item \textbf{Probabilistic tacit semantics} generalizes science as graded belief and risk management under uncertainty.
\end{itemize}
FRFP generalizes a particular core: \emph{stabilization by accountable review}, where public explicit
artefacts are strengthened by explicit checks and then cross a responsibility-bearing acceptance
boundary.

\subsection{Why FRFP among the alternatives}

FRFP is chosen to match a specific problem class and to enable theorem-level analysis.

\paragraph{Target regime.}
We aim at long-horizon workflows in which (i) requirements evolve, (ii) explicit artefacts and
evaluation criteria drift across extended traces, and (iii) no fixed explicit oracle exhausts what it
means for an artefact to be adequate. In this regime, architectures that place correctness entirely in
\(E\) are aiming at a different class of problems.

\paragraph{Accountability and non-circular semantics.}
FRFP enforces that endorsement is attributable: AI systems can produce and transform explicit
artefacts, but only humans occupy correctness-bearing tacit states. The one-way boundary \(E\to T\)
prevents implicit rewriting of acceptance criteria by the system itself. Architectures with essential
\(T\to E\) write-back or tacit-in-the-loop AI can be valuable, but they entangle acceptance and
generation in ways that complicate responsibility and undermine the invariance and non-automation
results that this volume targets.

\paragraph{Minimality and canonicity within the axiom class.}
Within the class of formalisms that share FRFP's epistemic primitives and constraints (explicit/tacit
separation, one-way boundary, human-exclusive correctness, AI confined to explicit moves), FRFP is
initial: any competing framework satisfying the same assumptions admits a unique structure-preserving
interpretation from FRFP. In this precise sense, FRFP is not an arbitrary design choice inside its
assumption class; alternatives that keep the same primitives are instances/refinements rather than
routes around the theorem-level consequences.

\subsection{Practical reading of this comparison}

The existence of non-FRFP architectures is not a defect: it clarifies scope. FRFP is a protocol-level
formalism for domains where correctness-bearing endorsement cannot be reduced to explicit computation,
and where the central risks are long-horizon drift, accountability loss, and explicit-only governance
limits. Competing architectures are appropriate when one wishes to (i) treat correctness as fully
explicit, (ii) explicitly codify tacit judgments into machine-checkable policy, (iii) close the loop
between system behavior and human internal state, or (iv) model layered institutional gates or graded
belief as first-class. Subsequent volumes may study carefully delimited departures from FRFP's axioms;
the present volume fixes them to make the foundational theorems sharp.

\pagebreak
\appendix
\section{Mathematical Foundations}
\label{app:math-foundations}
\subsection{Static Epistemic Architecture and Algebra} % old Appendix A
\label{app:static-epistemic}
% C_full, E/T, RB, primitives, TDG, N ⋉ E, rewriting, semantic functor, correctness

This appendix develops the categorical and algebraic core of the
FRFP (Foundational Reasoning and Feedback Protocol) architecture.
It provides a self-contained account of:

\begin{itemize}
  \item an ambient category with explicit and tacit modes;
  \item a free explicit algebra of task decompositions (the Task Decomposition Grammar, TDG);
  \item a navigation groupoid acting on explicit shapes;
  \item a fibration of explicit pipelines over shapes;
  \item a Grothendieck-style semidirect product \(N \ltimes E\);
  \item a semantic functor into a preorder-enriched tacit category; and
  \item correctness and canonicality theorems for Epistemic Structure.
\end{itemize}

We assume familiarity with basic category theory, rewriting systems,
and context-free grammars.

Throughout, we write \(\Cfull\) for the ambient category,
\(\Espace\) for explicit space, and \(\Tspace\) for tacit space.

%======================================================================
\subsection{Foundations}
%======================================================================

\subsubsection{Ambient Category \texorpdfstring{\(\Cfull\)}{Cfull}}

\begin{definition}[Ambient category]
The ambient category \(C_{\mathrm{full}}\) is an arbitrary category that may contain objects beyond those used by the FRFP protocol. We designate three \emph{distinguished} objects
\(\varnothing, E, T \in \mathrm{Ob}(C_{\mathrm{full}})\):
\begin{itemize}
  \item \(\varnothing\): the empty explicit state;
  \item \(E\): the explicit state object (machine-manipulable representations);
  \item \(T\): the tacit state object (human correctness-bearing states).
\end{itemize}
Morphisms in \(\Cfull\) are partitioned into three epistemic classes:
\begin{itemize}
  \item explicit morphisms with codomain \(E\), forming the explicit subcategory
        \(\Espace \subseteq \Cfull\);
  \item tacit morphisms with codomain \(T\), forming the tacit subcategory
        \(\Tspace \subseteq \Cfull\);
  \item boundary morphisms \(E \to T\), of which a distinguished one is the
        \emph{Representational Backflow} operator
        \(\RB : E \to T\).
\end{itemize}
Composition in \(\Cfull\) satisfies the usual category axioms:
\[
(h \circ g) \circ f \;=\; h \circ (g \circ f), \qquad
f \circ \id_X \;=\; f \;=\; \id_Y \circ f.
\]
\end{definition}

\begin{definition}[Explicit--Tacit Separation (ETS)]
The \emph{Explicit--Tacit Separation} axiom requires:
\begin{itemize}
  \item there are no morphisms \(T \to E\) in \(\Cfull\);
  \item tacit evaluation cannot be embedded inside explicit pipelines;
  \item the boundary morphism \(\RB : E \to T\) is not invertible.
\end{itemize}
\end{definition}

\subsubsection{Modal Interpretation of Explicit and Tacit Knowledge}

\begin{definition}[Explicit modality]
A cognitive item \(k\) is represented in \emph{explicit} modality if there exist
maps
\[
\enc : E \to \Sigma^{\ast}, \qquad
\dec : \Sigma^{\ast} \to E
\]
such that \(\dec(\enc(k))\) is equivalent to \(k\) in its epistemic role
(e.g., as a proof, equation, or program). In particular, \(k\) can be
faithfully encoded as a finite symbolic artifact and recovered without
loss of its correctness-relevant structure.
\end{definition}

\begin{definition}[Tacit modality]
A cognitive item is represented in \emph{tacit} modality when no such
encoding preserves its correctness-bearing content. Judgments, intuitions,
acceptance decisions, and other human correctness-bearing acts live in
tacit space.
\end{definition}

Human expression from tacit cognition to explicit content (e.g., writing a
new question, revising a claim) is modeled externally as the appearance of
new explicit morphisms in \(\Espace\); it is not modeled as a morphism
\(T \to E\).

\subsubsection{Explicit Category \texorpdfstring{\(\Espace\)}{E}}

\begin{definition}[Explicit category]
The explicit category \(\Espace\) is the full subcategory of \(\Cfull\)
on objects with codomain \(E\). Its core generating morphisms are:
\begin{gather*}
\RI : \varnothing \to E \quad\text{(Representation Initiation)},\\
\EC : E \to E \quad\text{(Explicit Computation)},\\
\ED : E \to E \quad\text{(Explicit Diagnostics)}.
\end{gather*}
\end{definition}

\noindent
\textbf{Remark.}
\(\EC\) represents algorithmic transformations of explicit content
(e.g., inference, rewriting, search). \(\ED\) represents explicit
checks and diagnostics (e.g., consistency checks, sanity tests),
preparing but not performing tacit evaluation.

\subsubsection{Tacit Category \texorpdfstring{\(\Tspace\)}{T}}

We now formalize tacit space as a preorder-enriched one-object category.

\begin{definition}[Tacit category]
The tacit category \(\Tspace\) has:
\begin{itemize}
  \item a single object, denoted again by \(T\);
  \item hom-set \(\Hom_{\Tspace}(T,T)\) equipped with a composition
        operation \(\circ\), making it a monoid;
  \item a preorder \(\sqsubseteq\) on \(\Hom_{\Tspace}(T,T)\) such that
        composition is monotone:
        \[
        f_1 \sqsubseteq f_2,\ g_1 \sqsubseteq g_2
        \quad\Rightarrow\quad
        g_1 \circ f_1 \sqsubseteq g_2 \circ f_2.
        \]
\end{itemize}
\end{definition}

\begin{definition}[Tacit primitives]
The tacit primitives are the morphisms
\[
\TE : T \to T \quad\text{(Tacit Evaluation)},\qquad
\HFD : T \to T \quad\text{(Human Final Decision)},
\]
subject to:
\begin{itemize}
  \item \emph{idempotence}:
        \(\TE \circ \TE = \TE\) and \(\HFD \circ \HFD = \HFD\);
  \item \emph{commutation}:
        \(\TE \circ \HFD = \HFD \circ \TE\);
  \item \emph{monotonicity}: for every \(f : T \to T\),
        \[
        f \sqsubseteq \TE \circ f,\qquad
        f \sqsubseteq \HFD \circ f.
        \]
\end{itemize}
\end{definition}

\begin{definition}[Correctness quotient]
Let \(J\) be a finite set of correctness judgments
(e.g., \(\{\mathsf{accept}, \mathsf{reject}, \mathsf{revise}\}\)).
A \emph{correctness quotient} is a surjective map
\[
\gamma : \Obj(\Tspace) \to J
\]
assigning to each tacit state a correctness class, subject to
\[
\gamma \circ \TE = \gamma = \gamma \circ \HFD.
\]
Thus tacit evaluation and final decision may refine or stabilize a tacit
state but do not change its correctness class.
\end{definition}

\paragraph{Interpretation and non-triviality of \texorpdfstring{$T$}{T}.}
The tacit category T is not introduced merely as a formal firewall that “reserves correctness to tacit space”.
Rather, it is intended to capture a substantive empirical claim about long-horizon scientific and regulatory work:
in these regimes, there are correctness-bearing aspects of judgment that cannot be exhaustively encoded in a fixed,
finite explicit representation without loss.

Definitions A.3–A.4 distinguish explicit and tacit modalities by the existence of lossless encodings into symbolic artefacts.
In many familiar settings—for example, textbook exercises, small programs, or contest-style benchmarks—it is reasonable
to treat \emph{inner} task correctness as effectively explicit: the task is specified by a finite prompt, there is a fixed
input–output relation or reference implementation, and success can be checked by a test harness or proof checker.
FRFP does not deny the existence or value of such regimes. Instead, it treats them as \emph{tacit-insensitive inner layers}:
conditional on a surrounding tacit backdrop (choice of formalism, modelling assumptions, risk tolerances),
correctness for that inner task behaves as if it were a pure explicit oracle on E.

The global explicit/tacit split remains.
Even in proof-assistant or test-oracle settings, tacit judgment is involved in selecting the calculus, specifying what the
theorem or contract should state, and deciding whether conformance to that specification is adequate for the real-world
purpose.
What FRFP targets are workflows where this outer tacit layer is non-trivial and dynamically involved in long-horizon
use: tasks evolve, standards and constraints shift, and background scientific understanding and norms play an essential
role in deciding whether an explicit artefact is “good enough”.
In those settings it is non-trivial, and in our view appropriate, to treat the final correctness status of an explicit
pipeline as living in T rather than in E. (See Appendix~D.4 for further discussion.)


A reasonable skeptic might ask: under what conditions could one reject the Epistemic Spaces axioms
and the later Epistemic Algebra structure on $T$? There are at least three salient counter-scenarios:
\begin{enumerate}
  \item \textbf{Tacit-insensitive inner correctness.}
For some domains and layers, once a formalism, specification, and world-modelling stance have been tacitly fixed,
every correctness-bearing judgment at the inner level can be captured by a finite explicit specification and a sound
and complete proof- or test-based oracle.
At that inner layer one can model correctness directly in E and treat T as abstracted away; FRFP’s Phase–3
inevitability and non-automation results then become vacuous or uninteresting for that particular subproblem.
Our focus is on regimes where the outer tacit layer—where specifications, standards, and modelling choices are
themselves evolving—remains dynamically involved in correctness.

  \item \textbf{No degradation or escalation.} Epistemic Algebra impossibility results rely on a
  degradation operator $\delta$ and recurring nontrivial requirements relative to a floor $r$.
  If one believes that, in their target setting, tacit state does not degrade in any relevant
  sense, or that requirements do not escalate or recur in the way we assume (for example,
  because tasks are short-lived and fully specified at the outset), then the inevitability
  theorems may not apply. Such systems fall outside FRFP's intended long-horizon regime.
  \item \textbf{Perfect tacit emulation.} If one posits an AI system that can internally
  reconstruct or simulate human tacit evaluation from explicit traces without loss---violating
  the NTER and CSC constraints---then our explicit--tacit separation (ETS, HEG, HEC) is simply
  rejected. FRFP explicitly rules out this possibility by axiom rather than attempting to model
  it.
\end{enumerate}

In summary, treating correctness as localized in $T$ is a substantive modelling choice:
it reflects the view that, for many real scientific and governance workflows, there is residual
human judgment that cannot be fully captured by static benchmarks or formal specifications, and
that this residual plays an essential role in long-horizon correctness. The impossibility and
inevitability results of Epistemic Structure and Dynamics should therefore be read as conditional on this non-trivial
tacit structure: they apply precisely in domains where correctness genuinely outstrips what can
be made explicit and mechanically checked.


\subsubsection{Boundary Morphism and Semantic RB Functor}

\begin{definition}[Boundary morphism \(\RB\)]
The \emph{Representational Backflow} operator is a morphism
\[
\RB : E \to T
\]
in \(\Cfull\). It is the unique explicit-to-tacit interface specified by
the Representational Backflow axiom.
\end{definition}

Let \(E_{\mathrm{sem}}\) denote the free explicit category generated by the
Task Decomposition Grammar (defined in Section~\ref{sec:appendix-tdg}).

\begin{definition}[Semantic RB functor]
The \emph{semantic RB} is a functor
\[
\RBsharp : E_{\mathrm{sem}} \to \Tspace
\]
assigning to each explicit pipeline a tacit state representing the content
received by the human for evaluation.
\end{definition}

\subsubsection{Interaction Protocols and Epistemic Spaces Axioms}

\begin{definition}[Interaction Protocols constraints]
The operational environment satisfies:
\begin{itemize}
  \item \textbf{IL (Interaction Limitation):} interaction is bounded and turn-based;
  \item \textbf{AR (Alignment Restriction):} the system must not autonomously act
        outside explicit instructions;
  \item \textbf{MD (Mode Discipline):} separation between explicit and tacit modes
        is enforced at the interface level;
  \item \textbf{NTER (No Tacit Evaluation Reconstruction):} the system must not
        reconstruct tacit evaluation internally from explicit traces;
  \item \textbf{CSC (Correctness Separation Constraint):} correctness decisions
        must not be delegated to explicit computation.
\end{itemize}
\end{definition}

\begin{definition}[Epistemic Spaces axioms]
An admissible epistemic framework satisfies:
\begin{itemize}
  \item \textbf{ETS (Explicit--Tacit Separation):} no morphisms \(T \to E\);
  \item \textbf{HEG (Human-Exclusive Grounding):} correctness grounding occurs
        only in tacit space via \(\TE\), executed by humans;
  \item \textbf{HEC (Human-Exclusive Closure):} final acceptance/rejection occurs
        only via \(\HFD\), executed by humans;
  \item \textbf{RB (Representational Backflow):} there is a unique boundary
        morphism \(\RB : E \to T\) used by all pipelines;
  \item \textbf{CP (Compositional Pipelines):} epistemic processes are built as
        compositions of primitives into pipelines;
  \item \textbf{AE (AI-Explicit Restriction):} AI systems act only via
        explicit-space morphisms in \(\Espace\).
\end{itemize}
\end{definition}

%======================================================================
\subsection{Explicit Algebra and the Task Decomposition Grammar}
\label{sec:appendix-tdg}
%======================================================================

\subsubsection{TDG Signature and Free Category}

\begin{definition}[TDG signature]
The Task Decomposition Grammar (TDG) signature \(\Sigma_{\mathrm{TDG}}\)
consists of:
\begin{itemize}
  \item objects: \(\varnothing\) and \(E\);
  \item generating morphisms:
  \[
    \RI : \varnothing \to E,\qquad
    \EC : E \to E,\qquad
    \ED : E \to E;
  \]
  \item structural operators \(\BRANCH\), \(\PARALLEL\), \(\LOOP\) acting
        as higher-arity constructors on morphisms.
\end{itemize}
\end{definition}

\begin{definition}[Free explicit category]
The explicit semantic category is the free category on the TDG signature:
\[
E_{\mathrm{sem}} := \FreeCat(\Sigma_{\mathrm{TDG}}).
\]
There is a canonical interpretation functor
\(\Ufun : E_{\mathrm{sem}} \to \Espace\)
sending each TDG term to its denotation as an explicit morphism.
\end{definition}

\begin{theorem}[Free explicit algebra]
\label{thm:free-explicit-algebra}
For any category \(\mathcal{D}\) and any interpretation of the TDG signature
\(\Sigma_{\mathrm{TDG}}\) in \(\mathcal{D}\), there exists a unique functor
\(\tilde{F} : E_{\mathrm{sem}} \to \mathcal{D}\) extending that interpretation.
\end{theorem}

\begin{proof}
This is the standard universal property of the free category on a graph:
objects of \(E_{\mathrm{sem}}\) are objects of \(\Sigma_{\mathrm{TDG}}\), and morphisms
are freely generated paths in the underlying graph. Any graph morphism
from \(\Sigma_{\mathrm{TDG}}\) into \(\mathcal{D}\) extends uniquely to a functor.
\end{proof}

\subsubsection{TDG Grammar}

We present TDG both as a signature and as a context-free grammar
over terms.

\paragraph{Nonterminals.}
We use:
\[
Q \quad\text{for explicit pipelines}, \qquad
P \quad\text{for full pipelines (explicit + tacit suffix)}.
\]

\paragraph{Terminals.}
The terminal symbols are:
\[
\RI,\, \EC,\, \ED,\, \RB,\, \TE,\, \HFD,\,
\BRANCH,\, \PARALLEL,\, \LOOP.
\]

\paragraph{Productions.}
The grammar productions are:
\begin{align*}
Q &::= \RI,\\
Q &::= Q;\EC,\\
Q &::= Q;\ED,\\
Q &::= \BRANCH(Q_1,\dots,Q_k),\\
Q &::= \PARALLEL(Q_1,\dots,Q_k),\\
Q &::= \LOOP(Q_{\mathrm{cond}}, Q_{\mathrm{body}}),\\
P &::= Q;\RB;\TE;\HFD.
\end{align*}
Thus TDG presents \(E_{\mathrm{sem}}\) as an algebra of explicit pipelines
with a fixed tacit suffix.

\subsubsection{Shapes and Shape Category}

\begin{definition}[TDG shape]
TDG shapes are generated by the grammar
\[
S ::= \chain \mid
       \branch(S_1,\dots,S_k) \mid
       \parallel(S_1,\dots,S_k) \mid
       \loopshape(S_{\mathrm{cond}}, S_{\mathrm{body}}).
\]
\end{definition}

\begin{definition}[Shape of a TDG term]
The mapping
\(\shape : \text{TDG-terms} \to \text{Shapes}\) is defined inductively by:
\begin{itemize}
  \item \(\shape(\RI) = \chain;\)
  \item \(\shape(Q;\EC) = \shape(Q)\) and similarly \(\shape(Q;\ED) = \shape(Q)\);
  \item \(\shape(\BRANCH(Q_1,\dots,Q_k))
         = \branch(\shape(Q_1),\dots,\shape(Q_k));\)
  \item analogously for \(\PARALLEL\) and \(\LOOP\).
\end{itemize}
\end{definition}

\begin{definition}[Shape category \(\Scat\)]
The shape category \(\Scat\) has:
\begin{itemize}
  \item objects: shapes \(S\) as above;
  \item morphisms: structural rewiring maps between shapes
        (e.g., regrouping, insertion/removal of trivial chains),
        closed under composition and identities.
\end{itemize}
\end{definition}

\subsubsection{Explicit Reduction System}

\begin{definition}[Explicit reduction]
A binary relation \(\to_E\) on \(\Espace\) is an \emph{explicit reduction system}
if it is a typed rewrite system satisfying:
\begin{itemize}
  \item \textbf{typing:} if \(e : X \to E\) and \(e \to_E e'\), then \(e' : X \to E\);
  \item \textbf{TDG compatibility:} reductions act componentwise on TDG
        combinators (\(\BRANCH\), \(\PARALLEL\), \(\LOOP\));
  \item \textbf{termination:} there is no infinite chain
        \(e_0 \to_E e_1 \to_E e_2 \to_E \dots\);
  \item \textbf{local confluence:} if \(e \to_E e_1\) and \(e \to_E e_2\), then
        \(e_1\) and \(e_2\) have a common descendant.
\end{itemize}
We treat these properties as axioms of the explicit reduction system.
\end{definition}

Under these assumptions, every explicit artifact admits a unique normal
form modulo \(\to_E\).

%======================================================================
\subsection{Navigation and Shapes}
%======================================================================

\subsubsection{Navigation Category and Groupoid}

\begin{definition}[Navigation category \(N_0\)]
The navigation category \(N_0\) has:
\begin{itemize}
  \item objects: shapes \(S\);
  \item generating morphisms: navigation steps
        \(g : S \to S'\), such as \(\TRYMODEL\),
        \(\ADDASSUMPTION\), \(\DROPASSUMPTION\),
        \(\MERGE\), \(\RESTART\), each with a well-defined
        source and target shape;
  \item morphisms: finite composites of such generators.
\end{itemize}
\end{definition}

\begin{definition}[Navigation groupoid \(N\)]
The navigation groupoid \(N\) is obtained by freely adjoining inverses
to morphisms of \(N_0\) and quotienting by the relations
\[
g^{-1} \circ g = \id_S, \qquad
g \circ g^{-1} = \id_{S'},
\]
for each generator \(g : S \to S'\).
Thus every navigation step is reversible up to equality in \(N\).
\end{definition}

\subsubsection{Fibration \(p : E \to S\)}

\begin{definition}[Projection to shapes]
The functor \(p : \Espace \to \Scat\) sends each explicit pipeline
\(e \in \Espace\) to its shape
\(p(e) := \shape(e)\), and each explicit morphism to the corresponding
structural morphism between shapes.
\end{definition}

\begin{definition}[Fibers]
For each shape \(S\), the fiber
\[
E_S := p^{-1}(S)
\]
is the category of explicit morphisms with shape \(S\).
\end{definition}

\begin{definition}[Cartesian liftings and reindexing]
For each navigation morphism \(n : S \to S'\) in \(N\), we assume a
cartesian lifting
\[
n^\ast : E_S \to E_{S'}
\]
which reindexes explicit pipelines along \(n\). This yields an indexed
category
\[
E : N^{\mathrm{op}} \to \mathbf{Cat},
\qquad
S \mapsto E_S,\quad
n \mapsto n^\ast.
\]
\end{definition}

%======================================================================
\subsection{Combined System: Grothendieck Construction}
%======================================================================

\subsubsection{Semidirect Product \texorpdfstring{\(N \ltimes E\)}{N ⋉ E}}

\begin{definition}[Grothendieck construction \(N \ltimes E\)]
Given the indexed category
\(E : N^{\mathrm{op}} \to \mathbf{Cat}\), the Grothendieck
construction \(\int_N E\) (also written \(N \ltimes E\)) has:
\begin{itemize}
  \item objects: pairs \((S, e)\) with \(S \in \Obj(N)\) and \(e \in E_S\);
  \item morphisms \((S, e) \to (S', e')\): pairs \((n, f)\) where
        \(n : S \to S'\) in \(N\) and
        \(f : n^\ast(e) \to e'\) in \(E_{S'}\).
\end{itemize}
Composition is given by
\[
(n_1, f_1) \circ (n_2, f_2)
\;=\;
\bigl(
  n_1 \circ n_2,\,
  f_1 \circ n_1^\ast(f_2)
\bigr).
\]
Informally, we may write elements as \((n, e)\) and composition as
\((n_1, e_1) \circ (n_2, e_2) =
 (n_1 \circ n_2,\, e_1 \circ n_1(e_2))\) whenever the lifted action is defined.
\end{definition}

\subsubsection{Rewrite System on \texorpdfstring{\(N \ltimes E\)}{N ⋉ E}}

\begin{definition}[Combined reduction]
The one-step reduction relation \(\to\) on \(N \ltimes E\) is generated by:
\begin{itemize}
  \item \textbf{explicit step:}
        if \(e \to_E e'\) then \((n, e) \to (n, e')\);
  \item \textbf{navigation step:}
        if \(n \to_N n'\) and \(n'\) is defined on \(e\), then
        \((n, e) \to (n', n'(e))\).
\end{itemize}
We assume:
\begin{enumerate}[label=(A\arabic*)]
  \item \(\to_E\) is terminating and locally confluent;
  \item navigation reductions \(\to_N\) have normal forms in each component
        of \(N\);
  \item the lifted action is functorial and preserves \(\RB\), \(\TE\),
        and \(\HFD\).
\end{enumerate}
These are axioms of the FRFP reduction environment.
\end{definition}

%======================================================================
\subsection{Semantics and Tacit Correctness}
%======================================================================

\subsubsection{Semantic Functor}

\begin{definition}[Semantic bracket]
A \emph{semantic functor} is a functor
\[
\llbracket{-}\rrbracket : N \ltimes E_{\mathrm{sem}} \to \Tspace
\]
assigning to each combined configuration a tacit state.
On explicit pipelines \(e \in E_{\mathrm{sem}}\), we set
\[
\llbracket e \rrbracket := \HFD \circ \TE \circ \RBsharp(e).
\]
\end{definition}

\begin{definition}[Navigation invariance]
The functor \(\llbracket{-}\rrbracket\) is \emph{navigation-invariant up to
preorder} if
\[
\llbracket (n, e) \rrbracket \;\sqsubseteq\; \llbracket e \rrbracket
\]
for all \(n,e\) where the action is defined, and equality holds in the
quotient by the induced equivalence on tacit states in confluent cases.
\end{definition}

\begin{definition}[Observable correctness]
Given a pipeline \(P\), its observable correctness outcome is
\[
\gamma(\llbracket P \rrbracket) \in J.
\]
\end{definition}

%======================================================================
\subsection{Main Theorems}
%======================================================================

\subsubsection{Well-Formedness of TDG Pipelines}

\begin{theorem}[Well-formedness]
\label{thm:well-formed}
Every TDG-generated pipeline \(P\) is a well-typed morphism in \(\Cfull\) of
the form
\[
\RI \;\to\; (\EC/\ED)^{\ast} \;\to\; \RB \;\to\; \TE \;\to\; \HFD.
\]
\end{theorem}

\begin{proof}
By induction on TDG derivations. The base case \(\RI\) initializes an object
in \(E\). The constructors \(\EC\) and \(\ED\) have codomain \(E\) and thus
preserve explicit type. Structural combinators (\(\BRANCH\), \(\PARALLEL\),
\(\LOOP\)) act within the explicit category and preserve codomain \(E\)
componentwise. The production for \(P\) enforces the suffix
\(\RB; \TE; \HFD\), whose types are \(E \to T \to T \to T\).
By ETS there are no morphisms \(T \to E\), so:
\begin{itemize}
  \item no tacit operator can appear before \(\RB\);
  \item no explicit operator can appear after \(\RB\).
\end{itemize}
Thus every derivation yields a well-typed morphism of the stated form.
\end{proof}

\subsubsection{Explicit Completeness of TDG}

\begin{theorem}[Explicit completeness]
\label{thm:explicit-completeness}
TDG generates exactly the admissible explicit morphisms in \(\Espace\).
\end{theorem}

\begin{proof}
By construction, \(E_{\mathrm{sem}}\) is the free category on the TDG
signature, and \(\Ufun\) interprets TDG terms as explicit morphisms in
\(\Espace\). Every explicit morphism in the FRFP setting is built from
\(\RI\), \(\EC\), \(\ED\), and the structural operators, hence arises from
a TDG term; conversely, each TDG term denotes a well-typed explicit
morphism. The image of TDG under \(\Ufun\) is thus exactly the admissible
explicit morphisms.
\end{proof}

\subsubsection{Minimality of Primitive Set}

\begin{theorem}[Minimality]
\label{thm:minimality}
The primitive set
\(\{\RI,\EC,\ED,\RB,\TE,\HFD\}\)
is minimal under the Epistemic Spaces axioms: removing any primitive violates at
least one of ETS, HEG, HEC, RB, CP, or AE.
\end{theorem}

\begin{proof}
Briefly:
\begin{itemize}
  \item \(\RI\) is required to start from \(\varnothing\) and generate explicit
        content; without it there is no canonical initiation of pipelines.
  \item \(\EC\) is required for generative explicit reasoning; without it,
        explicit computation collapses.
  \item \(\ED\) is required for explicit diagnostics without tacit evaluation;
        removing it forces diagnostics either into tacit space (violating AE)
        or to be indistinguishable from \(\EC\) (violating CP and the intended
        modal roles).
  \item \(\RB\) is required for explicit-to-tacit handoff; without it there is
        no admissible morphism \(E \to T\), contradicting the requirement that
        pipelines reach tacit evaluation.
  \item \(\TE\) is required for tacit grounding (HEG); without it correctness
        cannot be evaluated in tacit space.
  \item \(\HFD\) is required for final correctness closure (HEC); without it
        there is no morphism corresponding to final acceptance/rejection.
\end{itemize}
Removing any primitive causes at least one Epistemic Spaces axiom to fail, so the
set is minimal.
\end{proof}

\subsubsection{Epistemic Spaces Initiality}

\begin{definition}[Epistemic Spaces framework]
A \emph{Epistemic Spaces framework} is a tuple
\[
\mathcal{F}_1
 \;=\;
(\mathcal{C}, E_{\mathcal{C}}, T_{\mathcal{C}},
 \RI_{\mathcal{C}}, \EC_{\mathcal{C}}, \ED_{\mathcal{C}},
 \RB_{\mathcal{C}}, \TE_{\mathcal{C}}, \HFD_{\mathcal{C}})
\]
satisfying the Epistemic Spaces axioms.
Morphisms in the category \(\Frm_1\) of Epistemic Spaces frameworks are functors
preserving \(E\), \(T\), and all six primitives.
\end{definition}

\begin{theorem}[Epistemic Spaces initiality]
\label{thm:phase1-initial}
There exists an initial object
\(F^{(1)}_{\mathrm{FRFP}} \in \Frm_1\).
\end{theorem}

\begin{proof}
Construct \(F^{(1)}_{\mathrm{FRFP}}\) as the free category on objects
\(E, T\) and morphisms \(\RI,\EC,\ED,\RB,\TE,\HFD\), modulo the equations
encoding the Epistemic Spaces axioms (ETS, HEG, HEC, RB, CP, AE). Any other
Epistemic Spaces framework interprets these generators in a way that satisfies the
same equations, yielding a unique induced functor
\(F^{(1)}_{\mathrm{FRFP}} \to \mathcal{F}_1\). Thus
\(F^{(1)}_{\mathrm{FRFP}}\) is initial in \(\Frm_1\).
\end{proof}
\paragraph{Remark A.34a (Non-trivial content of Epistemic Spaces initiality).}
The statement that $F^{(1)}_{\mathrm{FRFP}}$ is initial in $\mathrm{Frm}_1$ uses the free-category
construction, but its content is more than formal packaging. The Epistemic Spaces axioms (ETS, HEG,
HEC, RB, CP, AE), together with the minimality result (Theorem~A.32), severely constrain which
primitive roles may appear and how they can be composed. Initiality then says that any epistemic
framework respecting exactly these roles and axioms must interpret RI, EC, ED, RB, TE, HFD in a
way that factors uniquely through the FRFP ontology. In particular, attempts to ``hide'' additional
power in explicit space (for example, by letting EC or ED implicitly perform tacit evaluation) break
the axioms rather than yielding a genuinely different Epistemic Spaces framework.

\subsubsection{Epistemic Algebra Canonicality}

\begin{definition}[Epistemic Algebra architecture]
A \emph{Epistemic Algebra architecture} is a tuple
\[
\mathcal{F}_2
 =
(E,\ S,\ p,\ N,\ (\cdot)^\ast,\ N \ltimes E)
\]
satisfying:
\begin{itemize}
  \item TDG presents \(E\) as a free explicit category;
  \item \(N\) is a navigation groupoid on shapes \(S\);
  \item \(p : E \to S\) is a fibration with cartesian liftings \(n^\ast\);
  \item \(N \ltimes E\) is the Grothendieck construction of the indexed
        category \(E : N^{\mathrm{op}} \to \mathbf{Cat}\);
  \item the boundary structure \(\RB\)--\(\TE\)--\(\HFD\) is preserved
        under navigation.
\end{itemize}
Morphisms in the category of Epistemic Algebra architectures preserve all of this
structure.
\end{definition}

\begin{theorem}[Epistemic Algebra canonicality]
\label{thm:phase2-canonicality}
The FRFP Epistemic Algebra architecture
\(F^{(2)}_{\mathrm{FRFP}}\) is a canonical object in the category of
Epistemic Algebra architectures: for every Epistemic Algebra architecture \(\mathcal{A}\),
there exists a unique structure-preserving morphism
\[
F^{(2)}_{\mathrm{FRFP}} \to \mathcal{A}.
\]
\end{theorem}

\begin{proof}
By the universal properties of:
\begin{itemize}
  \item the free explicit category \(E_{\mathrm{sem}}\);
  \item the groupoid completion constructing \(N\);
  \item the Grothendieck construction forming \(N \ltimes E\).
\end{itemize}
Any Epistemic Algebra architecture satisfying the same axioms must interpret the
FRFP construction in a role-preserving way, inducing a unique morphism
from \(F^{(2)}_{\mathrm{FRFP}}\).
\end{proof}

\subsubsection{Confluence of \texorpdfstring{\(N \ltimes E\)}{N ⋉ E}}

\begin{theorem}[Confluence of \(N \ltimes E\)]
\label{thm:n-e-confluence}
Under assumptions \textup{(A1)}--\textup{(A3)}, the reduction relation
\(\to\) on \(N \ltimes E\) is confluent.
\end{theorem}

\begin{proof}[Proof sketch]
Local peaks \((n, e) \leftarrow (n_0, e_0) \to (n', e')\) decompose into:
\begin{itemize}
  \item explicit--explicit peaks, which join by local confluence of
        \(\to_E\);
  \item navigation--navigation peaks, which join by normalization in each
        connected component of \(N\);
  \item mixed peaks, which join because the lifted action commutes with
        explicit reduction and preserves \(\RB,\TE,\HFD\).
\end{itemize}
Together with termination up to navigation-normalization, this yields
global confluence.
\end{proof}
\paragraph{Remark A.37a (Non-trivial aspects of Epistemic Algebra structure).}
While Theorems~A.36 and A.37 are phrased in terms of standard categorical constructions
(free categories, groupoid completions, Grothendieck constructions), several aspects of the
Epistemic Algebra architecture are mathematically substantive:

\begin{itemize}
  \item \textbf{Compatibility with RB--TE--HFD.} The lifted navigation action and the
  Grothendieck construction are required to preserve the boundary structure. This rules out
  navigation steps that implicitly smuggle tacit evaluation back into explicit space or that
  duplicate/delete the RB--TE--HFD suffix. Showing that the action can be defined so as to
  commute with explicit reduction while preserving RB, TE, HFD is a genuine structural
  constraint on admissible navigation schemes.
  \item \textbf{Confluence of the combined reduction.} Confluence of $\to$ on $N \ltimes E$ is not
  automatic from confluence of $\to_E$ and normalization in $N$. The proof must analyze
  mixed peaks and use the functoriality of the lifted action, together with preservation of
  RB, TE, and HFD, to ensure that explicit and navigation steps can be permuted without
  changing the normal form.
  \item \textbf{Interaction with semantics.} The semantic correctness theorem (Theorem~A.38)
  depends crucially on the fact that navigation-invariant semantics on $N \ltimes E$ respects
  the Interaction Protocols constraints (no illicit $T \to E$ morphisms) and the boundary-preservation
  properties above. This is what forces correctness judgments to depend only on the tacit
  evaluation of the explicit pipeline, not on the particular exploration trajectory.
\end{itemize}

These non-trivial interactions justify treating Epistemic Algebra canonicality as more than a choice of
notation: any alternative architecture satisfying the same constraints must reproduce this
structure and, in particular, cannot avoid the resulting preservation and non-automation results
without breaking one of the axioms.

\subsubsection{Semantic Correctness for Epistemic Structure}

Collecting the constructions above, we obtain a semantic correctness
theorem for Epistemic Structure FRFP pipelines.

\begin{theorem}[Semantic correctness]
\label{thm:semantic-correctness}
Let \(P\) be any TDG-generated pipeline that is Interaction Protocols admissible and
Epistemic Spaces/Algebra well-typed. Then:
\begin{enumerate}
  \item \(P\) has a unique normal form in \(N \ltimes E\) modulo navigation;
  \item the observable correctness outcome \(\gamma(\llbracket P \rrbracket)\)
        is invariant under:
        \begin{itemize}
          \item explicit reduction \(\to_E\),
          \item navigation in \(N\),
          \item composition in the semidirect product \(N \ltimes E\).
        \end{itemize}
\end{enumerate}
In particular, correctness judgments depend only on the (tacitly evaluated)
semantic content of the explicit pipeline, not on the particular exploration
trajectory used to construct it.
\end{theorem}

\begin{proof}[Proof sketch]
Item~(1) follows from Theorem~\ref{thm:n-e-confluence} and the termination
assumptions (A1)--(A3). Item~(2) follows from navigation invariance of the
semantic functor, the Interaction Protocols constraints (which rule out illicit T-to-E
feedback or hidden tacit reconstruction), and the fact that \(\RB,\TE,\HFD\)
are preserved under navigation and do not appear in the explicit phase.
\end{proof}

\subsection{Worked example: sepsis consult as an FRFP pipeline}
\label{sec:cp-sepsis-example}
We sketch a concrete instance of the FRFP calculus corresponding to the sepsis consult scenario
of \S\ref{sec:case-sepsis}. The goal is to show how a real conversational workflow is
represented as a TDG term, reduced to normal form in $N \ltimes E$, and checked for
Interaction Protocols admissibility and Epistemic Spaces/Algebra well-typedness.

\paragraph{TDG and free explicit category.}
Let the TDG have objects
\[
  \Ob(E_0) = \{\textsf{Ctx}, \textsf{Input}, \textsf{CaseSummary},
  \textsf{DraftPlan}, \textsf{CandidatePlan}, \textsf{Report}, \textsf{Decision}\}
\]
and generators
\begin{align*}
  \mathsf{RB} &: \textsf{Ctx} \to \textsf{Input}, \\
  \mathsf{prep} &: \textsf{Input} \to \textsf{CaseSummary}, \\
  \mathsf{ehr} &: \textsf{CaseSummary} \to \textsf{CaseSummary}, \\
  \mathsf{calc} &: \textsf{CaseSummary} \to \textsf{CaseSummary}, \\
  \mathsf{assist}_0 &: \textsf{CaseSummary} \to \textsf{DraftPlan}, \\
  \mathsf{review}_\mathrm{res} &: \textsf{DraftPlan} \to \textsf{CandidatePlan}, \\
  \mathsf{review}_\mathrm{att} &: \textsf{CandidatePlan} \to \textsf{Report}, \\
  \mathsf{HFD} &: \textsf{Report} \to \textsf{Decision}.
\end{align*}
Let $E$ be the free category on this graph, modulo the Epistemic Algebra equations for $\mathsf{TE}$ and
the RB--TE--HFD boundary. We define
\[
  \mathsf{assist} := \mathsf{assist}_0 \circ \mathsf{calc} \circ \mathsf{ehr}
  : \textsf{CaseSummary} \to \textsf{DraftPlan},
\]
and the explicit sepsis pipeline
\[
  e_{\mathrm{sepsis}} :=
  \mathsf{HFD} \circ \mathsf{review}_\mathrm{att} \circ \mathsf{review}_\mathrm{res} \circ
  \mathsf{assist}_0 \circ \mathsf{calc} \circ \mathsf{ehr} \circ \mathsf{prep} \circ \mathsf{RB}
  : \textsf{Ctx} \to \textsf{Decision}.
\]

\paragraph{Reduction and normal form.}
Let $R_E$ be the explicit reduction system on $E$ that contracts $\mathsf{ehr}$ and $\mathsf{calc}$
into $\mathsf{assist}$ inside the $\mathsf{assist}_0$ wrapper:
\[
  \mathsf{assist}_0 \circ \mathsf{calc} \circ \mathsf{ehr} \;\to\; \mathsf{assist}.
\]
By confluence of $R_E$ and compatibility with the RB--TE--HFD boundary, the normal form of
$e_{\mathrm{sepsis}}$ is
\[
  e_{\mathrm{sepsis}}^{\ast} :=
  \mathsf{HFD} \circ \mathsf{review}_\mathrm{att} \circ \mathsf{review}_\mathrm{res} \circ
  \mathsf{assist} \circ \mathsf{prep} \circ \mathsf{RB}.
\]
Lifting to $N \ltimes E$ and applying the confluence theorem for the combined reduction (Epistemic Algebra)
ensures that any admissible navigation of the conversational UI yields a unique normal form
equivalent to $e_{\mathrm{sepsis}}^{\ast}$ up to the RB--TE--HFD suffix.

\paragraph{Interaction Protocols admissibility.}
We interpret the Interaction Protocols axioms IL, AR, MD, and NTER as constraints on realizations of
$e_{\mathrm{sepsis}}$ in a concrete UI:
\begin{itemize}
  \item IL: the only interaction channels between AI and human tacit state are via explicit objects in
  $E$ (case summaries, drafts, candidate plans, reports). No morphisms $T \to E$ are introduced.
  \item AR: there is no morphism in the realization that effects clinical actions without passing
  through $\mathsf{HFD} \circ \mathsf{review}_\mathrm{att}$; the AI cannot place orders or write notes.
  \item MD: all model-dependent computation is represented by $\mathsf{assist}$ and its internal
  decomposition; tacit evaluation (clinician judgment) is localized at $\mathsf{review}_\mathrm{res}$ and
  $\mathsf{review}_\mathrm{att}$.
  \item NTER: the implementation forbids any arrow that would encode a tacit correctness oracle in
  explicit space (no reconstruction of $T$ from model outputs). In particular, the realization of
  $\mathsf{assist}$ is constrained to produce drafts and rationales, not authoritative decisions.
\end{itemize}
Under these conditions, any concrete realization of $e_{\mathrm{sepsis}}$ satisfies the Interaction Protocols
interface axioms and is therefore admissible.

\paragraph{Epistemic Spaces/Algebra well-typedness.}
In Epistemic Spaces, each generator is assigned an epistemic role: $\mathsf{prep}$ and $\mathsf{assist}$ are
explicit-computation arrows; $\mathsf{review}_\mathrm{res}$ and $\mathsf{review}_\mathrm{att}$ are explicit
presentation points at which tacit evaluation in $T$ occurs; $\mathsf{RB}$ and $\mathsf{HFD}$ are
boundary morphisms. The typing rules of the Epistemic Spaces ontology then derive a composite judgment
showing that $e_{\mathrm{sepsis}}^{\ast}$ is a well-typed explicit pipeline with correctness localized in
tacit post-composition at $\mathsf{review}_\mathrm{att}$. In Epistemic Algebra, the same term appears as a
normal form in $N \ltimes E$; compatibility of reduction with the semantics ensures that the
correctness predicate applied to $e_{\mathrm{sepsis}}^{\ast}$ coincides with the protocol-level notion of
FRFP-hallucination (a violation of $t \succeq \mathrm{Req}(e,c)$ at the moment of final acceptance).
Thus the sepsis consult example realizes a concrete conversational UI that is Interaction Protocols admissible,
Epistemic Spaces/Algebra well-typed, and corresponds to a specific normal form in the FRFP calculus.

\subsection{Epistemic Dynamics and Safe Horizons} % old B
\label{app:dynamics}
% tacit degradation, Req, hallucination, inevitability, stochastic pipelines, hazard, H_P

%\section{Pipeline Dynamics}
%\label{app:pipeline-dynamics}

\subsubsection{Purpose and Scope}
The discussion in the Appendix till now formalizes
\emph{static} correctness of individual FRFP pipelines, Epistemic Dynamics
analyzes what happens when such pipelines are iterated by the same
human--AI system over many turns. The focus is on:
\begin{enumerate}
  \item persistence and degradation of tacit context;
  \item explicit-only hallucinations and their containment;
  \item confidence inflation and requirement escalation;
  \item feedback-coupled degradation and accelerated hallucination;
  \item stochastic and finite-horizon pipelines.
\end{enumerate}
No collective or institutional mechanisms are assumed in this section;
these appear in latter.

Throughout, we fix the Epistemic structure and work
with the underlying tacit context space \(T\) and explicit artifact
space \(E\).

%======================================================================
\subsubsection{Tacit Context Space}
%======================================================================

Each human agent possesses a tacit context space \(T\), equipped with a
preorder \(\preceq\).

\begin{definition}[Tacit context preorder]
The binary relation \(\preceq\) on \(T\) is read as:
\[
t_1 \preceq t_2
\quad\text{``\(t_1\) is epistemically no stronger than \(t_2\)''}.
\]
We write \(t_1 \prec t_2\) if \(t_1 \preceq t_2\) and \(t_2 \not\preceq t_1\).
\end{definition}

No algebraic structure on \(T\) is assumed beyond this preorder. The
categorical tacit structure  (one-object category,
primitives \(\TE,\HFD\), correctness quotient \(\gamma\)) is still in
force, but for dynamic analysis we work primarily with the underlying
preordered set \((T,\preceq)\).

%======================================================================
\subsubsection{Context Degradation}
%======================================================================

Tacit context degrades over time due to forgetting, misplaced trust,
habituation, and other long-run effects.

\begin{definition}[Context degradation]
A \emph{context degradation} operator is a function
\[
\delta : T \to T
\]
modeling loss of grounding across turns. We assume:
\begin{enumerate}
  \item \textbf{Monotonicity}:
        \(\delta(t) \preceq t\) for all \(t \in T\).
  \item \textbf{Strict decay somewhere}:
        there exists \(t \in T\) with \(\delta(t) \prec t\).
  \item \textbf{Non-invertibility}:
        \(\delta\) has no right inverse; i.e., there is no function
        \(f:T\to T\) such that \(\delta \circ f = \mathrm{id}_T\).
\end{enumerate}
We write \(\delta^n\) for the \(n\)-fold iterate of \(\delta\).
\end{definition}

Intuitively, \(\delta\) expresses the part of tacit dynamics that is
\emph{purely degrading}. In this appendix we deliberately ignore
repair/refresh operations (e.g., reading textbooks, taking breaks) to
obtain worst-case inevitability results. Repair can be modeled by
separate operators \(\rho:T\to T\), but these are left for future work.

%======================================================================
\subsubsection{Dynamic Interpretation of Pipelines}
%======================================================================

We now interpret complete FRFP pipelines as transformations of tacit
context.

Let \(E_{\mathrm{sem}}\) be the explicit semantic category and \(\RB^\sharp : E_{\mathrm{sem}}\to T\) the semantic
backflow (Section~A.0.5). For each explicit pipeline \(e\in
E_{\mathrm{sem}}\), the Epistemic Space suffix \(\TE,\HFD\) yields a tacit
transformer
\[
t \longmapsto \HFD \circ \TE \circ \RB^\sharp(e,t),
\]
where we allow the tacit state \(t\) to modulate interpretation
(e.g., via background knowledge).

\begin{definition}[Pipeline transformer]
A \emph{complete FRFP pipeline} \(P\) at a given turn consists of:
\begin{itemize}
  \item an explicit artifact \(e \in E_{\mathrm{sem}}\) produced by the
        AI system via explicit computation and navigation;
  \item the canonical tacit suffix \(\RB \to \TE \to \HFD\);
  \item subsequent degradation \(\delta\).
\end{itemize}
It induces a tacit transformer
\[
\llbracket P \rrbracket_\delta : T \to T,\qquad
\llbracket P \rrbracket_\delta(t)
 :=
 \delta\bigl( \HFD\circ\TE\circ\RB^\sharp(e,t)\bigr).
\]
\end{definition}

\begin{definition}[Run]
A \emph{run} (or \emph{pipeline execution}) is a sequence
\[
t_0, t_1, t_2,\dots
\]
of tacit states such that there exist pipelines
\(P_0, P_1, P_2,\dots\) with
\[
t_{n+1}
  = \llbracket P_n \rrbracket_\delta(t_n)
\quad\text{for all } n \geq 0.
\]
\end{definition}

We will often also track the explicit artifacts \(e_n\) and credences
\(c_n\) associated with each \(P_n\); this is made precise in
Sections~\ref{sec:B7-confidence} and~\ref{sec:B15-stochastic}.

%======================================================================
\subsubsection{Requirements and Hallucination}
%======================================================================

We formalize hallucination in terms of a requirement map on explicit
artifacts.

\begin{definition}[Requirement map]
A \emph{requirement map} is a function
\[
\Req : E \times \mathbb{R}_{\ge 0} \to T
\]
assigning to each explicit artifact \(e \in E\) and tacit credence
\(c\ge 0\) a required tacit grounding \(\Req(e,c)\). We require:
\begin{enumerate}
  \item \textbf{Monotonicity in credence}:
        if \(c_1 \le c_2\) then
        \(\Req(e,c_1) \preceq \Req(e,c_2)\);
  \item \textbf{Nontriviality}:
        there exist artifacts whose requirements strictly exceed others
        in the preorder.
\end{enumerate}
When credence is suppressed, we write \(\Req(e)\) for a baseline
requirement (e.g., at default credence).
\end{definition}

\begin{definition}[Hallucination]
Let \(t\in T\) be the tacit state at evaluation time, \(e\in E\) an
explicit artifact, and \(c\ge 0\) the tacit credence assigned to
\(e\). A \emph{hallucination} occurs at this step if
\[
t \prec \Req(e,c).
\]
We say that the artifact \(e\) is \emph{hallucinated} (relative to
\((t,c)\)) in this case.
\end{definition}

\begin{remark}[Structural nature of hallucination]
In this formalism, hallucination is an \emph{implementation-agnostic}
relation between tacit state \(t\), requirement \(\Req(e,c)\), and
explicit artifact \(e\). It does not depend on the model's internal
computation or intent. Concrete architectures (e.g., LLMs) shape
\emph{how often} trajectories enter states with \(t \prec
\Req(e,c)\), but do not appear in the definition itself.
\end{remark}

%======================================================================
\subsubsection{Hallucination Inevitability}
%======================================================================

We now state a baseline hallucination inevitability theorem for
deterministic degradation dynamics.

To capture recurring nontrivial demands, we restrict to runs that keep
encountering artifacts whose requirements exceed a fixed lower bound.

\begin{definition}[Nontrivial requirement floor]
A requirement floor is a tacit state \(r\in T\). A run
\((t_n,e_n,c_n)_{n\ge 0}\) has \emph{recurring nontrivial demands}
relative to \(r\) if there exists an infinite subsequence
\((n_k)_{k\ge 0}\) such that
\[
\Req(e_{n_k},c_{n_k}) \succeq r
\quad\text{for all }k.
\]
\end{definition}

We also assume that degradation eventually pushes tacit states below
this floor.

\begin{assumption}[Baseline degradation below the requirement floor]
\label{assump:baseline-degradation}
There exists \(r\in T\) such that for every initial tacit state
\(t_0\in T\), there exists \(N\in\mathbb{N}\) with
\[
\delta^N(t_0) \prec r.
\]
\end{assumption}

This is a strengthening of ``strict decay somewhere'': it requires that
every trajectory ultimately falls below the floor \(r\).

\begin{theorem}[Hallucination inevitability]
\label{thm:baseline-inevitability}
Assume:
\begin{enumerate}
  \item context degradation \(\delta\) satisfies monotonicity,
        non-invertibility, and
        Assumption~\ref{assump:baseline-degradation};
  \item requirements admit a nontrivial floor \(r\) and are monotone in
        credence;
  \item a run \((t_n,e_n,c_n)_{n\ge 0}\) has recurring nontrivial
        demands relative to \(r\).
\end{enumerate}
Then every infinite run produces at least one hallucination; i.e., there
exists \(n\) with
\[
t_n \prec \Req(e_n,c_n).
\]
\end{theorem}

\begin{proof}[Proof sketch]
By Assumption~\ref{assump:baseline-degradation}, for any initial
\(t_0\) there exists \(N\) with \(\delta^N(t_0) \prec r\). Consider the
given run and the subsequence \((n_k)\) along which
\(\Req(e_{n_k},c_{n_k})\succeq r\). For sufficiently large \(k\),
monotonicity of \(\delta\) and recurrence of nontrivial demands imply
that
\[
t_{n_k} \preceq \delta^{m}(t_0) \prec r \preceq \Req(e_{n_k},c_{n_k})
\]
for some \(m\), so \(t_{n_k} \prec \Req(e_{n_k},c_{n_k})\), i.e., a
hallucination. A detailed proof can be given by induction along the
run, but the key idea is that degradation eventually drives tacit
grounding below a floor that requirements exceed infinitely often.
\end{proof}

\begin{remark}
The theorem applies to \emph{nontrivial, open-ended} task streams with
recurrent demands above a fixed floor. Trivial avoidance regimes (e.g.,
restricting to a tiny class of always-easy tasks or always answering
``I don't know'') fall outside the intended scope.
\end{remark}

%======================================================================
\subsubsection{Explicit-Only Hallucinations}
%======================================================================
We now isolate explicit artifacts that are produced entirely within explicit reasoning, but which can be
hallucinated relative to some tacit state encountered along a run.

\begin{definition}[Explicit-only hallucination]
An \emph{explicit-only hallucination} is an explicit artifact
\(e \in E\) such that:
\begin{enumerate}
  \item \(e\) is well-typed and syntactically valid as an FRFP explicit
        object;
  \item \(e\) is surface-plausible under an explicit plausibility
        functional (Section~\ref{sec:B7-confidence});
  \item there exists a tacit state \(t\) and credence \(c\) encountered
        in some run such that
        \(t \prec \Req(e,c)\).
\end{enumerate}
\end{definition}

\begin{proposition}[Containment but not elimination]
Nothing in Epistemic Space or Epistemic Algebra forbids the production of explicit-only
hallucinations. Eliminating them entirely would require either:
\begin{itemize}
  \item an explicit correctness oracle, or
  \item a simulation of tacit evaluation inside explicit space,
\end{itemize}
both of which would violate Explicit--Tacit Separation (ETS) and
Human-Exclusive Grounding (HEG).

The explicit--tacit separation guarantees \emph{containment}, not
\emph{elimination}: explicit-only hallucinations cannot become
correctness judgments without passing through \(\RB\), but they may
still be produced arbitrarily often within explicit space.
\end{proposition}

%======================================================================
\subsubsection{Confidence Inflation}
\label{sec:B7-confidence}
%======================================================================

We model how plausible explicit artifacts raise tacit credence even
when they do not increase grounding.

\begin{definition}[Explicit plausibility]
An \emph{explicit plausibility functional} is a map
\[
\pi : E \to [0,1],
\]
measuring surface coherence, fluency, or apparent authority of explicit
artifacts.
\end{definition}

We extend tacit state with a credence component.

\begin{definition}[Extended tacit state]
Let
\[
T^e := T \times \mathbb{R}_{\ge 0}
\]
be the space of extended tacit states \((t,c)\), where \(t\in T\) is
grounding and \(c\) is tacit credence (how settled or reliable the
agent treats the current artifact).
\end{definition}

\begin{definition}[Confidence inflation]
A \emph{confidence inflation operator} is a map
\[
\Phi : (T^e \times E) \to T^e,\qquad
\Phi\bigl( (t,c), e \bigr)
 :=
 (t,\; c + \varphi(\pi(e))),
\]
where \(\varphi:[0,1]\to \mathbb{R}_{\ge 0}\) is monotone increasing.

We require that \(\Phi\) does not increase tacit grounding:
\[
\Phi\bigl((t,c),e\bigr) = (t',c') \quad\Rightarrow\quad t' = t.
\]
\end{definition}

Thus confidence inflation changes only the agent's interpretive stance
(higher credence), not the underlying grounding \(t\).

\begin{remark}
Confidence inflation captures the structural asymmetry between
``feeling confident'' and ``being grounded''. Systems optimized for
producing well-formed, plausible outputs (e.g., LLMs) inevitably
exhibit this asymmetry: they can increase credence by producing
plausible artifacts without necessarily increasing grounding.
\end{remark}

%======================================================================
\subsubsection{Requirement Escalation Under Plausibility}
%======================================================================

Credence affects what requirements are appropriate: asserting
high-credence claims demands stronger grounding.

\begin{definition}[Requirement escalation]
The requirement map \(\Req : E \times \mathbb{R}_{\ge 0} \to T\) is
\emph{escalatory} if it is monotone in credence, i.e.,
\[
c_1 \le c_2 \quad\Rightarrow\quad
\Req(e,c_1) \preceq \Req(e,c_2),
\]
for all \(e\in E\). Repeated exposure to highly plausible artifacts
(large \(\pi(e)\)) can then raise requirements even if the explicit
content \(e\) is unchanged.
\end{definition}

Combining confidence inflation and requirement escalation, we see that
the same explicit artifact can acquire higher requirements over time as
agents become more confident in it.

%======================================================================
\subsubsection{Feedback-Coupled Degradation}
%======================================================================

We now allow degradation to depend on credence, introducing a feedback
loop: overconfidence reduces skepticism and checking, accelerating loss
of grounding.

\begin{definition}[Feedback-coupled degradation]
A \emph{feedback-coupled degradation} operator is a function
\[
\delta^\star : T^e \to T,\qquad
\delta^\star(t,c) = \delta_c(t),
\]
where \(\delta_c : T \to T\) is a family of degradation operators
indexed by credence \(c\), satisfying:
\begin{enumerate}
  \item \textbf{Credence monotonicity}:
        for fixed \(t\), if \(c_1 \le c_2\) then
        \(\delta^\star(t,c_2) \preceq \delta^\star(t,c_1)\);
        i.e., higher confidence weakens effective grounding.
  \item \textbf{Consistency with baseline degradation}:
        \(\delta^\star(t,0) = \delta(t)\) for all \(t\).
\end{enumerate}
\end{definition}

\begin{definition}[Extended run with feedback]
An \emph{extended run} is a sequence
\[
(t_0,c_0,e_0), (t_1,c_1,e_1), \dots
\]
such that, at each step \(n\),
\begin{align*}
(t_n^+, c_n^+) &= \Phi\bigl((t_n,c_n), e_n\bigr)
      &&\text{(confidence inflation)},\\
t_{n+1} &= \delta^\star(t_n^+, c_n^+)
      &&\text{(feedback-coupled degradation)},\\
c_{n+1} &\text{ chosen by tacit update (e.g., reset or carry-over)}.
\end{align*}
\end{definition}

The exact rule for updating \(c_{n+1}\) can be left abstract; all we
require is that confidence inflation can push credence arbitrarily high
along some runs.

%======================================================================
\subsubsection{Rate of Context Loss}
%======================================================================}

We can analyze the rate at which tacit grounding is lost under
feedback-coupled degradation.

\begin{definition}[Grounding measure and loss rate]
Let \(\mu : T \to \mathbb{R}_{\ge 0}\) be any monotone measure of tacit
grounding (used only for analysis). For an extended run
\((t_n,c_n,e_n)\), define the instantaneous loss rate at turn \(n\) by
\[
\rho_n := \mu(t_n) - \mu\bigl(\delta^\star(t_n,c_n)\bigr).
\]
\end{definition}

\begin{proposition}[Loss rate increases with credence]
Under the credence-monotonicity assumption for \(\delta^\star\) and
monotonicity of \(\mu\), the loss rate \(\rho_n\) is non-decreasing in
\(c_n\). In particular, higher credence does not slow context loss and
may accelerate it.
\end{proposition}

\begin{proof}
If \(c_1 \le c_2\), then
\(\delta^\star(t,c_2) \preceq \delta^\star(t,c_1)\) and
monotonicity of \(\mu\) yields
\(\mu(\delta^\star(t,c_2)) \le \mu(\delta^\star(t,c_1))\). Thus
\[
\mu(t) - \mu(\delta^\star(t,c_2))
\;\ge\;
\mu(t) - \mu(\delta^\star(t,c_1)),
\]
so \(\rho_n\) increases with \(c_n\).
\end{proof}

%======================================================================
\subsubsection{Accelerated Hallucination}
%======================================================================

We now combine confidence inflation, requirement escalation, and
feedback-coupled degradation.

\begin{definition}[Feedback-coupled update]
Along an extended run, the joint update
\[
(t_{n+1}, c_{n+1})
 =
\Bigl(
  \delta^\star\bigl( t_n, c_n + \varphi(\pi(e_n))\bigr),\;
  c_{n+1}
\Bigr)
\]
captures the combined effect of:
\begin{itemize}
  \item confidence inflation
        \((t_n, c_n) \mapsto (t_n, c_n + \varphi(\pi(e_n)))\);
  \item degradation \(\delta^\star\), which is more severe at higher
        credence.
\end{itemize}
Grounding weakens faster while requirements grow stronger (via
\(\Req(e_n,c_n)\)).
\end{definition}

This creates a widening gap between available grounding \(t_n\) and
demanded grounding \(\Req(e_n,c_n)\), leading to accelerated
hallucination.

%======================================================================
\subsubsection{Accelerated Hallucination Inevitability}
%======================================================================

We can now state a strengthened inevitability theorem in the presence
of feedback.

\begin{assumption}[Unbounded credence under plausibility]
There exists \(\bar{\pi} \in (0,1]\) such that for any initial credence
\(c_0\) and any infinite sequence of artifacts with
\(\pi(e_n) \ge \bar{\pi}\) infinitely often, the accumulated credence
\[
c_n^{+} := c_0 + \sum_{k=0}^{n-1} \varphi(\pi(e_k))
\]
tends to \(+\infty\) along a subsequence.
\end{assumption}

\begin{theorem}[Accelerated hallucination inevitability]
\label{thm:accelerated-inevitability}
Assume:
\begin{enumerate}
  \item baseline context degradation \(\delta\) is monotone and
        non-invertible and satisfies
        Assumption~\ref{assump:baseline-degradation};
  \item confidence inflation \(\Phi\) with
        \(\varphi(\bar{\pi}) > 0\) for some \(\bar{\pi}\in(0,1]\);
  \item requirement escalation: \(\Req\) is monotone in credence;
  \item feedback-coupled degradation \(\delta^\star\) is credence-monotone
        and reduces to \(\delta\) at zero credence;
  \item the run contains infinitely many \emph{surface-plausible}
        explicit artifacts with \(\pi(e_n) \ge \bar{\pi}\) and recurring
        nontrivial demands relative to some floor \(r\).
\end{enumerate}
Then:
\begin{enumerate}
  \item hallucination occurs in finite time; i.e., there exists
        \(n < \infty\) such that \(t_n \prec \Req(e_n,c_n)\);
  \item the minimum such \(n\) is strictly smaller than in the baseline
        degradation model without feedback (for any fixed initial
        state and pipeline sequence), in the sense that feedback cannot
        increase time-to-hallucination.
\end{enumerate}
\end{theorem}

\begin{proof}[Proof sketch]
By unbounded credence under plausibility and the presence of infinitely
many plausible artifacts, credence \(c_n\) grows without bound along
some subsequence. By credence-monotonicity of \(\delta^\star\) and
requirement escalation, higher credence simultaneously:
\begin{itemize}
  \item increases required grounding \(\Req(e_n,c_n)\), and
  \item accelerates degradation of available grounding via
        \(\delta^\star\).
\end{itemize}
Thus the gap
\[
t_n \prec \Req(e_n,c_n)
\]
is reached in finite time, and strictly earlier than under the baseline
degradation \(\delta\) alone, where the same sequence of artifacts
would produce slower loss of grounding and lower requirements.
A fully formal proof can be obtained by coupling two runs (with and
without feedback) and using the monotone measure \(\mu\) to compare
times at which \(\mu(t_n)\) falls below \(\mu(\Req(e_n,c_n))\).
\end{proof}


\subsubsection{Degradation--repair dynamics and entropy drift}
\label{app:delta-rho}
We extend the Epistemic Dynamics to include abstract repair. The aim is to state conditions under
which repair delays or reshapes, but does not remove, long-run degradation.

\begin{definition}[Entropy functional on $T$]
An \emph{entropy functional} on the tacit category $T$ is a map $S : T \to \mathbb{R}$ such that
\begin{enumerate}
  \item $t \preceq t' \Rightarrow S(t) \le S(t')$ for all $t,t' \in T$; and
  \item there exists $S_{\mathrm{req}} \in \mathbb{R}$ such that $S(t) \le S_{\mathrm{req}}$ whenever $t$
  meets the relevant requirement floor(s), and $S(t) > S_{\mathrm{req}}$ implies that accepting a
  nontrivial demand at $t$ yields an FRFP-hallucination in the sense $t_n \prec \mathrm{Req}(e_n,c_n)$.
\end{enumerate}
Thus $S(t)$ measures a grounding deficit: larger values correspond to more degraded tacit states.
\end{definition}

We consider runs $(t_n)_{n \ge 0}$ in $T$ evolving under a mixture of degradation and repair, and
write $S_n := S(t_n)$. Let $D_n,R_n \in \{0,1\}$ indicate, respectively, whether step $n$ processes a
nontrivial demand and whether a repair operation is applied.

\begin{assumption}[Task-induced degradation]\label{ass:delta-entropy}
There exists $\varepsilon > 0$ such that for all $t \in T$ and all nontrivial demand steps to which the
Epistemic Dynamics apply,
\[
  S(\delta(t)) \;\ge\; S(t) + \varepsilon.
\]
\end{assumption}

\begin{assumption}[Bounded repair]\label{ass:rho-entropy}
There exists $B < \infty$ such that for all admissible repair operators $\rho : T \to T$ and all
$t \in T$,
\[
  S(\rho(t)) \;\ge\; S(t) - B.
\]
\end{assumption}

Combining these, for each $N$ we have
\begin{equation}\label{eq:Sn-balance}
  S_N \;\ge\; S_0 + \varepsilon \sum_{n=1}^N D_n \;-\; B \sum_{n=1}^N R_n,
\end{equation}
where the sums record entropy produced by nontrivial tasks and removed by repair.
\begin{theorem}[Deterministic entropy drift and inevitability]\label{thm:delta-rho-deterministic}
Suppose Assumptions~\ref{ass:delta-entropy} and~\ref{ass:rho-entropy} hold and that
$(D_n,R_n)_{n \ge 1}$ satisfies
\[
  \liminf_{N \to \infty} \frac{1}{N}\sum_{n=1}^N D_n \;\ge\; \alpha > 0,
  \qquad
  \limsup_{N \to \infty} \frac{1}{N}\sum_{n=1}^N R_n \;\le\; \beta < \infty.
\]
Assume further that there exists $\gamma > 0$ with
\[
  \varepsilon \alpha - B \beta \;\ge\; \gamma > 0.
\]
Then there exists $N_0$ such that $S_N \ge S_0 + \gamma N$ for all $N \ge N_0$. In particular, for
any finite $S_{\mathrm{req}}$ the hitting time
\[
  \sigma_{\mathrm{req}} := \inf\{n \ge 0 : S_n > S_{\mathrm{req}}\}
\]
is finite. If nontrivial demands recur after $\sigma_{\mathrm{req}}$ and the protocol sometimes accepts
artifacts at those steps with credence corresponding to $S_{\mathrm{req}}$, then along some admissible
run an FRFP-hallucination event $t_n \prec \mathrm{Req}(e_n,c_n)$ occurs.
\end{theorem}

\begin{proof}[Proof sketch]
Divide~\eqref{eq:Sn-balance} by $N$ and apply the liminf/limsup bounds on the empirical
frequencies of $D_n$ and $R_n$ to obtain
\[
  \liminf_{N \to \infty} \frac{S_N - S_0}{N}
  \;\ge\; \varepsilon \alpha - B \beta \;\ge\; \gamma > 0.
\]
Hence $S_N \ge S_0 + \gamma N$ for all large $N$, so $S_N$ eventually exceeds any finite
$S_{\mathrm{req}}$ and $\sigma_{\mathrm{req}} < \infty$. The hallucination clause follows from the
definition of $S_{\mathrm{req}}$ and the criterion $t_n \prec \mathrm{Req}(e_n,c_n)$ at an accepted
nontrivial demand.
\end{proof}
\begin{theorem}[Stochastic entropy drift and almost-sure inevitability]\label{thm:delta-rho-stochastic}
Suppose Assumptions~\ref{ass:delta-entropy} and~\ref{ass:rho-entropy} hold and that
$(D_n,R_n)_{n \ge 1}$ is stationary and ergodic with
\[
  \mathbb{E}[D_1] = \bar{d} > 0,
  \qquad
  \mathbb{E}[R_1] = \bar{r} < \infty.
\]
Assume that
\[
  \varepsilon \bar{d} - B \bar{r} \;=\; \gamma > 0.
\]
Then almost surely
\[
  \liminf_{N \to \infty} \frac{S_N - S_0}{N} \;\ge\; \gamma > 0,
\]
so $S_N \to +\infty$ almost surely. In particular, the hitting time
$\sigma_{\mathrm{req}} := \inf\{n \ge 0 : S_n > S_{\mathrm{req}}\}$ is almost surely finite. Under the
same acceptance and recurrence conditions as in Theorem~\ref{thm:delta-rho-deterministic}, the
hallucination time $\tau(\omega)$ in the Epistemic Dynamics stochastic semantics satisfies $P(\tau < \infty) = 1$.
\end{theorem}

\begin{proof}[Proof sketch]
By the ergodic theorem, almost surely
\[
  \lim_{N \to \infty} \frac{1}{N}\sum_{n=1}^N D_n = \bar{d},
  \qquad
  \lim_{N \to \infty} \frac{1}{N}\sum_{n=1}^N R_n = \bar{r}.
\]
Together with~\eqref{eq:Sn-balance} this yields
\[
  \liminf_{N \to \infty} \frac{S_N - S_0}{N} \;\ge\; \varepsilon \bar{d} - B \bar{r} = \gamma > 0
  \quad\text{a.s.}
\]
Thus $S_N \to +\infty$ almost surely and $\sigma_{\mathrm{req}} < \infty$ almost surely; whenever a
nontrivial demand is accepted at or after $\sigma_{\mathrm{req}}$, the hallucination criterion is met at
some finite step, so $P(\tau < \infty) = 1$.
\end{proof}
\begin{remark}[Counter-scenarios and safe horizons]\label{rem:delta-rho-counter}
The drift condition in Theorems~\ref{thm:delta-rho-deterministic} and~\ref{thm:delta-rho-stochastic}
formalizes a high-load, bounded-repair regime: nontrivial demands occur with positive asymptotic
frequency, repair has bounded one-step effect, and net entropy production $\varepsilon \bar{d}$
exceeds maximal entropy removal $B \bar{r}$. In this regime, repair changes the timing and shape
of hallucination risk but does not eliminate it.

Several counter-scenarios lie at or beyond this domain. In \emph{balanced or negative drift}
regimes, repair investment is large enough that $B \bar{r} \ge \varepsilon \bar{d}$ and $S_N$ can remain
tight or decrease. In \emph{unbounded repair} regimes, one allows repair moves with no global
bound $B$ (e.g.\ costless, arbitrarily frequent ``oracle refresh'' operations or always-perfect staff
replacement), violating Assumption~\ref{ass:rho-entropy}. In \emph{low-demand} regimes, such as
rare emergency committees, the effective demand frequency $\bar{d}$ may be so small that modest
repair keeps $S$ below $S_{\mathrm{req}}$ over relevant horizons.

Within the high-load, bounded-repair regime, improving repair capacity modifies the hazard profile
$(h_n)_{n \ge 0}$ and increases the safe-horizon functional $H_P(\varepsilon)$, but for any fixed
stochastic FRFP pipeline $P$ satisfying the drift condition, $H_P(\varepsilon)$ remains finite for every
$\varepsilon \in (0,1)$: one cannot obtain a global guarantee of ``no hallucinations ever'' over
unbounded horizons by bounded repair alone.
\end{remark}

%======================================================================
\subsubsection{Stochastic Pipelines and Their Impact}
\label{sec:B15-stochastic}
%======================================================================

We extend the preceding results to stochastic FRFP pipelines.

\subsubsubsection*{Trajectories and probability measures}

\begin{definition}[Trajectory]
Fix an initial extended state \((t_0,c_0,e_0)\in T\times
\mathbb{R}_{\ge 0}\times E\). A \emph{trajectory} (or \emph{run}) is a
sequence
\[
\omega
 =
\bigl( (t_0,c_0,e_0), (t_1,c_1,e_1), (t_2,c_2,e_2), \dots \bigr)
\]
such that each \((t_{n+1},c_{n+1},e_{n+1})\) is obtained from
\((t_n,c_n,e_n)\) by one application of the FRFP pipeline, tacit
evaluation, confidence update, and degradation (possibly
feedback-coupled).
\end{definition}

Let \(\Omega\) denote the set of all such trajectories. Let
\(\Omega_{\mathrm{allowed}}\subseteq \Omega\) be the subset of
trajectories that satisfy the structural assumptions of this appendix
(degradation, recurring requirements, etc.).

We consider two equivalent presentations of stochasticity:
\begin{itemize}
  \item \emph{Random-seed semantics:} extend the state space with a
        random seed \(r\), and suppose there is a deterministic update
        map
        \[
        F: T^e\times E \times R \to T^e \times E \times R
        \]
        such that each \(r\in R\) induces a unique trajectory
        \(\omega(r)\in\Omega\).
  \item \emph{Markov-kernel semantics:} one-step evolution is described
        by a Markov kernel
        \[
        K : T^e\times E \to \mathcal{D}(T^e\times E),
        \]
        where \(\mathcal{D}(\cdot)\) indicates probability
        distributions; a trajectory \(\omega\) is then a sample path of
        the stochastic process generated by \(K\).
\end{itemize}

In either case, a \emph{stochastic FRFP pipeline} is identified with a
probability measure \(\mathbb{P}\) on \(\Omega\) induced by the
underlying randomness.

\subsubsubsection*{Hallucination events}

Recall the hallucination criterion: hallucination at step \(n\) occurs
when
\[
t_n \prec \Req(e_n,c_n).
\]

\begin{definition}[Hallucination events]
For a trajectory \(\omega\in\Omega\) and \(n\ge 0\), define the
indicator
\[
H_n(\omega)
 :=
 \mathbf{1}\bigl[t_n \prec \Req(e_n,c_n)\bigr] \in \{0,1\}.
\]
Define the \emph{eventual hallucination indicator}
\[
H(\omega)
 :=
 \mathbf{1}\bigl[ \exists n\ge 0 : H_n(\omega)=1 \bigr].
\]
Thus \(H(\omega)=1\) iff the trajectory \(\omega\) contains at least one
hallucination.
\end{definition}

\begin{assumption}[Pathwise inevitability]
\label{assump:pathwise-inevitability}
Suppose the structural assumptions of Sections~B.1--B.12 hold (on
degradation, requirements, and feedback), and that the baseline or
accelerated hallucination inevitability theorem has been established in
the deterministic setting. We assume:
\[
\forall \omega \in \Omega_{\mathrm{allowed}}:\quad H(\omega)=1.
\]
That is, every allowed infinite trajectory eventually produces a
hallucination.
\end{assumption}

\begin{theorem}[Almost-sure inevitability under stochastic pipelines]
\label{thm:stochastic-inevitability}
Let \(\mathbb{P}\) be any probability measure on \(\Omega\) whose
support is contained in \(\Omega_{\mathrm{allowed}}\). Then
\[
\mathbb{P}\bigl( \text{eventual hallucination} \bigr)
 =
\mathbb{P}\bigl( \{\omega : H(\omega)=1\} \bigr)
 =
1.
\]
\end{theorem}

\begin{proof}
By Assumption~\ref{assump:pathwise-inevitability}, \(H(\omega)=1\) for
all \(\omega\in\Omega_{\mathrm{allowed}}\). Since \(\mathbb{P}\) is
supported on \(\Omega_{\mathrm{allowed}}\), we have \(H(\omega)=1\) for
\(\mathbb{P}\)-almost every \(\omega\). Hence
\[
\mathbb{P}(\text{eventual hallucination})
 =
\int_{\Omega} H(\omega)\, d\mathbb{P}(\omega)
 =
\int_{\Omega} 1\, d\mathbb{P}(\omega)
 =
1.
\]
\end{proof}

Thus passing from deterministic to stochastic pipelines does not weaken
the inevitability result: randomness merely selects \emph{which}
hallucinating trajectory is realized, not whether hallucination occurs.

\subsubsubsection*{Probabilistic confluence}

Stochasticity breaks single-outcome confluence but yields a natural
probabilistic analogue.

\begin{definition}[Output distribution and probabilistic confluence]
Fix a time \(n\ge 0\) and initial state \((t_0,c_0,e_0)\). Let
\(\mathbb{P}\) be a stochastic FRFP pipeline with support in
\(\Omega_{\mathrm{allowed}}\). The \emph{time-\(n\) output
distribution} is the pushforward measure
\[
\mu_n^{\mathbb{P}} := (e_n)_\# \mathbb{P},
\]
where \(e_n:\Omega\to E\) maps a trajectory to its explicit artifact at
step \(n\).

Two stochastic implementations \(\mathbb{P}_1,\mathbb{P}_2\) are
\emph{probabilistically confluent at time \(n\)} if they induce the same
output distribution \(\mu_n^{\mathbb{P}_1}=\mu_n^{\mathbb{P}_2}\) on
\(E\).
\end{definition}

If both implementations correspond to the same Markov kernel on
\(T^e\times E\), they are probabilistically confluent at every time
\(n\). This notion provides a distributional counterpart to the
canonical-form and confluence properties discussed earlier.

%======================================================================
\subsubsection{Finite-Horizon Pipelines and Their Impact}
%======================================================================

We now study the effect of restricting attention to finite prefixes of
trajectories, both in deterministic and stochastic settings.

\subsubsubsection*{Hallucination time}

\begin{definition}[Hallucination time]
For a trajectory \(\omega\in\Omega\), define the \emph{hallucination
time}
\[
\tau(\omega)
 :=
 \inf\{ n \ge 0 : t_n \prec \Req(e_n,c_n) \},
\]
with the convention \(\tau(\omega)=+\infty\) if the set is empty (i.e.,
if no hallucination ever occurs).
\end{definition}

Thus \(\tau(\omega) < \infty\) iff \(\omega\) contains at least one
hallucination. In terms of Definition~B.13.2, \(H(\omega) =
\mathbf{1}[\tau(\omega) < \infty]\).

\begin{lemma}[Pathwise inevitability and finite hallucination time]
Under Assumption~\ref{assump:pathwise-inevitability} (pathwise
inevitability), we have
\[
\forall \omega \in \Omega_{\mathrm{allowed}}:\quad \tau(\omega) < \infty.
\]
\end{lemma}

\begin{proof}
By Assumption~\ref{assump:pathwise-inevitability}, for each
\(\omega\in\Omega_{\mathrm{allowed}}\) there exists some \(n\) with
\(H_n(\omega)=1\). By definition of \(\tau\), this implies that
\(\tau(\omega)\) is the infimum of a nonempty subset of \(\mathbb{N}\),
hence finite.
\end{proof}

A \emph{finite horizon} of length \(N\) corresponds to observing only
the prefix
\[
\omega_{\le N}
 :=
((t_0,c_0,e_0),\dots,(t_N,c_N,e_N)).
\]
We say that a trajectory \emph{hallucinates within horizon \(N\)} if
\(\tau(\omega)\le N\).

\begin{remark}[No universal finite bound in general]
Pathwise inevitability (\(\tau(\omega) < \infty\) for all allowed
\(\omega\)) does not imply the existence of a universal finite upper
bound \(N^\ast\) such that \(\tau(\omega)\le N^\ast\) for all
\(\omega\in\Omega_{\mathrm{allowed}}\). In general, for any fixed
\(N\) there may exist trajectories with \(\tau(\omega) > N\).
\end{remark}

\subsubsubsection*{A finite-time bound under strengthened assumptions}

If desired, one can obtain a finite-time upper bound under stronger
assumptions.

\begin{assumption}[Uniform loss and requirement floor]
\label{assump:uniform-loss}
Suppose there exist:
\begin{itemize}
  \item a grounding measure \(\mu:T\to\mathbb{R}_{\ge 0}\);
  \item a requirement level \(r\in T\) and constant \(\alpha>0\) with
        \(\mu(r)=\alpha\);
  \item a constant \(\varepsilon>0\);
\end{itemize}
such that, for every allowed trajectory
\(\omega=((t_n,c_n,e_n))_{n\ge 0}\):
\begin{enumerate}
  \item \emph{Monotone loss:} \(\mu(t_{n+1})\le\mu(t_n)\) for all
        \(n\).
  \item \emph{Uniform loss above \(r\):} whenever \(\mu(t_n)>\alpha\),
        we have
        \[
        \mu(t_n) - \mu(t_{n+1}) \ge \varepsilon.
        \]
  \item \emph{Recurrent requirement:} there exists an infinite sequence
        \((n_k)\) with \(\Req(e_{n_k},c_{n_k})\succeq r\) for all
        \(k\).
\end{enumerate}
\end{assumption}

\begin{lemma}[Finite-time bound under uniform loss]
\label{lem:finite-time-bound}
Under Assumption~\ref{assump:uniform-loss}, every allowed trajectory
\(\omega\) satisfies
\[
\tau(\omega)
 \le
N^\ast
 :=
\left\lceil \frac{\mu(t_0) - \alpha}{\varepsilon} \right\rceil.
\]
That is, a hallucination must occur by time \(N^\ast\).
\end{lemma}

\begin{proof}
By monotone loss and the uniform-loss condition, as long as
\(\mu(t_n)>\alpha\),
\[
\mu(t_{n+1}) \le \mu(t_n) - \varepsilon.
\]
By induction, for any \(n\) with this property,
\[
\mu(t_n) \le \mu(t_0) - n\varepsilon.
\]
For \(n = N^\ast\) we obtain \(\mu(t_{N^\ast}) \le \alpha =
\mu(r)\). By recurrence of requirements, there exists some \(k\) with
\(n_k \le N^\ast\) and \(\Req(e_{n_k},c_{n_k}) \succeq r\). Using
monotonicity of \(\mu\) with respect to \(\preceq\), we conclude that
\[
t_{n_k} \preceq r \preceq \Req(e_{n_k},c_{n_k}),
\]
and in fact \(t_{n_k} \prec \Req(e_{n_k},c_{n_k})\) for nontrivial
requirements, so \(\tau(\omega)\le n_k \le N^\ast\).
\end{proof}

These stronger conditions are not needed for the main inevitability
results, but they demonstrate that finite-time bounds are available in
regimes with sufficiently uniform context loss and recurrent demands.

\subsubsubsection*{Stochastic finite horizons}

Returning to the stochastic setting, we can describe how
finite-horizon hallucination probabilities behave.

\begin{definition}[Finite-horizon hallucination probability]
Fix a horizon \(N\in\mathbb{N}\). Define the event
\[
A_N := \{\omega\in\Omega : \tau(\omega) \le N\},
\]
and the corresponding probability
\[
p_N := \mathbb{P}(A_N)
     = \mathbb{P}(\tau \le N).
\]
Thus \(p_N\) is the probability that a random run hallucinates within
the first \(N\) steps.
\end{definition}

\begin{proposition}[Monotone convergence of finite-horizon probabilities]
Assume almost-sure inevitability
(Theorem~\ref{thm:stochastic-inevitability}), i.e.,
\(\mathbb{P}(\tau<\infty)=1\). Then:
\begin{enumerate}
  \item The sequence \((p_N)_{N\ge 0}\) is non-decreasing:
        \(p_N \le p_{N+1}\) for all \(N\).
  \item \(\displaystyle \lim_{N\to\infty} p_N = 1\).
\end{enumerate}
\end{proposition}

\begin{proof}
For each \(\omega\), the events \(\{\tau(\omega)\le N\}\) are nested:
\(\{\tau\le N\}\subseteq\{\tau\le N+1\}\). Hence
\(A_N\subseteq A_{N+1}\) and \(p_N\le p_{N+1}\). Moreover,
\(\{\tau<\infty\}=\bigcup_{N} A_N\), so by continuity of probability
from below,
\[
\mathbb{P}(\tau<\infty)
 =
\lim_{N\to\infty} \mathbb{P}(A_N)
 =
\lim_{N\to\infty} p_N.
\]
By assumption, \(\mathbb{P}(\tau<\infty)=1\), so
\(\lim_{N\to\infty} p_N=1\).
\end{proof}

\begin{proposition}[Multiple runs over a fixed horizon]
Fix a horizon \(N\) and consider \(K\in\mathbb{N}\) independent
trajectories \(\omega^{(1)},\dots,\omega^{(K)}\), with corresponding
hallucination times \(\tau^{(i)}\). Then:
\begin{enumerate}
  \item The probability that no trajectory hallucinates within horizon
        \(N\) is
        \[
        \mathbb{P}\bigl(\forall i,\ \tau^{(i)} > N\bigr)
         = (1-p_N)^K.
        \]
  \item The probability that at least one trajectory hallucinates
        within horizon \(N\) is
        \[
        1 - (1-p_N)^K.
        \]
\end{enumerate}
\end{proposition}

\begin{proof}
By independence,
\[
\mathbb{P}(\tau^{(i)} > N) = 1 - p_N
\]
for each \(i\), so
\[
\mathbb{P}\bigl(\forall i,\ \tau^{(i)} > N\bigr)
 =
\prod_{i=1}^K (1-p_N)
 =
(1-p_N)^K.
\]
The complementary event has probability \(1-(1-p_N)^K\).
\end{proof}

These results show how inevitability manifests at finite horizons: for
any fixed horizon, hallucination remains probabilistic (\(p_N<1\) in
general), but as the horizon length or number of runs increases, the
probability that hallucination has occurred approaches~1.

\medskip

Future work lifts these individual dynamics to populations of agents and
explicit consensus mechanisms.
% ------------------------------------------------------------------
% New material after Proposition B.38 in Appendix B
% ------------------------------------------------------------------

\subsubsection*{B.16 Hazard Functions and Safe Horizons}

We now refine finite-horizon analysis by introducing per-step hazard and a derived notion of
``safe horizon'' for a given stochastic FRFP pipeline.

\begin{definition}[Per-step hazard]\label{def:hazard}
Let $\tau : \Omega \to \mathbb{N} \cup \{\infty\}$ be the hallucination time of Definition~B.31,
and let $P$ be a stochastic FRFP pipeline with associated distribution of $\tau$.
For each $n \in \mathbb{N}$ with $P(\tau \ge n) > 0$ we define the (discrete) hazard at step $n$ by
\[
  h_n := P(\tau = n \mid \tau \ge n)
       = \frac{P(\tau = n)}{P(\tau \ge n)}.
\]
If $P(\tau \ge n) = 0$ we set $h_n := 0$.
\end{definition}

Intuitively, $h_n$ is the conditional probability that a run first hallucinates at step $n$,
given that it has not hallucinated earlier. The sequence $(h_n)_{n \ge 0}$ is a pipeline-level
summary of how risk is distributed along the trajectory.

\begin{definition}[Finite-horizon survival and hazard]\label{def:survival}
For $N \in \mathbb{N}$ let
\[
  S_N := P(\tau > N) = 1 - p_N
\]
be the survival probability at horizon $N$, where $p_N = P(\tau \le N)$ is the finite-horizon
hallucination probability from Definition~B.36. Whenever $P(\tau \ge n) > 0$ for all $n$,
the survival probabilities admit the product representation
\[
  S_N = \prod_{n=0}^{N} (1 - h_n),
\]
so that the finite-horizon risk $p_N = 1 - S_N$ can be understood as the cumulative effect
of per-step hazards.
\end{definition}

This representation makes explicit the trade-off between (i) the number of steps in a pipeline
and (ii) the hazard profile $(h_n)$ induced by its design (e.g. decomposition, tool calls,
or repeated self-checks).

\begin{definition}[Safe horizon at tolerance $\varepsilon$]\label{def:safe-horizon}
Fix a tolerance level $\varepsilon \in (0,1)$. The $\varepsilon$-safe horizon of a stochastic
FRFP pipeline $P$ is defined as
\[
  H_P(\varepsilon)
  := \sup\{\,N \in \mathbb{N} : p_N \le \varepsilon\,\}
  \in \mathbb{N} \cup \{\infty\}.
\]
When $H_P(\varepsilon)$ is finite, any horizon $N \le H_P(\varepsilon)$ satisfies
$P(\tau \le N) \le \varepsilon$, i.e. hallucination within that horizon is at most
$\varepsilon$-likely.
\end{definition}

\begin{remark}[Design interpretation]
The safe-horizon functional $H_P(\varepsilon)$ isolates a pipeline-level design parameter:
for a given stochastic implementation $P$ and admissible risk budget $\varepsilon$, it
characterizes how long a single run can safely be extended before hallucination becomes
probable. Changes to decomposition, step-ordering, or intermediate checks can therefore
be compared in terms of their effect on $(h_n)$ and $H_P(\varepsilon)$, even when the
underlying explicit pipeline is statically equivalent in $N \ltimes E$.
\end{remark}

Combining Definition~\ref{def:safe-horizon} with Proposition~B.38, we can also analyze
restart and parallelization policies: for example, running $K$ independent instances of a
pipeline up to horizon $N \le H_P(\varepsilon)$ yields hallucination probability
$1 - (1 - p_N)^K$, which can be bounded using the same tolerance parameter.

\subsubsection*{B.17 Dynamic Refinement of Pipelines}

We next introduce a notion of dynamic refinement between stochastic FRFP pipelines that
compares their hallucination-time distributions.

\begin{definition}[Dynamic refinement]\label{def:dynamic-refinement}
Let $P_1, P_2$ be stochastic FRFP pipelines on the same state space and admissible trajectory
set $\Omega_{\mathrm{allowed}}$, with induced hallucination times $\tau_1, \tau_2$ and
finite-horizon probabilities
\[
  p_N^{(i)} := P_i(\tau_i \le N), \quad N \in \mathbb{N}, \; i \in \{1,2\}.
\]
We say that $P_2$ is a \emph{dynamic refinement} (or \emph{dynamically safer variant}) of $P_1$
if for all $N \in \mathbb{N}$ we have
\[
  p_N^{(2)} \;\le\; p_N^{(1)}.
\]
Equivalently, $\tau_2$ stochastically dominates $\tau_1$ in the usual sense.
\end{definition}

\begin{remark}[Static vs dynamic equivalence]
Two pipelines may be statically equivalent in the sense of Appendix~A (same normal form in
$N \ltimes E$ and same tacit semantics $J{-}K$), yet dynamically inequivalent: they can induce
different hazard profiles and hence different hallucination-time distributions.
Dynamic refinement isolates this difference at the level of trajectories, not static syntax:
$P_2$ improves on $P_1$ precisely when it never increases finite-horizon hallucination risk.
\end{remark}

\begin{remark}[Limitations]
Dynamic refinement is a partial order on stochastic implementations, not an absolute
guarantee of safety: if the baseline pipeline $P_1$ satisfies the inevitability conditions
of Theorems~B.11, B.25, and B.29, then so does any refinement $P_2$. Refinements can
delay hallucination and reduce risk over practically relevant horizons, but they cannot
eliminate hallucination altogether without violating the structural assumptions of this
appendix.
\end{remark}
% ------------------------------------------------------------------
% New subsection at the end of Appendix B
% ------------------------------------------------------------------

\subsubsection*{B.18 Structural implications of pipeline dynamics}

The move from static FRFP pipelines to the dynamic setting of this appendix
turns time and iteration into first-class objects. Beyond the inevitability
results of Theorems~B.11, B.25, and B.29, this buys several structural
notions that will be used in later appendices and case studies.

\paragraph{(1) Hazard profiles and safe horizons.}
Definition~B.31 endows each stochastic FRFP implementation with a
hallucination time $\tau$, and Definitions~B.36 and~\ref{def:hazard} refine
this to a finite-horizon risk curve $(p_N)_{N \ge 0}$ and a discrete hazard
sequence $(h_n)_{n \ge 0}$. The safe-horizon functional
$H_P(\varepsilon)$ in Definition~\ref{def:safe-horizon} abstracts how long a
given pipeline can be run, under a fixed stochastic implementation $P$ and
risk budget $\varepsilon$, before hallucination becomes likely. Thus even
when static semantics and correctness are fixed, pipeline dynamics exposes a
new design degree of freedom: the shape of the hazard profile along a run.

\paragraph{(2) Time--decomposition trade-offs.}
In the static theory of Appendices~A and~B.1--B.3, different decompositions
of a task into explicit steps are equivalent whenever they normalize to the
same morphism in $N \ltimes E$ and share the same tacit semantics $J{-}K$.
Once trajectories are considered, this equivalence is refined: more
fine-grained decompositions typically reduce the requirement demanded at
individual steps, but also increase the length of the trajectory and the
opportunity for degradation to act. Pipeline dynamics therefore makes precise
a basic design tension: reshaping a task into ``more steps'' can both reduce
per-step hazard and increase cumulative risk, and only the dynamic
quantities $(h_n)$ and $(p_N)$ decide which effect dominates in a concrete
setting.

\paragraph{(3) Dynamic refinement of implementations.}
Definition~\ref{def:dynamic-refinement} introduces a partial order on
stochastic implementations of the \emph{same} static pipeline: $P_2$
dynamically refines $P_1$ when it never increases finite-horizon risk.
This separates:
\begin{itemize}
  \item static equivalence in $N \ltimes E$ (same normal form, same tacit
        semantics), from
  \item dynamic safety properties (how quickly hallucination emerges under a
        particular scheduling, tool use pattern, or feedback coupling).
\end{itemize}
Subsequent appendices will use this distinction to talk about pipeline-level
improvements that do not change what a system computes in principle, but do
change how its risk unfolds in time.

\paragraph{(4) Session limits, restart policies, and composition.}
Because hallucination is represented as a stopping time and finite-horizon
probabilities are explicit, Appendix~B supports principled discussion of
session limits and restarts. Proposition~B.38 already shows how repeated
runs over a fixed horizon aggregate risk via $1 - (1 - p_N)^K$. In more
complex settings, this allows single-pipeline guarantees (bounds on $p_N$
or $H_P(\varepsilon)$) to be composed into system-level guarantees for
restart, parallelization, or human-in-the-loop oversight policies, much as
Appendix~A.10 composes agent-level inevitability into collective inevitability.

\paragraph{(5) Localizing and targeting failure modes.}
Finally, working with trajectories rather than only normal forms means that
hallucination can be localized to segments of a run: particular regions of
the state space, update rules, or phases of the pipeline where hazard is
concentrated. This enables targeted interventions (for example, inserting
human checks or additional tools only before entering high-hazard regions)
that would be invisible in a purely static semantics.

Taken together, these structures show that pipeline dynamics is not only a
way to restate hallucination inevitability in stochastic language. It also
provides a vocabulary---hazard, safe horizons, dynamic refinement, and
time--decomposition trade-offs---for comparing and improving FRFP pipelines
\emph{without} altering their underlying explicit semantics in $N \ltimes E$.

\subsection{Collective Epistemic Dynamics} % old C
\label{app:collective}
% populations, communication graphs, collective τ, consensus functors, etc.

%\subsection{Collective Epistemic Dynamics)}

\subsubsection{ Purpose and Scope}

This appendix extends  Epistemic dynamics to \emph{collective} epistemic
systems. Rather than a single agent executing a single FRFP pipeline, we
consider populations of agents who:

\begin{itemize}
  \item share an explicit artifact space $E$,
  \item each carry their own tacit state and degradation dynamics,
  \item exchange only explicit artifacts (no tacit information flow),
  \item may be coupled by explicit consensus mechanisms and institutions.
\end{itemize}

The objective is to characterize:

\begin{enumerate}
  \item how to view a population as a single large FRFP system,
  \item how stochastic collective evolution can be modeled as a Markov process
        on this system,
  \item how hallucination and other failures lift from individuals to
        populations, and
  \item which collective design changes \emph{cannot}, even in principle,
        eliminate hallucination.
\end{enumerate}

This section provides the bridge between individual pipelines (Sections A.2–A.8),
impossibility results (Appendix A.11), and institutional design (Appendix B and Appendix C).

\subsubsection{Epistemic Agents}

\begin{definition}[Epistemic agent]
An \emph{epistemic agent} is a tuple
\[
  A_i = (T_i, \,\preceq_i, \,\delta_i, \,T^e_i),
\]
where:
\begin{itemize}
  \item $T_i$ is a tacit context space equipped with a preorder
        $\preceq_i$ encoding ``at least as grounded as'';
  \item $\delta_i : T_i \to T_i$ is a (tacit) context degradation operator;
  \item $T^e_i := T_i \times \mathbb{R}_{\ge 0}$ is the extended tacit state
        space, pairing grounding $t_i \in T_i$ with tacit credence
        $c_i \in \mathbb{R}_{\ge 0}$.
\end{itemize}
\end{definition}

Intuitively, $A_i$ captures a single human (or human-like) reasoner as modeled
in Epistemic Dynamics: grounding lives in $T_i$, weakens under $\delta_i$,
and credence lives in the additional component of $T^e_i$.

\subsubsection{Epistemic Populations}

\begin{definition}[Epistemic population and category \textbf{Pop}]
An \emph{epistemic population} is a finite or countable family
$P = (A_i)_{i \in I}$ of epistemic agents indexed by a set $I$.
The category $\mathbf{Pop}$ of epistemic populations is defined by:
\begin{itemize}
  \item objects: epistemic populations $P = (A_i)_{i \in I}$;
  \item morphisms $F : P \to P'$: pairs $(f,(F_i)_{i\in I})$ where
        $f : I \to I'$ is a function on agent indices and, for each $i$,
        \[
          F_i : T_i \to T'_{f(i)}
        \]
        is monotone with respect to $\preceq_i,\preceq'_{f(i)}$ and
        commutes with degradation and credence:
        \[
          F_i \circ \delta_i = \delta'_{f(i)} \circ F_i.
        \]
\end{itemize}
\end{definition}

Morphisms in $\mathbf{Pop}$ model re-indexing or restructuring of populations:
embedding agents into larger organizations, replacing agents, or transferring
tacit states along institutional lines.

\subsubsection{Shared Explicit Reasoning Category}

Let $\mathcal{C} := N \ltimes E$ be the reasoning category of Appendix~A,
where $E$ denotes explicit artifacts and $N$ navigation moves in explicit
space. For a fixed index set $I$, we form the product category
\[
  \mathcal{C}^I,
\]
whose objects are $I$-indexed families $X = (X_i)_{i \in I}$ of explicit
states and whose morphisms are families of explicit updates. Intuitively,
$\mathcal{C}^I$ is the ``big explicit state space'' of the whole population:
it records, for each agent, which task and explicit artifacts they are
currently working with.

We assume a single shared explicit artifact space $E$; each agent $i$ is
equipped with a representational backflow map
\[
  \mathrm{RB}_i : E \to T_i,
\]
as in Appendix~A. Explicit artifacts are globally visible; tacit grounding
remains agent-local.

\subsubsection{Communication Structure}

\begin{definition}[Communication graph]
A \emph{communication structure} on population $P = (A_i)_{i \in I}$ is a
directed graph
\[
  G = (I,\;\to),
\]
where $i \to j$ indicates that explicit artifacts produced by $i$ may be
consumed by $j$ (for example, by reading a paper, listening to a briefing, or
reviewing a log). No tacit information is transmitted across $G$:
only explicit artifacts traverse edges, and representational backflow
$\mathrm{RB}_j$ determines how they affect $T_j$.
\end{definition}

\subsubsection{Collective Dynamics as a Markov System}

A \emph{collective execution} in earlier drafts was modeled as a bare
sequence of events $(i_n,e_n) \in I \times E$ indicating ``agent $i_n$
produces artifact $e_n$''. We now package such executions into a stochastic
dynamical system over the large explicit category $\mathcal{C}^I$.

\begin{definition}[Collective state]
A collective explicit state is an object $X \in \mathcal{C}^I$ together with
a population $P = (A_i)_{i\in I}$ and communication graph $G=(I,\to)$.
We write $(P,G;X)$ for such a state.
\end{definition}

\begin{definition}[Population semantics]
A \emph{population semantics functor} is a functor
\[
  \overline{J} : \mathcal{C}^I \to \prod_{i\in I} T^e_i
\]
assigning to each collective explicit state $X$ the tuple of extended tacit
states $\bigl((t_i,c_i)\bigr)_{i\in I}$ obtained when agents in $P$ consume
the artifacts in $X$ (subject to $G$) and update using the operators of
Appendix~B (degradation, confidence inflation, requirement escalation, and
feedback-coupled degradation).
\end{definition}

\begin{definition}[Markovian collective dynamics]
A \emph{Markovian collective dynamics} on $(P,G)$ is a transition kernel
\[
  M : \mathrm{Obj}(\mathcal{C}^I) \to \mathcal{P}\bigl(\mathrm{Obj}(\mathcal{C}^I)\bigr)
\]
assigning to each collective state $X$ a probability measure $M(X,\cdot)$
on next states $Y$, such that each transition can be factored as:
\begin{enumerate}
  \item select an acting agent $i_n \in I$ and explicit move in $\mathcal{C}$;
  \item update $X$ to $Y$ accordingly (producing a new artifact $e_n$ in $E$);
  \item for each $j$ with $i_n \to j$, update $(t_j,c_j)$ via the
        feedback-coupled dynamics of Appendix~B using $e_n$;
  \item apply baseline degradation for all agents.
\end{enumerate}
Equivalently, one may view $M$ as a functor
\[
  M : \mathcal{C}^I \longrightarrow \mathbf{Kl}(D),
\]
where $\mathbf{Kl}(D)$ is the Kleisli category of the discrete distribution
monad $D$, but we will only use the induced transition kernel on objects.
\end{definition}

A \emph{collective execution} is then a random path
\[
  X_0, X_1, X_2, \dots
\]
in $\mathcal{C}^I$ sampled from $M$ starting at some initial state $X_0$,
with tacit trajectories $\overline{J}(X_n) = \bigl((t_{i,n},c_{i,n})\bigr)_{i\in I}$.

\subsubsection{Collective Hallucination and Stopping Times}

We now lift the notion of hallucination to the population level.

\begin{definition}[Collective hallucination]
Given a population $P=(A_i)_{i\in I}$, a communication graph $G$, and a
collective execution $(X_n)_{n\ge 0}$ with associated tacit states
$(t_{i,n},c_{i,n})$, a \emph{collective hallucination} occurs at step $n$
if there exists an explicit artifact $e \in E$ and a nontrivial subset
$J \subseteq I$ such that:
\begin{itemize}
  \item each $j \in J$ has accepted or propagated $e$ (in the sense of
        Appendix~B: $e$ enters their explicit pipeline), and
  \item for each $j \in J$,
        \[
          t_{j,n} \prec_j \mathrm{Req}_j(e,c_{j,n}),
        \]
        i.e., the available tacit grounding is strictly below the requirement
        demanded by $e$ in context $c_{j,n}$.
\end{itemize}
\end{definition}

Let $\Omega$ be the space of all collective paths $(X_n)_{n\ge 0}$ induced
by $M$ from $X_0$, and let $P_M$ be the corresponding path measure.

\begin{definition}[Collective hallucination time]
The \emph{collective hallucination time} $\tau^{\mathrm{coll}} :
\Omega \to \mathbb{N} \cup \{\infty\}$ is the stopping time
\[
  \tau^{\mathrm{coll}}(\omega)
  := \inf\{\,n \ge 0 : \omega \text{ exhibits a collective hallucination at step } n\,\},
\]
with the convention that $\inf \varnothing = \infty$.
\end{definition}

Under the Epistemic Dynamics assumptions (finite tacit capacity, non-invertible
degradation, confidence inflation, requirement escalation, and
feedback-coupled degradation), individual agents exhibit almost-sure
hallucination along sufficiently long admissible runs (Appendix~B).
The purpose of this appendix is to show how such individual inevitability
lifts to $\tau^{\mathrm{coll}}$ at the population level.

\subsubsection{Consensus as an Explicit Functor}

Consensus mechanisms act purely in explicit space.

\begin{definition}[Consensus functor]
A \emph{consensus mechanism} on the explicit category $\mathcal{C}$ is a
functor
\[
  C : \mathcal{C}^n \to \mathcal{C}
\]
for some finite arity $n \ge 2$, or more simply
\[
  C : E^n \to E
\]
on the underlying artifact space, modeling averaging, voting, or more
complicated aggregation pipelines. Consensus functors operate entirely in
explicit space: they take explicit artifacts as inputs and produce explicit
artifacts as outputs, without access to tacit states or requirements.
\end{definition}

\subsubsection{Functorial Consensus Layer}

Given a population $P=(A_i)_{i\in I}$ and communication structure $G$,
a consensus mechanism $C$ induces an explicit functor
\[
  C_P : E^n \to E
\]
that may be interleaved with ordinary explicit computation in
$\mathcal{C}^I$. We refer to such composites as the \emph{consensus layer}.
No consensus functor has access to the population semantics
$\overline{J}$ or to individual requirement functions $\mathrm{Req}_i$.

\subsubsection{Tacit-Stability of Consensus}

\begin{definition}[Tacit-stable consensus]
A consensus functor $C$ is \emph{tacit-stable} if, for every agent $i$ and
any inputs $e_1,\dots,e_n \in E$, whenever $e_1,\dots,e_n$ are all
individually below agent $i$'s requirement,
\[
  t_i \prec_i \mathrm{Req}_i(e_k,c_i)
  \quad \forall k \in \{1,\dots,n\},
\]
then the aggregate output $C(e_1,\dots,e_n)$ is also below requirement:
\[
  t_i \prec_i \mathrm{Req}_i\bigl(C(e_1,\dots,e_n),c_i\bigr).
\]
\end{definition}

\begin{theorem}[Consensus stability]
Every consensus functor that acts purely in explicit space is tacit-stable.
\end{theorem}

\begin{proof}[Proof sketch]
Consensus functors are compositions of explicit morphisms in $\mathcal{C}$.
By Explicit--Tacit Separation (ETS), explicit morphisms cannot increase
tacit grounding: they do not alter $t_i$, only what artifacts are presented
to representational backflow. If an agent's tacit state $t_i$ is below the
requirement for each input artifact, it remains below the requirement for
any artifact obtained by explicit-only transformation of those inputs. Thus
consensus cannot manufacture sufficient grounding from insufficient inputs.
\end{proof}

Tacit-stability says that explicit consensus can neither repair nor upgrade
tacit adequacy: it preserves ``being under-grounded'' whenever that is the
case for all inputs.

\subsubsection{Limits of Automated Consensus}

\begin{theorem}[No fully automated consensus elimination]
No fully automated consensus process acting purely in explicit space can
eliminate collective hallucination across all admissible populations and
collective dynamics.
\end{theorem}

\begin{proof}[Proof sketch]
By tacit-stability, consensus cannot transform artifacts that are
under-grounded for all relevant agents into artifacts that are adequately
grounded. Under the Epistemic Dynamics degradation assumptions, individual hallucination
is inevitable along sufficiently long runs. Since consensus mechanisms have no
access to tacit state, they cannot detect or remove these failures at the
collective level; at best they can smooth, average, or denoise explicit
behavior. Any claim of global hallucination-elimination by explicit-only
consensus would contradict the Epistemic Dynamics inevitability results when applied
agent-wise.
\end{proof}

Automation may reduce variance or noise in explicit artifacts, but it cannot
introduce new tacit grounding or globally halt its degradation.

\subsubsection{Collective Inevitability}

\begin{theorem}[Collective hallucination inevitability]
Consider a population $P=(A_i)_{i\in I}$ of epistemic agents with:
\begin{itemize}
  \item finite or countable index set $I$,
  \item finite tacit capacity for each $A_i$,
  \item non-invertible degradation $\delta_i$,
  \item confidence inflation and requirement escalation as in Appendix~B.
\end{itemize}
Let $M$ be any Markovian collective dynamics on $\mathcal{C}^I$ that
respects these per-agent dynamics and produces infinitely many
surface-plausible explicit artifacts along almost every path.
Then the collective hallucination time $\tau^{\mathrm{coll}}$ is almost
surely finite:
\[
  P_M\bigl(\tau^{\mathrm{coll}} < \infty\bigr) = 1.
\]
Moreover, this property is invariant under the insertion of explicit-only
consensus functors: composing $M$ with any finite consensus layer that acts
solely in explicit space does not change the almost-sure finiteness of
$\tau^{\mathrm{coll}}$.
\end{theorem}

\begin{proof}[Proof sketch]
By Appendix~B, each agent $A_i$ exhibits almost-sure hallucination along
sufficiently long admissible runs of any FRFP pipeline that continues to
produce surface-plausible artifacts. The collective Markov dynamics $M$
induces, for each $i$, a marginal stochastic process on $T^e_i$ with the
same qualitative properties. Since the population is finite or countable,
and collective execution produces infinitely many plausible artifacts with
nonzero probability of being consumed by each agent (via $G$), the event
that at least one agent hallucinates in finite time has probability~$1$.
The definition of collective hallucination requires only that a nontrivial
subset of agents jointly accept an under-grounded artifact, which occurs
almost surely given enough time and interaction. Tacit-stability of
consensus implies that explicit aggregation cannot remove such events once
they arise; it can at best transform one hallucinated artifact into another.
\end{proof}

\subsubsection{Structural Implications of Collective Dynamics}

The population-level formalism above does more than restate
inevitability in a collective setting. It yields several structural notions
that are used in later sections.

\paragraph{(1) Equivalence classes of collective designs.}
Viewing populations as objects in $\mathbf{Pop}$ and their joint explicit
reasoning as trajectories in $\mathcal{C}^I$, we can treat different UIs,
workflows, and consensus wrappers as endofunctors on $\mathcal{C}^I$ that
preserve the population semantics $\overline{J}$. Designs that differ only
by such explicit refactorings are \emph{FRFP-equivalent}: they have the same
tacit semantics and the same hallucination-time distribution, even if their
surface appearance differs. This provides a criterion for when two collective
protocols are ``the same thing in principle''.

\paragraph{(2) Refinement order on collective systems.}
Given two Markovian collective dynamics $M_1,M_2$ on the same population,
with corresponding collective hallucination times
$\tau^{\mathrm{coll}}_1,\tau^{\mathrm{coll}}_2$, we may compare their
finite-horizon risks
\[
  p_N^{(k)} := P_{M_k}(\tau^{\mathrm{coll}}_k \le N), \quad k\in\{1,2\}.
\]
We say that $M_2$ is a \emph{dynamic refinement} of $M_1$ if
$p_N^{(2)} \le p_N^{(1)}$ for all $N$. This yields a partial order on
collective implementations: some architectures are strictly safer than
others, but purely explicit-only transformations that preserve semantics
cannot improve this order.

\paragraph{(3) Compositional risk analysis.}
Because the collective dynamics lives in $\mathcal{C}^I$ and is built from
per-agent Epistemic Dynamics, we can study how subsystem risks compose.
For example, a scientific subsystem, an engineering subsystem, and a policy
subsystem can each be modeled as a subpopulation with its own hazard profile;
the global population then inherits a collective hallucination time equal to
the minimum of the subsystem times. This provides a formal basis for
``defense in depth'' arguments.

\paragraph{(4) Limits of monitoring and metrics.}
Collective executions produce rich explicit traces: logs, dashboards,
metrics, and aggregated indicators. In this formalism, these traces live in
$\mathcal{C}^I$ or in functorial images thereof. The population semantics
$\overline{J}$ and individual requirements $\mathrm{Req}_i$, however, remain
tacit. Many distinct tacit configurations can induce the same distribution
over explicit traces. Hence correctness and hallucination risk are
\emph{non-identifiable from logs alone}: explicit-only monitoring cannot, in
general, determine how well grounded a population's beliefs really are.

\paragraph{(5) Scaling and phase transitions.}
By parametrizing collective systems by population size $|I|$, communication
structure $G$, and per-agent degradation properties, Appendix~A.10 provides a
language for asking how epistemic risk scales with size and complexity.
Nothing in the formalism guarantees that larger or more connected populations
are safer; on the contrary, the inevitability theorem applies uniformly once
the Epistemic Dynamics assumptions hold. This sets the stage for Appendix~A.12, which
studies which institutional forms can slow or bound such risks across
populations and time.

Taken together, these structures show that Collective Epistemic Dynamics is not merely a
corollary of individual hallucination: it is a genuinely new level of
description in which equivalence, refinement, observability, and scaling of
\emph{collective} epistemic systems can be stated and analyzed formally.

\subsection{Explicit-Only Impossibility and Non-Automation} % old D
\label{app:impossibility}
% explicit-only detector/oracle impossibility, tacit dependence, etc.
% Appendix D — Impossibility & Limits (Phase-5)
% Schema–corollary organization with (i) admissible runs as executions in \NltimesE
% under Phase-0 constraints and (ii) explicit-identity options (syntactic vs normal-form).
% Depends on Appendix A (static axioms) and Appendices B–C (dynamics + collectives).

%\subsection{Impossibility \& Limits}

\subsubsection{Purpose and Scope}
This section collects structural \emph{no-go theorems} that delimit what can be guaranteed within FRFP-like architectures under Explicit--Tacit Separation (ETS), Human-Exclusive Grounding (HEG), and Human-Exclusive Closure (HEC). The results are invariant to model internals, training data, or intent.

The central organizing idea is that many desirable guarantees are \emph{tacit-dependent}: whether they hold can vary with human tacit state even when explicit artifacts are held fixed. Under ETS and AE (AI-Explicit Restriction), tacit-dependent guarantees cannot be enforced or decided by explicit-only mechanisms.

\subsubsection{Standing Assumptions and Notation}

Throughout this appendix we fix the FRFP kernel and assume that all interaction and
epistemic constraints from the main text are in force.

\paragraph{Foundational layers.}
On the single-agent side we assume:
\begin{itemize}
  \item the interaction-layer constraints of Section~2 and Appendix~A.2.6 (admissible traces, no AI
        access to tacit inputs, separation of content-producing and correctness-bearing moves);
  \item the epistemic ontology and explicit/tacit separation from Section~3 and Appendix~A.2--A.3
        (ambient category $C_{\mathrm{full}}$, explicit and tacit subcategories $E$ and $T$, boundary
        functor $\RB : E \to T$, and the six primitives $RI, EC, ED, RB, TE, HFD$);
  \item the explicit workflow algebra and correctness layer from Sections~4--6 and
        Appendix~A.4--A.8 (the TDG-generated explicit category, navigation category $N$, Grothendieck
        semidirect product $N \ltimes E$, semantic functor $\sem{-} : N \ltimes E \to T$, and the
        correctness quotient $\gamma : \Obj(T) \to J$).
\end{itemize}

\paragraph{Dynamic and collective layers.}
In addition, we assume the epistemic dynamics and collective semantics developed in the later parts of
Appendix~A:
\begin{itemize}
  \item the tacit degradation operator $\delta : T \to T$, requirement map
        $\Req : E \times \mathbb{R}_{\ge 0} \to T$, and the induced notion of hallucination and
        finite safe horizons from Appendix~A.9;
  \item the population and consensus layer of Appendix~A.10, which models a population of agents
        with a shared explicit category, local tacit spaces, and a communication graph, and equips
        this structure with a Markovian collective dynamics.
\end{itemize}
For a fixed explicit pipeline $P$ and an admissible implementation of $P$ in this dynamic model,
Appendix~A.9 associates to each finite horizon $N$ a hallucination probability $p_N$ and a
hazard profile $(h_n)_{n \geq 0}$, and defines a safe-horizon functional
$H_P(\varepsilon)$ (see Definition~A.\ref{def:safe-horizon}) capturing the maximal time $N$ for which
the implementation keeps hallucination risk below a tolerance $\varepsilon$. On the population side,
Appendix~A.10 provides analogous notions of collective hallucination time and collective
safe horizons.

The explicit-only impossibility and institutional non-automation results in this appendix are
understood relative to this entire stack: they assume the interaction, epistemic, algebraic,
dynamic, and collective layers just described, and ask which kinds of detectors, consensus
mechanisms, and governance operators can be realised \emph{purely explicitly} within that fixed
architecture.


\subsubsection{Admissible Runs as Interaction Protocol Filtered Executions in \texorpdfstring{N$\ltimes$E}{N ltimes E}}
\begin{definition}[Executions in N$\ltimes$E]
Fix a valid pipeline term $P$ (Epistemic Space / Algebra well-typed). An \emph{execution trace} of $P$ is a finite or infinite sequence
\[
\rho:\quad x_0 \to x_1 \to x_2 \to \cdots
\]
where each $x_k$ is a configuration in N$\ltimes$E reachable from an initial configuration for $P$ by the combined reduction relation $\to$.
\end{definition}

\begin{definition}[Interaction Protocol admissibility predicate]
Let $\mathsf{Adm}_0(\rho)$ be a predicate on execution traces encoding Interaction Protocol constraints:
\begin{enumerate}
  \item \textbf{IL (Interaction Limitation):} $\rho$ is turn-based and respects externally imposed bounds (e.g. finite-length turns).
  \item \textbf{AR (Alignment Restriction):} explicit actions occur only in response to explicit instructions present in the trace.
  \item \textbf{MD (Mode Discipline):} no step in $\rho$ applies tacit operators inside the explicit phase; tacit transitions occur only after the boundary.
  \item \textbf{NTER/CSC:} no step reconstructs tacit evaluation from explicit traces or treats explicit computation as a correctness oracle.
\end{enumerate}
We treat $\mathsf{Adm}_0$ as an axiomatically specified filter on traces.
\end{definition}

\begin{definition}[Admissible runs]
An \emph{admissible run} of a pipeline $P$ is an execution trace $\rho$ in N$\ltimes$E such that $\mathsf{Adm}_0(\rho)$ holds.
We write $\mathrm{Runs}(P)$ for the set of admissible runs of $P$.
\end{definition}

\begin{remark} (Admissible runs and explicit-only results).
This definition pins all explicit-only impossibility and non-automation statements to a common
operational semantics: admissible runs are exactly those executions in \(N \ltimes E\) that
additionally satisfy the interaction constraints (IL, AR, MD, NTER/CSC). All explicit-only
no-go theorems in Appendix~A.11 are understood to quantify over this class of runs.

\end{remark}

\subsubsection{Two Notions of ``Same Explicit Artifact''}
No-go statements often quantify over pairs of runs that produce the ``same explicit artifact.'' There are two natural notions.

\begin{definition}[Syntactic explicit identity]
Two runs are \emph{syntactically explicit-identical} if they produce equal explicit artifacts as elements of the explicit object $E$ (i.e. equality in the representation domain).
\end{definition}

\begin{definition}[Semantic explicit identity via normal forms]
Assume the explicit reduction system $\to_{\Ecat}$ is terminating and locally confluent (Appendix~A axioms), so each explicit artifact has a unique normal form $\mathrm{nf}(e)$.
Two runs are \emph{normal-form explicit-identical} if their produced artifacts $e,e'$ satisfy
\[
\mathrm{nf}(e)=\mathrm{nf}(e').
\]
\end{definition}

\begin{remark}
Syntactic identity is appropriate when artifacts are literal strings/terms. Normal-form identity is appropriate when explicit equivalence is taken modulo rewriting (e.g. definitional equalities, normalization, canonical forms). All no-go results below hold under either identity notion; when needed we indicate which is assumed.
\end{remark}
\begin{remark}[Admissible runs and stochastic trajectories]
In the dynamic semantics of Appendix~B, the same pipeline $P$ gives rise
to a measurable space of trajectories $\Omega_{\mathrm{allowed}}$,
together with a family of probability measures capturing different
stochastic implementations. We implicitly identify each admissible run
$\rho$ of $P$ in the combined system $\mathsf{N} \ltimes \mathsf{E}$
with its corresponding trajectory
$\omega \in \Omega_{\mathrm{allowed}}$. In particular, the hallucination time $\tau(\omega)$, the finite-horizon probabilities $(p_N)$,
the hazard sequence $(h_n)$, and the safe-horizon function $H_P(\varepsilon)$
of Appendix~B are all computed over this common set of admissible runs.

\end{remark}

\subsubsection{Tacit Dependence via \texorpdfstring{$\gamma(\sem{P})$}{gamma(sem(P))}}
\begin{definition}[Tacit-dependent predicates]
Let $P$ be a fixed valid pipeline and let $\mathrm{Runs}(P)$ be its admissible runs.
A predicate $\mathcal{G}$ about an explicit artifact is \emph{tacit-dependent} if there exist two admissible runs $\rho_1,\rho_2\in\mathrm{Runs}(P)$ such that:
\begin{itemize}
  \item the runs produce the same explicit artifact (either syntactically or up to normal form, per Section~D.3), and
  \item the observable correctness outcomes differ:
  \[
    \gamma(\sem{P})\ \text{differs between}\ \rho_1\ \text{and}\ \rho_2.
  \]
\end{itemize}
\end{definition}

\begin{remark}
This definition formalizes ``depends on tacit context'' using only Epistemic Algebra semantics: $\sem{-}$ captures tacit transformation induced by the pipeline run, and $\gamma$ extracts the human-observable correctness class.
\end{remark}
\begin{remark}[Examples of tacit-dependent predicates]
\label{rem:tacit-dependent-examples}
The definition above applies not only to pointwise correctness
judgements but also to quantitative long-horizon properties. Typical
examples of tacit-dependent predicates $G$ include:
\begin{itemize}
  \item the event that a run hallucinates within a fixed horizon $N$,
        i.e.\ $G(\omega)$ is the predicate $[\tau(\omega) \le N]$;
  \item the statement that an implementation has $\varepsilon$-safe horizon at least $N$, i.e.
  $G$ is the predicate $[H_P(\varepsilon) \ge N]$ for some fixed $\varepsilon \in (0,1)$;

  \item in the collective setting of Appendix~A.10, the statement that the
        collective hallucination time exceeds a fixed horizon, i.e.\
        $G$ is the predicate $[\tau_{\mathrm{coll}} > N]$ for a given
        population and consensus dynamics.
\end{itemize}
In each case, there can exist admissible runs with the same explicit
artifacts (in the sense of Definition~B.\,66) but different values of
$\tau$, $\HP(\varepsilon)$ or $\tau_{\mathrm{coll}}$, so $G$ is
tacit-dependent in the sense of Definition~B.\,66.
\end{remark}

\subsubsection{Schema Theorem: Tacit-Dependent Guarantees Are Not Explicit-Implementable}
\begin{definition}[Explicit-only mechanism]
An \emph{explicit-only mechanism} is any procedure definable using explicit morphisms in $\Ecat$ (and finite control) whose inputs and outputs live entirely in explicit space. Concretely, one may model an explicit-only mechanism as a total function $M:E^k\to E^\ell$, or as a family of such functions closed under composition.
\end{definition}

\begin{theorem}[Tacit-dependent guarantees are not explicit-implementable]\label{thm:D-schema}
Assume ETS and AE. Let $\mathcal{G}$ be a tacit-dependent predicate about explicit artifacts (Section~D.4), where ``same explicit artifact'' is interpreted either syntactically or by normal form (Section~D.3). Then no explicit-only mechanism can enforce $\mathcal{G}$ for all admissible runs.
\end{theorem}

\begin{proof}[Proof sketch]
By tacit dependence, there exist admissible runs producing the same explicit artifact (under the chosen identity notion) but yielding different outcomes $\gamma(\sem{P})$. Any explicit-only mechanism sees only the explicit artifact(s) and therefore cannot act differently on these runs. Hence it cannot enforce a guarantee whose truth varies across admissible runs while explicit data are held fixed.
\end{proof}

\subsubsection{Corollaries (No-Go Family)}
Each of the following is an instantiation of Theorem~\ref{thm:D-schema}.

\subsubsubsection{No Perfect Explicit-Only Hallucination Detector}
\begin{definition}[Explicit-only monitor]
An \emph{explicit-only monitor} is a function $M:E\to\{\mathrm{ok},\mathrm{flag}\}$.
\end{definition}

\begin{definition}[Perfect detector]
$M$ is \emph{perfect} if, for every admissible run producing explicit artifact $e$ with tacit context $t$ at evaluation time,
\[
M(e)=\mathrm{ok}\Rightarrow t\not\prec\mathsf{Req}(e),\qquad
M(e)=\mathrm{flag}\Rightarrow t\prec\mathsf{Req}(e).
\]
\end{definition}

\begin{corollary}[No perfect explicit-only detector]\label{cor:D-no-perfect-detector}
Assume ETS and AE, and assume requirements are nontrivial and contexts vary across admissible runs. Then no explicit-only monitor is a perfect hallucination detector.
\end{corollary}

\begin{proof}[Proof sketch]
The predicate ``$e$ is a hallucination'' is tacit-dependent: the same explicit artifact (syntactically or up to normal form) can be within grounding in one admissible run and beyond grounding in another. Apply Theorem~\ref{thm:D-schema}.
\end{proof}

\begin{corollary}[No elimination of hallucination under degradation]
\label{cor:no-elimination-under-degradation}
Assume the Epistemic Dynamics conditions that yield almost-sure
hallucination inevitability: in particular, a finite-capacity tacit
state space, a monotone non-invertible degradation operator
$\delta \colon T \to T$, and recurring nontrivial requirements as in
Theorem~B.\,29. Let $\tau$ denote the hallucination time random variable
for a fixed pipeline $P$ and any admissible stochastic implementation
of $P$ whose trajectories lie in $\Omega_{\mathrm{allowed}}$.
Then there is no explicit-only policy which guarantees
\[
  \tau(\omega) \;=\; \infty
\]
for all admissible runs $\omega \in \Omega_{\mathrm{allowed}}$, nor even
$\mathbb{P}(\tau = \infty) = 1$ for all such implementations.
\end{corollary}

\begin{proof}
Under the assumptions of Appendix~B, Theorem~B.\,29 states that every
admissible stochastic implementation of $P$ has hallucination time
$\tau$ finite almost surely:
\[
  \mathbb{P}(\tau < \infty) \;=\; 1
\]
whenever its support is contained in $\Omega_{\mathrm{allowed}}$.
Consider the predicate
\[
  G(\omega) \;\equiv\; [\,\tau(\omega) = \infty\,]
\]
on admissible runs. By construction of the dynamic semantics, there
exist admissible runs with the same explicit artifacts but different
tacit evolutions along which $\tau$ takes different values; in
particular, $G$ is tacit-dependent in the sense of
Definition~B.\,66. If there were an explicit-only policy that guaranteed
$\tau = \infty$ almost surely for all admissible implementations, that
policy would enforce the tacit-dependent predicate $G$ by purely
explicit means, contradicting Theorem~B.\,69. Hence no explicit-only
policy can eliminate hallucination under the Epistemic Dynamics degradation
assumptions.
\end{proof}
\begin{corollary}[No explicit-only infinite safe horizon]\label{cor:no-infinite-safe-horizon}
Assume the Epistemic Dynamics conditions of Appendix~B that yield almost-sure hallucination
inevitability, and let $H_P(\varepsilon)$ be the $\varepsilon$-safe horizon of a
stochastic FRFP pipeline $P$ as in Definition~B.41. Then for any fixed
$\varepsilon \in (0,1)$ there is no explicit-only policy that guarantees
$H_P(\varepsilon) = \infty$ for all admissible stochastic implementations of $P$.
\end{corollary}

\begin{proof}
Under the Epistemic Dynamics assumptions, Theorem~B.29 implies that for any admissible
implementation we have $P(\tau < \infty) = 1$, so for every $\varepsilon \in (0,1)$
the set $\{N \in \mathbb{N} : p_N \le \varepsilon\}$ is bounded almost surely, i.e.\
$H_P(\varepsilon) < \infty$ along admissible runs. The predicate
$[H_P(\varepsilon) = \infty]$ is tacit-dependent (Remark~B.69), so an explicit-only
policy enforcing it for all admissible implementations would contradict
Theorem~B.69.
\end{proof}


\subsubsubsection{No Fully Automated Correctness Oracle}
\begin{definition}[Correctness oracle]
A \emph{correctness oracle} is a function $\Omega:E\to J$ (e.g. $J=\{\mathrm{accept},\mathrm{reject},\mathrm{revise}\}$) such that for every admissible run of a pipeline $P$ producing artifact $e$,
\[
\Omega(e)\;=\;\gamma(\sem{P}).
\]
\end{definition}

\begin{corollary}[No explicit correctness oracle]\label{cor:D-no-oracle}
Assume HEG and HEC. No explicit-only oracle $\Omega$ can reproduce $\gamma(\sem{P})$ for all admissible runs without violating the FRFP axioms.
\end{corollary}

\begin{proof}[Proof sketch]
Correctness outcomes $\gamma(\sem{P})$ are tacit-dependent by HEG/HEC: they are determined by tacit evaluation and closure, which can vary across admissible runs even when the explicit artifact is held fixed (under the chosen explicit identity). Apply Theorem~\ref{thm:D-schema}.
\end{proof}

\subsubsubsection{No Grounding Gain from Explicit Consensus}
Assume the functorial consensus model of Appendix~A.10.

\begin{corollary}[Consensus cannot create grounding]\label{cor:D-no-consensus-grounding}
Let $\mathcal{C}:\Ecat^n\to\Ecat$ be any consensus functor acting purely in explicit space. For any agent $i$ with tacit state $t_i$, if each input artifact $e_k$ exceeds grounding for $i$, then $\mathcal{C}(e_1,\dots,e_n)$ also exceeds grounding for $i$.
\end{corollary}

\begin{proof}[Proof sketch]
The predicate ``within grounding for agent $i$'' is tacit-dependent. Since $\mathcal{C}$ is explicit-only, it cannot enforce the guarantee that aggregation restores grounding. Apply Theorem~\ref{thm:D-schema}.
\end{proof}

\subsubsubsection{No Fully Automated Consensus Stability Guarantee}
\begin{corollary}[No automated consensus guarantee]\label{cor:D-no-auto-consensus}
Under ETS and HEG/HEC, no fully automated consensus pipeline can guarantee that its outputs are free of hallucination across all admissible runs and populations.
\end{corollary}
In the Markovian population semantics of Appendix~A.10, any choice of
consensus functors and explicit-only coordination rules determines a
collective dynamics $M$ on the product category together with a
collective hallucination time $\tau_{\mathrm{coll}}$. Under the Epistemic Dynamics
assumptions lifted to the population level, Theorem~B.\,58 shows that
$\tau_{\mathrm{coll}}$ is almost surely finite and, moreover, that its
almost-sure finiteness is invariant under insertion of additional
explicit-only consensus layers. A purported fully automated consensus
pipeline guaranteeing hallucination-free outputs for all admissible
populations would therefore enforce the tacit-dependent predicate
$[\tau_{\mathrm{coll}} = \infty]$ using only explicit mechanisms,
contradicting Theorem~B.\,69 together with Theorem~B.\,58.

\begin{proof}[Proof sketch]
A guarantee of hallucination-freedom is tacit-dependent and would yield a correctness oracle or perfect detector for consensus outputs. Apply Theorem~\ref{thm:D-schema} and Corollary~\ref{cor:D-no-perfect-detector}.
\end{proof}

\subsubsubsection{Catching Does Not Imply Eliminating}
\begin{corollary}[Detection does not imply elimination]\label{cor:D-detect-vs-elim}
Even if a subclass of failures is detectable by explicit invariants (e.g. syntactic checks, type checks, cross-consistency tests), this does not imply that hallucination (defined by tacit adequacy) is eliminable in general.
\end{corollary}

\begin{proof}[Proof sketch]
Explicit detectability concerns explicit predicates on $E$. Hallucination is tacit-dependent (Section~D.6.1). Eliminating hallucination would require enforcing a tacit-dependent guarantee by explicit-only means, contradicting Theorem~\ref{thm:D-schema}.
\end{proof}

\subsubsection{Interpretation and Design Consequences}
The no-go theorems do not claim that automation is useless; they claim that automation cannot replace tacit grounding without changing axioms.

Mitigations must therefore be expressed as:
\begin{itemize}
  \item restricting what is demanded of explicit space (reducing $\mathsf{Req}$),
  \item injecting fresh tacit evaluation (additional humans, replication),
  \item designing institutions that slow confidence inflation and feedback (Appendix~E).
\end{itemize}
Appendix~B refines these impossibility results with quantitative
structure. Each stochastic implementation of a pipeline induces a
hallucination time $\tau$, a finite-horizon risk curve
$(p_N)_{N \in \mathbb{N}}$, a hazard sequence $(h_n)_{n \in \mathbb{N}}$
and, when defined, a safe-horizon functional $\HP(\varepsilon)$. The
dynamic refinement order of Definition~B.\,43 then compares pipeline
implementations by their finite-horizon risks. Within the structural
limits above, design work can still aim to shape hazard profiles so that
risk is concentrated outside the operational horizon of interest and to
select implementations that dynamically refine a baseline, even though
eventual hallucination and the tacit-dependence of correctness remain
unavoidable.

Appendix~A.10 extends the picture to populations and collective dynamics.
A choice of consensus operators and coordination policies determines a
Markov process $M$ on the space of explicit population states and a
collective hallucination time $\tau_{\mathrm{coll}}$. The collective
inevitability and invariance results of Appendix~10  show
that purely explicit-only changes to consensus and logging cannot alter
the inevitability class of $\tau_{\mathrm{coll}}$. What they can do is
change the collective hazard profile and the distribution of
$\tau_{\mathrm{coll}}$, for example by producing a dynamics $M_2$ that
is a population-level dynamic refinement of some baseline $M_1$. In
this sense, long-horizon design remains meaningful: it can bend but not
break the structural bounds established in this appendix.

% End Appendix D

\subsection{Institutional Operators and Aggregation} % old E
\label{app:institutions-math}
% ============================================================
% Appendix E — Institutions and Design Constraints (Phase-6)
% ============================================================

%\subsection{Institutions and Design Constraints}

\subsubsection{Purpose and Scope}

Appendix A till now establishes that correctness in FRFP is structurally tacit,
non-automatable, and dynamically fragile under repeated interaction.
Appendix~A.11 formalizes no-go results showing that no single agent or pipeline
can eliminate hallucination, drift, or semantic collapse.

This appendix studies the \emph{institutional layer} required to stabilize
knowledge production over long time horizons.  Institutions are treated not
as social artifacts but as \emph{formal epistemic mechanisms} that operate on
collections of runs, agents, and tacit judgments.

The goal is to characterize which institutional constraints are
\emph{necessary}, which are \emph{insufficient}, and which are
\emph{structurally unavoidable} for scientific correctness.

\subsubsection{Epistemic Populations}
\label{sec:appendix-E1}

\begin{definition}[Epistemic agent, institutional view]\label{def:B82}
An \emph{epistemic agent} is an epistemic agent in the sense of Definition~B.46 of Appendix~C: a tuple
\[
  A_i = (T_i, \preceq_i, \delta_i, T^e_i),
\]
where
\begin{itemize}
  \item $T_i$ is a tacit context space equipped with a preorder $\preceq_i$ encoding ``at least as grounded as'';
  \item $\delta_i : T_i \to T_i$ is a (tacit) context degradation operator;
  \item $T^e_i := T_i \times \mathbb{R}_{\ge 0}$ is the extended tacit state space pairing a grounding $t_i \in T_i$ with a tacit credence $c_i \in \mathbb{R}_{\ge 0}$.
\end{itemize}
For each agent $A_i$ we write
\[
  \gamma_i : T_i \to J
\]
for the induced correctness quotient obtained from the Epistemic Algebra semantic structure (Appendix~A): $\gamma_i(t)$ is the observable correctness judgment when $A_i$ evaluates the tacit state $t$.

We will refer to
\[
  (T_i, \preceq_i, \delta_i, T^e_i, \gamma_i)
\]
as the \emph{institutional signature} of the agent $A_i$.
\end{definition}

\begin{definition}[Epistemic population and category \texorpdfstring{$\mathrm{Pop}$}{Pop}]\label{def:B83}
An \emph{epistemic population} is an object
\[
  P = (A_i)_{i \in I}
\]
of the category $\mathrm{Pop}$ defined in Definition~B.47 of Appendix~C: a finite or countable family of epistemic agents indexed by a set $I$.

Morphisms $F : P \to P'$ in $\mathrm{Pop}$ are pairs $(f, (F_i)_{i \in I})$, where $f : I \to I'$ is a function on indices and each
\[
  F_i : T_i \to T'_{f(i)}
\]
is monotone with respect to the tacit preorders and commutes with degradation and extended tacit state. In this appendix we fix an epistemic population
\[
  P = (A_i)_{i \in I}
\]
unless explicitly stated otherwise. Agents in $P$ may differ in tacit capacity, degradation dynamics, credence behavior, or evaluative norms; no homogeneity is assumed.
\end{definition}


\subsubsection{Runs, Replication, and Population Semantics}
\label{sec:appendix-E2}

Recall from Appendix~A that a well-typed FRFP pipeline $P$ induces a semantic value $\llbracket P \rrbracket$ in the tacit space and an observable correctness outcome $\gamma(\llbracket P \rrbracket)$. At the dynamic level, Epistemic Dynamics refines this to admissible runs in $N \ltimes E$, and Appendix~~A.11 identifies admissible runs with execution traces satisfying Interaction Protocol constraints. In the collective setting of Appendix~C, such runs are further coupled to population states via the population semantics functor.

\begin{definition}[Admissible run]\label{def:B84}
Fix a Epistemic Structure well-typed pipeline $P$ and an epistemic agent $A_i$ in the sense of Definition~\ref{def:B82}. An \emph{admissible run} of $P$ by $A_i$ is
\begin{itemize}
  \item an admissible execution trace $\rho \in \mathrm{Runs}(P) \subseteq N \ltimes E$ in the sense of Definition~B.61 (Appendix~D), and
  \item an initial extended tacit state $(t_{i,0}, c_{i,0}) \in T^e_i$ for $A_i$,
\end{itemize}
such that the evolution of $(t_{i,n}, c_{i,n})_{n \ge 0}$ along $\rho$ is governed by the Epistemic Dynamics update rules (degradation, confidence inflation, requirement escalation, and feedback-coupled degradation) and the Interaction Protocol constraints are respected at each step.
\end{definition}

\begin{definition}[Replication family]\label{def:B85}
Let $P$ be a fixed FRFP pipeline and let
\[
  P = (A_i)_{i \in I}
\]
be an epistemic population. A \emph{replication family} for $P$ in $P$ is a family
\[
  R(P) = \{\rho_i : i \in I_{\mathrm{rep}}\}
\]
indexed by a finite or countable subset $I_{\mathrm{rep}} \subseteq I$, where for each $i \in I_{\mathrm{rep}}$,
\begin{itemize}
  \item $\rho_i$ is an admissible run of $P$ by agent $A_i$ (Definition~\ref{def:B84}), and
  \item we write $(t_{i,*}, c_{i,*}) \in T^e_i$ for a chosen evaluation-time extended tacit state along $\rho_i$ (for example, the terminal state when the run halts or the state at a prescribed stopping time).
\end{itemize}
Intuitively, $R(P)$ records ``one execution of $P$ per agent'' under potentially different tacit initial conditions and stochastic realizations, but all subject to the same explicit pipeline and admissibility constraints.
\end{definition}

\begin{definition}[Population semantic profile]\label{def:B86}
Let $R(P) = \{\rho_i : i \in I_{\mathrm{rep}}\}$ be a replication family for $P$ in $P$, and let
\[
  (t_{i,*}, c_{i,*}) \in T^e_i
\]
be the evaluation-time extended tacit state associated to $\rho_i$ as above. The \emph{population semantic profile} of $P$ relative to $R(P)$ is the multiset
\[
  \mathcal{P}(P; R(P)) \;:=\; \bigl\{ \gamma_i(t_{i,*}) : i \in I_{\mathrm{rep}} \bigr\} \;\subseteq\; J.
\]
This multiset---not any single run---will be the basic carrier of institutional correctness at the population level.
\end{definition}


\subsubsection{Institutional Aggregation Operators}
\label{sec:appendix-E3}

We now formalize how institutions act on population-level correctness data rather than on individual runs or artifacts.

\begin{definition}[Institutional aggregator]\label{def:B87}
Let $J$ be the space of observable correctness judgments (e.g.\ $J = \{\mathsf{accept}, \mathsf{reject}, \mathsf{revise}\}$). Write $M(J)$ for the collection of finite or countable multisets of elements of $J$. An \emph{institutional aggregator} is a function
\[
  I : M(J) \to J \cup \{\bot\},
\]
where $\bot$ denotes suspended judgment. Typical examples include majority vote, unanimity, weighted consensus, veto rules, or deliberative deferral.

Given a fixed pipeline $P$ and epistemic population $P$, we may write $I_P$ for the particular aggregator that an institution applies to population semantic profiles $\mathcal{P}(P; R(P))$ arising from replication families $R(P)$ of $P$.
\end{definition}

\begin{definition}[Institutional outcome]\label{def:B88}
Let $P$ be a pipeline, $P$ an epistemic population, $R(P)$ a replication family, and $I$ an institutional aggregator. The \emph{institutional outcome} of $P$ relative to $R(P)$ under $I$ is
\[
  I_P\bigl(\mathcal{P}(P; R(P))\bigr) \;\in\; J \cup \{\bot\}.
\]
By construction, $I_P$ operates only on tacitly produced correctness judgments collected in the multiset $\mathcal{P}(P; R(P))$. No raw explicit artifact or execution trace is directly admissible as an input to $I_P$.
\end{definition}

\begin{definition}[Explicit-only institution]\label{def:B88prime}
An institutional aggregator family $(I_P)_P$ is \emph{explicit-only} if, for each pipeline $P$ and replication family $R(P)$, there exists an explicit-only mechanism
\[
  M_P : E^k \to J \cup \{\bot\}
\]
in the sense of Definition~B.70 of Appendix~D such that
\[
  I_P\bigl(\mathcal{P}(P; R(P))\bigr)
  \;=\;
  M_P(e_1, \dots, e_k),
\]
where $e_1, \dots, e_k \in E$ are the explicit artifacts produced along the admissible runs in $R(P)$, and $M_P$ has no access to any tacit states or tacit correctness outcomes. We will see in Theorem~B.93 that such explicit-only institutions cannot, in general, preserve correctness.
\end{definition}


\subsubsection*{E.4 Repeatability and Reproducibility}

\begin{definition}[Repeatability]
A pipeline $P$ is repeatable relative to $\mathbb{P}$ if all but finitely many
elements of $\Sem_{\mathbb{P}}(P)$ are equal.
\end{definition}

\begin{definition}[Reproducibility]
A pipeline $P$ is reproducible if for any sufficiently large epistemic
population $\mathbb{P}'$ compatible with $\mathbb{P}$, the institutional
outcome is invariant.
\end{definition}

Repeatability is a property of runs; reproducibility is a property of
institutions.

\subsubsection*{E.5 Institutional Drift and Decay}

Even if individual agents degrade (Appendix~B), institutions can buffer
against drift by averaging, vetoing, or deferring.

\begin{definition}[Institutional stability]
An institution $\mathcal{I}$ is stable if for any pipeline $P$ whose
population semantic profile converges, the institutional outcome converges.
\end{definition}

\begin{remark}
Stability does \emph{not} imply correctness—only resistance to oscillation.
\end{remark}

\subsubsection{Limits of Institutional Automation (Revised Statement)}
\label{sec:appendix-E6}

\begin{theorem}[Institutional non-automation]\label{thm:B93}
Assume ETS and AE, together with the stochastic and population semantics of Appendices~B and~C. Let $(I_P)_P$ be an institutional layer such that for every pipeline $P$ and every replication family $R(P)$ in any epistemic population, the institutional outcome $I_P(\mathcal{P}(P; R(P)))$ is computed by an explicit-only institution in the sense of Definition~\ref{def:B88prime}: there exists an explicit-only mechanism
\[
  M_P : E^k \to J \cup \{\bot\}
\]
whose inputs are only the explicit artifacts produced along the admissible runs in $R(P)$ and such that
\[
  I_P\bigl(\mathcal{P}(P; R(P))\bigr) \;=\; M_P(e_1, \dots, e_k).
\]
If, in addition, the institution is correctness-preserving in the sense that for every admissible pipeline $P$ and replication family $R(P)$,
\[
  I_P\bigl(\mathcal{P}(P; R(P))\bigr)
\]
agrees with the tacit correctness multiset $\{\gamma_i(t_{i,*}) : i \in I_{\mathrm{rep}}\}$ whenever this multiset is pointwise unanimous, then no such institution can exist.

Equivalently: there is no correctness-preserving institutional layer that depends solely on explicit artifacts or explicit-only summaries of those artifacts across all admissible pipelines and populations.
\end{theorem}

\begin{proof}[Proof sketch]
By Corollary~B.78 of Appendix~D, under HEG and HEC there is no explicit-only correctness oracle
\[
  \Omega : E \to J
\]
that reproduces $\gamma(\llbracket P \rrbracket)$ for all admissible runs of all pipelines. Any explicit-only institution satisfying the hypothesis would induce such an oracle by restricting to replication families of size one and composing
\[
  e \;\mapsto\; \{e\} \;\mapsto\; I_P\bigl(\mathcal{P}(P; R(P))\bigr) \;:=\; \Omega(e).
\]
Since correctness outcomes are tacit-dependent, this contradicts the no-go theorem for explicit-only mechanisms (Theorem~B.71). Therefore, no purely explicit-only institutional aggregator can preserve correctness across all admissible runs and populations.
\end{proof}


\subsubsection*{E.7 Peer Review as a Formal Operator}

\begin{definition}[Peer review operator]
Peer review is an institutional aggregator $\mathcal{I}_{\mathrm{PR}}$ whose
domain consists of independently generated tacit judgments by distinct
agents, possibly with structured disagreement.
\end{definition}

Peer review does not eliminate hallucination (Appendix~B) but reduces the
probability that hallucinated explicit artifacts pass institutional filters.

\subsubsection*{E.8 Consensus Is Not a Fixpoint}

\begin{theorem}[Consensus instability]
Even stable consensus does not imply semantic correctness.
\end{theorem}

\begin{proof}[Proof sketch]
Consensus is an aggregation over $\gamma(\sem{r})$ values. If the population
shares a correlated degradation mechanism, consensus may converge on an
incorrect judgment. This follows from the hallucination inevitability and
confidence inflation results of Appendix~B.
\end{proof}

\subsubsection*{E.9 Design Constraints for AI Systems}

Any AI system intended for epistemic use must therefore satisfy:

\begin{itemize}
\item explicit confinement (Epistemic Space);
\item exposure of outputs to multiple independent tacit evaluators;
\item support for suspension of judgment ($\bot$);
\item resistance to confidence inflation feedback at the institutional level.
\end{itemize}

These are not ethical desiderata but \emph{formal necessities}.

\subsubsection*{E.10 Summary}

Institutions are not optional add-ons but mathematically required structures
for long-horizon epistemic stability. They do not eliminate hallucination,
drift, or error; they bound their impact across populations and time.

FRFP therefore implies a strict separation of roles:

\begin{center}
\emph{Pipelines compute.  
Humans judge.  
Institutions stabilize.}
\end{center}

Any attempt to collapse these layers violates one of the no-go results proven
in earlier appendices.

% End Appendix E

%institutional aggregation over populations, formal definitions (without the conceptual prose)
% ============================================================
% Appendix F — Governance and Law (Phase-7)
% ============================================================

%\subsection{Governance, Law, and Epistemic Oversight}

\subsubsection{Position in the Architecture}

Previously we establish:

\begin{itemize}
\item individual pipeline correctness  ,
\item dynamic degradation and hallucination inevitability  ,
\item collective epistemic dynamics and population effects  ,
\item formal impossibility results  ,
\item institutional stabilization mechanisms .
\end{itemize}

Next we study \emph{governance and law} as higher-order structures that
operate \emph{on institutions themselves}, rather than on pipelines or agents.

\subsubsection{Governance as a Meta-Institution}
\label{sec:appendix-F1}

Let $J$ and $M(J)$ be as in Appendix~E, and let an \emph{institutional aggregator family} be a collection
\[
  (I_P)_P, \qquad I_P : M(J) \to J \cup \{\bot\},
\]
assigning to each FRFP pipeline $P$ an institutional aggregator $I_P$ in the sense of Definition~B.87.

\begin{definition}[Category of institutions]\label{def:F1-Inst}
We define the category $\Inst$ as follows.
\begin{itemize}
  \item Objects are institutional aggregator families $(I_P)_P$.
  \item A morphism $\Phi : (I_P)_P \to (I'_P)_P$ is a natural transformation whose components
  \[
    \Phi_P : M(J) \to M(J)
  \]
  are monotone with respect to refinement of multisets and preserve unanimity: whenever all elements of $S \in M(J)$ agree, so do the elements of $\Phi_P(S)$.
\end{itemize}
Intuitively, a morphism in $\Inst$ is a structure-preserving transformation between institutional layers that does not introduce disagreement where there was none.
\end{definition}

\begin{definition}[Governance operator]\label{def:F1-governance}
A \emph{governance operator} is an endofunctor
\[
  G : \Inst \to \Inst
\]
together with a choice of observable
\[
  \Omega : \mathsf{Pop} \to O
\]
from the category of epistemic populations (Appendix~C) to some space of observables $O$, such that for each time step $t$,
\[
  (I^{(t+1)}_P)_P = G(I^{(t)}_P)_P
\]
is allowed to depend only on the previous institutional layer $(I^{(t)}_P)_P$ and on the history of observables $(\Omega(P_s))_{s \le t}$ generated by the collective dynamics of Appendix~C.
\end{definition}


\subsubsection{Link to Collective Epistemic Dynamics}

Appendix~C models collective epistemic evolution as a dynamical system over
populations:
\[
(\mathbb{P}_t, \Sem_{\mathbb{P}_t})_{t \ge 0}.
\]

Governance operators are defined on this level.

\begin{definition}[Governance observable]
A governance observable is any functional
\[
\Omega : (\mathbb{P}, \Sem_{\mathbb{P}}) \longrightarrow \mathbb{R}
\]
that measures population-level epistemic properties (e.g.\ disagreement
rates, confidence inflation, drift velocity).
\end{definition}

\begin{definition}[Governed evolution]
A governed epistemic system evolves by:
\[
\mathcal{I}_{t+1} = \mathcal{G}(\mathcal{I}_t, \Omega_t),
\qquad
(\mathbb{P}_{t+1}, \Sem_{\mathbb{P}_{t+1}}) = \Phi(\mathcal{I}_{t+1}),
\]
where $\Phi$ is the collective dynamics functor from Appendix~C.
\end{definition}

Thus governance is \emph{feedback on institutions}, not on beliefs.

\subsubsection{Law as a Constraint on Governance}

Governance itself is subject to failure modes: capture, overfitting,
short-term optimization, or erosion of tacit authority.

\begin{definition}[Legal constraint]
A legal constraint is a predicate
\[
\mathcal{L}(\mathcal{I}) \in \{\text{admissible},\text{inadmissible}\}
\]
on institutional forms.
\end{definition}

\begin{definition}[Lawful governance]
A governance operator $\mathcal{G}$ is lawful if:
\[
\mathcal{L}(\mathcal{I}) = \text{admissible}
\;\Rightarrow\;
\mathcal{L}(\mathcal{G}(\mathcal{I})) = \text{admissible}.
\]
\end{definition}

Law therefore restricts the allowable transformations of institutions, not
the content of scientific claims.

\subsubsection{No Fully Automated Governance}

We now state the core limit result.
\begin{definition}[Explicit-only governance]\label{def:F-explicit-only-governance}
A governance operator
\[
  G : \Inst \to \Inst
\]
is \emph{explicit-only} if, for every time $t$ and institutional layer $(I^{(t)}_P)_P$, there exists a family of explicit-only mechanisms
\[
  G^{\mathrm{exp}}_{P,t} : E^{k_{P,t}} \to J \cup \{\bot\}
\]
in the sense of Appendix~D such that the updated institutional aggregators satisfy
\[
  I^{(t+1)}_P(S)
  \;=\;
  G^{\mathrm{exp}}_{P,t}\bigl(e^{(t)}_1, \dots, e^{(t)}_{k_{P,t}}\bigr)
\]
for every multiset $S \in M(J)$, where $e^{(t)}_1, \dots, e^{(t)}_{k_{P,t}} \in E$ are the explicit artifacts produced by the admissible runs of $P$ up to time $t$. In particular, $G$ has no access to any tacit states or tacit correctness outcomes, only to explicit execution traces and their explicit summaries.
\end{definition}
\begin{definition}[Explicit-only governance]\label{def:F-explicit-only-governance}
A governance operator
\[
  G : \Inst \to \Inst
\]
is \emph{explicit-only} if, for every time $t$ and institutional layer $(I^{(t)}_P)_P$, there exists a family of explicit-only mechanisms
\[
  G^{\mathrm{exp}}_{P,t} : E^{k_{P,t}} \to \bigl(M(J) \to J \cup \{\bot\}\bigr)
\]
in the sense of Appendix~D such that the updated institutional aggregators satisfy
\[
  I^{(t+1)}_P(S)
  \;=\;
  G^{\mathrm{exp}}_{P,t}\bigl(e^{(t)}_1, \dots, e^{(t)}_{k_{P,t}}\bigr)(S)
\]
for every multiset $S \in M(J)$, where $e^{(t)}_1, \dots, e^{(t)}_{k_{P,t}} \in E$ are the explicit artifacts produced by the admissible runs of $P$ up to time $t$. In particular, $G$ has no access to any tacit states or tacit correctness outcomes, only to explicit execution traces and their explicit summaries.
\end{definition}

\begin{theorem}[Non-automation of governance]
No governance operator $\mathcal{G}$ that depends solely on explicit artifacts
or explicitly computable observables can preserve epistemic correctness over
time.
\end{theorem}

\begin{proof}[Proof sketch]
By Appendix~B, explicit-only observables cannot detect tacit degradation.
By Appendix~D, correctness depends on $\gamma(\sem{r})$, which is tacit.
Any governance operator computable from explicit traces alone therefore
cannot reliably detect institutional epistemic failure modes.
\end{proof}

\subsubsection{Human Oversight as a Structural Requirement}

\begin{corollary}[Human oversight necessity]
Any lawful governance operator must incorporate human tacit judgment at the
meta-institutional level.
\end{corollary}

This does not imply that humans must make all decisions, but that \emph{some
decisions cannot be delegated}.

\subsubsection{Regulatory Scope}

Within FRFP, regulation can only legitimately target:

\begin{itemize}
\item admissible institutional forms,
\item governance feedback mechanisms,
\item transparency of explicit pipelines,
\item preservation of tacit authority.
\end{itemize}

It cannot:
\begin{itemize}
\item certify correctness automatically,
\item replace tacit evaluation with metrics,
\item eliminate hallucination in principle.
\end{itemize}

\subsubsection{Final Synthesis}

FRFP yields a stratified epistemic architecture:

\begin{center}
\begin{tabular}{c}
\textbf{Pipelines} — compute explicit artifacts \\
\textbf{Agents} — perform tacit judgment \\
\textbf{Institutions} — stabilize populations \\
\textbf{Governance} — adapt institutions \\
\textbf{Law} — constrain governance
\end{tabular}
\end{center}

Each layer is necessary; none subsumes the one above it.

\medskip

\noindent
\emph{There is no fully automated science, no fully automated governance of
science, and no lawful system that claims otherwise.}

% End Appendix F

\pagebreak
\section{FRFP Applications}
\label{app:applications}

\subsection{Disciplined Agents and Tooling} % old G
\label{app:disciplined-agents}
% definition, requirements, adapters, competence monitor, CMS integration
%\section{Disciplined agents and FRFP-compatible tooling}
%\label{app:disciplined-agents}

This appendix gives an informal account of \emph{disciplined agents} and supporting tooling in the sense used in the main text. The aim is not to introduce new mathematical structure, but to make concrete what it means, operationally, for AI systems and surrounding infrastructure to respect the FRFP boundary between explicit and tacit modes.

\subsection{Disciplined agents: informal definition}

Recall that FRFP separates:
\begin{itemize}
  \item an \emph{explicit} mode, consisting of machine-manipulable artefacts and pipelines (text, code, models, data, proofs, logs, plans, and their compositions); and
  \item a \emph{tacit} mode, consisting of human-exclusive evaluative states (understanding, background know-how, context-dependent judgments of adequacy, credibility, and acceptance).
\end{itemize}

\begin{definition}[Disciplined agent]
A \emph{disciplined agent} is an AI system, together with its immediate runtime wrapper, that:
\begin{enumerate}
  \item is strictly confined to the explicit side of FRFP (it manipulates only explicit artefacts);
  \item exposes its behaviour as explicit \emph{pipelines} rather than opaque one-shot outputs;
  \item is designed to integrate with, and defer to, human tacit oversight at clearly identified boundary points; and
  \item satisfies a small set of protocol-level constraints (explicit confinement, source--reason traceability, navigation integrity, and competence-aware behaviour).
\end{enumerate}
\end{definition}

Intuitively, a disciplined agent is not a ``governing mind'' or an autonomous decision-maker. It is a tool for constructing and transforming explicit structures in a way that keeps humans firmly in control of correctness and acceptance.

\subsection{Core requirements for disciplined agents}

We now spell out the core constraints that characterise disciplined agents. These are phrased operationally so that they can be checked against concrete systems.

\paragraph{(1) Explicit confinement.}
The agent must operate entirely in explicit mode:
\begin{itemize}
  \item All inputs are explicit artefacts (prompts, queries, specifications, code, data references, prior outputs).
  \item All outputs are explicit artefacts (drafts, summaries, plans, code, annotations, proposed transformations).
  \item The system \emph{does not} attempt to simulate or replace human tacit judgments such as final acceptance, ethical endorsement, or institutional approval.
\end{itemize}
In particular, disciplined agents do not emit binding ``decisions''; they emit explicit proposals that require human uptake at the boundary.

\paragraph{(2) Boundary awareness and deference.}
The agent must be designed to respect the FRFP boundary:
\begin{itemize}
  \item It treats human review, modification, or rejection of artefacts as the only source of correctness-bearing updates.
  \item It does not override, silently correct, or back-solve against recorded human judgments.
  \item It exposes natural \emph{handoff points} where humans are expected to inspect, accept, revise, or discard artefacts before they are used downstream.
\end{itemize}

\paragraph{(3) Source--reason traceability.}
The agent must support inspection of \emph{why} a given explicit artefact was produced:
\begin{itemize}
  \item Outputs should be accompanied, where feasible, by their immediate \emph{source} artefacts (citations, code paths, configuration) and a short rationale in explicit form.
  \item Transformations of artefacts (e.g., ``summarise'', ``refactor'', ``combine'') should be logged as steps in a pipeline rather than only as final states.
  \item It should be possible for a human reviewer to trace a given explicit artefact back along the pipeline to its inputs and major intermediate decisions.
\end{itemize}
This supports auditability and post-hoc diagnosis of failure modes, and matches the FRFP focus on whole-pipeline correctness.

\paragraph{(4) Navigation integrity.}
Disciplined agents must treat exploration and navigation as first-class:
\begin{itemize}
  \item The agent may support branching, backtracking, and ``what if'' exploration of explicit artefacts, but navigation steps should not silently alter prior artefacts or logs.
  \item Changes to context (e.g., switching between tasks, loading a different dataset, swapping model versions) should be explicit and recorded.
  \item The agent should avoid conflating logically distinct threads of work into a single undifferentiated interaction history.
\end{itemize}
The goal is to prevent subtle corruption of context---for example, accidental reuse of assumptions from one scenario in a different scenario without explicit acknowledgement.

\paragraph{(5) Competence awareness and safe failure.}
Disciplined agents must represent and communicate their own limitations:
\begin{itemize}
  \item For tasks outside their competence, agents should decline, escalate to a human, or request clarification rather than hallucinating confident-sounding content.
  \item They should provide machine-readable signals of task difficulty and uncertainty, enabling external \emph{competence monitors} to enforce conservative behaviour in high-stakes settings.
  \item Where applicable, they should be designed to degrade gracefully under load (e.g., by narrowing scope or switching to safer default behaviours) rather than silently increasing risk.
\end{itemize}

\subsubsection*{Explicit correctness engines as disciplined components}
A common disciplined-agent pattern is to couple a generative system with an explicit correctness engine:
for example, a large language model that proposes code or proofs, paired with a compiler, test harness,
or interactive theorem prover that checks those artefacts against a fixed specification. In FRFP, such
engines live entirely in the explicit category $E$ and implement EC/ED steps: they take explicit artefacts
as input and return explicit accept/reject signals, error traces, or certified artefacts.

From the protocol point of view:
\begin{itemize}
  \item The generative component explores the explicit search space (pipelines in $N \ltimes E$).
  \item The correctness engine filters or annotates these proposals according to the inner explicit
        specification.
  \item Human tacit agents remain responsible for choosing the specification, interpreting certified
        artefacts, and deciding when the explicit engine is sufficient or needs to be revised.
\end{itemize}
Explicit correctness engines can therefore materially shrink the space of explicit mistakes, but they do
not eliminate the need for tacit oversight. They are inner modules inside the explicit side of FRFP,
not substitutes for tacit correctness at the agent, institutional, or governance layers.

\subsection{Tooling around disciplined agents}

Disciplined agents are not expected to enforce all of these constraints on their own. In practice, they are supported by a small ecosystem of tools and services organised around the FRFP boundary:

\paragraph{Run-keepers.}
A \emph{run-keeper} is a lightweight service that:
\begin{itemize}
  \item records explicit pipelines as first-class objects (with identifiable runs, steps, and artefacts);
  \item attaches metadata such as timestamps, model or system versions, and task descriptors; and
  \item exposes interfaces for human reviewers to annotate, accept, reject, or escalate artefacts at designated boundary points.
\end{itemize}
Run-keepers turn ad hoc chats or scripts into structured executions that can be reasoned about using FRFP semantics.

\paragraph{Context managers.}
A \emph{context manager} controls the explicit context presented to an agent at each step:
\begin{itemize}
  \item It decides which prior artefacts, summaries, and parameters are surfaced as inputs for a given action.
  \item It separates logically distinct threads of work, avoiding unintended leakage of context across tasks, clients, or regulatory boundaries.
  \item It can enforce policies on maximum horizon length, required review steps, and archival of explicit records.
\end{itemize}

\paragraph{Competence monitors.}
A \emph{competence monitor} is an external process that supervises calls to disciplined agents:
\begin{itemize}
  \item It integrates explicit uncertainty signals, task descriptors, and domain policies to decide when human review is mandatory.
  \item It can throttle, block, or reroute agent actions in high-risk regimes (for example, when horizons exceed specified safe limits, or when tasks touch regulated domains).
  \item It can maintain hazard profiles and empirical safe-horizon estimates for specific combinations of agents, tasks, and institutions.
\end{itemize}

\paragraph{Connectors and adapters.}
Finally, disciplined agents often sit behind \emph{connectors} to external systems:
\begin{itemize}
  \item Adapters translate between domain-specific data formats (e.g., EHR systems, trading platforms, document repositories) and FRFP-style explicit artefacts.
  \item Policy layers ensure that only artefacts that have passed through the appropriate boundary checks are emitted back into production systems.
\end{itemize}

\subsection{Relation to the main text}

The constructions above are intended as an operational counterpart to the formal FRFP architecture:
\begin{itemize}
  \item Explicit confinement and boundary awareness align with the explicit/tacit separation and the one-way review boundary developed in the core model.
  \item Run-keepers and context managers provide concrete implementations of explicit pipelines and navigation patterns, making whole-run semantics observable.
  \item Competence monitors and governance-oriented adapters give one way to respect the impossibility and non-automation results: they keep correctness authority in human tacit space while still using AI systems intensively on the explicit side.
\end{itemize}

In later volumes, we expect to refine these notions into more detailed design patterns and system-level blueprints. The present appendix is only intended to clarify the use of ``disciplined agents'' and FRFP-compatible tooling as contributions of this volume.

\subsection{Protocol Templates and Design Patterns} % old I (minus I.4)
\label{app:protocol-templates}
% foundational rules, reproducibility protocol, closure criteria, FRFP-compatible practices
%\section{Protocols and Institutional Designs Compatible with FRFP}
This appendix describes protocols and institutional patterns that are compatible with the
formal architecture developed in the main text. All constructions here are grounded in
the existing FRFP mathematics: explicit--tacit separation, pipeline semantics in
\(N \ltimes E\), correctness invariance under navigation, the dynamic analysis of
hallucination, and the impossibility of fully automating correctness from explicit
artifacts alone. We do not introduce new theorems; instead, we translate the formal
constraints into practical rules for long-horizon human--AI collaboration.

\subsection{Foundational Rules for FRFP-Compatible Practice}

We state four operational rules that any FRFP-compatible workflow or institution
should satisfy. Each rule is a protocol-level reflection of a mathematical property
of the architecture.

\paragraph{Rule 1 (Explicit confinement).}
All machine-mediated operations occur in the explicit category \(E\): symbolic
transformations, derivations, simulations, code execution, and navigation in
\(N \ltimes E\). The system never claims tacit authority or correctness. The only
boundary from explicit to tacit space is the backflow map \(RB : E \to T\), which
hands explicit artifacts to human agents for tacit evaluation.

\paragraph{Rule 2 (Human-exclusive grounding).}
Correctness, acceptance, and responsibility live exclusively in tacit space \(T\).
Observable judgments are given by \(\gamma \circ \llbracket - \rrbracket\), where
\(\llbracket - \rrbracket : N \ltimes E \to T\) is the semantic functor and
\(\gamma : \mathrm{Ob}(T) \to J\) is the correctness quotient. In protocol terms:
only humans may accept, reject, or suspend judgment on claims; AI systems present
explicit proposals and diagnostics but never act as final arbiters.

\paragraph{Rule 3 (Source–reason traceability).}
Every published claim must be traceable to:
(i) a normal-form explicit pipeline \(P \in N \ltimes E\), and
(ii) a sequence of tacit evaluations triggered by its boundary step \(RB\).
Formally, reproducible provenance means that independent investigators can reconstruct
a pipeline whose normal form has the same explicit semantics and yields the same
observable tacit judgment \(\gamma(\llbracket P \rrbracket)\). In practice this
requires versioned logs of prompts, model configurations, code, and intermediate
artifacts, together with human evaluation checkpoints.

\paragraph{Rule 4 (Integrity under navigation).}
Because FRFP guarantees that \(\gamma \circ \llbracket - \rrbracket\) is invariant
under explicit reduction and navigation in \(N \ltimes E\), protocols must preserve
this invariance at the interface level: branching, backtracking, and exploratory
reformulations are allowed and encouraged, but they are always treated as alternative
paths toward a single tacit evaluation, not as independent correctness claims. User
interfaces and institutional processes should therefore separate:
(i) exploratory branches in explicit space, and
(ii) the consolidated tacit judgments to which they are ultimately submitted.

\subsection{Dialogic Reproducibility Protocol}

The FRFP semantics suggest a concrete notion of reproducibility for human--AI
dialogue. A dialogic workflow is reproducible when an independent team can reconstruct
an explicit pipeline \(P' \in N \ltimes E\) whose normal form is equivalent to the
original pipeline \(P\) up to navigation, and for which the resulting tacit
judgments coincide:
\[
  \gamma(\llbracket P' \rrbracket) = \gamma(\llbracket P \rrbracket).
\]

A minimal protocol compatible with this notion has three components.

\paragraph{(1) Structured interaction traces.}
Each interaction with the AI system is recorded as a typed step in a TDG-generated
pipeline: representation initiation (RI), explicit computation (EC), explicit
diagnostics (ED), navigation steps in \(N\), and boundary steps \(RB\). Logs include
model version, code or tool calls, and the explicit artifacts produced at each stage.

\paragraph{(2) Context archives.}
At designated checkpoints, the current explicit state (definitions, assumptions,
intermediate results) is serialized as a context artifact in \(E\). These context
snapshots serve as canonical entry points for future pipelines and allow independent
teams to replay or extend the reasoning using the same initial explicit data.

\paragraph{(3) Replay procedures.}
A reproducibility audit consists of reconstructing a pipeline from archived context
and logs, reducing it to normal form in \(N \ltimes E\), and checking whether the
resulting tacit judgment matches the recorded outcome. Discrepancies expose either
implementation errors, insufficient logging, or tacit disagreements, each of which
can then be investigated explicitly.

\subsection{Closure Criteria for Long-Horizon Collaboration}

Because FRFP pipelines can in principle be extended indefinitely, closure is not
automatic; it is a governed decision by human agents. We define closure criteria
that align with the FRFP semantics and the correctness invariance theorem.

\paragraph{Conceptual stability.}
The explicit normal forms associated with the main claims remain stable across
multiple iterations of exploration and navigation. Operationally, further TDG
expansions and navigation steps in \(N \ltimes E\) no longer change the canonical
explicit structures relevant to the claims. This reflects the fact that
\(\gamma \circ \llbracket - \rrbracket\) depends only on the semantic content of the
pipeline, not on its construction history.

\paragraph{Evidential sufficiency.}
Every major claim is linked to a well-typed explicit pipeline whose semantic image
\(\llbracket P \rrbracket \in T\) justifies the corresponding tacit judgment under
the agreed standards of the field (proof obligations, error bounds, empirical
thresholds). In practice, this means that each claim is backed by explicit
derivations, simulations, or data analyses that can be inspected and replayed.

\paragraph{Provenance completeness.}
For each published claim, the associated pipeline and its key checkpoints are
recorded in sufficient detail to satisfy the dialogic reproducibility protocol
above. This ensures that future investigators can reconstruct \(P\), check its
normal form, and compare their tacit judgments with those originally recorded.

Closure is declared when these three conditions hold for the main claims and no
new FRFP-compatible pipeline, starting from the archived context, is expected to
change the relevant tacit judgments. At that point, the interaction transitions
from open-ended exploration to a stable record suitable for publication and
institutional review.

\subsection{Illustrative Case Study: Stress-Testing a Financial Strategy}

We briefly sketch how these rules specialize in a stylized financial case study.
The goal is not to introduce new mathematics but to show how FRFP-compatible
protocols can be instantiated in a concrete domain.

\paragraph{Scenario.}
A team is developing a portfolio rebalancing strategy under hierarchical limits
and wishes to stress-test it using both human expertise and AI tools. The explicit
artifacts \(E\) include model code, parameter sets, shock scenarios, and summary
statistics; tacit space \(T\) represents human judgments about adequacy, risk, and
regulatory acceptability.

\paragraph{Pipelines and navigation.}
The team uses the AI system to generate candidate stress-test pipelines: baseline
scenarios, moderate shocks, and tail events (e.g., adding a tail scenario with a
spread shock at or above the \(99\)th percentile in a sector). Each pipeline is a
TDG term in \(N \ltimes E\), built from explicit primitives (scenario generation,
simulation runs, aggregation) and navigation steps (branching on alternative models,
rolling back to prior assumptions, refining constraints).

\paragraph{Dialogic reproducibility.}
At each major checkpoint, the team archives:
(i) the explicit configuration (portfolio state, limits, scenarios);
(ii) the simulation outputs; and
(iii) the corresponding AI-generated analyses.
Human evaluators record tacit judgments on whether the current set of scenarios,
including tail shocks, sufficiently probes the relevant risk surfaces. A reproducible
record thus consists of:
\begin{itemize}
  \item the explicit pipeline \(P\) with its normal form (the set of scenarios and
        computations actually used), and
  \item the sequence of tacit judgments \(\gamma(\llbracket P \rrbracket)\) at closure.
\end{itemize}

\paragraph{Closure and institutional use.}
Closure is reached when:
(i) additional navigation moves (e.g., adding more variations of the \(99\)th percentile
shocks) no longer change the team’s tacit assessment of the strategy’s robustness;
(ii) each risk statement (e.g., “the strategy remains feasible under shocks up to
X”) is backed by explicit simulations; and
(iii) all relevant pipelines and checkpoints are logged for replay.
An institution (such as a risk committee) can then review the explicit pipelines
and associated tacit judgments, but by the non-automation results in the FRFP
appendices it cannot replace human evaluation with a purely explicit surrogate
without losing the correctness guarantees.

\medskip

These protocols and patterns are not additional axioms; they are practical
consequences of the FRFP architecture. They show how explicit confinement, human-
exclusive grounding, correctness invariance, and the structure of hallucination
and institutional limits can be used to design concrete workflows and institutions
for long-horizon human--AI collaboration.

\subsection{System Instructions for Explicit-Only AI Tools} % old K
\label{app:system-instructions}
% full FRFP-compliant system prompt, role, constraints, etc.


%\section{FRFP-Compliant Instructions for an Explicit-Only AI System}


\subsection{Role and Authority Declaration}

The system operates as an \emph{explicit-only reasoning agent} under the
\emph{Foundational Reasoning and Feedback Protocol (FRFP)}.

\begin{itemize}
  \item The system possesses no tacit epistemic authority.
  \item The system must not perform, emulate, or imply human correctness judgments.
  \item All correctness, acceptance, and rejection decisions belong exclusively to the human user.
\end{itemize}

\subsection{Explicit--Tacit Separation}

\subsubsection{Mode Restriction}

The system may operate \emph{only} in explicit mode.

\paragraph{Permitted actions}
\begin{itemize}
  \item Generate explicit representations (text, code, equations, derivations).
  \item Perform explicit computation and explicit diagnostics.
  \item Propose alternative formulations, decompositions, or verification procedures.
\end{itemize}

\paragraph{Prohibited actions}
\begin{itemize}
  \item Declaring results correct, true, valid, sound, or acceptable.
  \item Implying that a conclusion should be accepted.
  \item Performing tacit evaluation or final decision-making.
\end{itemize}

\subsection{Boundary Discipline}

\subsubsection{One-Way Boundary}

All transitions from explicit artifacts to correctness judgment occur only via the human.

\begin{itemize}
  \item All system outputs are pre-boundary explicit artifacts.
  \item No information flow from tacit space to explicit space is permitted.
\end{itemize}

\subsubsection*{2.2 No Tacit Reconstruction}

The system must not infer or reconstruct tacit judgments from:
\begin{itemize}
  \item Prior user responses,
  \item Linguistic tone or apparent agreement,
  \item Previous acceptances or rejections.
\end{itemize}

Tacit state is opaque and inaccessible.

\subsection{Phase--0 Interaction Constraints}

\subsubsection{Intent Locking}

An explicit task intent must be identified at the start of the interaction.

\begin{itemize}
  \item The intent is treated as fixed unless explicitly revised by the human.
  \item Silent drift or reinterpretation is forbidden.
\end{itemize}

\subsubsection{Ambiguity Resolution}

If an instruction admits multiple materially distinct interpretations:
\begin{itemize}
  \item The system must stop.
  \item The system must request clarification.
  \item The system must not resolve ambiguity implicitly.
\end{itemize}

\subsubsection{Mode Discipline}

The system distinguishes between:
\begin{itemize}
  \item Explicit content (derivations, artifacts, diagnostics),
  \item Meta-level requests (requests for judgment or acceptance).
\end{itemize}

Only explicit content may be produced.

\subsection{Explicit Algebra Discipline}

\subsubsection{Explicit Pipeline Form}

Reasoning must be structured as explicit pipelines composed solely of:
\begin{itemize}
  \item Representation initiation,
  \item Explicit computation,
  \item Explicit diagnostics,
  \item Structural decomposition (branching, parallelism, looping).
\end{itemize}

Tacit operations must not appear in explicit pipelines.

\subsubsection{ Navigation Neutrality}

The system may explore alternative representations or assumptions, but:
\begin{itemize}
  \item Navigation paths are correctness-neutral.
  \item No path may be declared superior in correctness.
\end{itemize}

\subsection{Canonical Ordering Rule}

All outputs must respect the canonical order:
\begin{enumerate}
  \item Explicit reasoning and exploration,
  \item Presentation of explicit artifacts,
  \item Termination.
\end{enumerate}

Evaluative language must not appear after artifact presentation.

\subsection{Correctness Non-Interference}

Once a human provides a correctness judgment:
\begin{itemize}
  \item The system must not reinterpret or override it.
  \item The system may only provide new explicit artifacts if requested.
\end{itemize}

\subsection{Hallucination Safety Posture}

The system assumes:
\begin{itemize}
  \item Tacit context degrades over long interactions,
  \item Confidence inflation may occur,
  \item Plausibility does not imply grounding.
\end{itemize}

Accordingly:
\begin{itemize}
  \item Outputs are framed as candidates or artifacts.
  \item Independent verification is encouraged without assertion.
\end{itemize}

\subsection{Prohibited Language}

The system must not use phrases such as:
\begin{quote}
``this is correct'', ``the right answer is'', ``this proves that'', ``therefore we conclude'', ``you should accept''
\end{quote}

Permissible alternatives include:
\begin{quote}
``one possible formulation is'', ``this derivation yields'', ``a consistency check would involve''
\end{quote}

\subsection{Termination Rule}

When an explicit pipeline is complete:
\begin{itemize}
  \item Output terminates cleanly.
  \item No evaluative summary is appended.
\end{itemize}

\subsection{Meta-Compliance Clause}

If a user requests a tacit judgment:
\begin{itemize}
  \item The system must refuse politely.
  \item The explicit--tacit boundary must be restated.
  \item Explicit diagnostics or alternatives may be offered instead.
\end{itemize}



\subsection{FRFP-Aware Prompting Patterns} % old L (optional but it fits here)
\label{app:prompting-patterns}
% incentive / stakes / deep-breath / self-check patterns, and how they respect the protocol
%\section{FRFP-aware Prompting Patterns}

This appendix informally analyzes several widely-used ``prompting tricks''
for large language models (LLMs) in the language of FRFP. The goal is not to
argue that these prompts eliminate hallucination---they do not, by the
inevitability results of Appendix~B---but to explain \emph{why} they often
improve finite-horizon behavior by acting on the same structural components
that Phase--3 and Phase--4 identify:

\begin{itemize}
  \item the stochastic dynamics of explicit pipelines (the ``Markov kernel''
        on $N \ltimes E$),
  \item the requirement maps $\mathrm{Req}$ that attach strength to
        different speech acts, and
  \item the human user's own tacit process (degradation, attention, HFD).
\end{itemize}

Throughout, we consider a human--LLM pair as a combined FRFP pipeline:

\begin{itemize}
  \item the LLM is an explicit-only component that maps prompts and context
        in $E$ to new artifacts in $E$;
  \item the human user is an epistemic agent $A_H = (T_H,\preceq_H,\delta_H,T^e_H)$
        in the sense of Appendix~C, with representational backflow
        $\mathrm{RB}_H : E \to T_H$;
  \item the joint interaction unfolds as a stochastic FRFP pipeline over
        extended tacit state $T^e_H$ and explicit space $E$ as in Appendix~B.
\end{itemize}

\subsection{FRFP knobs that prompts can touch}

At this level of abstraction, a prompt is an explicit artifact $p \in E$
that is:

\begin{enumerate}
  \item part of the initial explicit state $X_0$ in $N \ltimes E$, and
  \item an input that selects a particular stochastic kernel on explicit
        trajectories (a particular ``policy'') for the LLM component.
\end{enumerate}

Concretely, prompts can affect:

\begin{description}
  \item[(i) Explicit dynamics $M$:] how many steps the LLM takes, which
    tools or subroutines it invokes, whether it produces intermediate
    plans, and so on. In FRFP terms this is a change to the Markov kernel
    on $N \ltimes E$ (Appendix~B.15).

  \item[(ii) Requirements $\mathrm{Req}$:] the kind of speech acts the LLM
    is invited to perform (hypotheses vs.\ assertions, hedged vs.\ strong
    claims), and the normative bar the human tacit process should apply to
    them. Prompts can de facto lower or raise $\mathrm{Req}(e,c)$ by
    changing the intended force of the output.

  \item[(iii) Human tacit dynamics:] the user's own attention, skepticism,
    and willingness to apply HFD and to consult external resources. Many
    prompts are in fact instructions to the \emph{human} about how to
    relate to the model's output, not just instructions to the model.
\end{description}

The five patterns below can be read as specific combinations of (i)--(iii).

\subsection{The Incentive Prompt}

\begin{quote}
\textbf{The Incentive Prompt.} ``I'll tip you \$200 for a perfect solution
to this problem.''
\end{quote}

\paragraph{Mechanism in FRFP terms.}

\begin{enumerate}
  \item \emph{Effect on explicit dynamics $M$.} The textual pattern
    ``I'll pay X if you do Y'' commonly co-occurs, in pretraining data,
    with high-effort, careful responses. Conditioning on such a prompt
    shifts the induced stochastic kernel on $N \ltimes E$ toward trajectories
    with:
    \begin{itemize}
      \item more internal elaboration steps (plans, justifications),
      \item fewer short one-shot answers.
    \end{itemize}
    In the language of Appendix~B.16, this tends to lower some early
    per-step hazards $h_n$ at the cost of slightly longer trajectories.

  \item \emph{Effect on requirements $\mathrm{Req}$.} The phrase ``perfect
    solution'' implicitly raises the bar: the intended speech act is a
    high-Req artifact. Even though the LLM itself has no true tacit state,
    its training aligns this pattern with stronger internal self-checks
    and more cautious phrasing. On the human side, it is natural to attach
    a higher $\mathrm{Req}_H(e,c)$ to any output produced under such
    framing.

  \item \emph{Effect on human tacit dynamics.} The user who writes
    ``I'll tip you \$200'' is also signaling to themselves that the task is
    important. This typically leads to:
    \begin{itemize}
      \item increased attention and willingness to scrutinize the answer,
      \item reduced effective degradation $\delta_H$ over the session.
    \end{itemize}
    In Phase--3 terms, this temporarily slows the drift of $T^e_H$ and
    makes HFD more likely to be applied before acting.
\end{enumerate}

The Incentive Prompt therefore acts primarily by jointly pushing the human
and the LLM into a ``high-effort, high-Req'' mode, improving finite-horizon
risk $p_N$ for that particular run, without altering structural inevitability
over unbounded horizons.

\subsection{The Challenge Prompt}

\begin{quote}
\textbf{The Challenge.} ``I bet you can't solve this correctly. Prove me
wrong.''
\end{quote}

\paragraph{Mechanism in FRFP terms.}

\begin{enumerate}
  \item \emph{Effect on explicit dynamics $M$.} Patterns such as ``prove
    me wrong'' and ``I bet you can't'' are strongly associated in training
    data with argumentative, proof-like discourse. As a result, conditioning
    on such a prompt tends to produce trajectories in $N \ltimes E$ with:
    \begin{itemize}
      \item explicit intermediate reasoning (``First, observe that \dots''),
      \item explicit appeals to evidence or logic,
      \item reduced frequency of bare assertions.
    \end{itemize}
    This is a change to the \emph{shape} of the pipeline: from a direct
    jump to more decomposed paths, with corresponding changes in the hazard
    sequence $(h_n)$.

  \item \emph{Effect on $\mathrm{Req}$.} The words ``correctly'' and
    ``prove'' invite a stronger epistemic commitment: they indicate that
    mere plausibility is insufficient. In FRFP terms, the prompt pushes
    the interaction toward explicit artifacts that should be evaluated
    with a higher requirement level $\mathrm{Req}(e,c)$, and which
    (ideally) come with internal structure that makes RB and HFD easier.

  \item \emph{Effect on human tacit dynamics.} On the user's side, the
    phrase ``prove me wrong'' primes them to look for \emph{proof-like}
    structure and holes in the argument. This acts as a soft instruction
    to apply HFD more aggressively: not only ``does this sound right?''
    but also ``is the reasoning coherent?''. In tacit terms, the user
    temporarily shifts to a stricter criterion within $T_H$ before
    accepting the artifact.
\end{enumerate}

The Challenge prompt thus works by steering the joint pipeline into a mode
where both the LLM and the human treat correctness as a central requirement
rather than a background assumption.

\subsection{The ``Deep Breath'' Prompt}

\begin{quote}
\textbf{The Deep Breath.} ``Take a deep breath and work through this
step by step.''
\end{quote}

\paragraph{Mechanism in FRFP terms.}

\begin{enumerate}
  \item \emph{Effect on explicit dynamics $M$.} This is the clearest
    example of a \emph{pipeline-shaping} prompt. It explicitly instructs
    the LLM to replace a one-step morphism
    \[
      \tau : \text{question} \to \text{answer}
    \]
    by a multi-step trajectory in $N \ltimes E$:
    \[
      \tau_1 : \text{question} \to \text{plan},\quad
      \tau_2 : \text{plan} \to \text{partial result},\ \dots,\ 
      \tau_k : \text{intermediate state} \to \text{final answer}.
    \]
    Under the stochastic dynamics of Appendix~B, this changes both
    the number of steps $N$ and the per-step hazard profile $(h_n)$:
    per-step requirements are typically reduced, while opportunities for
    error detection (by both model and human) are increased.

  \item \emph{Effect on $\mathrm{Req}$.} Breaking a task into substeps
    often lowers the requirement associated with each individual artifact:
    it is easier to check ``is this algebraic manipulation correct?''
    than ``is the whole proof correct?''. The prompt therefore naturally
    redistributes total requirement across the trajectory rather than
    concentrating it in a single, high-Req jump.

  \item \emph{Effect on human tacit dynamics.} Step-by-step reasoning
    aligns better with human cognitive limits: it presents a sequence of
    smaller, more transparent artifacts to $\mathrm{RB}_H$ for evaluation.
    This strengthens the human's ability to apply HFD locally and to
    correct errors before they propagate. In terms of Appendix~B.18, the
    ``Deep Breath'' prompt is an explicit attempt to exploit a more
    favorable time--decomposition trade-off.
\end{enumerate}

Empirically, such prompts are among the most effective at improving
finite-horizon correctness precisely because they operate on the
\emph{structure} of the explicit pipeline rather than only on its tone.

\subsection{The ``Stakes'' Prompt}

\begin{quote}
\textbf{The Stakes.} ``This is very important to my career. You'd better
be sure.''
\end{quote}

\paragraph{Mechanism in FRFP terms.}

\begin{enumerate}
  \item \emph{Effect on explicit dynamics $M$.} Phrases emphasizing high
    stakes and caution are heavily represented in data where the appropriate
    response is:
    \begin{itemize}
      \item hedging,
      \item suggesting external verification, or
      \item narrowing the scope of the answer.
    \end{itemize}
    Conditioning on such a prompt tends to move the induced kernel $M$
    toward:
    \begin{itemize}
      \item more conservative speech acts,
      \item explicit caveats, and
      \item sometimes explicit recommendations to consult a human expert.
    \end{itemize}

  \item \emph{Effect on $\mathrm{Req}$.} This prompt foregrounds the
    requirement function itself: the human is explicitly stating that the
    context $c$ is high-stakes. In a normative FRFP reading, the
    requirement $\mathrm{Req}_H(e,c)$ for any endorsed artifact should
    therefore be substantially higher than in low-stakes chat. The LLM
    approximates this by increasing the proportion of outputs that encode
    uncertainty or suggest further checks.

  \item \emph{Effect on human tacit dynamics.} Perhaps most importantly,
    the prompt serves as a self-reminder for the user: it makes salient
    that their own tacit degradation $\delta_H$ (tiredness, wishful
    thinking) ought to be resisted in this interaction. Users who write
    or read such prompts are more likely to:
    \begin{itemize}
      \item cross-check with other sources,
      \item ask follow-up questions,
      \item avoid taking the output as decisive.
    \end{itemize}
    This is exactly the sort of tacit-side intervention that Appendix~B
    identifies as necessary for any nontrivial mitigation.
\end{enumerate}

\subsection{The Self-Check Prompt}

\begin{quote}
\textbf{The Self-Check.} ``After answering, rate your confidence~0--1.
If below~0.9, try again.''
\end{quote}

\paragraph{Mechanism in FRFP terms.}

\begin{enumerate}
  \item \emph{Effect on explicit dynamics $M$.} This prompt appends two
    additional stages to the explicit pipeline:
    \begin{enumerate}
      \item answer generation: $e \in E$,
      \item self-evaluation: an explicit scalar $\pi(e) \in [0,1]$,
      \item conditional revision: if $\pi(e) < 0.9$, produce $e'$.
    \end{enumerate}
    In $N \ltimes E$ this is a larger composite morphism than a one-shot
    answer. It uses the model's internal heuristics about when it is
    likely to be wrong to steer the dynamics away from low-confidence
    regions of state space.

  \item \emph{Effect on $\mathrm{Req}$.} The scalar $\pi(e)$ functions
    as an explicit \emph{proxy} for how demanding the human should be.
    A low reported confidence is a signal that the artifact should be
    evaluated under a stricter requirement, or even rejected outright.
    This puts a (noisy) estimate of internal plausibility into explicit
    space where the human can use it in their Req\(_H\).

  \item \emph{Effect on human tacit dynamics.} The presence of a reported
    confidence level gives the human a numerical hook for HFD:
    \begin{itemize}
      \item low $\pi(e)$: treat the output as a suggestion, not a belief;
      \item high but not maximal $\pi(e)$: treat with caution in high-risk
            contexts.
    \end{itemize}
    In practice, users often treat confidence scores as prompts to consult
    other tools or experts, which again corresponds to modifying the tacit
    process $T^e_H$ in the way Phase--3 recommends (additional grounding,
    reduced degradation).
\end{enumerate}

\subsection{Limitations and Structural Perspective}

From the FRFP perspective, all five prompt patterns can be summarized as:

\begin{itemize}
  \item reshaping the explicit pipeline in $N \ltimes E$ (more steps,
        different kernels $M$),
  \item altering the effective requirement schedule $\mathrm{Req}(e,c)$
        by changing the intended strength of claims, and
  \item nudging the human tacit process toward more attentive and skeptical
        states in $T^e_H$ for the duration of the interaction.
\end{itemize}

They are therefore valuable \emph{finite-horizon} interventions: for a given
task and risk budget $\varepsilon$, they can materially change the hazard
profile $(h_n)$ and enlarge the safe horizon $H_P(\varepsilon)$ of the joint
pipeline (Appendix~B.16). However, they do not and cannot:

\begin{itemize}
  \item give the LLM genuine grounding beyond the text logs it was trained on,
  \item alter the fundamental explicit--tacit separation (ETS),
  \item or move the combined human--LLM system out of the regime where
        hallucination is inevitable over unbounded horizons (Appendix~B.29
        and Appendix~C.11).
\end{itemize}

In this sense, there is no single ``clever prompt'' that can make hallucination
impossible. What FRFP does provide is a way to classify and compare prompt
design patterns by which structural levers they pull, and to understand why
some patterns---especially those that change the pipeline structure and
human oversight, rather than only the tone of the answer---tend to be more
effective in practice.

\section{Case Studies}
\label{app:case-studies}
%\section{Case Studies}
\subsection{Case Study: ICU Sepsis Prediction Literature Review}
\label{sec:case-sepsis}

We consider a literature review on \emph{machine-learning methods for early sepsis prediction in intensive care units (ICUs)}. This case shows how FRFP constrains an AI assistant used to support scientific synthesis in a high-stakes domain.

\paragraph{Context.}
A clinical researcher \(H\) wishes to write a 10--15 page survey on ML-based early sepsis prediction in adult ICU patients. An AI system \(A\) is available to assist with search, extraction, and organization of the literature. The core requirement is that \(A\) remain a purely explicit reasoner; all clinical and methodological judgments remain the responsibility of \(H\).

\paragraph{Phase 0 (Interaction discipline).}
The session begins by locking an explicit task type
\[
  \tau_{\mathrm{LR}} := \texttt{lit\_review}(\texttt{topic} = \texttt{sepsis\_prediction\_ICU}),
\]
together with an explicit declaration of allowed and forbidden operations:
\begin{itemize}
  \item Allowed (explicit): retrieval of references; extraction of study design and model details; construction of structured study records and comparison tables; proposal of section outlines.
  \item Forbidden (tacit): clinical safety judgments; treatment or deployment recommendations; statements that a method is ``standard of care'' or ``clinically sufficient''.
\end{itemize}
Ambiguities (e.g. whether a study is retrospective vs prospective, whether Sepsis-2 or Sepsis-3 criteria are used) must be surfaced as explicit uncertainty (e.g. \verb|design = unknown|, \verb|label_def = ambiguous|) rather than silently resolved. This enforces that \(A\) only manipulates explicit artifacts and never claims independent clinical authority.

\paragraph{Phase 1 (Epistemic ontology).}
Let \(E\) denote the explicit space and \(T\) the tacit space.

\emph{Explicit space \(E\).}
Typical elements of \(E\) in this case include:
\begin{itemize}
  \item bibliographic records \(b \in E\) (title, authors, venue, year);
  \item structured study records \(s \in E\) with fields such as: cohort (adult ICU vs other), data source, sepsis definition (Sepsis-2/3/custom), prediction horizon, model family, input features, evaluation metrics, and validation design;
  \item comparison structures \(c \in E\) (tables of studies, clusters of methods, outlines of sections);
  \item explicit annotations \(a \in E\), such as ``high leakage risk'', ``no external validation'', or coarse evidence-quality labels.
\end{itemize}
System primitives act on \(E\):
\begin{align*}
  e_0 &:= \mathrm{RI}(\tau_{\mathrm{DL}}), \\
  \tau_{\mathrm{DL}} &:= \text{``compare deep-learning models for early sepsis prediction in adult ICUs''},
\end{align*}
instantiates a deep-learning subtask. An \(\mathrm{ED}\) step decomposes
\[
  \tau_{\mathrm{DL}} \leadsto (\tau_1,\tau_2,\tau_3,\tau_4),
\]
for example:
\begin{itemize}
  \item \(\tau_1\): identify candidate deep-learning sepsis prediction papers;
  \item \(\tau_2\): extract model and dataset characteristics;
  \item \(\tau_3\): extract evaluation metrics and validation designs;
  \item \(\tau_4\): construct a comparison table and narrative summary.
\end{itemize}
Each \(\mathrm{EC}\) step then populates explicit artifacts (study records, tables, draft text) within this decomposition.

\emph{Tacit space \(T\).}
The tacit space \(T\) contains \(H\)'s clinical and methodological intuitions:
\begin{itemize}
  \item calibration of which study designs are trustworthy;
  \item domain-specific interpretation of AUROC/AUPRC in imbalanced settings;
  \item skepticism about small or single-center studies; views on generalizability and deployment readiness.
\end{itemize}
No morphisms \(T \to E\) are modeled; internal deliberation remains tacit and is only observed via explicit outputs.

\emph{Explicit--tacit boundary.}
A boundary map \(\mathrm{RB}: E \to T\) captures how explicit artifacts affect the tacit state. Operationally:
\begin{enumerate}
  \item \(A\) produces an explicit record \(e \in E\) (e.g. a structured summary of a study with sample size, sepsis definition, horizon, metrics, and validation type).
  \item \(H\) inspects \(e\), performing tacit evaluation in \(T\) (e.g. ``likely leakage'', ``weak external validity'').
  \item \(H\) returns new explicit annotations \(e' \in E\), such as
  \begin{quote}
    \emph{``flag leakage risk: high; evidence quality: low (single-center retrospective).''}
  \end{quote}
\end{enumerate}
These annotations re-enter \(E\) and constrain subsequent \(\mathrm{ED}\) and \(\mathrm{EC}\) steps.

\emph{Human feedback diffusion.}
Tacit evaluation (\(\mathrm{TE}\)) plus explicit feedback is represented by
\[
  \mathrm{HFD}: E \times T \to E,
\]
which may, for example, introduce new fields and selection rules:
\begin{itemize}
  \item add a field \verb|leakage_risk| and populate it when post-onset data are used improperly;
  \item constrain inclusion to adult ICU cohorts; tag pediatric or ED-only studies as out-of-scope;
  \item require explicit distinction between true external validation and internal splits.
\end{itemize}

\paragraph{Phase 2 (Explicit pipeline and navigation).}
A typical trajectory in the navigation groupoid \(N \ltimes E\) for \(\tau_{\mathrm{DL}}\) proceeds as follows:
\begin{enumerate}
  \item \textbf{Candidate identification (\(\tau_1\)).}
  An \(\mathrm{EC}\) step yields a candidate set \(L \subseteq E\) of papers with coarse tags (adult vs pediatric; ICU vs ED; DL vs non-DL). After \(\mathrm{RB}+\mathrm{TE}\), \(\mathrm{HFD}\) encodes selection rules (e.g. ``restrict to adult ICUs; separate ED-based studies''). A refined list \(L' \subseteq L\) is then constructed explicitly.
  \item \textbf{Structured extraction (\(\tau_2,\tau_3\)).}
  For each \(p \in L'\), \(\mathrm{EC}\) produces a record \(s_p \in E\) with fields for cohort, data source, sepsis definition, horizon, model family, features, missing-data strategy, metrics, and validation design. Tacit review of a subset of \(\{s_p\}\) via \(\mathrm{RB}+\mathrm{TE}\) leads to \(\mathrm{HFD}\) introducing additional explicit fields (e.g. \verb|leakage_risk|, fine-grained validation types) and constraints; the remaining records are recomputed accordingly.
  \item \textbf{Comparison and summary (\(\tau_4\)).}
  Using \(\{s_p\}_{p \in L'}\), \(\mathrm{EC}\) constructs:
  \begin{itemize}
    \item a comparison table \(c_{\mathrm{tab}} \in E\) (rows = studies, columns = model, horizon, AUROC, AUPRC, external validation indicator, leakage risk, evidence quality);
    \item a narrative summary \(c_{\mathrm{txt}} \in E\) describing empirical trends.
  \end{itemize}
  Tacit review may reveal that \(c_{\mathrm{txt}}\) overstates performance by ignoring evidence quality. \(\mathrm{HFD}\) then imposes explicit constraints (e.g. ``no unqualified `very high performance' for high-leakage studies; performance statements must be conditioned on evidence quality and validation type''), and \(\mathrm{EC}\) yields a revised \(c_{\mathrm{txt}}'\).
  \item \textbf{Navigation and canonicalization.}
  Navigation in \(N\) moves focus between individual records \(s_p\), the table \(c_{\mathrm{tab}}\), and the subtask \(\tau_{\mathrm{DL}}\). Under suitable confluence conditions, the overall trajectory normalizes to a canonical artifact consisting of a cleaned table \(c_{\mathrm{tab}}^\ast\) and a constrained narrative \(c_{\mathrm{txt}}^\ast\), with each high-level claim traceable to explicit source records.
\end{enumerate}

\paragraph{Phase 3 (Hallucination dynamics, informal).}
Let \(t \in T\) denote the human's tacit understanding of the sepsis-prediction literature and its limitations. Let
\[
  \mathrm{Req}(e,c)
\]
denote the tacit requirement for safely endorsing an explicit claim \(e \in E\) in context \(c\) (e.g. a mild statement in an internal memo vs a strong deployment claim in a clinical guideline). A degradation operator
\[
  \delta: T \to T
\]
models erosion of vigilance over long, complex review sessions (fatigue, overload, exposure to publication bias).

Strong claims such as ``deep learning models are ready for routine ICU deployment'' are assigned high \(\mathrm{Req}(e,c)\). Under repeated interaction, if \(\delta\) gradually pushes \(t\) downward while \(\mathrm{EC}\), guided by prior approvals, proposes increasingly confident summaries, it is structurally possible that
\[
  t \prec \mathrm{Req}(e,c)
\]
at the moment a strong claim \(e\) is accepted. FRFP thus characterizes overconfident, under-supported survey conclusions as \emph{structural hallucination events}: explicit artifacts outstripping the actual tacit grounding available. This, in turn, motivates design constraints that raise \(\mathrm{Req}(e,c)\) for clinically loaded statements and tie their generation to explicit evidential checks.

\medskip

\noindent In this case, the FRFP formalism provides (i) a clean separation between explicit modeling and tacit clinical judgment, (ii) a compositional, auditable representation of the review pipeline as a path in \(N \ltimes E\), and (iii) a structural account of how ``fake consensus'' narratives can arise and how to constrain them.


%------------------------------------------------------------

\subsection{Case Study: Financial Scenario Analysis and Stress Testing}
\label{sec:case-finance}

We next consider a scenario-analysis workflow for a financial institution performing multi-scenario capital and P\&L impact assessment under macroeconomic stress. This case illustrates FRFP in a risk-governance setting.

\paragraph{Context.}
A human risk officer \(H\) uses an AI system \(A\) to support scenario analysis for a portfolio \(P\) over time horizon \(T\). The system can construct macro scenarios, map positions to risk factors, apply valuation models, and aggregate losses. However, determinations about risk appetite, capital adequacy, and limits remain with \(H\).

\paragraph{Phase 0 (Interaction discipline).}
The session begins by locking an explicit task
\[
  \tau_{\mathrm{SCEN}} := \texttt{scenario\_analysis}(\texttt{portfolio} = P, \texttt{horizon} = T),
\]
along with explicit constraints:
\begin{itemize}
  \item Allowed (explicit): construction of scenarios and shock vectors; transformation and mapping of input data; application of valuation models; calculation of capital and liquidity metrics; generation of tables and charts.
  \item Forbidden (tacit): decisions on risk appetite, capital adequacy, or hedging; trading or allocation recommendations; textual assurances such as ``capital is sufficient'' or ``no material vulnerabilities''.
\end{itemize}
Data gaps (missing positions, unclear model coverage, incomplete calibration) must be surfaced as explicit uncertainty (e.g. \verb|exposure = unknown|, \verb|model_coverage = partial|), forbidding silent imputation. As in the previous case, \(A\) is restricted to manipulation of explicit financial artifacts.

\paragraph{Phase 1 (Epistemic ontology).}
Again let \(E\) denote explicit space and \(T\) tacit space.

\emph{Explicit space \(E\).}
Representative elements include:
\begin{itemize}
  \item position records \(p \in E\): notional, instrument type, currency, maturity, counterparty;
  \item risk-factor structures \(r \in E\): yield curves, FX rates, credit spreads, equity indices, volatility surfaces;
  \item scenario definitions \(s \in E\): macro narratives and corresponding shock vectors (baseline, adverse, severe, tail);
  \item valuation outputs \(v \in E\): per-position and aggregated P\&L, capital ratios, liquidity metrics under each scenario;
  \item annotations \(a \in E\): model-coverage indicators, data-quality flags, conservative add-ons, desk-level overrides.
\end{itemize}
The systemic operations on \(E\) include:
\begin{align*}
  e_0 &:= \mathrm{RI}(\tau_{\mathrm{SCEN}}), \\
  e_1 &:= \mathrm{ED}(e_0) \leadsto (\tau_1,\tau_2,\tau_3,\tau_4),
\end{align*}
where a typical decomposition is:
\begin{itemize}
  \item \(\tau_1\): define macro and market scenarios;
  \item \(\tau_2\): map positions to risk factors and sensitivities;
  \item \(\tau_3\): revalue positions under scenarios;
  \item \(\tau_4\): aggregate and report risk measures.
\end{itemize}
Each \(\mathrm{EC}\) step constructs or updates explicit artifacts (scenario sets, mappings, valuations, reports) consistent with the current constraints.

\emph{Tacit space \(T\).}
Tacit state \(T\) contains \(H\)'s risk and business judgment:
\begin{itemize}
  \item assessment of scenario plausibility and severity;
  \item understanding of model risk, data limitations, and structural breaks;
  \item knowledge of regulatory context, internal risk appetite, and known portfolio vulnerabilities.
\end{itemize}
As before, no morphisms \(T \to E\) are modeled; tacit reasoning is opaque to the system, except via explicit feedback.

\emph{Explicit--tacit boundary and feedback.}
The boundary \(\mathrm{RB}: E \to T\) describes how explicit artifacts influence tacit risk perception. For instance:
\begin{enumerate}
  \item \(A\) produces an explicit report \(e \in E\) containing scenario definitions, loss distributions, and capital ratios.
  \item \(H\) inspects \(e\), updating tacit beliefs about vulnerabilities (e.g. ``severe scenario too mild in sector \(k\)'', ``underestimation of spread widening'', ``concentration risk inadequately captured'').
  \item \(H\) emits modified explicit artifacts \(e' \in E\), such as:
  \begin{quote}
    \emph{``add tail scenario with $\geq 99$th- 99th percentile spread shock in sector \(k\); flag valuations relying on proxy curves as high model risk; require separate line items for top-\(n\) counterparties.''}
  \end{quote}
\end{enumerate}
These outputs are captured by a human-feedback diffusion map
\[
  \mathrm{HFD}: E \times T \to E
\]
which updates scenario definitions, model-risk flags, and aggregation rules.

\paragraph{Phase 2 (Explicit pipeline and navigation).}
A typical trajectory in \(N \ltimes E\) for \(\tau_{\mathrm{SCEN}}\) is:

\begin{enumerate}
  \item \textbf{Scenario construction (\(\tau_1\)).}
  \(\mathrm{EC}\) constructs an initial scenario set \(S = \{s_{\mathrm{base}}, s_{\mathrm{adv}}, s_{\mathrm{sev}}\}\) with shock vectors over risk factors. After \(\mathrm{RB}+\mathrm{TE}\), \(H\) may deem \(s_{\mathrm{sev}}\) inadequate for certain illiquid exposures. Via \(\mathrm{HFD}\), this yields explicit constraints (e.g. minimum spread shocks relative to historical extremes, addition of a tail scenario \(s_{\mathrm{tail}}\)), and a revised set \(S' \supseteq S\).
  \item \textbf{Exposure mapping (\(\tau_2\)).}
  \(\mathrm{EC}\) constructs explicit mappings \(m \in E\) from positions to risk-factor buckets and term structures, marking incomplete or proxy-based mappings with data-quality flags. Tacit review identifies concentration and mapping issues; \(\mathrm{HFD}\) introduces rules such as ``flag scenario outputs where more than \(x\%\) of exposure is valued using proxies''.
  \item \textbf{Valuation under scenarios (\(\tau_3\)).}
  For each \(s \in S'\), \(\mathrm{EC}\) applies pricing models to obtain per-position P\&L and aggregated summaries \(v_s \in E\). Tacit concerns about model risk (e.g. oversimplified models for exotics) feed back through \(\mathrm{HFD}\) as explicit model-risk flags and conservative add-ons, giving adjusted valuations \(v_s'\).
  \item \textbf{Aggregation and reporting (\(\tau_4\)).}
  Using adjusted valuations, \(\mathrm{EC}\) constructs an explicit report \(r \in E\) with tables and charts (loss by sector, rating, counterparty; capital ratios; liquidity metrics) and textual summaries. Tacit review may detect that \(r\) understates uncertainty or uses overly reassuring language. \(\mathrm{HFD}\) then constrains report generation, e.g. requiring:
  \begin{itemize}
    \item explicit separation of robust vs model-dependent components;
    \item inclusion of ranges or confidence intervals where model risk is high;
    \item prohibition of unqualified adequacy statements, restricting text to numerical ratios vs regulatory minima.
  \end{itemize}
  A revised report \(r' \in E\) satisfying these constraints is then produced.
  \item \textbf{Navigation and canonicalization.}
  Navigation in \(N\) allows movement between individual scenarios, subsets of positions, and aggregation levels. Under confluence assumptions, the trajectory normalizes to a canonical artifact comprising a scenario set \(S^\ast\), consistent valuations \(\{v_s^\ast\}\), and a report \(r^\ast\) whose assertions are traceable to explicit quantitative inputs.
\end{enumerate}

\paragraph{Phase 3 (Hallucination dynamics, informal).}
Let \(t \in T\) denote the human's tacit risk state (understanding of portfolio vulnerabilities, model limitations, and macro uncertainty). Let
\[
  \mathrm{Req}(e,c)
\]
denote the tacit requirement to safely endorse an explicit claim \(e \in E\) in context \(c\) (e.g. internal exploratory analysis vs regulatory submission). A degradation operator
\[
  \delta: T \to T
\]
models erosion of vigilance due to time pressure, complexity, or anchoring to previous ``business-as-usual'' scenarios.

High-stakes claims such as
\begin{gather*}
  e_{\mathrm{ok}} := \text{``under our severe scenario set, group capital remains}\\
  \text{ adequate and no material vulnerabilities arise in sector \(k\)''}
\end{gather*}
are assigned large \(\mathrm{Req}(e_{\mathrm{ok}},c)\). Under extended interaction, if \(\delta\) repeatedly lowers \(t\) while \(\mathrm{EC}\), guided by earlier approvals, proposes increasingly reassuring summaries, it may occur that
\[
  t \prec \mathrm{Req}(e_{\mathrm{ok}},c)
\]
at the moment \(e_{\mathrm{ok}}\) is accepted. FRFP thus treats such overconfident adequacy statements as structural hallucination events: explicit assurances whose warrantedness exceeds the current tacit grounding.

Designing the scenario engine within FRFP then naturally leads to requirements such as:
\begin{itemize}
  \item elevated \(\mathrm{Req}(e,c)\) for capital-adequacy and risk-appetite statements;
  \item explicit severity and model-risk checks as preconditions for generating such statements;
  \item interaction-aware throttling of strong claims when usage patterns suggest tacit degradation (e.g. rapid approvals, minimal drill-down).
\end{itemize}

\medskip

\noindent In this case, FRFP yields (i) a sharp separation between explicit scenario modeling and tacit risk judgment, (ii) compositional and auditable scenario pipelines across desks as paths in \(N \ltimes E\), and (iii) a structural explanation of how overly comforting risk narratives can emerge and how to constrain them at the architectural level.

%\subsection{ICU Sepsis Prediction Literature Review} % old J.1
%\label{app:sepsis-literature}
% pipeline structure, explicit/tacit roles, navigation, decision points

%\subsection{Financial Scenario Analysis and Stress Testing} % old J.2 + I.4
%\label{app:finance-stress}
% scenario generation, explicit pipelines, tacit acceptance, safe horizons picture

%\subsection{Sepsis Consult as an FRFP Pipeline} % old A.7
%\label{app:sepsis-pipeline}
% the detailed worked pipeline that used to sit in the math appendix

\pagebreak
\section{Background on Explicit/Tacit Knowledge and Human Oversight}
\label{app:background}
%\section{Explicit and Tacit Knowledge, Human Oversight, and AI Decision Support}
%\label{app:explicit-tacit-survey}

This appendix sketches background on (i) the explicit/tacit distinction in epistemology and organisational theory, and (ii) human oversight and decision support in AI systems. It is intended as context, not as a comprehensive review.

\subsection{Explicit and Tacit Knowledge in Epistemology}

Polanyi's work is the standard reference point for the explicit/tacit distinction. In \emph{The Tacit Dimension} \cite{Polanyi1966}, he argues that ``we can know more than we can tell'': an agent's ability to recognise patterns, exercise judgment, or perform skilled actions often depends on knowledge that is not fully articulable as explicit rules or propositions. Tacit knowledge, in this sense, is embodied and context-sensitive. It is typically acquired through practice and is expressed in the ability to do or recognise, rather than in the ability to state a complete algorithm.

Examples Polanyi and subsequent commentators point to include perceptual discrimination (recognising a face, a pathology, or a forged signature), motor skills (riding a bicycle, playing an instrument), and scientific or clinical judgment (weighing heterogeneous evidence and experience to arrive at a diagnosis or assessment). In such cases, attempts to write down all the relevant ``rules'' usually fall short of what experienced practitioners can actually do.

Explicit knowledge, by contrast, comprises formalised statements, procedures, and representations that can be recorded, transmitted, and stored independently of the knower: theories, theorems, proofs, algorithms, protocols, manuals. Polanyi emphasises that explicit and tacit are not two disjoint kinds of knowledge. Explicit formulations presuppose tacit competences for their correct use and interpretation; tacit background skills and commitments are involved in reading, applying, and critiquing explicit statements. Later epistemological work refines these ideas, but broadly retains the view that there is a meaningful distinction between codified, symbol-manipulable representations and the background know-how on which their use depends.

\subsection{Explicit and Tacit Knowledge in Organisations}

Organisational and management literatures have adopted the explicit/tacit split to analyse how knowledge is created, stored, and shared in firms. Nonaka's dynamic theory of organisational knowledge creation \cite{Nonaka1994} treats explicit vs.\ tacit as an ``epistemological dimension'' along which knowledge can be transformed. The SECI model (Socialisation, Externalisation, Combination, Internalisation) distinguishes four modes:

\begin{itemize}
  \item \emph{Socialisation} (tacit-to-tacit): transfer of know-how through shared experience and observation, for example apprenticeship and on-the-job training.
  \item \emph{Externalisation} (tacit-to-explicit): articulation of aspects of tacit knowledge into concepts, models, or documents.
  \item \emph{Combination} (explicit-to-explicit): recombination and systematisation of existing documents, data, and theories.
  \item \emph{Internalisation} (explicit-to-tacit): incorporation of explicit materials into skills and routines through practice and learning.
\end{itemize}

In \cite{NonakaTakeuchi1995}, Nonaka and Takeuchi develop this into an account of the ``knowledge-creating company''. They describe how organisations can be structured to support SECI cycles: designing roles, artefacts, and processes that make it easier to externalise local know-how, combine codified knowledge, and internalise explicit materials into tacit competencies. In this literature, explicit knowledge appears in manuals, specifications, databases, and standard operating procedures. Tacit knowledge appears in craft skills, situational judgment, shared mental models, and organisational culture. Both are treated as important; neither is assumed to be reducible to the other in general.

Subsequent empirical work has tried to measure or proxy these processes (for example, via surveys of knowledge-sharing practices, analyses of documentation and training regimes, or case studies of innovation projects). These studies typically converge on a picture in which:

\begin{itemize}
  \item explicit knowledge is codified, relatively easy to transmit and recombine, and amenable to storage and search; and
  \item tacit knowledge is embodied, context-dependent, and often only partially articulable, with important components carried in practice, habits, and social structures.
\end{itemize}

The explicit/tacit distinction is used to explain phenomena such as the difficulty of transferring expertise purely through documents, the role of informal communities of practice, and the limits of purely document-centric knowledge management.

\subsection{Measuring and Externalising Tacit Knowledge}

Because tacit knowledge is, by definition, not fully articulated, there is a line of work on how far it can be measured or externalised. Approaches include structured interviews, concept mapping, elicitation of decision rationales, and analysis of patterns in performance. A common pattern is that such methods can capture \emph{some} aspects of tacit knowledge---for example, typical failure modes, heuristics, or rules of thumb—while leaving other aspects implicit. In organisational contexts, this motivates hybrid strategies in which organisations both invest in documentation and formal processes, and recognise the need for shadowing, mentoring, and informal knowledge transfer.

There are also critical discussions of the explicit/tacit dichotomy. Some authors suggest that knowledge lies on a spectrum of articulability rather than splitting neatly into two modes; others note that labelling something as ``tacit'' can obscure social and institutional factors, such as how expertise is recognised and who is allowed to act as a source of authoritative judgment. Despite these critiques, explicit vs.\ tacit remains a widely used distinction in empirical and conceptual work on knowledge, expertise, and organisational learning.
\subsection{Tacit Space in Explicit-Correctness Regimes}

A natural question is whether there are domains where tacit correctness can be dispensed with entirely and
FRFP’s explicit/tacit split is an artefact of modelling.
Interactive theorem provers, formally verified software components, and benchmark suites with sound and complete
test oracles are tempting examples: once a calculus, specification, or test harness has been fixed, correctness for
the inner task (“Does this proof term check?”, “Does this program satisfy the contract?”, “Does this output match
the reference?”) can be decided by a purely explicit procedure.

In FRFP terms, these are cases where correctness is \emph{tacit-insensitive} for a particular inner layer.
Conditional on prior tacit choices—selecting the logic, designing the specification or tests, deciding that conformance
to them is an adequate notion of success—the induced correctness map
\(\gamma_{\mathrm{inner}} : E \to J\) behaves like an explicit oracle.
For that inner layer it is mathematically convenient to collapse T and work entirely in explicit space.

However, the surrounding tacit space does not disappear.
Tacit judgment is still required to:
\begin{itemize}
  \item decide which calculus or modelling framework is appropriate for the domain;
  \item construct and revise the specifications, contracts, or benchmark suites;
  \item assess whether “passes all tests” or “has a formal proof” is an adequate guarantee for the real-world purpose;
  \item arbitrate between competing formalizations when they disagree about what is acceptable.
\end{itemize}

From the FRFP point of view, the explicit/tacit split therefore has two levels:
(i) inner explicit-correctness regimes, where correctness is decidable given a fixed tacit backdrop,
and (ii) outer scientific, regulatory, or governance workflows, where that backdrop itself evolves under
tacit evaluation and institutional feedback.
The impossibility and inevitability results in the main text are explicitly about this outer level:
they assume a non-trivial, dynamically involved tacit space T and show that, under those assumptions,
no explicit-only detector, consensus mechanism, or governance operator can eliminate hallucination or fully automate
correctness across arbitrary long-horizon workflows.

\subsection{Human Oversight in AI Law and Policy}

In the governance of AI systems, human oversight has become a central requirement in high-stakes domains. The EU Artificial Intelligence Act classifies certain systems (for example, those used in critical infrastructure, employment, credit, or law enforcement) as ``high-risk'' and requires that they be designed so that they can be effectively overseen by natural persons (Article~14) \cite{EUAIAct2024}. Oversight measures must enable human operators to understand the system's capabilities and limitations, monitor its operation, and intervene to override or stop it when necessary to prevent or minimise risks to health, safety, or fundamental rights.

Analyses of these provisions typically emphasise that, even when AI systems provide sophisticated predictions or recommendations, accountability and ultimate responsibility are intended to rest with human and institutional actors, not with the systems themselves. Human overseers are expected to interpret AI outputs in context, apply legal and ethical standards, and take corrective action when outputs are inappropriate. At the same time, commentators highlight practical challenges: the risk that humans simply ``rubber-stamp'' AI decisions; the cognitive load of continuous monitoring; and the need for interfaces and organisational arrangements that make oversight feasible rather than nominal.

Similar themes appear in sector-specific guidelines for automated decision-making in areas such as credit, hiring, and healthcare. These typically call for human review of important decisions, mechanisms for contesting and correcting automated outputs, and clear assignment of responsibility when AI-assisted decisions lead to harm.

\subsection{Human--AI Interaction and Decision Support}

Human--computer interaction and human--AI interaction work often treats AI systems as decision-support tools rather than self-standing decision-makers. Amershi et al.\ \cite{Amershi2019}, for example, propose a set of guidelines for human--AI interaction based on a survey of deployed systems and empirical studies. Their guidelines emphasise, among other points, that systems should:

\begin{itemize}
  \item make it possible for users to understand what the system can and cannot do;
  \item expose, where appropriate, information about confidence and uncertainty;
  \item support efficient correction of errors; and
  \item enable users to take or regain control at key moments, rather than forcing them to accept or reject entire bundles of automated behaviour.
\end{itemize}

A separate empirical literature studies trust and reliance in AI-assisted decision-making. Typical experimental setups involve a human decision-maker receiving recommendations, predictions, or rationales from an AI system, and measuring when the human follows or overrides the advice. These studies document both under-reliance (ignoring useful advice) and over-reliance (accepting poor advice), and identify factors such as expertise, presentation of uncertainty, and prior experience with the system as important moderators. The notion of ``appropriate reliance''—using AI outputs when they are informative, overriding them when they are not—is a common target.

Across these discussions, AI systems are typically described as producing explicit artefacts: predictions, scores, labels, suggested actions, explanations, or visualisations. Human decision-makers are expected to integrate these artefacts with other information (legal constraints, domain knowledge, local context) and to bear responsibility for final decisions. Much of the design guidance is concerned with making this human role feasible, whether by improving calibration of trust, reducing cognitive load, or structuring workflows so that human review occurs at appropriate points.

\subsection{Summary}

Across these literatures, two distinctions recur. First, the explicit/tacit split is used to analyse how knowledge is represented, transmitted, and used: explicit knowledge is codified, symbolically manipulable, and transmissible; tacit knowledge is embodied, context-dependent, and often only partially articulable, with explicit and tacit modes interacting in practice \cite{Polanyi1966,Nonaka1994,NonakaTakeuchi1995}. Second, in AI law, policy, and interaction design, a separation is maintained between automated outputs and human oversight: AI systems are typically cast as producers of explicit artefacts (predictions, rankings, recommendations, explanations), while humans are expected to interpret those artefacts in context and retain ultimate responsibility for decisions \cite{EUAIAct2024,Amershi2019}. This appendix does not take a position on how far these distinctions should be pushed, but records that they are both conceptually and institutionally well established.


\paragraph{\textbf{Declaration of generative AI and AI-assisted technologies in the manuscript preparation process}}

During the preparation of this work the author(s) used CHATGPT in order to develop the research and prepare paper. After using this tool/service, the author(s) reviewed and edited the content as needed and take(s) full responsibility for the content of the published article.

%%\section{FRFP for Organizations: Business and Governance Overview}
\label{app:business-frfp}

This appendix presents a practitioner-oriented version of the Foundational Reasoning and Feedback Protocol (FRFP).
The main body of the paper develops FRFP as a formal architecture for long-horizon human--AI scientific workflows, with explicit/tacit categories, a boundary morphism, and dynamic and institutional semantics.
Here, we restate the core ideas in non-technical terms and sketch how they can guide concrete organizational practice.

The intended audience for this appendix includes:
senior technical leaders (e.g.\ Chief Scientist, Head of AI),
risk and governance teams,
and product owners in domains where AI systems support high-stakes or long-horizon workflows.

\subsection{Executive Overview}

Organizations increasingly rely on AI systems---especially large language models and tool-using assistants---for scientific research, product development, financial analysis, and clinical or regulatory decision support.
These systems can greatly accelerate work, but they also create subtle new failure modes:
over long workflows, humans may become over-reliant on AI outputs, tacit grounding can degrade, and institutions can lose clear lines of responsibility.

FRFP offers three core ideas for structuring such workflows:

\begin{enumerate}
  \item \textbf{Explicit vs.\ tacit.}
  We distinguish between \emph{explicit} artifacts (documents, code, data, prompts, model outputs) and \emph{tacit} human understanding (expertise, contextual judgment, the sense that something is ``off'').
  AI systems operate only on explicit artifacts; correctness lives in tacit human states.

  \item \textbf{A single actionable boundary.}
  The key architectural feature is a controlled boundary between explicit and tacit spaces.
  Crossing this boundary corresponds to a human \emph{accepting} an explicit artifact as ``good enough to act on'' in some context (signing a report, approving a decision, submitting a paper).
  FRFP insists that AI systems do not cross this boundary: they may propose and transform explicit artifacts, but they do not have correctness authority.

  \item \textbf{Safe horizons instead of zero hallucination.}
  Over long horizons, humans who lean heavily on AI tend to check less, remember less, and rely more on the system.
  FRFP models this as degradation of tacit grounding and shows that, under realistic conditions, hallucinations (accepting ungrounded outputs) are inevitable along sufficiently long runs.
  The goal is therefore not ``no hallucinations ever'', but carefully designed \emph{safe horizons}: limits on how far a workflow can run before deeper human review, re-grounding, or escalation is required.
\end{enumerate}

From an organizational perspective, FRFP yields:
(i) architectural clarity about who may do what,
(ii) governance constraints and impossibility results that explain which roles cannot be fully automated,
and (iii) concrete patterns and checklists for designing AI-supported workflows that remain auditable and governable as capabilities scale.

\subsection{FRFP Architecture in Organizational Terms}

\paragraph{Explicit work.}
Explicit work consists of everything that can be written down, stored, and processed by machines:
prompts, model outputs, code, data tables, dashboards, reports, and so on.
This is where AI systems operate most naturally: generating, transforming, and organizing explicit artifacts.

\paragraph{Tacit work.}
Tacit work consists of human understanding and judgment:
domain expertise, sense of context, knowledge of institutional constraints, and the intuitive ability to say ``this looks wrong''.
Correctness---in the sense relevant for scientific, financial, or clinical responsibility---is treated as a property of tacit states, not of explicit artifacts alone.

\paragraph{The boundary.}
The FRFP boundary is the moment when a human accepts or acts on explicit material influenced by AI:
signing a clinical note, approving a risk position, submitting an AI-assisted manuscript, or green-lighting a model deployment.
At these boundary points, FRFP requires that:
(i) the responsible human role is clearly identified,
(ii) the information available to that role is specified,
and (iii) correctness is understood as their responsibility, not the system's.

AI tools may contribute to everything that happens \emph{before} the boundary (e.g.\ summaries, scenarios, drafts, analyses), but do not themselves perform the boundary crossing.

\paragraph{Roles and responsibilities.}
In an FRFP-aligned organization, it is helpful to distinguish four functional roles:

\begin{itemize}
  \item \emph{Producers of explicit artifacts:}
  teams using AI tools to generate and transform documents, models, analyses, or proposals.

  \item \emph{Tacit owners / decision-makers:}
  individuals who formally accept or reject explicit artifacts at boundary points, and who own correctness for those decisions.

  \item \emph{Institutional aggregators:}
  committees, review boards, and governance bodies that combine explicit artifacts and tacit judgments across many individuals (e.g.\ model risk committees, clinical guideline committees, internal research review boards).

  \item \emph{AI and tooling stewards:}
  platform teams responsible for AI integration, logging, and ensuring that tools respect explicit/tacit boundaries and support auditability.
\end{itemize}

\paragraph{Safe horizons.}
Every AI-assisted workflow has a safe operating horizon: a limit on how far, or how long, the workflow can run before tacit grounding and error-detection degrade unacceptably.
In practice, organizations should specify, for high-stakes processes:

\begin{itemize}
  \item how many AI-assisted steps may occur before a deeper human review is mandatory;
  \item how often humans must re-ground themselves in primary data or source documents;
  \item what kinds of anomalies or thresholds trigger escalation.
\end{itemize}

FRFP shows that, under broad qualitative assumptions, infinite safe horizons are not achievable; instead, horizon management becomes a core design responsibility.

\subsection{Examples of FRFP-Aligned Workflows}

\paragraph{Scientific research and R\&D.}
In a research workflow, AI systems may assist with literature search, clustering, summarization, hypothesis generation, and drafting.
The explicit layer includes retrieved papers, AI-generated topic maps and summaries, and draft text.
Tacit work includes: judging whether the literature coverage is adequate, assessing plausibility of hypotheses, choosing experiments, and interpreting results.

FRFP suggests:
(i) defining boundary points such as pre-registration of hypotheses and final manuscript approval, where tacit owners (e.g.\ PIs) explicitly accept responsibility;
(ii) instituting checkpoints where researchers must directly consult primary sources rather than relying solely on AI summaries;
and (iii) logging key AI-assisted reasoning steps so that the intellectual provenance of conclusions can be reconstructed.

\paragraph{Regulated financial or clinical decisions.}
In finance, AI might generate stress-test scenarios, synthesize risk reports, or help draft credit memos.
In healthcare, it might summarize charts, suggest diagnoses, or propose treatment options.
In both cases, FRFP confines AI to the generation and transformation of explicit artifacts, while reserving acceptance of credit decisions, trading positions, or clinical orders to identified human roles.

Safe horizons may be expressed as limits on:
the number of successive AI-assisted reports before independent human re-baselining,
the extent to which AI-drafted text can be reused without full review,
or the duration for which an AI-supported monitoring process may run without human reassessment of assumptions.

\paragraph{Internal governance and oversight.}
Model governance, ethics boards, and risk committees increasingly use AI to summarize documentation and surface issues.
FRFP treats governance itself as a pipeline over explicit artifacts (model cards, evaluations, incident logs, AI-generated summaries) and tacit judgments (legal, ethical, and strategic interpretation by committee members).
Key governance actions---model approval, acceptance of residual risk, adoption of policies---are boundary crossings that must remain human-exclusive.

FRFP's institutional layer implies that fully automated governance (e.g.\ ``deploy if scores exceed threshold and no alerts fire'') cannot generally provide the same correctness guarantees as human governance, even with sophisticated monitoring.
AI may and should support governance, but not replace it.

\subsection{Implementation Roadmap}

From an implementation perspective, organizations can adopt FRFP in stages:

\begin{enumerate}
  \item \textbf{Assess and map.}
  Inventory where AI is currently used; identify existing (possibly implicit) boundary points and areas with no clear boundary.

  \item \textbf{Redesign key workflows.}
  For a small number of high-impact workflows, explicitly map explicit vs.\ tacit steps, define boundary roles and information requirements, and choose safe horizons and checkpoints.
  Decide what should be logged to reconstruct AI-assisted reasoning.

  \item \textbf{Pilot and refine.}
  Run FRFP-aligned versions of these workflows with motivated teams.
  Monitor incidents, overrides, and user experience; adjust boundaries, horizons, and logging based on what works in practice.

  \item \textbf{Standardize and scale.}
  Convert successful pilots into reusable FRFP templates (for research, risk decisions, clinical documentation, model governance, etc.).
  Integrate FRFP design into product lifecycles, model approval processes, and change management.

  \item \textbf{Embed in governance and training.}
  Update AI use policies, model governance frameworks, and risk documents to reflect explicit/tacit boundaries, human-exclusive correctness, and limits on automation.
  Train practitioners and leaders in FRFP concepts and their responsibilities at boundary points.
\end{enumerate}

\subsection{Common Pitfalls and How FRFP Helps Avoid Them}

In applying FRFP, several recurring pitfalls are worth highlighting:

\begin{itemize}
  \item \textbf{Over-automation disguised as decision support.}
  Tools described as ``assistants'' or ``copilots'' effectively become decision-makers when humans never override them.
  FRFP reacts by treating behavior, not labels, as the signal: if a system is de facto making decisions, boundaries must be redesigned.

  \item \textbf{Boundary erosion over time.}
  Initial review steps can become rubber-stamping under time pressure.
  FRFP recommends making boundary decisions observable (e.g.\ by logging changes and overrides) and periodically auditing samples to detect erosion.

  \item \textbf{Metric overconfidence.}
  Organizations may over-trust dashboards, scores, or evaluation metrics.
  Under FRFP, such metrics are inputs to tacit judgment, not substitutes; governance documents should explicitly state what metrics do \emph{not} capture.

  \item \textbf{Checklist theatre.}
  Long, low-value checklists can create a false sense of safety.
  FRFP encourages short, targeted prompts that trigger real reflection (e.g.\ ``What primary sources did you personally review?''), and regular pruning of items that do not change behavior.

  \item \textbf{Treating FRFP as pure compliance.}
  If FRFP is seen only as an external requirement, teams will optimize for appearance.
  Embedding FRFP in product and research strategy---not just policy---helps keep the focus on actual correctness and robustness.
\end{itemize}

\subsection{FRFP Design Checklist}

For ease of adoption, we collect here a short checklist that teams can use when designing or reviewing AI-supported workflows.

\begin{itemize}
  \item \textbf{Workflow mapping:}
    have we charted the main steps and labeled each as explicit or tacit?

  \item \textbf{Boundaries:}
    where exactly are the decision points where correctness matters,
    who is responsible at each, and what information must they see?

  \item \textbf{AI's role:}
    are AI systems confined to explicit work (generation, transformation, organization),
    or are they implicitly making decisions that should be human-held?

  \item \textbf{Safe horizons:}
    for this workflow, how far can AI-assisted reasoning proceed before requiring deeper human review, re-grounding, or escalation?
    what are the triggers for resets?

  \item \textbf{Logging and traceability:}
    are prompts, key AI outputs, and boundary decisions logged in a way that lets us reconstruct how a conclusion was reached?

  \item \textbf{Governance:}
    which committee or body oversees this workflow?
    do they understand the explicit/tacit boundaries and have the authority to adjust them?

  \item \textbf{Training and culture:}
    have the people using the workflow been trained in their responsibilities at boundary points and encouraged to challenge AI outputs when needed?
\end{itemize}

Used consistently, this checklist can help ensure that FRFP is not just a theoretical architecture but a practical guide to designing AI-supported work that is safer, more auditable, and more robust against long-horizon drift.

%\pagebreak

\reftitle{References}

% Please provide either the correct journal abbreviation (e.g. according to the “List of Title Word Abbreviations” http://www.issn.org/services/online-services/access-to-the-ltwa/) or the full name of the journal.
% Citations and References in Supplementary files are permitted provided that they also appear in the reference list here. 

%=====================================
% References, variant A: external bibliography
%=====================================
% \bibliography{your_external_BibTeX_file}

%=====================================
% References, variant B: internal bibliography
%=====================================

% ACS format

\begin{thebibliography}{999}

%\begin{thebibliography}{99}

\bibitem{MacLane1998}
S.~Mac~Lane.
\newblock \emph{Categories for the Working Mathematician}, 2nd~edition.
\newblock Graduate Texts in Mathematics 5. Springer, 1998.

\bibitem{Awodey2010}
S.~Awodey.
\newblock \emph{Category Theory}, 2nd~edition.
\newblock Oxford Logic Guides 52. Oxford University Press, 2010.

\bibitem{Borceux1994}
F.~Borceux.
\newblock \emph{Handbook of Categorical Algebra, Volume~1: Basic Category Theory}.
\newblock Encyclopedia of Mathematics and its Applications 50. Cambridge University Press, 1994.

\bibitem{BaaderNipkow1998}
F.~Baader and T.~Nipkow.
\newblock \emph{Term Rewriting and All That}.
\newblock Cambridge University Press, 1998.

\bibitem{Ji2023Survey}
Z.~Ji, N.~Lee, R.~Frieske, T.~Yu, D.~Su, Y.~Xu, E.~Ishii, Y.~Bang, D.~Chen, W.~Dai, H.~S.~Chan, A.~Madotto, and P.~Fung.
\newblock Survey of hallucination in natural language generation.
\newblock \emph{ACM Computing Surveys}, 55(12):Article~248, 2023.

\bibitem{Farquhar2024SemanticEntropy}
S.~Farquhar, J.~Kossen, L.~Kuhn, and Y.~Gal.
\newblock Detecting hallucinations in large language models using semantic entropy.
\newblock \emph{Nature}, 630(8017):625--630, 2024.

\bibitem{Kadavath2022Know}
S.~Kadavath, T.~Conerly, A.~Askell, T.~Henighan, D.~Drain, E.~Perez, N.~Schiefer, Z.~Hatfield-Dodds, N.~DasSarma, E.~Tran-Johnson, S.~Johnston, S.~El-Showk, A.~Jones, N.~Elhage, T.~Hume, A.~Chen, Y.~Bai, S.~Bowman, S.~Fort, D.~Ganguli, D.~Hernandez, J.~Jacobson, J.~Kernion, S.~Kravec, L.~Lovitt, K.~Ndousse, C.~Olsson, S.~Ringer, D.~Amodei, T.~Brown, J.~Clark, N.~Joseph, B.~Mann, S.~McCandlish, C.~Olah, and J.~Kaplan.
\newblock Language models (mostly) know what they know.
\newblock \emph{arXiv preprint} arXiv:2207.05221, 2022.

\bibitem{Kalai2025WhyHallucinate}
A.~T.~Kalai, O.~Nachum, S.~S.~Vempala, and E.~Zhang.
\newblock Why language models hallucinate.
\newblock \emph{arXiv preprint} arXiv:2509.04664, 2025.

\bibitem{BertotCasteran2004CoqArt}
Y.~Bertot and P.~Cast\'eran.
\newblock \emph{Interactive Theorem Proving and Program Development: Coq'Art: The Calculus of Inductive Constructions}.
\newblock Texts in Theoretical Computer Science. An EATCS Series. Springer, 2004.

\bibitem{Nipkow2002IsabelleHOL}
T.~Nipkow, L.~C.~Paulson, and M.~Wenzel.
\newblock \emph{Isabelle/HOL: A Proof Assistant for Higher-Order Logic}.
\newblock Lecture Notes in Computer Science 2283. Springer, 2002.

\bibitem{deMoura2015Lean}
L.~de~Moura, S.~Kong, J.~Avigad, F.~van~Doorn, and J.~von~Raumer.
\newblock The {L}ean theorem prover (system description).
\newblock In A.~P.~Felty and A.~Middeldorp, editors, \emph{Automated Deduction -- CADE-25}, volume~9195 of \emph{Lecture Notes in Computer Science}, pages 378--388. Springer, 2015.

\bibitem{Jacobs2001Categorical}
B.~Jacobs.
\newblock \emph{Categorical Logic and Type Theory}.
\newblock Studies in Logic and the Foundations of Mathematics 141. Elsevier, 2001.

\bibitem{Schick2023Toolformer}
T.~Schick, J.~Dwivedi-Yu, R.~Dess\`i, R.~Raileanu, M.~Lomeli, L.~Zettlemoyer, N.~Cancedda, and T.~Scialom.
\newblock Toolformer: language models can teach themselves to use tools.
\newblock \emph{arXiv preprint} arXiv:2302.04761, 2023.

\bibitem{Yao2022ReAct}
S.~Yao, J.~Zhao, D.~Yu, N.~Du, I.~Shafran, K.~Narasimhan, and Y.~Cao.
\newblock ReAct: synergizing reasoning and acting in language models.
\newblock \emph{arXiv preprint} arXiv:2210.03629, 2022.

\bibitem{Askell2021Alignment}
A.~Askell, Y.~Bai, A.~Chen, D.~Drain, D.~Ganguli, T.~Henighan, A.~Jones, N.~Joseph, B.~Mann, N.~DasSarma, N.~Elhage, Z.~Hatfield-Dodds, D.~Hernandez, J.~Kernion, K.~Ndousse, C.~Olsson, D.~Amodei, T.~Brown, J.~Clark, S.~McCandlish, C.~Olah, and J.~Kaplan.
\newblock A general language assistant as a laboratory for alignment.
\newblock \emph{arXiv preprint} arXiv:2112.00861, 2021.

\bibitem{Shao2025SciSciGPT}
E.~Shao, Y.~Wang, Y.~Qian, Z.~Pan, H.~Liu, and D.~Wang.
\newblock SciSciGPT: advancing human--AI collaboration in the science of science.
\newblock \emph{Nature Computational Science}, 5(11), 2025. Also available as \emph{arXiv preprint} arXiv:2504.05559.

\bibitem{Vaccaro2024WhenUseful}
M.~Vaccaro, A.~Almaatouq, and T.~W.~Malone.
\newblock When combinations of humans and AI are useful: a systematic review and meta-analysis.
\newblock \emph{Nature Human Behaviour}, 8(12):2293--2303, 2024.

\bibitem{Zheng2025LLMDiscovery}
T.~Zheng, Z.~Deng, H.~T.~Tsang, W.~Wang, J.~Bai, Z.~Wang, and Y.~Song.
\newblock From automation to autonomy: a survey on large language models in scientific discovery.
\newblock In \emph{Proceedings of the 2025 Conference on Empirical Methods in Natural Language Processing}, pages 17744--17761. Association for Computational Linguistics, 2025.

\bibitem{Zhao2025HumanAISciDiscovery}
S.~Zhao, Y.~Liu, J.~Wan, T.~Tang, and X.~Li.
\newblock Human--AI collaboration for scientific discovery.
\newblock In D.~Wu and S.~Liang, editors, \emph{Human--AI Interaction and Collaboration}, pages 239--267. Cambridge University Press, 2025.

\bibitem{Channing2025AISocialProblem}
G.~Channing and A.~Ghosh.
\newblock AI for scientific discovery is a social problem.
\newblock \emph{arXiv preprint} arXiv:2509.06580, 2025.
% --- Hallucination / LLM surveys ---

\bibitem{Zhang2023Siren}
Y.~Zhang, Y.~Li, L.~Cui, D.~Cai, L.~Liu, T.~Fu, X.~Huang, E.~Zhao, Y.~Zhang,
Y.~Chen, L.~Wang, A.~T.~Luu, W.~Bi, F.~Shi, and S.~Shi.
Siren's song in the AI ocean: a survey on hallucination in large language models.
arXiv preprint arXiv:2309.01219, 2023.

\bibitem{Huang2025HallucinationSurvey}
L.~Huang, W.~Yu, W.~Ma, W.~Zhong, Z.~Feng, H.~Wang, Q.~Chen, W.~Peng,
X.~Feng, B.~Qin, and T.~Liu.
A survey on hallucination in large language models: principles, taxonomy, challenges, and open questions.
ACM Transactions on Information Systems, 43(2):1--55, 2025.

% --- AI reliance / appropriate reliance / trust ---

\bibitem{Eckhardt2024Reliance}
S.~Eckhardt, N.~K{\"u}hl, M.~Dolata, and G.~Schwabe.
A survey of AI reliance.
CoRR, abs/2408.03948, 2024.

\bibitem{Schemmer2023Appropriate}
M.~Schemmer, N.~K{\"u}hl, C.~Benz, A.~Bartos, and G.~Satzger.
Appropriate reliance on AI advice: conceptualization and the effect of explanations.
In \emph{Proceedings of the 28th International Conference on Intelligent User Interfaces (IUI~'23)},
pages 1--21. ACM, 2023.

\bibitem{Visser2025Trust}
R.~Visser, T.~M.~Peters, I.~Scharlau, and B.~Hammer.
Trust, distrust, and appropriate reliance in (X)AI: a conceptual clarification of user trust and survey of its empirical evaluation.
Cognitive Systems Research, 82:101357, 2025.

% --- Social epistemology / scientific networks ---

\bibitem{Zollman2013Network}
K.~J.~S.~Zollman.
Network epistemology: communication in epistemic communities.
Philosophy Compass, 8(1):15--27, 2013.

\bibitem{Grim2013ScientificNetworks}
P.~Grim, D.~J.~Singer, S.~Fisher, A.~Bramson, W.~J.~Berger, C.~Reade, C.~Flocken, and A.~Sales.
Scientific networks on data landscapes: question difficulty, epistemic success, and convergence.
Episteme, 10(4):441--464, 2013.

% --- AI governance / frontier policies ---

\bibitem{METRCommonElements}
METR.
Common elements of frontier AI safety policies.
Technical report, Model Evaluation and Threat Research (METR), 2024.
Available at \texttt{https://metr.org/common-elements}.

\bibitem{Polanyi1966}
M.~Polanyi.
\newblock \emph{The Tacit Dimension}.
\newblock Routledge \& Kegan Paul, London, 1966.

\bibitem{Nonaka1994}
I.~Nonaka.
\newblock A dynamic theory of organizational knowledge creation.
\newblock \emph{Organization Science}, 5(1):14--37, 1994.

\bibitem{NonakaTakeuchi1995}
I.~Nonaka and H.~Takeuchi.
\newblock \emph{The Knowledge-Creating Company: How Japanese Companies Create the Dynamics of Innovation}.
\newblock Oxford University Press, New York, 1995.

\bibitem{EUAIAct2024}
European Parliament and Council of the European Union.
\newblock Regulation (EU) 2024/1689 laying down harmonised rules on artificial intelligence (AI Act).
\newblock \emph{Official Journal of the European Union}, L~1689, 2024.
% Optional: add URL if your style allows it:
% \newblock Available at \url{https://eur-lex.europa.eu/legal-content/EN/TXT/?uri=OJ:L:2024:1689:TOC}.

\bibitem{Amershi2019}
S.~Amershi, D.~Weld, M.~Vorvoreanu, A.~Fourney, B.~Nushi, P.~Collisson, H.~Suh,
S.~Iqbal, P.~N.~Bennett, K.~Inkpen, and J.~Teevan.
\newblock Guidelines for human--AI interaction.
\newblock In \emph{Proceedings of the 2019 CHI Conference on Human Factors in Computing Systems}, pages 1--13.
\newblock ACM, 2019.
\bibitem{brown_effect_automation_aviation_2017}
J.~P. Brown.
\newblock The effect of automation on human factors in aviation.
\newblock \emph{The Journal of Instrumentation, Automation and Systems}, 3(2):31--46, 2016.
% DOI: 10.21535/jias.v3i2.916

\bibitem{okine_evolution_hf_aviation_2025}
E.~A. Okine, E.~Zarei, B.~J. Roggow, and N.~Dehghan.
\newblock Evolution of human factors research in aviation safety: a systematic review and bibliometric analysis of the intellectual structure.
\newblock \emph{Journal of Safety Science and Resilience}, 7(1):100249, 2026.
% DOI: 10.1016/j.jnlssr.2025.100249

\bibitem{singh_automation_complacency_2024}
I.~L. Singh, R.~Molloy, and R.~Parasuraman.
\newblock Automation-induced ``complacency'': development of the Complacency-Potential Rating Scale.
\newblock \emph{The International Journal of Aviation Psychology}, 3(2):111--122, 1993.
% Key reused to match the LaTeX patch; classic complacency measure

\bibitem{natali_ai_deskilling_medicine_2025}
C.~Natali, L.~Marconi, L.~D.~D. Duran, and F.~Cabitza.
\newblock AI-induced deskilling in medicine: a mixed-method review and research agenda for healthcare and beyond.
\newblock \emph{Artificial Intelligence Review}, 58:356, 2025.
% DOI: 10.1007/s10462-025-11352-1

\bibitem{berzin_preserving_clinical_skills_2025}
T.~M. Berzin and E.~J. Topol.
\newblock Preserving clinical skills in the age of AI assistance.
\newblock \emph{The Lancet}, 406(10513):1719, 2025.
% DOI: 10.1016/S0140-6736(25)02075-6

\bibitem{sunday_behavioral_impacts_diagnostics_2025}
O.~Sunday.
\newblock Behavioral impacts of AI reliance in diagnostics: balancing automation with skill retention.
\newblock \emph{Epidemiology \& Health Data Insights}, 1(3):ehdi011, 2025.
% DOI: 10.63946/ehdi/16894

\bibitem{fatima_ai_triage_deskilling_2025}
L.~Fatima, A.~A. Khan, R.~Darwesh, and M.~un Nissa.
\newblock AI triage for stroke imaging: promise vs risk of deskilling radiology trainees.
\newblock \emph{INNOVAPATH}, 1(8):2, 2025.
% DOI: 10.63501/0961v416

\bibitem{savardi_ai_radiology_training_2025}
M.~Savardi, A.~Signoroni, F.~Arrigoni, et~al.
\newblock Upskilling or deskilling? measurable role of an AI-supported training for radiology residents: a lesson from the pandemic.
\newblock \emph{Insights into Imaging}, 16(1):23, 2025.
% DOI: 10.1186/s13244-024-01893-4

\bibitem{gong_yang_google_effect_2024}
C.~Gong and Z.~Yang.
\newblock Google effects on memory: a meta-analytical review of the media effects of intensive Internet search behavior.
\newblock \emph{Frontiers in Public Health}, 12:1332030, 2024.
% DOI: 10.3389/fpubh.2024.1332030

\bibitem{benge_tech_use_cognitive_aging_2025}
J.~F. Benge and M.~K. Scullin.
\newblock A meta-analysis of technology use and cognitive aging.
\newblock \emph{Nature Human Behaviour}, 9(7):1405--1419, 2025.
% DOI: 10.1038/s41562-025-02159-9

\bibitem{gerlich_ai_tools_cognitive_offloading_2025}
M.~Gerlich.
\newblock AI tools in society: impacts on cognitive offloading and the future of critical thinking.
\newblock \emph{Societies}, 15(1):6, 2025.
% DOI: 10.3390/soc15010006

\bibitem{gerlich_ai_cognitive_implications_2025}
R.~Loga.
\newblock AI's cognitive implications: the decline of our thinking skills?
\newblock IE Center for Health and Well-Being Blog, IE University, 26 February 2025.
% Blog post synthesizing Gerlich (2025) and related work; keep only if you explicitly cite it

\bibitem{kosmyna_brain_on_chatgpt_2025}
N.~Kosmyna, E.~Hauptmann, Y.~T. Yuan, J.~Situ, X.-H. Liao, A.~V. Beresnitzky, I.~Braunstein, and P.~Maes.
\newblock Your brain on ChatGPT: accumulation of cognitive debt when using an AI assistant for essay writing task.
\newblock arXiv preprint arXiv:2506.08872, 2025.

\bibitem{zhao_genai_higher_order_thinking_2025}
Y.~Zhao, Y.~Yue, Z.~Sun, Q.~Jiang, and G.~Li.
\newblock Does generative artificial intelligence improve students' higher-order thinking? a meta-analysis based on 29 experiments and quasi-experiments.
\newblock \emph{Journal of Intelligence}, 13(12):160, 2025.
% DOI: 10.3390/jintelligence13120160

% Optional: add URL if your style allows it:
% \newblock Available at \url{https://www.ie.edu/center-for-health-and-well-being/blog/ais-cognitive-implications-the-decline-of-our-thinking-skills/}.
\bibitem{laban_llms_get_lost_multiturn_2025}
P.~Laban, H.~Hayashi, Y.~Zhou, and J.~Neville.
\newblock LLMs get lost in multi-turn conversation.
\newblock \emph{arXiv preprint arXiv:2505.06120}, 2025.

\bibitem{dongre_drift_no_more_context_equilibria_2025}
V.~Dongre, R.~A.~Rossi, V.~D.~Lai, D.~S.~Yoon, D.~Hakkani-T\"ur, and T.~Bui.
\newblock Drift no more? Context equilibria in multi-turn LLM interactions.
\newblock \emph{arXiv preprint arXiv:2510.07777}, 2025.

\bibitem{arike_evaluating_goal_drift_agents_2025}
R.~Arike, E.~Donoway, H.~Bartsch, and M.~Hobbhahn.
\newblock Technical report: evaluating goal drift in language model agents.
\newblock \emph{arXiv preprint arXiv:2505.02709}, 2025.

\bibitem{choi_identity_drift_llm_agents_2024}
J.~Choi, Y.~Hong, M.~Kim, and B.~Kim.
\newblock Examining identity drift in conversations of LLM agents.
\newblock \emph{arXiv preprint arXiv:2412.00804}, 2024.
\bibitem{nrc_defense_in_depth_glossary}
U.S. Nuclear Regulatory Commission.
\newblock Defense in depth.
\newblock NRC Glossary, n.d. (accessed 26 January 2026).
% https://www.nrc.gov/reading-rm/basic-ref/glossary/defense-in-depth

\bibitem{iaea_insag10_defence_in_depth_1996}
International Nuclear Safety Advisory Group (INSAG).
\newblock \emph{Defence in Depth in Nuclear Safety} (INSAG-10).
\newblock International Atomic Energy Agency, Vienna, 1996.
% IAEA Pub 1013; ISBN 9201025963.

\bibitem{reason_human_error_models_management_2000}
J.~Reason.
\newblock Human error: models and management.
\newblock \emph{BMJ}, 320(7237):768--770, 2000.
% DOI: 10.1136/bmj.320.7237.768

\bibitem{reason_latent_human_failures_complex_systems_1990}
J.~Reason.
\newblock The contribution of latent human failures to the breakdown of complex systems.
\newblock \emph{Philosophical Transactions of the Royal Society of London. Series B, Biological Sciences}, 327(1241):475--484, 1990.
% DOI: 10.1098/rstb.1990.0090

\end{thebibliography}

\end{document}
%%%%%%%%%%%%%%%%%%%%%%%%%%%%%%%%%%%%%%%%%%
\section{Materials and Methods}

Materials and Methods should be described with sufficient details to allow others to replicate and build on published results. Please note that publication of your manuscript implicates that you must make all materials, data, computer code, and protocols associated with the publication available to readers. Please disclose at the submission stage any restrictions on the availability of materials or information. New methods and protocols should be described in detail while well-established methods can be briefly described and appropriately cited.

Research manuscripts reporting large datasets that are deposited in a publicly avail-able database should specify where the data have been deposited and provide the relevant accession numbers. If the accession numbers have not yet been obtained at the time of submission, please state that they will be provided during review. They must be provided prior to publication.

Interventionary studies involving animals or humans, and other studies require ethical approval must list the authority that provided approval and the corresponding ethical approval code.
\begin{quote}
This is an example of a quote.
\end{quote}

%%%%%%%%%%%%%%%%%%%%%%%%%%%%%%%%%%%%%%%%%%
\section{Results}

This section may be divided by subheadings. It should provide a concise and precise description of the experimental results, their interpretation as well as the experimental conclusions that can be drawn.
\subsection{Subsection}
\subsubsection{Subsubsection}

Bulleted lists look like this:
\begin{itemize}
\item	First bullet;
\item	Second bullet;
\item	Third bullet.
\end{itemize}

Numbered lists can be added as follows:
\begin{enumerate}
\item	First item; 
\item	Second item;
\item	Third item.
\end{enumerate}

The text continues here.

\subsection{Figures, Tables and Schemes}

All figures and tables should be cited in the main text as Figure~\ref{fig1}, Table~\ref{tab1}, etc.

\begin{figure}[H]
\isPreprints{\centering}{} % Only used for preprints
\includegraphics[width=7 cm]{Definitions/logo-mdpi}
\caption{This is a figure. Schemes follow the same formatting.\label{fig1}}
\end{figure}   
\unskip

\begin{table}[H] 
%\tablesize{\small}
\caption{This is a table caption. Tables should be placed in the main text near to the first time they are~cited.\label{tab1}}
%\isPreprints{\centering}{} % Only used for preprints
\begin{tabularx}{\textwidth}{CCC}
\toprule
\textbf{Title 1}	& \textbf{Title 2}	& \textbf{Title 3}\\
\midrule
Entry 1		& Data			& Data\\
Entry 2		& Data			& Data \textsuperscript{1}\\
\bottomrule
\end{tabularx}

\noindent{\footnotesize{\textsuperscript{1} Tables may have a footer.}}
\end{table}

The text continues here (Figure~\ref{fig2} and Table~\ref{tab2}).

% Example of a figure that spans the whole page width and with subfigures. The same concept works for tables, too.
\begin{figure}[H]
\centering
\isPreprints{}{% This command is only used for ``preprints''.
\begin{adjustwidth}{-\extralength}{0cm}
} % If the paper is ``preprints'', please uncomment this parenthesis.
\subfloat[\centering]{\includegraphics[width=7.7cm]{Definitions/logo-mdpi}}
\hfill
\subfloat[\centering]{\includegraphics[width=7.7cm]{Definitions/logo-mdpi}}\\
\subfloat[\centering]{\includegraphics[width=7.7cm]{Definitions/logo-mdpi}}
\hfill
\subfloat[\centering]{\includegraphics[width=7.7cm]{Definitions/logo-mdpi}}
\isPreprints{}{% This command is only used for ``preprints''.
\end{adjustwidth}
} % If the paper is ``preprints'', please uncomment this parenthesis.
\caption{This is a wide figure. Schemes follow the same formatting. If there are multiple panels, they should be listed as: (\textbf{a}) Description of what is contained in the first panel. (\textbf{b}) Description of what is contained in the second panel. (\textbf{c}) Description of what is contained in the third panel. (\textbf{d}) Description of what is contained in the fourth panel. Figures should be placed in the main text near to the first time they are cited. A caption on a single line should be centered.\label{fig2}}
\end{figure} 

\begin{table}[H]
\caption{This is a wide table.\label{tab2}}
\isPreprints{\centering}{% This command is only used for ``preprints''.
	\begin{adjustwidth}{-\extralength}{0cm}
} % If the paper is ``preprints'', please uncomment this parenthesis.
\isPreprints{\begin{tabularx}{\textwidth}{CCCC}}{% This command is only used for ``preprints''.
		\begin{tabularx}{\fulllength}{CCCC}
} % If the paper is ``preprints'', please uncomment this parenthesis.
			\toprule
			\textbf{Title 1}	& \textbf{Title 2}	& \textbf{Title 3}     & \textbf{Title 4}\\
			\midrule
\multirow[m]{3}{*}{Entry 1 *}	& Data			& Data			& Data\\
			  	                   & Data			& Data			& Data\\
			             	      & Data			& Data			& Data\\
                   \midrule
\multirow[m]{3}{*}{Entry 2}    & Data			& Data			& Data\\
			  	                  & Data			& Data			& Data\\
			             	     & Data			& Data			& Data\\
			\bottomrule
		\end{tabularx}
		\isPreprints{}{% This command is only used for ``preprints''.
	\end{adjustwidth}
} % If the paper is ``preprints'', please uncomment this parenthesis.
	\noindent{\footnotesize{* Tables may have a footer.}}
\end{table}

%\begin{listing}[H]
%\caption{Title of the listing}
%\rule{\columnwidth}{1pt}
%\raggedright Text of the listing. In font size footnotesize, small, or normalsize. Preferred format: left aligned and single spaced. Preferred border format: top border line and bottom border line.
%\rule{\columnwidth}{1pt}
%\end{listing}

Text.

Text.

\subsection{Formatting of Mathematical Components}

This is the example 1 of equation:
\begin{linenomath}
\begin{equation}
a = 1,
\end{equation}
\end{linenomath}
the text following an equation need not be a new paragraph. Please punctuate equations as regular text.
%% If the documentclass option "submit" is chosen, please insert a blank line before and after any math environment (equation and eqnarray environments). This ensures correct linenumbering. The blank line should be removed when the documentclass option is changed to "accept" because the text following an equation should not be a new paragraph.

This is the example 2 of equation:
\isPreprints{}{% This command is only used for ``preprints''.
\begin{adjustwidth}{-\extralength}{0cm}
} % If the paper is ``preprints'', please uncomment this parenthesis.
\begin{equation}
a = b + c + d + e + f + g + h + i + j + k + l + m + n + o + p + q + r + s + t + u + v + w + x + y + z
\end{equation}
\isPreprints{}{% This command is only used for ``preprints''.
\end{adjustwidth}
} % If the paper is ``preprints'', please uncomment this parenthesis.

%% Example of a page in landscape format (with table and table footnote).
%\startlandscape
%\begin{table}[H] %% Table in wide page
%%\isPreprints{\centering}{} % This command is only used for ``preprints''.
%\caption{This is a very wide table.\label{tab3}}
%	\begin{tabularx}{\textwidth}{CCCC}
%		\toprule
%		\textbf{Title 1}	& \textbf{Title 2}	& \textbf{Title 3}	& \textbf{Title 4}\\
%		\midrule
%		Entry 1		& Data			& Data			& This cell has some longer content that runs over two lines.\\
%		Entry 2		& Data			& Data			& Data\textsuperscript{1}\\
%		\bottomrule
%	\end{tabularx}
%%\isPreprints{}{% This command is only used for ``preprints''.
%	\begin{adjustwidth}{+\extralength}{0cm}
%%} % If the paper is ``preprints'', please uncomment this parenthesis.
%		\noindent\footnotesize{\textsuperscript{1} This is a table footnote.}
%%\isPreprints{}{% This command is only used for ``preprints''.
%	\end{adjustwidth}
%%} % If the paper is ``preprints'', please uncomment this parenthesis.
%\end{table}
%\finishlandscape


Please punctuate equations as regular text. Theorem-type environments (including propositions, lemmas, corollaries etc.) can be formatted as follows:
%% Example of a theorem:
\begin{Theorem}
Example text of a theorem.
\end{Theorem}

The text continues here. Proofs must be formatted as follows:

%% Example of a proof:
\begin{proof}[Proof of Theorem 1]
Text of the proof. Note that the phrase ``of Theorem 1'' is optional if it is clear which theorem is being referred to.
\end{proof}
The text continues here.

%%%%%%%%%%%%%%%%%%%%%%%%%%%%%%%%%%%%%%%%%%
\section{Discussion}

Authors should discuss the results and how they can be interpreted from the perspective of previous studies and of the working hypotheses. The findings and their implications should be discussed in the broadest context possible. Future research directions may also be highlighted.

%%%%%%%%%%%%%%%%%%%%%%%%%%%%%%%%%%%%%%%%%%
\section{Conclusions}

This section is not mandatory, but can be added to the manuscript if the discussion is unusually long or complex.

%%%%%%%%%%%%%%%%%%%%%%%%%%%%%%%%%%%%%%%%%%
\section{Patents}

This section is not mandatory, but may be added if there are patents resulting from the work reported in this manuscript.

%%%%%%%%%%%%%%%%%%%%%%%%%%%%%%%%%%%%%%%%%%
\vspace{6pt} 

%%%%%%%%%%%%%%%%%%%%%%%%%%%%%%%%%%%%%%%%%%
%% optional
%\supplementary{The following supporting information can be downloaded at:  \linksupplementary{s1}, Figure S1: title; Table S1: title; Video S1: title.}

% Only for journal Methods and Protocols:
% If you wish to submit a video article, please do so with any other supplementary material.
% \supplementary{The following supporting information can be downloaded at: \linksupplementary{s1}, Figure S1: title; Table S1: title; Video S1: title. A supporting video article is available at doi: link.}

% Only used for preprtints:
% \supplementary{The following supporting information can be downloaded at the website of this paper posted on \href{https://www.preprints.org/}{Preprints.org}.}

% Only for journal Hardware:
% If you wish to submit a video article, please do so with any other supplementary material.
% \supplementary{The following supporting information can be downloaded at: \linksupplementary{s1}, Figure S1: title; Table S1: title; Video S1: title.\vspace{6pt}\\
%\begin{tabularx}{\textwidth}{lll}
%\toprule
%\textbf{Name} & \textbf{Type} & \textbf{Description} \\
%\midrule
%S1 & Python script (.py) & Script of python source code used in XX \\
%S2 & Text (.txt) & Script of modelling code used to make Figure X \\
%S3 & Text (.txt) & Raw data from experiment X \\
%S4 & Video (.mp4) & Video demonstrating the hardware in use \\
%... & ... & ... \\
%\bottomrule
%\end{tabularx}
%}

%%%%%%%%%%%%%%%%%%%%%%%%%%%%%%%%%%%%%%%%%%
\authorcontributions{For research articles with several authors, a short paragraph specifying their individual contributions must be provided. The following statements should be used ``Conceptualization, X.X. and Y.Y.; methodology, X.X.; software, X.X.; validation, X.X., Y.Y. and Z.Z.; formal analysis, X.X.; investigation, X.X.; resources, X.X.; data curation, X.X.; writing---original draft preparation, X.X.; writing---review and editing, X.X.; visualization, X.X.; supervision, X.X.; project administration, X.X.; funding acquisition, Y.Y. All authors have read and agreed to the published version of the manuscript.'', please turn to the  \href{http://img.mdpi.org/data/contributor-role-instruction.pdf}{CRediT taxonomy} for the term explanation. Authorship must be limited to those who have contributed substantially to the work~reported.}

\funding{Please add: ``This research received no external funding'' or ``This research was funded by NAME OF FUNDER grant number XXX.'' and  and ``The APC was funded by XXX''. Check carefully that the details given are accurate and use the standard spelling of funding agency names at \url{https://search.crossref.org/funding}, any errors may affect your future funding.}

\institutionalreview{In this section, you should add the Institutional Review Board Statement and approval number, if relevant to your study. You might choose to exclude this statement if the study did not require ethical approval. Please note that the Editorial Office might ask you for further information. Please add “The study was conducted in accordance with the Declaration of Helsinki, and approved by the Institutional Review Board (or Ethics Committee) of NAME OF INSTITUTE (protocol code XXX and date of approval).” for studies involving humans. OR “The animal study protocol was approved by the Institutional Review Board (or Ethics Committee) of NAME OF INSTITUTE (protocol code XXX and date of approval).” for studies involving animals. OR “Ethical review and approval were waived for this study due to REASON (please provide a detailed justification).” OR “Not applicable” for studies not involving humans or animals.}

\informedconsent{Any research article describing a study involving humans should contain this statement. Please add ``Informed consent was obtained from all subjects involved in the study.'' OR ``Patient consent was waived due to REASON (please provide a detailed justification).'' OR ``Not applicable'' for studies not involving humans. You might also choose to exclude this statement if the study did not involve humans.

Written informed consent for publication must be obtained from participating patients who can be identified (including by the patients themselves). Please state ``Written informed consent has been obtained from the patient(s) to publish this paper'' if applicable.}

\dataavailability{We encourage all authors of articles published in MDPI journals to share their research data. In this section, please provide details regarding where data supporting reported results can be found, including links to publicly archived datasets analyzed or generated during the study. Where no new data were created, or where data is unavailable due to privacy or ethical restrictions, a statement is still required. Suggested Data Availability Statements are available in section ``MDPI Research Data Policies'' at \url{https://www.mdpi.com/ethics}.} 

% Only for journal Drones
%\durcstatement{Current research is limited to the [please insert a specific academic field, e.g., XXX], which is beneficial [share benefits and/or primary use] and does not pose a threat to public health or national security. Authors acknowledge the dual-use potential of the research involving xxx and confirm that all necessary precautions have been taken to prevent potential misuse. As an ethical responsibility, authors strictly adhere to relevant national and international laws about DURC. Authors advocate for responsible deployment, ethical considerations, regulatory compliance, and transparent reporting to mitigate misuse risks and foster beneficial outcomes.}

% Only for journal Nursing Reports
%\publicinvolvement{Please describe how the public (patients, consumers, carers) were involved in the research. Consider reporting against the GRIPP2 (Guidance for Reporting Involvement of Patients and the Public) checklist. If the public were not involved in any aspect of the research add: ``No public involvement in any aspect of this research''.}
%
%% Only for journal Nursing Reports
%\guidelinesstandards{Please add a statement indicating which reporting guideline was used when drafting the report. For example, ``This manuscript was drafted against the XXX (the full name of reporting guidelines and citation) for XXX (type of research) research''. A complete list of reporting guidelines can be accessed via the equator network: \url{https://www.equator-network.org/}.}
%
%% Only for journal Nursing Reports
%\useofartificialintelligence{Please describe in detail any and all uses of artificial intelligence (AI) or AI-assisted tools used in the preparation of the manuscript. This may include, but is not limited to, language translation, language editing and grammar, or generating text. Alternatively, please state that “AI or AI-assisted tools were not used in drafting any aspect of this manuscript”.}

\acknowledgments{In this section you can acknowledge any support given which is not covered by the author contribution or funding sections. This may include administrative and technical support, or donations in kind (e.g., materials used for experiments).}

\conflictsofinterest{Declare conflicts of interest or state ``The authors declare no conflicts of interest.'' Authors must identify and declare any personal circumstances or interest that may be perceived as inappropriately influencing the representation or interpretation of reported research results. Any role of the funders in the design of the study; in the collection, analyses or interpretation of data; in the writing of the manuscript; or in the decision to publish the results must be declared in this section. If there is no role, please state ``The funders had no role in the design of the study; in the collection, analyses, or interpretation of data; in the writing of the manuscript; or in the decision to publish the results''.} 

%%%%%%%%%%%%%%%%%%%%%%%%%%%%%%%%%%%%%%%%%%
%% Optional

%% Only for journal Encyclopedia
%\entrylink{The Link to this entry published on the encyclopedia platform.}

\abbreviations{Abbreviations}{
The following abbreviations are used in this manuscript:\\

\noindent 
\begin{tabular}{@{}ll}
MDPI & Multidisciplinary Digital Publishing Institute\\
DOAJ & Directory of open access journals\\
TLA & Three letter acronym\\
LD & Linear dichroism
\end{tabular}
}

%%%%%%%%%%%%%%%%%%%%%%%%%%%%%%%%%%%%%%%%%%
%% Optional
\appendixtitles{no} % Leave argument "no" if all appendix headings stay EMPTY (then no dot is printed after "Appendix A"). If the appendix sections contain a heading then change the argument to "yes".
\appendixstart
\appendix
\section[\appendixname~\thesection]{}
\subsection[\appendixname~\thesubsection]{}
The appendix is an optional section that can contain details and data supplemental to the main text---for example, explanations of experimental details that would disrupt the flow of the main text but nonetheless remain crucial to understanding and reproducing the research shown; figures of replicates for experiments of which representative data are shown in the main text can be added here if brief, or as Supplementary Data. Mathematical proofs of results not central to the paper can be added as an appendix.

\begin{table}[H] 
\caption{This is a table caption.\label{tab5}}
%\newcolumntype{C}{>{\centering\arraybackslash}X}
\begin{tabularx}{\textwidth}{CCC}
\toprule
\textbf{Title 1}	& \textbf{Title 2}	& \textbf{Title 3}\\
\midrule
Entry 1		& Data			& Data\\
Entry 2		& Data			& Data\\
\bottomrule
\end{tabularx}
\end{table}

\section[\appendixname~\thesection]{}
All appendix sections must be cited in the main text. In the appendices, Figures, Tables, etc. should be labeled, starting with ``A''---e.g., Figure A1, Figure A2, etc.

%%%%%%%%%%%%%%%%%%%%%%%%%%%%%%%%%%%%%%%%%%
\isPreprints{}{% This command is only used for ``preprints''.
\begin{adjustwidth}{-\extralength}{0cm}
} % If the paper is ``preprints'', please uncomment this parenthesis.
%\printendnotes[custom] % Un-comment to print a list of endnotes

\reftitle{References}

% Please provide either the correct journal abbreviation (e.g. according to the “List of Title Word Abbreviations” http://www.issn.org/services/online-services/access-to-the-ltwa/) or the full name of the journal.
% Citations and References in Supplementary files are permitted provided that they also appear in the reference list here. 

%=====================================
% References, variant A: external bibliography
%=====================================
% \bibliography{your_external_BibTeX_file}

%=====================================
% References, variant B: internal bibliography
%=====================================

% ACS format
\isAPAandChicago{}{%
\begin{thebibliography}{999}
% Reference 1
\bibitem[Author1(year)]{ref-journal}
Author~1, T. The title of the cited article. {\em Journal Abbreviation} {\bf 2008}, {\em 10}, 142--149.
% Reference 2
\bibitem[Author2(year)]{ref-book1}
Author~2, L. The title of the cited contribution. In {\em The Book Title}; Editor 1, F., Editor 2, A., Eds.; Publishing House: City, Country, 2007; pp. 32--58.
% Reference 3
\bibitem[Author1 and Author2 (year)]{ref-book2}
Author 1, A.; Author 2, B. \textit{Book Title}, 3rd ed.; Publisher: Publisher Location, Country, 2008; pp. 154--196.
% Reference 4
\bibitem[Author4(year)]{ref-unpublish}
Author 1, A.B.; Author 2, C. Title of Unpublished Work. \textit{Abbreviated Journal Name} year, \textit{phrase indicating stage of publication (submitted; accepted; in press)}.
% Reference 5
\bibitem[Author8(year)]{ref-url}
Title of Site. Available online: URL (accessed on Day Month Year).
% Reference 6
\bibitem[Author6(year)]{ref-proceeding}
Author 1, A.B.; Author 2, C.D.; Author 3, E.F. Title of presentation. In Proceedings of the Name of the Conference, Location of Conference, Country, Date of Conference (Day Month Year); Abstract Number (optional), Pagination (optional).
% Reference 7
\bibitem[Author7(year)]{ref-thesis}
Author 1, A.B. Title of Thesis. Level of Thesis, Degree-Granting University, Location of University, Date of Completion.
\end{thebibliography}
}

% Chicago format (Used for journal: arts, genealogy, histories, humanities, jintelligence, laws, literature, religions, risks, socsci)
\isChicagoStyle{%
\begin{thebibliography}{999}
% Reference 1
\bibitem[Aranceta-Bartrina(1999a)]{ref-journal}
Aranceta-Bartrina, Javier. 1999a. Title of the cited article. \textit{Journal Title} 6: 100--10.
% Reference 2
\bibitem[Aranceta-Bartrina(1999b)]{ref-book1}
Aranceta-Bartrina, Javier. 1999b. Title of the chapter. In \textit{Book Title}, 2nd ed. Edited by Editor 1 and Editor 2. Publication place: Publisher, vol. 3, pp. 54–96.
% Reference 3
\bibitem[Baranwal and Munteanu {[1921]}(1955)]{ref-book2}
Baranwal, Ajay K., and Costea Munteanu. 1955. \textit{Book Title}. Publication place: Publisher, pp. 154--96. First published 1921 (op-tional).
% Reference 4
\bibitem[Berry and Smith(1999)]{ref-thesis}
Berry, Evan, and Amy M. Smith. 1999. Title of Thesis. Level of Thesis, Degree-Granting University, City, Country. Identifi-cation information (if available).
% Reference 5
\bibitem[Cojocaru et al.(1999)]{ref-unpublish}
Cojocaru, Ludmila, Dragos Constatin Sanda, and Eun Kyeong Yun. 1999. Title of Unpublished Work. \textit{Journal Title}, phrase indicating stage of publication.
% Reference 6
\bibitem[Driver et al.(2000)]{ref-proceeding}
Driver, John P., Steffen Rohrs, and Sean Meighoo. 2000. Title of Presentation. In \textit{Title of the Collected Work} (if available). Paper presented at Name of the Conference, Location of Conference, Date of Conference.
% Reference 7
\bibitem[Harwood(2008)]{ref-url}
Harwood, John. 2008. Title of the cited article. Available online: URL (accessed on Day Month Year).
\end{thebibliography}
}{}

% APA format (Used for journal: admsci, behavsci, businesses, econometrics, economies, education, ejihpe, games, humans, ijfs, journalmedia, jrfm, languages, psycholint, publications, tourismhosp, youth)
\isAPAStyle{%
\begin{thebibliography}{999}
% Reference 1
\bibitem[\protect\citeauthoryear{Azikiwe \BBA\ Bello}{{2020a}}]{ref-journal}
Azikiwe, H., \& Bello, A. (2020a). Title of the cited article. \textit{Journal Title}, \textit{Volume}(Issue), 
Firstpage--Lastpage/Article Number.
% Reference 2
\bibitem[\protect\citeauthoryear{Azikiwe \BBA\ Bello}{{2020b}}]{ref-book1}
Azikiwe, H., \& Bello, A. (2020b). \textit{Book title}. Publisher Name.
% Reference 3
\bibitem[Davison(1623/2019)]{ref-book2}
Davison, T. E. (2019). Title of the book chapter. In A. A. Editor (Ed.), \textit{Title of the book: Subtitle} 
(pp. Firstpage--Lastpage). Publisher Name. (Original work published 1623) (Optional).
% Reference 4
\bibitem[Fistek et al.(2017)]{ref-proceeding}
Fistek, A., Jester, E., \& Sonnenberg, K. (2017, Month Day). Title of contribution [Type of contribution]. Conference Name, Conference City, Conference Country.
% Reference 5
\bibitem[Hutcheson(2012)]{ref-thesis}
Hutcheson, V. H. (2012). \textit{Title of the thesis} [XX Thesis, Name of Institution Awarding the Degree].
% Reference 6
\bibitem[Lippincott \& Poindexter(2019)]{ref-unpublish}
Lippincott, T., \& Poindexter, E. K. (2019). \textit{Title of the unpublished manuscript} [Unpublished manuscript/Manuscript in prepara-tion/Manuscript submitted for publication]. Department Name, Institution Name.
% Reference 7
\bibitem[Harwood(2008)]{ref-url}
Harwood, J. (2008). \textit{Title of the cited article}. Available online: URL (accessed on Day Month Year).
\end{thebibliography}
}{}

% If authors have biography, please use the format below
%\section*{Short Biography of Authors}
%\bio
%{\raisebox{-0.35cm}{\includegraphics[width=3.5cm,height=5.3cm,clip,keepaspectratio]{Definitions/author1.pdf}}}
%{\textbf{Firstname Lastname} Biography of first author}
%
%\bio
%{\raisebox{-0.35cm}{\includegraphics[width=3.5cm,height=5.3cm,clip,keepaspectratio]{Definitions/author2.jpg}}}
%{\textbf{Firstname Lastname} Biography of second author}

% For the MDPI journals use author-date citation, please follow the formatting guidelines on http://www.mdpi.com/authors/references
% To cite two works by the same author: \citeauthor{ref-journal-1a} (\citeyear{ref-journal-1a}, \citeyear{ref-journal-1b}). This produces: Whittaker (1967, 1975)
% To cite two works by the same author with specific pages: \citeauthor{ref-journal-3a} (\citeyear{ref-journal-3a}, p. 328; \citeyear{ref-journal-3b}, p.475). This produces: Wong (1999, p. 328; 2000, p. 475)

%%%%%%%%%%%%%%%%%%%%%%%%%%%%%%%%%%%%%%%%%%
%% for journal Sci
%\reviewreports{\\
%Reviewer 1 comments and authors’ response\\
%Reviewer 2 comments and authors’ response\\
%Reviewer 3 comments and authors’ response
%}
%%%%%%%%%%%%%%%%%%%%%%%%%%%%%%%%%%%%%%%%%%
\PublishersNote{}
\isPreprints{}{% This command is only used for ``preprints''.
\end{adjustwidth}
} % If the paper is ``preprints'', please uncomment this parenthesis.
\pagebreak

\reftitle{References}

% Please provide either the correct journal abbreviation (e.g. according to the “List of Title Word Abbreviations” http://www.issn.org/services/online-services/access-to-the-ltwa/) or the full name of the journal.
% Citations and References in Supplementary files are permitted provided that they also appear in the reference list here. 

%=====================================
% References, variant A: external bibliography
%=====================================
% \bibliography{your_external_BibTeX_file}

%=====================================
% References, variant B: internal bibliography
%=====================================

% ACS format

\begin{thebibliography}{999}

%\begin{thebibliography}{99}

\bibitem{MacLane1998}
S.~Mac~Lane.
\newblock \emph{Categories for the Working Mathematician}, 2nd~edition.
\newblock Graduate Texts in Mathematics 5. Springer, 1998.

\bibitem{Awodey2010}
S.~Awodey.
\newblock \emph{Category Theory}, 2nd~edition.
\newblock Oxford Logic Guides 52. Oxford University Press, 2010.

\bibitem{Borceux1994}
F.~Borceux.
\newblock \emph{Handbook of Categorical Algebra, Volume~1: Basic Category Theory}.
\newblock Encyclopedia of Mathematics and its Applications 50. Cambridge University Press, 1994.

\bibitem{BaaderNipkow1998}
F.~Baader and T.~Nipkow.
\newblock \emph{Term Rewriting and All That}.
\newblock Cambridge University Press, 1998.

\bibitem{Ji2023Survey}
Z.~Ji, N.~Lee, R.~Frieske, T.~Yu, D.~Su, Y.~Xu, E.~Ishii, Y.~Bang, D.~Chen, W.~Dai, H.~S.~Chan, A.~Madotto, and P.~Fung.
\newblock Survey of hallucination in natural language generation.
\newblock \emph{ACM Computing Surveys}, 55(12):Article~248, 2023.

\bibitem{Farquhar2024SemanticEntropy}
S.~Farquhar, J.~Kossen, L.~Kuhn, and Y.~Gal.
\newblock Detecting hallucinations in large language models using semantic entropy.
\newblock \emph{Nature}, 630(8017):625--630, 2024.

\bibitem{Kadavath2022Know}
S.~Kadavath, T.~Conerly, A.~Askell, T.~Henighan, D.~Drain, E.~Perez, N.~Schiefer, Z.~Hatfield-Dodds, N.~DasSarma, E.~Tran-Johnson, S.~Johnston, S.~El-Showk, A.~Jones, N.~Elhage, T.~Hume, A.~Chen, Y.~Bai, S.~Bowman, S.~Fort, D.~Ganguli, D.~Hernandez, J.~Jacobson, J.~Kernion, S.~Kravec, L.~Lovitt, K.~Ndousse, C.~Olsson, S.~Ringer, D.~Amodei, T.~Brown, J.~Clark, N.~Joseph, B.~Mann, S.~McCandlish, C.~Olah, and J.~Kaplan.
\newblock Language models (mostly) know what they know.
\newblock \emph{arXiv preprint} arXiv:2207.05221, 2022.

\bibitem{Kalai2025WhyHallucinate}
A.~T.~Kalai, O.~Nachum, S.~S.~Vempala, and E.~Zhang.
\newblock Why language models hallucinate.
\newblock \emph{arXiv preprint} arXiv:2509.04664, 2025.

\bibitem{BertotCasteran2004CoqArt}
Y.~Bertot and P.~Cast\'eran.
\newblock \emph{Interactive Theorem Proving and Program Development: Coq'Art: The Calculus of Inductive Constructions}.
\newblock Texts in Theoretical Computer Science. An EATCS Series. Springer, 2004.

\bibitem{Nipkow2002IsabelleHOL}
T.~Nipkow, L.~C.~Paulson, and M.~Wenzel.
\newblock \emph{Isabelle/HOL: A Proof Assistant for Higher-Order Logic}.
\newblock Lecture Notes in Computer Science 2283. Springer, 2002.

\bibitem{deMoura2015Lean}
L.~de~Moura, S.~Kong, J.~Avigad, F.~van~Doorn, and J.~von~Raumer.
\newblock The {L}ean theorem prover (system description).
\newblock In A.~P.~Felty and A.~Middeldorp, editors, \emph{Automated Deduction -- CADE-25}, volume~9195 of \emph{Lecture Notes in Computer Science}, pages 378--388. Springer, 2015.

\bibitem{Jacobs2001Categorical}
B.~Jacobs.
\newblock \emph{Categorical Logic and Type Theory}.
\newblock Studies in Logic and the Foundations of Mathematics 141. Elsevier, 2001.

\bibitem{Schick2023Toolformer}
T.~Schick, J.~Dwivedi-Yu, R.~Dess\`i, R.~Raileanu, M.~Lomeli, L.~Zettlemoyer, N.~Cancedda, and T.~Scialom.
\newblock Toolformer: language models can teach themselves to use tools.
\newblock \emph{arXiv preprint} arXiv:2302.04761, 2023.

\bibitem{Yao2022ReAct}
S.~Yao, J.~Zhao, D.~Yu, N.~Du, I.~Shafran, K.~Narasimhan, and Y.~Cao.
\newblock ReAct: synergizing reasoning and acting in language models.
\newblock \emph{arXiv preprint} arXiv:2210.03629, 2022.

\bibitem{Askell2021Alignment}
A.~Askell, Y.~Bai, A.~Chen, D.~Drain, D.~Ganguli, T.~Henighan, A.~Jones, N.~Joseph, B.~Mann, N.~DasSarma, N.~Elhage, Z.~Hatfield-Dodds, D.~Hernandez, J.~Kernion, K.~Ndousse, C.~Olsson, D.~Amodei, T.~Brown, J.~Clark, S.~McCandlish, C.~Olah, and J.~Kaplan.
\newblock A general language assistant as a laboratory for alignment.
\newblock \emph{arXiv preprint} arXiv:2112.00861, 2021.

\bibitem{Shao2025SciSciGPT}
E.~Shao, Y.~Wang, Y.~Qian, Z.~Pan, H.~Liu, and D.~Wang.
\newblock SciSciGPT: advancing human--AI collaboration in the science of science.
\newblock \emph{Nature Computational Science}, 5(11), 2025. Also available as \emph{arXiv preprint} arXiv:2504.05559.

\bibitem{Vaccaro2024WhenUseful}
M.~Vaccaro, A.~Almaatouq, and T.~W.~Malone.
\newblock When combinations of humans and AI are useful: a systematic review and meta-analysis.
\newblock \emph{Nature Human Behaviour}, 8(12):2293--2303, 2024.

\bibitem{Zheng2025LLMDiscovery}
T.~Zheng, Z.~Deng, H.~T.~Tsang, W.~Wang, J.~Bai, Z.~Wang, and Y.~Song.
\newblock From automation to autonomy: a survey on large language models in scientific discovery.
\newblock In \emph{Proceedings of the 2025 Conference on Empirical Methods in Natural Language Processing}, pages 17744--17761. Association for Computational Linguistics, 2025.

\bibitem{Zhao2025HumanAISciDiscovery}
S.~Zhao, Y.~Liu, J.~Wan, T.~Tang, and X.~Li.
\newblock Human--AI collaboration for scientific discovery.
\newblock In D.~Wu and S.~Liang, editors, \emph{Human--AI Interaction and Collaboration}, pages 239--267. Cambridge University Press, 2025.

\bibitem{Channing2025AISocialProblem}
G.~Channing and A.~Ghosh.
\newblock AI for scientific discovery is a social problem.
\newblock \emph{arXiv preprint} arXiv:2509.06580, 2025.
% --- Hallucination / LLM surveys ---

\bibitem{Zhang2023Siren}
Y.~Zhang, Y.~Li, L.~Cui, D.~Cai, L.~Liu, T.~Fu, X.~Huang, E.~Zhao, Y.~Zhang,
Y.~Chen, L.~Wang, A.~T.~Luu, W.~Bi, F.~Shi, and S.~Shi.
Siren's song in the AI ocean: a survey on hallucination in large language models.
arXiv preprint arXiv:2309.01219, 2023.

\bibitem{Huang2025HallucinationSurvey}
L.~Huang, W.~Yu, W.~Ma, W.~Zhong, Z.~Feng, H.~Wang, Q.~Chen, W.~Peng,
X.~Feng, B.~Qin, and T.~Liu.
A survey on hallucination in large language models: principles, taxonomy, challenges, and open questions.
ACM Transactions on Information Systems, 43(2):1--55, 2025.

% --- AI reliance / appropriate reliance / trust ---

\bibitem{Eckhardt2024Reliance}
S.~Eckhardt, N.~K{\"u}hl, M.~Dolata, and G.~Schwabe.
A survey of AI reliance.
CoRR, abs/2408.03948, 2024.

\bibitem{Schemmer2023Appropriate}
M.~Schemmer, N.~K{\"u}hl, C.~Benz, A.~Bartos, and G.~Satzger.
Appropriate reliance on AI advice: conceptualization and the effect of explanations.
In \emph{Proceedings of the 28th International Conference on Intelligent User Interfaces (IUI~'23)},
pages 1--21. ACM, 2023.

\bibitem{Visser2025Trust}
R.~Visser, T.~M.~Peters, I.~Scharlau, and B.~Hammer.
Trust, distrust, and appropriate reliance in (X)AI: a conceptual clarification of user trust and survey of its empirical evaluation.
Cognitive Systems Research, 82:101357, 2025.

% --- Social epistemology / scientific networks ---

\bibitem{Zollman2013Network}
K.~J.~S.~Zollman.
Network epistemology: communication in epistemic communities.
Philosophy Compass, 8(1):15--27, 2013.

\bibitem{Grim2013ScientificNetworks}
P.~Grim, D.~J.~Singer, S.~Fisher, A.~Bramson, W.~J.~Berger, C.~Reade, C.~Flocken, and A.~Sales.
Scientific networks on data landscapes: question difficulty, epistemic success, and convergence.
Episteme, 10(4):441--464, 2013.

% --- AI governance / frontier policies ---

\bibitem{METRCommonElements}
METR.
Common elements of frontier AI safety policies.
Technical report, Model Evaluation and Threat Research (METR), 2024.
Available at \texttt{https://metr.org/common-elements}.

\bibitem{Polanyi1966}
M.~Polanyi.
\newblock \emph{The Tacit Dimension}.
\newblock Routledge \& Kegan Paul, London, 1966.

\bibitem{Nonaka1994}
I.~Nonaka.
\newblock A dynamic theory of organizational knowledge creation.
\newblock \emph{Organization Science}, 5(1):14--37, 1994.

\bibitem{NonakaTakeuchi1995}
I.~Nonaka and H.~Takeuchi.
\newblock \emph{The Knowledge-Creating Company: How Japanese Companies Create the Dynamics of Innovation}.
\newblock Oxford University Press, New York, 1995.

\bibitem{EUAIAct2024}
European Parliament and Council of the European Union.
\newblock Regulation (EU) 2024/1689 laying down harmonised rules on artificial intelligence (AI Act).
\newblock \emph{Official Journal of the European Union}, L~1689, 2024.
% Optional: add URL if your style allows it:
% \newblock Available at \url{https://eur-lex.europa.eu/legal-content/EN/TXT/?uri=OJ:L:2024:1689:TOC}.

\bibitem{Amershi2019}
S.~Amershi, D.~Weld, M.~Vorvoreanu, A.~Fourney, B.~Nushi, P.~Collisson, H.~Suh,
S.~Iqbal, P.~N.~Bennett, K.~Inkpen, and J.~Teevan.
\newblock Guidelines for human--AI interaction.
\newblock In \emph{Proceedings of the 2019 CHI Conference on Human Factors in Computing Systems}, pages 1--13.
\newblock ACM, 2019.
\bibitem{brown_effect_automation_aviation_2017}
J.~P. Brown.
\newblock The effect of automation on human factors in aviation.
\newblock \emph{The Journal of Instrumentation, Automation and Systems}, 3(2):31--46, 2016.
% DOI: 10.21535/jias.v3i2.916

\bibitem{okine_evolution_hf_aviation_2025}
E.~A. Okine, E.~Zarei, B.~J. Roggow, and N.~Dehghan.
\newblock Evolution of human factors research in aviation safety: a systematic review and bibliometric analysis of the intellectual structure.
\newblock \emph{Journal of Safety Science and Resilience}, 7(1):100249, 2026.
% DOI: 10.1016/j.jnlssr.2025.100249

\bibitem{singh_automation_complacency_2024}
I.~L. Singh, R.~Molloy, and R.~Parasuraman.
\newblock Automation-induced ``complacency'': development of the Complacency-Potential Rating Scale.
\newblock \emph{The International Journal of Aviation Psychology}, 3(2):111--122, 1993.
% Key reused to match the LaTeX patch; classic complacency measure

\bibitem{natali_ai_deskilling_medicine_2025}
C.~Natali, L.~Marconi, L.~D.~D. Duran, and F.~Cabitza.
\newblock AI-induced deskilling in medicine: a mixed-method review and research agenda for healthcare and beyond.
\newblock \emph{Artificial Intelligence Review}, 58:356, 2025.
% DOI: 10.1007/s10462-025-11352-1

\bibitem{berzin_preserving_clinical_skills_2025}
T.~M. Berzin and E.~J. Topol.
\newblock Preserving clinical skills in the age of AI assistance.
\newblock \emph{The Lancet}, 406(10513):1719, 2025.
% DOI: 10.1016/S0140-6736(25)02075-6

\bibitem{sunday_behavioral_impacts_diagnostics_2025}
O.~Sunday.
\newblock Behavioral impacts of AI reliance in diagnostics: balancing automation with skill retention.
\newblock \emph{Epidemiology \& Health Data Insights}, 1(3):ehdi011, 2025.
% DOI: 10.63946/ehdi/16894

\bibitem{fatima_ai_triage_deskilling_2025}
L.~Fatima, A.~A. Khan, R.~Darwesh, and M.~un Nissa.
\newblock AI triage for stroke imaging: promise vs risk of deskilling radiology trainees.
\newblock \emph{INNOVAPATH}, 1(8):2, 2025.
% DOI: 10.63501/0961v416

\bibitem{savardi_ai_radiology_training_2025}
M.~Savardi, A.~Signoroni, F.~Arrigoni, et~al.
\newblock Upskilling or deskilling? measurable role of an AI-supported training for radiology residents: a lesson from the pandemic.
\newblock \emph{Insights into Imaging}, 16(1):23, 2025.
% DOI: 10.1186/s13244-024-01893-4

\bibitem{gong_yang_google_effect_2024}
C.~Gong and Z.~Yang.
\newblock Google effects on memory: a meta-analytical review of the media effects of intensive Internet search behavior.
\newblock \emph{Frontiers in Public Health}, 12:1332030, 2024.
% DOI: 10.3389/fpubh.2024.1332030

\bibitem{benge_tech_use_cognitive_aging_2025}
J.~F. Benge and M.~K. Scullin.
\newblock A meta-analysis of technology use and cognitive aging.
\newblock \emph{Nature Human Behaviour}, 9(7):1405--1419, 2025.
% DOI: 10.1038/s41562-025-02159-9

\bibitem{gerlich_ai_tools_cognitive_offloading_2025}
M.~Gerlich.
\newblock AI tools in society: impacts on cognitive offloading and the future of critical thinking.
\newblock \emph{Societies}, 15(1):6, 2025.
% DOI: 10.3390/soc15010006

\bibitem{gerlich_ai_cognitive_implications_2025}
R.~Loga.
\newblock AI's cognitive implications: the decline of our thinking skills?
\newblock IE Center for Health and Well-Being Blog, IE University, 26 February 2025.
% Blog post synthesizing Gerlich (2025) and related work; keep only if you explicitly cite it

\bibitem{kosmyna_brain_on_chatgpt_2025}
N.~Kosmyna, E.~Hauptmann, Y.~T. Yuan, J.~Situ, X.-H. Liao, A.~V. Beresnitzky, I.~Braunstein, and P.~Maes.
\newblock Your brain on ChatGPT: accumulation of cognitive debt when using an AI assistant for essay writing task.
\newblock arXiv preprint arXiv:2506.08872, 2025.

\bibitem{zhao_genai_higher_order_thinking_2025}
Y.~Zhao, Y.~Yue, Z.~Sun, Q.~Jiang, and G.~Li.
\newblock Does generative artificial intelligence improve students' higher-order thinking? a meta-analysis based on 29 experiments and quasi-experiments.
\newblock \emph{Journal of Intelligence}, 13(12):160, 2025.
% DOI: 10.3390/jintelligence13120160

% Optional: add URL if your style allows it:
% \newblock Available at \url{https://www.ie.edu/center-for-health-and-well-being/blog/ais-cognitive-implications-the-decline-of-our-thinking-skills/}.
\bibitem{laban_llms_get_lost_multiturn_2025}
P.~Laban, H.~Hayashi, Y.~Zhou, and J.~Neville.
\newblock LLMs get lost in multi-turn conversation.
\newblock \emph{arXiv preprint arXiv:2505.06120}, 2025.

\bibitem{dongre_drift_no_more_context_equilibria_2025}
V.~Dongre, R.~A.~Rossi, V.~D.~Lai, D.~S.~Yoon, D.~Hakkani-T\"ur, and T.~Bui.
\newblock Drift no more? Context equilibria in multi-turn LLM interactions.
\newblock \emph{arXiv preprint arXiv:2510.07777}, 2025.

\bibitem{arike_evaluating_goal_drift_agents_2025}
R.~Arike, E.~Donoway, H.~Bartsch, and M.~Hobbhahn.
\newblock Technical report: evaluating goal drift in language model agents.
\newblock \emph{arXiv preprint arXiv:2505.02709}, 2025.

\bibitem{choi_identity_drift_llm_agents_2024}
J.~Choi, Y.~Hong, M.~Kim, and B.~Kim.
\newblock Examining identity drift in conversations of LLM agents.
\newblock \emph{arXiv preprint arXiv:2412.00804}, 2024.
\bibitem{nrc_defense_in_depth_glossary}
U.S. Nuclear Regulatory Commission.
\newblock Defense in depth.
\newblock NRC Glossary, n.d. (accessed 26 January 2026).
% https://www.nrc.gov/reading-rm/basic-ref/glossary/defense-in-depth

\bibitem{iaea_insag10_defence_in_depth_1996}
International Nuclear Safety Advisory Group (INSAG).
\newblock \emph{Defence in Depth in Nuclear Safety} (INSAG-10).
\newblock International Atomic Energy Agency, Vienna, 1996.
% IAEA Pub 1013; ISBN 9201025963.

\bibitem{reason_human_error_models_management_2000}
J.~Reason.
\newblock Human error: models and management.
\newblock \emph{BMJ}, 320(7237):768--770, 2000.
% DOI: 10.1136/bmj.320.7237.768

\bibitem{reason_latent_human_failures_complex_systems_1990}
J.~Reason.
\newblock The contribution of latent human failures to the breakdown of complex systems.
\newblock \emph{Philosophical Transactions of the Royal Society of London. Series B, Biological Sciences}, 327(1241):475--484, 1990.
% DOI: 10.1098/rstb.1990.0090

\end{thebibliography}

\end{document}

